% Generated display of complete source files; originals remain unchanged.
\begingroup
\lstset{commentstyle=\color[rgb]{0.12,0.34,0.24},stringstyle=\color[rgb]{0.15,0.30,0.35},keywordstyle=\color[rgb]{0.13,0.23,0.43}}
\subsection{\texttt{src/p1\_intersection.py}}
% SHA256 ef4e3d3c83f2cfbf878f9803b33bc8969e5e8fa172bcf384b106dab8e77cd0b5
\begin{lstlisting}[language=python,escapeinside={(*@}{@*)}]
"""B 题问题 1：示向度交会定位区域、区域直径与直径圆覆盖判定。

模型的严格表述
--------------
检测点 :math:`S_i=(x_i,y_i)`，该点测得示向度 :math:`\\theta_i`（度，x 轴正向逆时针，[0,360)），
示向度误差有界 :math:`|e_i|\\le\\varepsilon=1^\\circ`（题目附录 2 图 2：实线与相邻虚线夹角均为 1°）。
真源必在

    R = (*@\ensuremath{\cap}@*)_i W_i,   W_i = { X : |ang(X - S_i) - (*@\ensuremath{\theta}@*)_i| <= (*@\ensuremath{\epsilon}@*) }

W_i 是以 S_i 为顶点、张角 2(*@\ensuremath{\epsilon}@*) 的**角楔**，等价于两个半平面之交：

    W_i = { X : a(*@\ensuremath{{}^{-}}@*)_i·X + b(*@\ensuremath{{}^{-}}@*)_i <= 0 } (*@\ensuremath{\cap}@*) { X : a(*@\ensuremath{{}^{+}}@*)_i·X + b(*@\ensuremath{{}^{+}}@*)_i <= 0 }

因此 R 是 2k 个半平面的交 —— **凸多边形**，顶点数 (*@\ensuremath{\le}@*) 2k。

算法要点（含两处易错点的正确处理）
----------------------------------
1. **可行性 / 有界性**（先判空、再判无界，顺序不能反）
   * 空集：两两示向度可以相差 > 2(*@\ensuremath{\epsilon}@*)，但两个"错开"的角楔也可能没有公共点；
   * 无界：存在方向 d 使 |ang(d) - (*@\ensuremath{\theta}@*)_i| (*@\ensuremath{\le}@*) (*@\ensuremath{\epsilon}@*) 对所有 i 成立
     (*@\ensuremath{\Longleftrightarrow}@*) 所有示向度**两两夹角 (*@\ensuremath{\le}@*) 2(*@\ensuremath{\epsilon}@*)**；此时直径 = +∞，不能报有限值。
2. **区域顶点**：凸集的顶点必是两条起作用约束边界直线的交点，故枚举 2k 条边界直线的
   O(k²) 个交点、用**半平面判据** a·X + b (*@\ensuremath{\le}@*) 0 筛选即可（精确）。
   注意：**检测点本身也可能是顶点**（当它落在其它所有楔内时）。旧实现在这里用
   "候选点到检测点的距离 < 1e-9 则丢弃" 处理角度奇异，会把这种顶点误删 —— 见
   `region_by_vertices` 的半平面判据。
3. 提供两套**相互独立**的区域算法，供自检交叉验证：
   * `region_by_vertices`：顶点枚举 + 凸包（主算法，精确）；
   * `region_by_clipping`：Sutherland–Hodgman 逐个半平面裁剪大框（参考实现，可判无界）。
4. **直径**：顶点数 (*@\ensuremath{\le}@*) `DIAM_BRUTE_LIMIT` 时暴力枚举；否则旋转卡壳，两者在自检中互相校验。
5. **覆盖判定**：以直径 AB 为直径的圆 C，
   V (*@\ensuremath{\in}@*) C (*@\ensuremath{\Longleftrightarrow}@*) (*@\ensuremath{\angle}@*)AVB (*@\ensuremath{\ge}@*) 90°（Thales）(*@\ensuremath{\Longleftrightarrow}@*) |V - c| (*@\ensuremath{\le}@*) D/2（c 为 AB 中点）。
   一般**不能覆盖**；正确做法是用**最小包围圆**（半径 r* (*@\ensuremath{\in}@*) [D/2, D/√3]，Jung 定理）。
"""
from __future__ import annotations
import csv
import json
import math
import os
import sys

import numpy as np

BEARING_ERR = 1.0        # 示向度误差界（度）
ANG_TOL = 1e-9           # 角度判据容差（度）
DIST_TOL = 1e-6          # 半平面判据容差（米）
CLIP = 1e7               # 裁剪参考实现的外框半边长（米）
DIAM_BRUTE_LIMIT = 128   # 顶点数不超过该值时直径用暴力枚举


# --------------------------------------------------------------------------
# 基本几何
# --------------------------------------------------------------------------
def bearing(p, q):
    """p 处指向 q 的方位角（度，[0,360)）"""
    return math.degrees(math.atan2(q[1] - p[1], q[0] - p[0])) % 360.0


def ang_diff(a, b):
    """a-b 归一化到 (-180,180]"""
    return (a - b + 180.0) % 360.0 - 180.0


def wedge_halfplanes(p, theta, err=BEARING_ERR):
    """角楔 W = {X : |ang(X-p) - theta| <= err} 的两个半平面，形式 a·X + b <= 0。

    err < 90°，故 W 是两条边界直线所夹的尖楔。
    """
    p = np.asarray(p, float)
    tm = math.radians((theta - err) % 360.0)
    tp = math.radians((theta + err) % 360.0)
    dm = np.array([math.cos(tm), math.sin(tm)])
    dp = np.array([math.cos(tp), math.sin(tp)])
    # cross(dm, X-p) >= 0   <=>   a1·X + b1 <= 0
    a1 = np.array([dm[1], -dm[0]])
    b1 = float(-dm[1] * p[0] + dm[0] * p[1])
    # cross(dp, X-p) <= 0   <=>   a2·X + b2 <= 0
    a2 = np.array([-dp[1], dp[0]])
    b2 = float(dp[1] * p[0] - dp[0] * p[1])
    return [(a1, b1), (a2, b2)]


def region_is_unbounded(svds, err=BEARING_ERR):
    """区域无界判据：存在方向 d 与所有示向度夹角 <= err。

    返回一个这样的方向（度），有界时返回 None。
    仅在区域非空时有意义（先判空、再判无界）。
    """
    for th in svds:
        for cand in ((th - err) % 360.0, (th + err) % 360.0):
            if all(abs(ang_diff(cand, t)) <= err + ANG_TOL for t in svds):
                return cand
    return None


def _line_intersection(P, u, Q, v):
    """直线 P + t·u 与 Q + s·v 的交点，平行时返回 None"""
    det = u[0] * (-v[1]) - (-v[0]) * u[1]
    if abs(det) < 1e-14:
        return None
    r = Q - P
    t = (r[0] * (-v[1]) - (-v[0]) * r[1]) / det
    return P + t * u


def convex_hull(P):
    """Andrew monotone chain，逆时针，去掉共线点"""
    def half(pts):
        st = []
        for p in pts:
            while len(st) >= 2 and _cross2(st[-1] - st[-2], p - st[-2]) <= 0:
                st.pop()
            st.append(p)
        return st
    P = P[np.lexsort((P[:, 1], P[:, 0]))]
    lo = half(P)
    hi = half(P[::-1])
    return np.array(lo[:-1] + hi[:-1])


def _cross2(u, v):
    return float(u[0] * v[1] - u[1] * v[0])


def _dedup_hull(cand):
    """候选顶点去重 + 凸包，返回 (P, "empty"|"degenerate"|"polygon")"""
    if len(cand) == 0:
        return np.zeros((0, 2)), "empty"
    P = np.array(cand, float)
    P = P[np.lexsort((P[:, 1], P[:, 0]))]
    ded = [P[0]]
    for v in P[1:]:
        if np.linalg.norm(v - ded[-1]) > 1e-7:
            ded.append(v)
    P = np.array(ded)
    if len(P) < 3:
        return P, "degenerate"
    H = convex_hull(P)
    if len(H) < 3:
        return H, "degenerate"
    return H, "polygon"


# --------------------------------------------------------------------------
# 区域算法 1：顶点枚举 + 凸包（主算法，精确）
# --------------------------------------------------------------------------
def region_by_vertices(pts, svds, err=BEARING_ERR, tol=DIST_TOL):
    """R = (*@\ensuremath{\cap}@*) W_i 的顶点集合。

    R 是凸集，其任一极点必是两条起作用约束边界直线的交点；反之，落在所有半平面内的
    边界直线交点必为极点（若 X = (Y+Z)/2 且两个非平行约束在 X 处取等，则 Y、Z 也
    必须取等，故 Y = Z = X）。因此枚举 + 筛选是精确的。
    """
    if len(pts) == 0:
        return np.zeros((0, 2)), "empty"
    H, lines = [], []
    for p, th in zip(pts, svds):
        H += wedge_halfplanes(p, th, err)
        for s in (-1.0, 1.0):
            a = math.radians((th + s * err) % 360.0)
            lines.append((np.asarray(p, float), np.array([math.cos(a), math.sin(a)])))
    cand = []
    n = len(lines)
    for i in range(n):
        for j in range(i + 1, n):
            X = _line_intersection(lines[i][0], lines[i][1], lines[j][0], lines[j][1])
            if X is None or not np.all(np.isfinite(X)):
                continue
            # 半平面判据：在楔顶点处也良定义（无角度奇异），故检测点可作顶点保留
            if all(float(a @ X + b) <= tol for (a, b) in H):
                cand.append(X)
    return _dedup_hull(cand)


# --------------------------------------------------------------------------
# 区域算法 2：半平面裁剪（独立参考实现，可直接判无界）
# --------------------------------------------------------------------------
def clip_halfplane(poly, a, b, tol=DIST_TOL):
    """Sutherland–Hodgman：保留 {X : a·X + b <= 0} 的部分"""
    if not poly:
        return []
    out = []
    n = len(poly)
    for i in range(n):
        P = poly[i]
        Q = poly[(i + 1) % n]
        fp = float(a @ P + b)
        fq = float(a @ Q + b)
        if fp <= tol:
            out.append(P)
        if (fp < -tol and fq > tol) or (fp > tol and fq < -tol):
            t = fp / (fp - fq)
            out.append(P + t * (Q - P))
    ded = []
    for v in out:
        if not ded or np.linalg.norm(v - ded[-1]) > 1e-9:
            ded.append(v)
    if len(ded) > 1 and np.linalg.norm(ded[0] - ded[-1]) < 1e-9:
        ded.pop()
    return ded


def region_by_clipping(pts, svds, err=BEARING_ERR, box=CLIP, tol=DIST_TOL):
    """从大框出发逐个半平面裁剪。

    返回 (P, status)：status (*@\ensuremath{\in}@*) {"empty","degenerate","polygon","unbounded"}；
    多边形触到外框（|坐标| > box/2）即判定区域无界，P 是被截断的多边形。
    """
    poly = [np.array([-box, -box], float), np.array([box, -box], float),
            np.array([box, box], float), np.array([-box, box], float)]
    for p, th in zip(pts, svds):
        for (a, b) in wedge_halfplanes(p, th, err):
            poly = clip_halfplane(poly, a, b, tol)
            if not poly:
                return np.zeros((0, 2)), "empty"
    P = np.array(poly, float)
    if len(P) < 3:
        return P, "degenerate"
    if float(np.max(np.abs(P))) > box * 0.5:
        return P, "unbounded"
    H = convex_hull(P)
    return (H if len(H) >= 3 else P), "polygon"


def region_from_bearings(pts, svds, err=BEARING_ERR):
    """主入口（兼容旧接口）：返回 (顶点数组, status)。

    status (*@\ensuremath{\in}@*) {"empty", "unbounded", "degenerate", "bounded"}。
    """
    P, st = region_by_vertices(pts, svds, err)
    if st == "empty":
        return P, "empty"
    rec = region_is_unbounded(list(svds), err) if len(svds) else None
    if rec is not None:
        return P, "unbounded"
    if st == "degenerate":
        return P, "degenerate"
    return P, "bounded"


# --------------------------------------------------------------------------
# 直径与最小包围圆
# --------------------------------------------------------------------------
def _diameter_brute(P):
    n = len(P)
    best, pair = -1.0, (P[0], P[1])
    for i in range(n - 1):
        d = np.linalg.norm(P[i + 1:] - P[i], axis=1)
        k = int(np.argmax(d))
        if float(d[k]) > best:
            best, pair = float(d[k]), (P[i], P[i + 1 + k])
    return best, pair


def rotating_calipers(P):
    """凸多边形直径（旋转卡壳，O(n)）。P 需为逆时针凸序（允许共线点）。"""
    n = len(P)
    if n < 3:
        return diameter(P, brute_limit=2)
    k = 1
    best, pair = -1.0, (P[0], P[1])
    for i in range(n):
        i2 = (i + 1) % n
        e = P[i2] - P[i]
        while True:
            k2 = (k + 1) % n
            if abs(_cross2(e, P[k2] - P[i])) > abs(_cross2(e, P[k] - P[i])):
                k = k2
            else:
                break
        for j in (k, (k + 1) % n):
            d = float(np.linalg.norm(P[i] - P[j]))
            if d > best:
                best, pair = d, (P[i], P[j])
            d = float(np.linalg.norm(P[i2] - P[j]))
            if d > best:
                best, pair = d, (P[i2], P[j])
    return best, pair


def diameter(P, brute_limit=DIAM_BRUTE_LIMIT):
    """凸多边形直径，返回 (D, (A, B))"""
    n = len(P)
    if n == 0:
        return 0.0, None
    if n == 1:
        return 0.0, (P[0], P[0])
    if n == 2:
        return float(np.linalg.norm(P[0] - P[1])), (P[0], P[1])
    if n <= brute_limit:
        return _diameter_brute(P)
    return rotating_calipers(P)


def min_enclosing_circle(P):
    """最小包围圆 (center, radius)。顶点数很少，直接枚举 1/2/3 点候选圆。

    理论：r* (*@\ensuremath{\in}@*) [D/2, D/√3]（Jung 定理，平面），等号右端仅当区域为等边三角形时取到。
    """
    P = np.asarray(P, float)
    n = len(P)
    if n == 0:
        return None, 0.0
    if n == 1:
        return P[0].copy(), 0.0
    cands = [(P[i].copy(), 0.0) for i in range(n)]
    for i in range(n):
        for j in range(i + 1, n):
            c = 0.5 * (P[i] + P[j])
            cands.append((c, float(np.linalg.norm(P[i] - c))))
    for i in range(n):
        for j in range(i + 1, n):
            for k in range(j + 1, n):
                A, B, C = P[i], P[j], P[k]
                d = 2.0 * (A[0] * (B[1] - C[1]) + B[0] * (C[1] - A[1])
                           + C[0] * (A[1] - B[1]))
                if abs(d) < 1e-12:
                    continue
                aa, bb, cc = float(A @ A), float(B @ B), float(C @ C)
                ux = (aa * (B[1] - C[1]) + bb * (C[1] - A[1]) + cc * (A[1] - B[1])) / d
                uy = (aa * (C[0] - B[0]) + bb * (A[0] - C[0]) + cc * (B[0] - A[0])) / d
                c = np.array([ux, uy])
                cands.append((c, float(np.linalg.norm(A - c))))
    best = None
    for c, r in cands:
        if all(float(np.linalg.norm(P[i] - c)) <= r + 1e-9 for i in range(n)):
            if best is None or r < best[1]:
                best = (c, r)
    return best


def region_area(P):
    n = len(P)
    if n < 3:
        return 0.0
    return abs(0.5 * float(np.sum(P[:, 0] * np.roll(P[:, 1], -1)
                                   - np.roll(P[:, 0], -1) * P[:, 1])))


# --------------------------------------------------------------------------
# 汇总
# --------------------------------------------------------------------------
def analyse(pts, svds, err=BEARING_ERR):
    """给定检测点与示向度，输出定位区域的全套几何量。"""
    P, status = region_from_bearings(pts, svds, err)
    base = {"n_bearings": len(pts), "status": status}
    if status == "empty":
        base["reason"] = ("所有示向度锥无公共交集（数据不一致：两两夹角可以 > 2°，"
                          "但错开的两个楔也可能不相交）")
        return base
    if status == "unbounded":
        rec = region_is_unbounded(list(svds), err)
        base["reason"] = ("示向度锥近似平行（两两夹角 <= 2°），定位区域无界，"
                          "直径 = +inf；必须换检测点重测")
        base["recession_direction_deg"] = None if rec is None else round(float(rec), 4)
        base["n_vertices"] = int(len(P))
        base["vertices"] = [[round(float(v[0]), 4), round(float(v[1]), 4)] for v in P]
        return base
    if status == "degenerate" or len(P) < 3:
        base["status"] = "degenerate"
        base["reason"] = "公共交集退化为线段或点"
        base["vertices"] = [[round(float(v[0]), 4), round(float(v[1]), 4)] for v in P]
        return base

    D, (A, B) = diameter(P)
    c = 0.5 * (A + B)
    maxr = float(max(np.linalg.norm(v - c) for v in P))
    circle = min_enclosing_circle(P)
    if circle is None:
        # Very thin polygons far from the origin can lose circumcircle
        # precision. Translation preserves the geometric problem.
        origin = np.mean(P, axis=0)
        local = min_enclosing_circle(P - origin)
        mc = origin + local[0] if local is not None else origin
        mr = float(np.max(np.linalg.norm(P - mc, axis=1)))
    else:
        mc, mr = circle
    base.update({
        "n_vertices": int(len(P)),
        "vertices": [[round(float(v[0]), 4), round(float(v[1]), 4)] for v in P],
        "diameter_m": round(D, 4),
        "diameter_endpoints": [[round(float(A[0]), 4), round(float(A[1]), 4)],
                               [round(float(B[0]), 4), round(float(B[1]), 4)]],
        "circle_center": [round(float(c[0]), 4), round(float(c[1]), 4)],
        "circle_radius_m": round(D / 2.0, 4),
        "max_vertex_dist_m": round(maxr, 4),
        "diameter_circle_covers": bool(maxr <= D / 2.0 + 1e-9),
        "min_enclosing_center": [round(float(mc[0]), 4), round(float(mc[1]), 4)],
        "min_enclosing_radius_m": round(float(mr), 4),
        "min_circle_over_half_diameter": round(float(mr) / (D / 2.0), 6),
        "area_m2": round(region_area(P), 4),
    })
    return base


def run_case_csv(path, truth_path=None, err=BEARING_ERR):
    with open(path, encoding="utf-8") as f:
        rows = list(csv.DictReader(f))
    pts = [(float(r["x_m"]), float(r["y_m"])) for r in rows]
    svds = [float(r["svd_deg"]) for r in rows]
    res = analyse(pts, svds, err)
    res["n_bearings"] = len(rows)
    if truth_path and os.path.exists(truth_path):
        with open(truth_path, encoding="utf-8") as f:
            truth = json.load(f)
        G = truth["source_xy"]
        res["truth_xy"] = G
        if "circle_center" in res:
            res["center_error_vs_truth_m"] = round(math.dist(res["circle_center"], G), 3)
        if "min_enclosing_center" in res:
            res["mincircle_error_vs_truth_m"] = round(
                math.dist(res["min_enclosing_center"], G), 3)
        # 真值是否满足全部约束（等价于在半平面内）
        H = []
        for p, th in zip(pts, svds):
            H += wedge_halfplanes(p, th, err)
        res["truth_inside_region"] = bool(all(a @ np.asarray(G, float) + b <= 1e-6
                                             for (a, b) in H))
    return res


def main():
    d = sys.argv[1] if len(sys.argv) > 1 else os.path.join(
        os.path.dirname(os.path.abspath(__file__)), "..", "data")
    out = []
    for fn in sorted(os.listdir(d)):
        if not (fn.startswith("p1_case") and fn.endswith(".csv")):
            continue
        res = run_case_csv(os.path.join(d, fn),
                           os.path.join(d, fn.replace(".csv", ".truth.json")))
        res["case"] = fn
        out.append(res)
    jl = os.path.join(d, "p1_results.jsonl")
    with open(jl, "w", encoding="utf-8") as f:
        for r in out:
            f.write(json.dumps(r, ensure_ascii=False) + "\n")

    ok = [r for r in out if r["status"] == "bounded"]
    unlimited = [r for r in out if r["status"] == "unbounded"]
    summary = {
        "n_cases": len(out),
        "n_bounded": len(ok),
        "n_unbounded": len(unlimited),
        "n_degenerate": sum(1 for r in out if r["status"] == "degenerate"),
        "n_empty": sum(1 for r in out if r["status"] == "empty"),
        "diameter_m_range": [min((r["diameter_m"] for r in ok), default=None),
                             max((r["diameter_m"] for r in ok), default=None)],
        "n_diameter_circle_misses": sum(1 for r in ok if not r["diameter_circle_covers"]),
        "min_circle_over_half_diameter_max": max(
            (r["min_circle_over_half_diameter"] for r in ok), default=None),
        "n_truth_inside_region": sum(1 for r in ok if r.get("truth_inside_region")),
    }
    with open(os.path.join(d, "p1_summary.json"), "w", encoding="utf-8") as f:
        json.dump(summary, f, ensure_ascii=False, indent=2)
    for r in out:
        print("%-16s %-10s D=%-9s cover=%-5s n=%-2s r*/D*2=%-8s center_err=%s"
              % (r["case"], r["status"], r.get("diameter_m", "-"),
                 r.get("diameter_circle_covers", "-"), r.get("n_bearings", "-"),
                 r.get("min_circle_over_half_diameter", "-"),
                 r.get("center_error_vs_truth_m", "-")))
    print(json.dumps(summary, ensure_ascii=False))


if __name__ == "__main__":
    main()

\end{lstlisting}
\subsection{\texttt{src/p2\_grid\_expectation.py}}
% SHA256 5f08251fbd6384b898e89036cb90ff6d8e579bce7c938cf37606fc9e756b4424
\begin{lstlisting}[language=python,escapeinside={(*@}{@*)}]
"""B 题问题 2（模拟）：固定首次探测扇形后，第二检测点网格上的"交会区域直径期望"热力图。

问题设定（按需求方给定的简化模型）
----------------------------------
1. **固定当前探测得到的扇形区域**：探测器（机器狗）位于原点 ``S1 = (0, 0)``，
   探测方向为 x 轴正方向，示向度误差界 ``eps = 1 deg``。于是第一次探测给出的信息是角楔

       W1 = { X : |ang(X - S1) - theta1| <= eps }

   ``theta1`` 取自当前探测数据（默认 0 deg，即"探测方向为 x 轴正方向"）。
2. **干扰源均匀分布在扇形内，近似认为都在 x 轴上**：源位 ``G = (t, 0)``，
   权重由"扇形内面积均匀"退化得到：面积元 ``dA = r dr dtheta -> w(t) (*@\ensuremath{\propto}@*) t``
   （``--weight area``，默认）；另提供等权（``--weight length``）作对照。
3. **网格**：把每个网格点当作第二个检测点 ``S2``。默认网格框在右侧多留 0.2R
   （``x_range=(-1800, 2160)``），因为最优点在右半区，这样同样点数下右半区更密。
   注意**可行域仍是靶区圆域** ``R <= 1800 m``，``x > 1800`` 的部分是空框。
4. **交会区域与直径**：给定源位 ``G``，两个检测点各得一条带 ±eps 的示向度锥，
   两者之交是凸四边形 ``R = W1 (*@\ensuremath{\cap}@*) W2``，其**直径** D = 区域内任意两点距离的最大值。
5. **期望**：``E[D](S2) = Σ_i w_i D(S2, t_i)``，只在"可检测且有界"的样本上取平均，
   同时记录"可检测概率"，避免把不可用点伪装成好点。
6. **热力图**：``E[D](S2)`` 作为网格点 ``(x, y)`` 的函数，对数色标。

两套口径（主 / 辅）
-------------------
* ``measured``（主）：两个楔都用**测得值** ± 1°，即 ``R = W1(theta1) (*@\ensuremath{\cap}@*) W2(theta2)``，
  期望只对源位 ``t`` 取 —— 完全对应"固定当前探测数据"的字面含义。
* ``error_avg``（辅）：同一个区域，但期望**同时对 ±1° 测量误差取**，
  即 ``E_{t, e1, e2}[ D ]`` —— 更贴近"定位效果"的真实统计含义。

为什么网格必须是真正二维的（重要结论）
--------------------------------------
若把 ``S2`` 放在 x 轴上（``y = 0``），则两条示向度近似平行，夹角 ``|dtheta| <= 2*eps``，
交会区域**无界**、直径为 +∞（这正是 ``p1_intersection.region_by_vertices`` 会报
``unbounded`` 的退化几何，等价于"沿示向度方向逼近"）。所以 ``y = 0`` 必须在网格上被
剔除。更进一步，``|y|`` 很小但未退化的点，交会区域会被拉成极长的细条
（直径可达 10^4–10^5 m），数学上没错但工程上无用，程序用
``usable_threshold = 10 × 中位数`` 标注并掩码。

可检测性门控
------------
第二个检测点要能收到该频道信号，必须 ``|S2 G| <= 有效接收半径上界 1500 m``。
不加这个门控会选出几何有效但物理收不到信号的假最优点，所以它是模型的一部分
（默认开启，``--no-reach`` 可关闭对照）。

正确性验证
----------
直径由三条**互相独立**的实现给出，随机算例最大相对差 < 1e-13：
``p1_intersection.region_by_vertices``（半平面枚举）、``independent_diameter``
（世界坐标直接解边界交点）、``fast_diameters_vec``（角点枚举 + 楔定义筛选）。
``--selfcheck`` 另含蒙特卡洛与退化算例（T0–T10）。

用法
----
    python src/p2_grid_expectation.py --selfcheck
    python src/p2_grid_expectation.py --step 20 --n-t 500 --out out/p2_grid
    python src/p2_grid_expectation.py --theta1 30 --out out/p2_grid_t30
"""
from __future__ import annotations

import argparse
import csv
import json
import math
import os
import sys
import time

import numpy as np

_HERE = os.path.dirname(os.path.abspath(__file__))
if _HERE not in sys.path:
    sys.path.insert(0, _HERE)

from p1_intersection import BEARING_ERR, ang_diff, bearing, region_by_vertices  # noqa: E402

# --------------------------------------------------------------------------
# 物理常数（题目附录 1/2）
# --------------------------------------------------------------------------
EPS_DEG = BEARING_ERR          # 示向度误差界 1°
R_ARENA = 1800.0               # 目标区域半径（m）
R_RECV_MIN = 1000.0            # 有效接收半径下界（m，接口不返回）
R_RECV_MAX = 1500.0            # 有效接收半径上界（m，取信号的必要条件）
NEAR_R = 5.0                   # 近场阈值（m）：更近则没有示向度

# 默认计算参数
DEF_STEP = 20.0
DEF_X_RANGE = (-R_ARENA, 1.2 * R_ARENA)      # 右侧多留 0.2R（框内留白，便于看全）
DEF_Y_RANGE = (-R_ARENA, R_ARENA)
DEF_N_T = 500                  # x 轴上源位采样数
DEF_N_E = 400                  # 误差 MC 次数（error_avg 口径）
DEFAULT_SEED = 20260913


# --------------------------------------------------------------------------
# 基本几何
# --------------------------------------------------------------------------
def point_in_wedge(X, p, theta, err=EPS_DEG, tol=1e-9):
    """X 是否在角楔内（与楔顶点重合视为在楔内）"""
    X = np.asarray(X, float)
    p = np.asarray(p, float)
    if float(np.linalg.norm(X - p)) < tol:
        return True
    return abs(ang_diff(bearing(p, X), theta)) <= err + 1e-12


def _unbounded_bearings(theta1, theta2, err=EPS_DEG):
    """两楔交无界 <=> 存在方向与两条示向度夹角都 <= err。

    对两个示向度，等价于 ``min(|dth|, 180-|dth|) <= 2*err``。
    """
    d = abs(ang_diff(theta2, theta1))
    return min(d, 180.0 - d) <= 2.0 * err + 1e-12


def _corner_points(Pw, thP, Qw, thQ, eps):
    """两个角楔交的 6 个候选角点，**全部在世界坐标下求解**。

    每个角楔 ``|ang(X - P) - (*@\ensuremath{\theta}@*)| <= eps`` 的两条边界线方向为 ``(*@\ensuremath{\theta}@*) (*@\ensuremath{\mp}@*) eps``，
    交点由四线两两相交得到（6 个候选）；顶点判定直接用楔的定义
    （候选点 X 是顶点 (*@\ensuremath{\Longleftrightarrow}@*) 对每个楔都有 ``|ang(X-P) - (*@\ensuremath{\theta}@*)| <= eps`` 或 X 恰在顶点 P）。
    这个判据与坐标系定向、2x2 解法的行列式符号都无关，因此最稳健。

    返回 ``(corner_world, valid)``，形状 ``(..., 6, 2)`` 与 ``(..., 6)``。
    """
    Pw = np.asarray(Pw, float)
    Qw = np.asarray(Qw, float)
    thP = np.asarray(thP, float)
    thQ = np.asarray(thQ, float)

    def two_lines(P, th):
        out = []
        for s in (-eps, eps):
            a = np.radians(th + s)
            n = np.stack([np.sin(a), -np.cos(a)], axis=-1)
            out.append((n, np.einsum("...i,...i->...", n, P)))
        return out

    lines = two_lines(Pw, thP) + two_lines(Qw, thQ)

    def in_wedge(X, P, th):
        d = X - P
        dist = np.hypot(d[..., 0], d[..., 1])
        ang = np.degrees(np.arctan2(d[..., 1], d[..., 0]))
        diff = (ang - th + 180.0) % 360.0 - 180.0
        return (np.abs(diff) <= eps + 1e-9) | (dist <= 1e-9)

    corners, valid = [], []
    for ii in range(4):
        for jj in range(ii + 1, 4):
            (n1, c1), (n2, c2) = lines[ii], lines[jj]
            det = n1[..., 0] * n2[..., 1] - n1[..., 1] * n2[..., 0]
            safe = np.where(np.abs(det) < 1e-15, 1e-15, det)
            X = np.stack([(c1 * n2[..., 1] - c2 * n1[..., 1]) / safe,
                          (n1[..., 0] * c2 - n2[..., 0] * c1) / safe], axis=-1)
            ok = in_wedge(X, Pw, thP) & in_wedge(X, Qw, thQ)
            corners.append(X)
            valid.append(ok)
    return np.stack(corners, axis=-2), np.stack(valid, axis=-1)


def _diam_from_corners(cw, v, shp):
    """从候选角点与其可用性掩码求凸四边形直径（世界坐标下，距离即真实距离）"""
    nP = cw.shape[-2]
    best = np.zeros(shp)
    for ii in range(nP):
        for jj in range(ii + 1, nP):
            d2 = ((cw[..., ii, 0] - cw[..., jj, 0]) ** 2
                  + (cw[..., ii, 1] - cw[..., jj, 1]) ** 2)
            best = np.where(v[..., ii] & v[..., jj], np.maximum(best, d2), best)
    return best, v.sum(axis=-1)


def independent_diameter(t, S2, eps=EPS_DEG):
    """交会区域直径的**独立参考实现**（标量、世界坐标、直接解边界交点）。

    与 ``_corner_points`` 走的是不同代码路径（这里用射线参数式解交点）。
    无界（``min(|dth|,180-|dth|) <= 2*eps``）返回 ``inf``。
    """
    S2 = np.asarray(S2, float)
    t = float(t)
    th1 = 0.0                                                   # S1 -> G
    th2 = math.degrees(math.atan2(-S2[1], t - S2[0])) % 360.0   # S2 -> G
    if _unbounded_bearings(th1, th2, eps):
        return math.inf

    lines = []
    for apex, th in (((0.0, 0.0), th1), (tuple(S2), th2)):
        for s in (-eps, eps):
            a = math.radians(th + s)
            lines.append((apex, (math.cos(a), math.sin(a))))

    def in_wedge(X, apex, th):
        dx, dy = X[0] - apex[0], X[1] - apex[1]
        if math.hypot(dx, dy) < 1e-9:
            return True
        ad = math.degrees(math.atan2(dy, dx))
        return abs((ad - th + 180.0) % 360.0 - 180.0) <= eps + 1e-9

    verts = []
    for i in range(4):
        for j2 in range(i + 1, 4):
            (P1, u), (P2, v) = lines[i], lines[j2]
            det = u[0] * (-v[1]) - (-v[0]) * u[1]
            if abs(det) < 1e-14:
                continue
            rv = (P2[0] - P1[0], P2[1] - P1[1])
            s = (rv[0] * (-v[1]) - (-v[0]) * rv[1]) / det
            X = (P1[0] + s * u[0], P1[1] + s * u[1])
            if in_wedge(X, (0.0, 0.0), th1) and in_wedge(X, tuple(S2), th2):
                verts.append(X)
    if len(verts) < 2:
        return math.inf
    return max(math.dist(a, b) for i, a in enumerate(verts) for b in verts[i + 1:])


def fast_diameters_vec(S2s, ts, eps=EPS_DEG):
    """批量快路径：S2s (m,2) × ts (n,) -> D (m,n), status (m,n)。

    **源必须固定在 x 轴正半轴（S1 在原点）**，这是本题设定。

    status 编码：0 = 有界多边形，1 = 空，2 = 退化，3 = 无界。
    """
    S2s = np.asarray(S2s, float)
    ts = np.asarray(ts, float).ravel()
    m, n = S2s.shape[0], ts.size
    D = np.full((m, n), np.nan)
    status = np.ones((m, n), dtype=np.int8)
    if m == 0 or n == 0:
        return D, status
    r = np.hypot(S2s[:, 0], S2s[:, 1])
    ok_row = r >= 1e-9
    status[~ok_row, :] = 2
    th2 = np.degrees(np.arctan2(-S2s[:, 1][:, None],
                                ts[None, :] - S2s[:, 0][:, None])) % 360.0
    dth = np.abs((th2 - 0.0 + 180.0) % 360.0 - 180.0)
    sep = np.minimum(dth, 180.0 - dth)
    unb = (sep <= 2.0 * eps + 1e-12) | (~ok_row[:, None])
    status[unb] = np.where(np.broadcast_to(ok_row[:, None], unb.shape), 3, 2)[unb]
    ok = ~unb
    if not np.any(ok):
        return D, status

    S1w = np.zeros((m, n, 2))
    S2w = np.broadcast_to(S2s[:, None, :], (m, n, 2))
    th1g = np.zeros((m, n))
    corner, valid = _corner_points(S1w, th1g, S2w, th2, eps)     # 世界坐标
    cw = corner.reshape(m, n, corner.shape[-2], 2)
    v = valid.reshape(m, n, valid.shape[-1])
    best, n_valid = _diam_from_corners(cw, v, (m, n))
    status = np.where(ok & (n_valid < 2), 1, status).astype(np.int8)
    finite = ok & (n_valid >= 2) & np.isfinite(best)
    D = np.where(finite, np.sqrt(np.maximum(best, 0.0)), np.nan)
    status = np.where(finite, 0, status).astype(np.int8)
    return D, status


def fast_diameters(S2, ts, eps=EPS_DEG):
    """``fast_diameters_vec`` 的单点版"""
    D, st = fast_diameters_vec(np.asarray(S2, float)[None, :],
                               np.asarray(ts, float).ravel(), eps)
    return D[0], st[0]


def _fast_pair_grid(S1, th1, S2, th2, eps):
    """通用批量快路径：``(...,)`` 形状的 S1/th1/S2/th2 -> D 与 status。

    与 ``fast_diameters_vec`` 同一套做法，但两个角楔都任意
    （不要求 S1 在原点、源在 x 轴上）。用于 error_avg 口径。
    """
    P1 = np.asarray(S1, float)
    P2 = np.asarray(S2, float)
    th1 = np.asarray(th1, float)
    th2 = np.asarray(th2, float)
    if th2.ndim > th1.ndim:
        th1 = th1.reshape(th1.shape + (1,) * (th2.ndim - th1.ndim))
    shp = np.broadcast_shapes(th1.shape, th2.shape)
    th1 = np.broadcast_to(th1, shp)
    th2 = np.broadcast_to(th2, shp)
    P1 = np.broadcast_to(np.asarray(P1, float)[..., None, :], shp + (2,))
    P2 = np.broadcast_to(np.asarray(P2, float)[..., None, :], shp + (2,))

    rho = np.linalg.norm(P1 - P2, axis=-1)
    sep = np.abs((th2 - th1 + 180.0) % 360.0 - 180.0)
    sep = np.minimum(sep, 180.0 - sep)
    D = np.full(shp, np.nan)
    status = np.ones(shp, dtype=np.int8)
    status[sep <= 2.0 * eps + 1e-12] = 3
    status[rho < 1e-9] = 2
    ok = ~((sep <= 2.0 * eps + 1e-12) | (rho < 1e-9))
    if not np.any(ok):
        return D, status
    corner, valid = _corner_points(P1, th1, P2, th2, eps)
    cw = corner.reshape(shp + (corner.shape[-2], 2))
    v = valid.reshape(shp + (valid.shape[-1],))
    best, n_valid = _diam_from_corners(cw, v, shp)
    status = np.where(ok & (n_valid < 2), 1, status).astype(np.int8)
    finite = ok & (n_valid >= 2) & np.isfinite(best)
    D = np.where(finite, np.sqrt(np.maximum(best, 0.0)), np.nan)
    status = np.where(finite, 0, status).astype(np.int8)
    return D, status


# --------------------------------------------------------------------------
# 源位采样与权重
# --------------------------------------------------------------------------
def sector_ray_exit(S1, a_deg, R=R_ARENA):
    """从 S1 沿方位角 a 到半径 R 圆域边界的距离"""
    S1 = np.asarray(S1, float)
    u = np.array([math.cos(math.radians(a_deg)), math.sin(math.radians(a_deg))])
    b = float(S1 @ u)
    c = float(S1 @ S1) - R * R
    disc = b * b - c
    return 0.0 if disc <= 0 else float(-b + math.sqrt(disc))


def source_samples(S1, theta1, err=EPS_DEG, n_t=DEF_N_T, r_arena=R_ARENA,
                   r_recv_max=R_RECV_MAX, near=NEAR_R):
    """x 轴上的源位采样（面积均匀权重）。

    来源：干扰源在扇形内面积均匀 -> ``dA = r dr dtheta``，
    对角向积分后 ``w(r) (*@\ensuremath{\propto}@*) r``，故 ``w_i (*@\ensuremath{\propto}@*) t_i``（归一化后即离散先验）。
    ``t_hi = min(r_recv_max, 沿 theta1 到靶区边界的距离)``。
    """
    S1 = np.asarray(S1, float)
    t_hi = min(r_recv_max, sector_ray_exit(S1, theta1, r_arena))
    t_lo = near
    if t_hi <= t_lo:
        return np.zeros(0), np.zeros(0), {"t_lo": t_lo, "t_hi": t_hi, "n_t": 0}
    ts = np.linspace(t_lo, t_hi, int(n_t))
    w_area = ts / float(ts.sum())
    info = {"t_lo": float(t_lo), "t_hi": float(t_hi), "n_t": int(n_t),
            "weight": "area",
            "t_peak_note": "w(t) (*@\ensuremath{\propto}@*) t（面积均匀退化的严格结果）",
            "sector_half_angle_deg": float(err),
            "sector_full_width_at_t_hi_m": float(2.0 * t_hi * math.tan(math.radians(err)))}
    return ts, w_area, info


def length_weights(ts):
    w = np.ones_like(np.asarray(ts, float))
    return w / float(w.sum())


# --------------------------------------------------------------------------
# 网格与期望
# --------------------------------------------------------------------------
def make_grid(step=DEF_STEP, x_range=DEF_X_RANGE, y_range=DEF_Y_RANGE,
              arena_r=R_ARENA, domain="arena"):
    """生成网格点（仅保留靶区圆域内、且不贴 x 轴退化带的点）。

    默认网格框右侧多留 0.2R（``x_range=(-1800, 2160)``）：最优点在右半区
    （约 ``x (*@\ensuremath{\approx}@*) +1080`` m），把点堆在那里比均匀铺满 ±1800 更划算。
    **注意可行域仍是靶区圆域** ``R <= 1800 m``，``x > 1800`` 的部分不会进入结果。
    """
    xs = np.arange(x_range[0], x_range[1] + 1e-9, step)
    ys = np.arange(y_range[0], y_range[1] + 1e-9, step)
    DX, DY = np.meshgrid(xs, ys)             # shape (ny, nx)
    X, Y = DX.ravel(), DY.ravel()
    keep = (X * X + Y * Y <= arena_r * arena_r) & (np.abs(Y) > 1e-9)
    if domain == "reach":
        d = np.sqrt(np.minimum(X ** 2, (np.abs(X) - R_RECV_MAX) ** 2) + Y ** 2)
        keep &= (d <= R_RECV_MAX)
    return X[keep], Y[keep], xs, ys, DX.shape


def reach_ok(X, Y, ts, r_recv_max=R_RECV_MAX):
    """可检测性门控矩阵：``|S2 - G_i| <= r_recv_max``。

    源在 x 轴上，对固定 (x,y) 关于 t 的最小值在 ``t = clip(x, t_lo, t_hi)`` 处取到。
    """
    X = np.asarray(X, float)[:, None]
    Y = np.asarray(Y, float)[:, None]
    ts = np.asarray(ts, float)[None, :]
    t_clip = np.clip(X, float(np.min(ts)), float(np.max(ts)))
    return np.sqrt((X - t_clip) ** 2 + Y ** 2) <= r_recv_max + 1e-9


def expected_diameter_measured(X, Y, ts, w, eps=EPS_DEG, reach=True, chunk=4096):
    """主口径：区域用测得值 ±eps，期望只对源位取。

    返回 ``E``（``E[D*1{可检测}]``）、``E_det``（``E[D|可检测]``，主目标）、
    ``cov``/``p_det``、``Dmin``/``Dmax``/``n_ok``。
    """
    S2 = np.column_stack([X, Y])
    n = X.size
    E = np.zeros(n)
    E_det = np.full(n, np.nan)
    cov = np.zeros(n)
    p_det = np.zeros(n)
    Dmin = np.full(n, np.nan)
    Dmax = np.full(n, np.nan)
    n_ok = np.zeros(n, dtype=int)
    wsum = float(np.sum(w))
    for s in range(0, n, chunk):
        e = min(n, s + chunk)
        D, _ = fast_diameters_vec(S2[s:e], ts, eps)
        m = np.isfinite(D)
        if reach:
            m &= reach_ok(S2[s:e, 0], S2[s:e, 1], ts)
        sw = np.where(m, w[None, :], 0.0).sum(axis=1)
        contrib = (np.where(m, D, 0.0) * w[None, :]).sum(axis=1)
        with np.errstate(invalid="ignore"):
            E[s:e] = contrib
            E_det[s:e] = np.where(sw > 0, contrib / np.maximum(sw, 1e-300), np.nan)
            Dmin[s:e] = np.where(m.any(axis=1),
                                 np.min(np.where(m, D, np.inf), axis=1), np.nan)
            Dmax[s:e] = np.where(m.any(axis=1),
                                 np.max(np.where(m, D, -np.inf), axis=1), np.nan)
        cov[s:e] = sw
        p_det[s:e] = sw / wsum
        n_ok[s:e] = m.sum(axis=1)
    return {"E": E, "E_det": E_det, "cov": cov, "p_det": p_det,
            "Dmin": Dmin, "Dmax": Dmax, "n_ok": n_ok}


def expected_diameter_error_avg(X, Y, ts, w, eps=EPS_DEG, n_e=DEF_N_E,
                                chunk=1024, seed=DEFAULT_SEED, reach=True):
    """辅口径：区域仍取 ``W1(theta1) (*@\ensuremath{\cap}@*) W2(theta2)``，但期望**同时对 ±eps 误差取**。

    定义为 ``E_err(X) = E_{e1,e2}[ E_{t|e}[D * 1{可检测}] ]``。

    实现用**分块蒙特卡洛**：对每批网格点抽 ``n_e`` 组独立误差 ``(e1,e2)``，
    第 k 组只作用到该批第 k 个点上 —— 每个网格点得到 ``n_e`` 个独立无偏估计，取均值。
    代价是 ``n_e × n_t × 网格点数``，而不是误差笛卡尔积
    （后者在 1 万网格点上要十几 GB，实测会被拖死）。
    """
    S2 = np.column_stack([X, Y])
    n = X.size
    if n == 0:
        return {"E_err": np.zeros(0), "E_err_det": np.zeros(0),
                "cov_err": np.zeros(0), "n_pairs": 0}
    rng = np.random.default_rng(seed)
    E_err = np.zeros(n)
    E_err_det = np.full(n, np.nan)
    cov_err = np.zeros(n)
    for s in range(0, n, chunk):
        e = min(n, s + chunk)
        blk = e - s
        E1 = rng.uniform(-eps, eps, blk)
        E2 = rng.uniform(-eps, eps, blk)
        S2b = S2[s:e]
        S1b = np.zeros((blk, 2))
        base2 = np.degrees(np.arctan2(-S2b[:, 1][:, None],
                                      ts[None, :] - S2b[:, 0][:, None])) % 360.0
        D, _ = _fast_pair_grid(S1b, E1[:, None], S2b, base2 + E2[:, None], eps)
        m = np.isfinite(D)
        if reach:
            m &= reach_ok(S2b[:, 0], S2b[:, 1], ts)
        contrib = (np.where(m, D, 0.0) * w[None, :]).sum(axis=1)
        wsum = np.where(m, w[None, :], 0.0).sum(axis=1)
        with np.errstate(invalid="ignore"):
            E_err_det[s:e] = np.where(wsum > 0, contrib / np.maximum(wsum, 1e-300), np.nan)
        cov_err[s:e] = wsum
        E_err[s:e] = contrib
    return {"E_err": E_err, "E_err_det": E_err_det, "cov_err": cov_err,
            "n_pairs": int(n_e), "estimator": "逐点蒙特卡洛（无偏）"}


# --------------------------------------------------------------------------
# 自检
# --------------------------------------------------------------------------
def selfcheck(verbose=True, seed=7):
    """独立实现交叉校验 + 退化算例 + 期望一致性（T0–T10）"""
    out = []
    ok_all = True

    def rec(name, ok, detail=""):
        nonlocal ok_all
        ok_all = ok_all and ok
        out.append({"check": name, "ok": bool(ok), "detail": detail})
        if verbose:
            print(f"[{'OK ' if ok else 'FAIL'}] {name}  {detail}")

    def th2_of(S2, t):
        return float(np.degrees(np.arctan2(-S2[1], t - S2[0])) % 360.0)

    def mc_points(S2, t, n_pt, rng_):
        """在两个楔各自锥内采样，保留落在交会区域内的点（向量化）"""
        th2 = th2_of(S2, t)
        L = rng_.uniform(5.0, 1900.0, n_pt)
        a1 = np.radians(rng_.uniform(-EPS_DEG, EPS_DEG, n_pt))
        P1 = np.stack([L * np.cos(a1), L * np.sin(a1)], axis=1)
        a2 = math.radians(th2) + np.radians(rng_.uniform(-EPS_DEG, EPS_DEG, n_pt))
        P2 = np.asarray(S2)[None, :] + np.stack([L * np.cos(a2), L * np.sin(a2)], axis=1)
        K = np.vstack([P1, P2])

        def inw(P, apex, th):
            d = P - np.asarray(apex)[None, :]
            dist = np.hypot(d[:, 0], d[:, 1])
            angd = np.degrees(np.arctan2(d[:, 1], d[:, 0]))
            return ((np.abs((angd - th + 180.0) % 360.0 - 180.0) <= EPS_DEG + 1e-9)
                    & (dist > 1e-9))

        return K[inw(K, (0.0, 0.0), 0.0) & inw(K, S2, th2)]

    def max_pair_dist(K, blk=512):
        dmax = 0.0
        for s0 in range(0, len(K), blk):
            dmax = max(dmax, float(np.max(
                np.linalg.norm(K[s0:s0 + blk, None, :] - K[None, :, :], axis=2))))
        return dmax

    rng = np.random.default_rng(seed)

    # ---- T0: 顶点个数=4 且每个顶点真在两个楔内 ----
    n0, bad0 = 0, 0
    for _ in range(2500):
        S2 = np.array([rng.uniform(-1500, 1500), rng.uniform(-1400, 1400)])
        if np.hypot(*S2) < 200.0:
            continue
        t = float(rng.uniform(200.0, 1450.0))
        if np.hypot(S2[0] - t, S2[1]) > R_RECV_MAX:
            continue
        th2 = th2_of(S2, t)
        if min(abs(ang_diff(th2, 0.0)), 180 - abs(ang_diff(th2, 0.0))) <= 10.0:
            continue
        D, _ = fast_diameters(S2, np.array([t]))
        if not np.isfinite(D[0]):
            continue
        n0 += 1
        c, v = _corner_points(np.zeros((1, 1, 2)), np.zeros((1, 1)),
                              S2.reshape(1, 1, 2), np.array([[th2]]), EPS_DEG)
        c, v = c[0, 0], v[0, 0]
        if int(v.sum()) != 4:
            bad0 += 1
            continue
        for k in np.nonzero(v)[0]:
            if not (point_in_wedge(c[k], (0.0, 0.0), 0.0)
                    and point_in_wedge(c[k], S2, th2)):
                bad0 += 1
                break
    rec("T0 顶点个数=4 且每个顶点真在两个楔内", bad0 <= 2,
        f"检查 {n0} 组，异常 {bad0}（允许 (*@\ensuremath{\le}@*)2 组数值退化）")

    # ---- T1: 快路径 vs p1 参考实现（良态） ----
    n_bad, n_cmp, n_skip, max_dd = 0, 0, 0, 0.0
    for _ in range(6000):
        S2 = np.array([rng.uniform(-1500, 1500), rng.uniform(-1500, 1500)])
        if np.hypot(*S2) < 1.0:
            continue
        t = float(rng.uniform(30.0, 1500.0))
        th2 = th2_of(S2, t)
        if min(abs(ang_diff(th2, 0.0)), 180 - abs(ang_diff(th2, 0.0))) <= 2 * EPS_DEG + 0.3:
            n_skip += 1
            continue
        D, _ = fast_diameters(S2, np.array([t]))
        P, pst = region_by_vertices([(0.0, 0.0), tuple(S2)], [0.0, th2])
        if pst != "polygon" or len(P) < 3:
            continue
        n_cmp += 1
        Pa = np.asarray(P, float)
        d_ref = max(math.dist(a, b) for i, a in enumerate(Pa) for b in Pa[i + 1:])
        if not np.isfinite(D[0]):
            n_bad += 1
            continue
        rel = abs(float(D[0]) - d_ref) / max(d_ref, 1e-12)
        max_dd = max(max_dd, rel)
        if rel > 1e-7:
            n_bad += 1
    rec("T1 快路径 vs p1 参考实现（直径，良态）", n_bad == 0,
        f"比较 {n_cmp} 组，跳过近临界 {n_skip} 组，最大相对差 {max_dd:.3e}")

    # ---- T2: 快路径 vs 独立实现 ----
    worst, worst_cfg, n2 = 0.0, None, 0
    for _ in range(8000):
        S2 = np.array([rng.uniform(-1500, 1500), rng.uniform(-1400, 1400)])
        if np.hypot(*S2) < 5.0:
            continue
        t = float(rng.uniform(20.0, 1500.0))
        D, _ = fast_diameters(S2, np.array([t]))
        if not np.isfinite(D[0]):
            continue
        dc = independent_diameter(t, S2)
        if not math.isfinite(dc):
            continue
        n2 += 1
        rel = abs(float(D[0]) - dc) / max(dc, 1e-12)
        if rel > worst:
            worst, worst_cfg = rel, (S2.copy(), t, float(D[0]), dc)
    rec("T2 快路径 vs 独立实现（直接解边界交点）", worst < 1e-9,
        f"{n2} 组，最大相对差 {worst:.3e}"
        + ("" if worst_cfg is None else
           f"（最差 S2={np.round(worst_cfg[0],3)}, t={worst_cfg[1]:.1f}, "
           f"数值={worst_cfg[2]:.4f}, 独立={worst_cfg[3]:.4f}）"))

    # ---- T3: 蒙特卡洛真值——直径 >= 区域内采样点对最大距离 ----
    worst3, worst3_cfg, n3 = 0.0, None, 0
    for _ in range(40):
        S2 = np.array([rng.uniform(-1200, 1200), rng.uniform(-1200, 1200)])
        if np.hypot(*S2) < 50.0:
            continue
        t = float(rng.uniform(100.0, 1400.0))
        if np.hypot(S2[0] - t, S2[1]) > R_RECV_MAX:
            continue
        th2 = th2_of(S2, t)
        if _unbounded_bearings(0.0, th2):
            continue
        D, _ = fast_diameters(S2, np.array([t]))
        if not np.isfinite(D[0]):
            continue
        K = mc_points(S2, t, 20000, rng)
        if len(K) < 2:
            continue
        dmax = max_pair_dist(K)
        n3 += 1
        over = dmax / float(D[0]) - 1.0
        if over > worst3:
            worst3, worst3_cfg = over, (S2.copy(), t, float(D[0]), dmax)
    rec("T3 直径 >= 区域内 MC 采样点对最大距离", worst3 <= 1e-9,
        f"{n3} 组，最大超出 {worst3:.3e}"
        + ("" if worst3_cfg is None else
           f"（S2={np.round(worst3_cfg[0],2)}, t={worst3_cfg[1]:.1f}, "
           f"数值={worst3_cfg[2]:.4f}, MC={worst3_cfg[3]:.4f}）"))

    # ---- T4: MC 逼近性 ----
    n4, ratios = 0, []
    for _ in range(40):
        S2 = np.array([rng.uniform(-800, 800), rng.uniform(-800, 800)])
        if np.hypot(*S2) < 50.0:
            continue
        t = float(rng.uniform(100.0, 1400.0))
        if np.hypot(S2[0] - t, S2[1]) > R_RECV_MAX:
            continue
        th2 = th2_of(S2, t)
        if _unbounded_bearings(0.0, th2):
            continue
        D, _ = fast_diameters(S2, np.array([t]))
        if not np.isfinite(D[0]):
            continue
        K = mc_points(S2, t, 40000, rng)
        if len(K) < 2:
            continue
        n4 += 1
        ratios.append(max_pair_dist(K) / float(D[0]))
    ratios = np.array(ratios)
    ok4 = bool(ratios.size and np.all(ratios >= 0.6) and np.all(ratios <= 1.0 + 1e-9))
    rec("T4 MC 最大距离 / 数值直径 (*@\ensuremath{\in}@*) [0.6, 1]", ok4,
        f"{n4} 组，最小比值 {ratios.min() if ratios.size else float('nan'):.4f}，"
        f"最大比值 {ratios.max() if ratios.size else float('nan'):.4f}")

    # ---- T5: 贴 x 轴退化必须被识别 ----
    bad5, det5 = 0, []
    for y in (0.0, 1e-6, 1.0):
        S2 = np.array([600.0, y])
        D, _ = fast_diameters(S2, np.array([300.0, 900.0]))
        det5.append(f"y={y}: D={np.round(D, 3).tolist()}")
        if not np.all(np.isnan(D)):
            bad5 += 1
    rec("T5 贴 x 轴退化被识别（无界/退化）", bad5 == 0, " | ".join(det5))

    # ---- T6: y=0 且源比 S2 远 => 无界 ----
    D, st = fast_diameters(np.array([600.0, 0.0]), np.array([900.0]))
    rec("T6 y=0 且源比 S2 远 -> 无界（必须剔除）",
        np.isnan(D[0]) and int(st[0]) == 3, f"D={D[0]}, status={int(st[0])}")

    # ---- T7: 交会角 90°（S2 在源正上方）时 D (*@\ensuremath{\approx}@*) 2 tan(eps)|S1S2| ----
    t = 1000.0
    t7_bad, t7_worst = [], 0.0
    for y in (30.0, 60.0, 120.0, 1000.0):
        S2 = np.array([t, y])
        Dv, _ = fast_diameters(S2, np.array([t]))
        dc = independent_diameter(t, S2)
        target = 2 * math.tan(math.radians(EPS_DEG)) * float(np.hypot(*S2))
        rel = abs(float(Dv[0]) - target) / target
        t7_worst = max(t7_worst, rel)
        if abs(float(Dv[0]) - dc) / dc > 1e-9:
            t7_bad.append(f"y={y}: 数值={float(Dv[0]):.6f} 独立={dc:.6f} 不一致")
        if rel > 2e-3:
            t7_bad.append(f"y={y}: 相对差={rel:.2e}")
    if _unbounded_bearings(0.0, 270.0):
        t7_bad.append("gamma*=90° 应判定为良态")
    rec("T7 交会角=90° 时 D (*@\ensuremath{\approx}@*) 2 tan(eps)|S1S2|（实测误差 < 2e-3）", not t7_bad,
        "；".join(t7_bad) if t7_bad else f"四组，最大相对差 {t7_worst:.2e}")

    # ---- T8: 源位采样权重归一 + 上界 ----
    ts, w, info = source_samples((0.0, 0.0), 0.0, n_t=200)
    rec("T8 源位采样权重归一 + 上界正确",
        abs(w.sum() - 1.0) < 1e-12 and abs(info["t_hi"] - min(R_RECV_MAX, R_ARENA)) < 1e-9,
        f"t(*@\ensuremath{\in}@*)[{info['t_lo']:.1f},{info['t_hi']:.1f}], Σw-1={w.sum()-1:.2e}")

    # ---- T9: E[D|可检测] 与逐点加权平均一致 ----
    X = np.array([300.0, 447.2, -200.0, 900.0])
    Y = np.array([600.0, 800.0, 300.0, 0.0])
    tt, ww, _ = source_samples((0.0, 0.0), 0.0, n_t=60)
    res9 = expected_diameter_measured(X, Y, tt, ww)
    manual = []
    for i in range(X.size):
        num = den = 0.0
        for j in range(tt.size):
            D, _ = fast_diameters(np.array([X[i], Y[i]]), np.array([tt[j]]))
            if np.isfinite(D[0]):
                num += ww[j] * float(D[0])
                den += ww[j]
        manual.append(num / den if den > 0 else np.nan)
    ok9 = np.allclose(np.array(manual), res9["E_det"], rtol=1e-10, equal_nan=True)
    rec("T9 E[D|可检测] 与逐点加权平均一致", ok9,
        f"manual={np.round(manual,4).tolist()}, batch={np.round(res9['E_det'],4).tolist()}")

    # ---- T10: 网格框扩展的语义（诚实版）----
    #   靶区是半径 1800 m 的圆域，格点是"圆域 (*@\ensuremath{\cap}@*) 网格框"；因此**平移/放大框边界
    #   不会改变圆域内的点数**（只改变框内空白的多少）。要真正加密只能缩小 step。
    Xa, Ya, xsa, _, _ = make_grid(step=50.0, x_range=(-R_ARENA, R_ARENA))
    Xb, Yb, xsb, _, _ = make_grid(step=50.0)
    inside = (Xb ** 2 + Yb ** 2) <= R_ARENA ** 2 + 1e-9
    same_count = (Xa.size == Xb.size)
    rec("T10 网格框扩展不改变圆域内点数（加密只能靠缩小 step）",
        bool(np.all(inside)) and same_count and xsb[-1] > xsa[-1],
        f"x_max {xsa[-1]:.0f}->{xsb[-1]:.0f} m：圆域内点数均为 {Xa.size}（不变），"
        f"全部满足 R<=1800；右半区点数与左半区相同（格点本已对称）")

    if verbose:
        print(f"\n自检总结：{'全部通过' if ok_all else '存在失败项'}")
    return {"ok": ok_all, "checks": out}


# --------------------------------------------------------------------------
# 输出：CSV / JSON
# --------------------------------------------------------------------------
def _json_default(o):
    """把 numpy 标量/布尔/数组转成可序列化对象（自检结果里会混进 np.bool_）"""
    if isinstance(o, (np.bool_,)):
        return bool(o)
    if isinstance(o, np.integer):
        return int(o)
    if isinstance(o, np.floating):
        return float(o)
    if isinstance(o, np.ndarray):
        return o.tolist()
    raise TypeError(f"Object of type {type(o).__name__} is not JSON serializable")


def _fmt(v, nd=6):
    if v is None:
        return ""
    if isinstance(v, float) and (math.isnan(v) or math.isinf(v)):
        return "nan" if math.isnan(v) else "inf"
    return round(float(v), nd)


def save_outputs(res, outdir, tag=""):
    os.makedirs(outdir, exist_ok=True)
    X, Y = res["X"], res["Y"]
    path = os.path.join(outdir, f"p2_grid_map{tag}.csv")
    thr = res.get("usable_threshold_m", float("inf"))
    with open(path, "w", newline="", encoding="utf-8-sig") as f:
        wr = csv.writer(f)
        wr.writerow(["x_m", "y_m", "r_m", "phi_deg",
                     "E_diam_given_detectable_m", "E_diam_mixed_mm",
                     "p_detectable", "cov_weight", "Dmin_m", "Dmax_m",
                     "E_diam_err_given_detectable_m", "E_diam_err_mixed_m",
                     "cov_err", "n_ok", "usable", "near_axis_unusable_band"])
        for i in range(X.size):
            r = math.hypot(X[i], Y[i])
            phi = math.degrees(math.atan2(Y[i], X[i])) % 360.0
            ed = res["E_det"][i]
            usable = int(np.isfinite(ed) and ed <= thr and res["p_det"][i] > 0)
            wr.writerow([_fmt(X[i], 4), _fmt(Y[i], 4), _fmt(r, 4), _fmt(phi, 4),
                         _fmt(ed), _fmt(res["E"][i]),
                         _fmt(res["p_det"][i]), _fmt(res["cov"][i]),
                         _fmt(res["Dmin"][i]), _fmt(res["Dmax"][i]),
                         _fmt(res["E_err_det"][i]), _fmt(res["E_err"][i]),
                         _fmt(res["cov_err"][i]), int(res["n_ok"][i]), usable,
                         int(np.isfinite(ed) and ed > thr)])
    rep = {k: v for k, v in res.items()
           if k not in ("X", "Y", "E", "E_det", "cov", "p_det", "Dmin", "Dmax", "n_ok",
                        "E_err", "E_err_det", "cov_err", "usable")}
    rep["csv"] = os.path.basename(path)
    with open(os.path.join(outdir, f"p2_grid_report{tag}.json"), "w", encoding="utf-8") as f:
        json.dump(rep, f, ensure_ascii=False, indent=2, default=_json_default)
    return path


# --------------------------------------------------------------------------
# 热力图
# --------------------------------------------------------------------------
def _savefig(fig, path, dpi=150, tries=6):
    """保存图片并重试（Windows 上杀毒/索引偶尔短暂锁住刚写出的 PNG）。"""
    import time as _t
    last = None
    for k in range(tries):
        try:
            fig.savefig(path, dpi=dpi)
            return path
        except OSError as ex:
            last = ex
            _t.sleep(0.5 * (k + 1))
    raise last


def make_figures(res, outdir, tag=""):
    import matplotlib
    matplotlib.use("Agg")
    import matplotlib.pyplot as plt
    from matplotlib.colors import LogNorm
    from matplotlib.ticker import FixedLocator, FuncFormatter

    os.makedirs(outdir, exist_ok=True)
    X, Y, xs, ys, shape = res["X"], res["Y"], res["xs"], res["ys"], res["shape"]
    ny, nx = shape
    figs = []
    t_hi = float(res["source_prior"]["t_hi"])       # 扇形径向上界 = 有效接收半径上界

    def to_grid(vals):
        g = np.full((ny, nx), np.nan)
        g[np.searchsorted(ys, Y), np.searchsorted(xs, X)] = vals
        return g

    XLIM = (-1500.0, 2200.0)      # x 负方向小一点；正方向盖住右侧扩展后的网格框
    YLIM = (-2000.0, 2000.0)

    # 关注区（E ~ 50–1500 m）用对数色标，低值段的差异才看得出来
    ZMIN, ZMAX = 50.0, 2000.0
    lnorm = LogNorm(vmin=ZMIN, vmax=ZMAX)
    MAJOR = [50, 70, 100, 150, 200, 300, 500, 700, 1000, 1500, 2000]
    MINOR = [60, 80, 120, 250, 400, 600, 800, 1200]

    ok_use = res.get("usable", np.ones(X.size, bool))
    GX = to_grid(X)
    GY = to_grid(Y)
    Gm = to_grid(np.where(ok_use, res["E_det"], np.nan))
    Ge = to_grid(np.where(ok_use, res["E_err_det"], np.nan))
    Gc = to_grid(res["p_det"])
    th = np.linspace(0, 2 * math.pi, 2000)

    def best_xy(G):
        r, c = np.unravel_index(int(np.nanargmin(G)), G.shape)
        return float(GX[r, c]), float(GY[r, c]), float(G[r, c])

    def show(ax, G, cmap, vmin=None, vmax=None, norm=None, ttl=None):
        kw = {"norm": norm} if norm is not None else {"vmin": vmin, "vmax": vmax}
        im = ax.imshow(G, origin="lower", extent=[xs[0], xs[-1], ys[0], ys[-1]],
                       cmap=cmap, interpolation="nearest", aspect="equal", **kw)
        ax.set_xlim(XLIM[0], XLIM[1])
        ax.set_ylim(YLIM[0], YLIM[1])
        ax.set_xlabel("x (m)"); ax.set_ylabel("y (m)")
        if ttl:
            ax.set_title(ttl, fontsize=9)
        return im

    def ring(ax, style="w--", lw=1.0, alpha=0.85):
        """只画落在显示范围内的靶区圆弧（R=1800 m）"""
        cx, cy = R_ARENA * np.cos(th), R_ARENA * np.sin(th)
        m = ((cx >= XLIM[0]) & (cx <= XLIM[1]) & (cy >= YLIM[0]) & (cy <= YLIM[1]))
        ax.plot(np.where(m, cx, np.nan), np.where(m, cy, np.nan),
                style, lw=lw, alpha=alpha)

    def draw_sector(ax, zorder=6, label=True):
        """画出**首次探测的扇形**（源位先验所在区域），半透明。

        角楔半张角 ``eps = 1°``，径向范围 ``[near, t_hi]``。``eps`` 只有 1°，
        所以这个扇形在整张图上很窄（1500 m 处全宽仅约 52 m）—— 为让它看得见，
        填充用半透明 + 亮色描边。**图内不放任何说明文字/框**（会遮挡内容），
        相关信息一律写进图例、标题或 `report.json`。

        注意 ``ax.fill`` 的 label 会按**网格点数**产生成百上千条图例项，
        所以这里不传 label，图例项改由下面的 ``sector_patch`` 提供。
        """
        t1 = t_hi
        ang_f = np.linspace(-EPS_DEG, EPS_DEG, 400)
        ax.plot(t1 * np.cos(np.radians(ang_f)), t1 * np.sin(np.radians(ang_f)),
                "-", color="darkorange", lw=1.6, alpha=0.95, zorder=zorder)
        n = 240
        rr = np.linspace(NEAR_R, t1, n)
        aa = np.linspace(-EPS_DEG, EPS_DEG, n)
        A, R_ = np.meshgrid(np.radians(aa), rr)
        ax.fill(R_ * np.cos(A), R_ * np.sin(A), color="darkorange", alpha=0.6,
                lw=0, zorder=zorder)
        for s in (-EPS_DEG, EPS_DEG):
            a = math.radians(s)
            ax.plot([NEAR_R * math.cos(a), t1 * math.cos(a)],
                    [NEAR_R * math.sin(a), t1 * math.sin(a)],
                    "-", color="darkorange", lw=1.6, alpha=0.95, zorder=zorder)

    def sector_patch():
        from matplotlib.patches import Patch
        return Patch(fc="darkorange", alpha=0.6, ec="darkorange",
                     label=f"1st-detection sector (±{EPS_DEG:g}°, r (*@\ensuremath{\le}@*) {t_hi:.0f} m)")

    def set_ticks(cb, size=9):
        """色标刻度：主/次刻度都带数字，且**字号统一**（不再大小不一）"""
        cb.set_ticks(MAJOR)
        cb.ax.set_yticklabels([f"{t:g}" for t in MAJOR])
        cb.ax.yaxis.set_minor_locator(FixedLocator(MINOR))
        cb.ax.yaxis.set_minor_formatter(FuncFormatter(lambda v, pos: f"{v:g}"))
        cb.ax.tick_params(which="both", labelsize=size, length=3.5)

    def origin_marker(ax, ms=13, label=True):
        """第一检测点（加粗圆点）。返回 artist，便于显式组织图例。"""
        ln, = ax.plot([0], [0], marker="o", ms=ms, mfc="red", mec="k", mew=1.5,
                      ls="none", zorder=9,
                      label=("S1 = origin (detector)" if label else None))
        return ln

    def best_marker(ax, G, ms=10, label=True):
        """最优第二检测点及关于 x 轴对称的那一点（图内不加文字，信息进图例）"""
        bx, by, vb = best_xy(G)
        pts = [by, -by] if abs(by) > 1e-9 else [by]
        first = None
        for n_, yy in enumerate(pts):
            ln, = ax.plot([bx], [yy], "wo", ms=ms, mfc="none", mew=2, ls="none",
                          zorder=7,
                          label=("best S2 (symmetric pair)" if (label and n_ == 0) else None))
            if first is None:
                first = ln
        return bx, by, vb, first

    # ---------- 图 1：主口径热力图 ----------
    fig, ax = plt.subplots(figsize=(10.6, 10.6))
    im = show(ax, Gm, "viridis_r", norm=lnorm)
    fig.subplots_adjust(top=0.88, bottom=0.07, left=0.09, right=0.98)
    ax.set_title("B-Problem 2: expected intersection-region diameter vs 2nd detection point",
                 fontsize=10, pad=10)
    fig.text(0.5, 0.965,
             "log colour scale;  measured bearings;  source prior: area-uniform inside the"
             " 1st-detection sector;\n"
             f"$\\varepsilon$={res['eps_deg']}°,  $\\theta_1$={res['theta1']}°,  "
             f"reach gate $|S_2G|\\leq${int(R_RECV_MAX)} m,  arena $R\\leq${int(R_ARENA)} m,  "
             f"grid step {res['step']:g} m ({res['grid_points']} points)",
             ha="center", va="bottom", fontsize=9)
    cb = fig.colorbar(im, ax=ax, shrink=0.8, extend="both")
    cb.set_label(f"E[intersection diameter | detectable]  (m, log scale),  "
                 f"colourbar range {ZMIN:g}–{ZMAX:g} m")
    set_ticks(cb)
    ring(ax)
    draw_sector(ax)
    h_origin = origin_marker(ax)
    bx, by, vb, h_best = best_marker(ax, Gm)
    ax.legend(handles=[h_origin, h_best, sector_patch()],
              loc="lower left", fontsize=8, framealpha=0.9)
    f1 = os.path.join(outdir, f"p2_grid_E_heatmap{tag}.png")
    _savefig(fig, f1); plt.close(fig); figs.append(f1)

    # ---------- 图 2：两口径并排 ----------
    fig, axes = plt.subplots(1, 2, figsize=(16.4, 9.4))
    for ax2, G, ttl in (
            (axes[0], Gm, "(a) measured bearings fixed\n"
                          "expectation over source position only"),
            (axes[1], Ge, "(b) error-averaged\n"
                          "expectation over source position AND ±1° reading errors")):
        im = show(ax2, G, "viridis_r", norm=lnorm)
        ring(ax2)
        draw_sector(ax2, label=False)
        origin_marker(ax2, ms=11, label=False)
        ibx, iby, ibv, _ = best_marker(ax2, G, ms=8, label=False)
        ax2.set_title(f"{ttl}\nmin E = {ibv:.1f} m at ({ibx:.0f}, {iby:.0f}) m", fontsize=9)
        cb2 = fig.colorbar(im, ax=ax2, shrink=0.75, extend="both")
        cb2.set_label("E[D | detectable]  (m, log scale)")
        set_ticks(cb2, size=8)
    f2 = os.path.join(outdir, f"p2_grid_E_heatmap_two{tag}.png")
    fig.tight_layout(); _savefig(fig, f2); plt.close(fig); figs.append(f2)

    # ---------- 图 3：可检测概率 ----------
    fig, ax = plt.subplots(figsize=(10.6, 10.0))
    im = show(ax, Gc, "magma", 0.0, 1.0,
              ttl="Detectability: probability the 2nd detector can even hear the channel")
    cbp = fig.colorbar(im, ax=ax, shrink=0.8)
    cbp.set_label("P(|S2G| (*@\ensuremath{\le}@*) 1500 m) under source prior")
    cbp.ax.tick_params(which="both", labelsize=9)
    ring(ax, "c--", 1.0, 0.9)
    draw_sector(ax)
    h_origin = origin_marker(ax)
    ax.legend(handles=[h_origin, sector_patch()],
              loc="lower left", fontsize=8, framealpha=0.9)
    f3 = os.path.join(outdir, f"p2_grid_coverage{tag}.png")
    _savefig(fig, f3); plt.close(fig); figs.append(f3)

    # ---------- 图 4：E(r) 与 E(phi) 剖面 ----------
    fig, axes = plt.subplots(1, 2, figsize=(14.0, 5.4))
    r_all = np.hypot(X, Y)
    phi_all = np.degrees(np.arctan2(Y, X)) % 360.0
    for phi0, style in ((90.0, "-"), (80.0, "--"), (70.0, "-."), (45.0, ":")):
        m = np.abs(((phi_all - phi0 + 180) % 360) - 180) < \
            res["step"] / (2 * np.maximum(r_all, 1)) * 1.2
        if m.sum() < 3:
            continue
        o = np.argsort(r_all[m])
        axes[0].plot(r_all[m][o], res["E_det"][m][o], style, lw=1.4, label=f"(*@\ensuremath{\phi}@*)={phi0:.0f}°")
    axes[0].set_xlabel("baseline r = |S1S2| (m)"); axes[0].set_ylabel("E[D | detectable] (m)")
    axes[0].set_title("radial profile (measured bearings)", fontsize=10)
    axes[0].legend(fontsize=8); axes[0].grid(alpha=0.3)
    for r0, style in ((200.0, "-"), (400.0, "--"), (800.0, "-."), (1200.0, ":")):
        m = np.abs(r_all - r0) < res["step"] * 1.2
        if m.sum() < 3:
            continue
        o = np.argsort(phi_all[m])
        axes[1].plot(phi_all[m][o], res["E_det"][m][o], style, lw=1.4, label=f"r={r0:.0f} m")
    axes[1].set_xlabel("(*@\ensuremath{\phi}@*) = arg(S2) (deg)"); axes[1].set_ylabel("E[D | detectable] (m)")
    axes[1].set_title("angular profile: minima near the 90°-crossing geometry", fontsize=10)
    axes[1].legend(fontsize=8); axes[1].grid(alpha=0.3)
    f4 = os.path.join(outdir, f"p2_grid_profiles{tag}.png")
    fig.tight_layout(); _savefig(fig, f4); plt.close(fig); figs.append(f4)

    # ---------- 图 5：交会区域随 S2 变化的示意 ----------
    from p1_intersection import convex_hull
    fig, axes = plt.subplots(1, 3, figsize=(15.0, 5.2))
    t_demo = 900.0
    for ax2, (rr_, pp_) in zip(axes, ((900.0, 90.0), (180.0, 80.0), (900.0, 20.0))):
        S2 = np.array([rr_ * math.cos(math.radians(pp_)), rr_ * math.sin(math.radians(pp_))])
        th2 = float(np.degrees(np.arctan2(-S2[1], t_demo - S2[0])) % 360.0)
        P, stt = region_by_vertices([(0.0, 0.0), tuple(S2)], [0.0, th2])
        D, stf = fast_diameters(S2, np.array([t_demo]))
        for base, colr, anch in ((0.0, "tab:blue", np.zeros(2)), (th2, "tab:orange", S2)):
            for s in (-1, 1):
                a = math.radians(base + s * EPS_DEG)
                L = 2200.0
                ax2.plot([anch[0], anch[0] + L * math.cos(a)],
                         [anch[1], anch[1] + L * math.sin(a)],
                         color=colr, lw=1.0, alpha=0.75)
        if len(P) >= 3 and not _unbounded_bearings(0.0, th2):
            H = convex_hull(np.asarray(P, float))
            H = np.vstack([H, H[:1]])
            ax2.fill(H[:, 0], H[:, 1], color="crimson", alpha=0.55, zorder=5)
        ax2.plot([0], [0], marker="o", ms=11, mfc="red", mec="k", mew=1.4, zorder=6)
        ax2.plot([S2[0]], [S2[1]], "o", color="tab:orange", ms=7, zorder=6)
        ax2.plot([t_demo], [0], "k^", ms=8, zorder=6)
        ax2.set_xlim(-1500, 2400); ax2.set_ylim(-1400, 2100)
        dv = float(D[0]) if np.isfinite(D[0]) else float("nan")
        ax2.set_title(f"source t=900 m,  r={rr_:.0f} m, (*@\ensuremath{\phi}@*)={pp_:.0f}°\n"
                      f"intersection D = {dv:.1f} m", fontsize=9)
        ax2.set_xlabel("x (m)"); ax2.set_ylabel("y (m)"); ax2.grid(alpha=0.25)
    f5 = os.path.join(outdir, f"p2_grid_region_demo{tag}.png")
    fig.tight_layout(); _savefig(fig, f5); plt.close(fig); figs.append(f5)
    return figs


# --------------------------------------------------------------------------
# 主流程
# --------------------------------------------------------------------------
def run(theta1=0.0, step=DEF_STEP, n_t=DEF_N_T, n_e=DEF_N_E, weight="area",
        domain="arena", eps=EPS_DEG, outdir=None, figures=True,
        seed=DEFAULT_SEED, reach=True, checks=True, verbose=True,
        x_range=DEF_X_RANGE, y_range=DEF_Y_RANGE):
    t0 = time.time()
    S1 = (0.0, 0.0)
    ts, w_area, sinfo = source_samples(S1, theta1, eps, n_t)
    w = w_area if weight == "area" else length_weights(ts)
    if ts.size == 0:
        raise SystemExit("源位采样为空：检查 theta1 与靶区/接收半径设置")
    X, Y, xs, ys, shape = make_grid(step=step, x_range=x_range, y_range=y_range,
                                    arena_r=R_ARENA, domain=domain)
    if verbose:
        print(f"网格点 {X.size} 个（step={step} m, domain={domain}, "
              f"x(*@\ensuremath{\in}@*)[{xs[0]:.0f},{xs[-1]:.0f}], y(*@\ensuremath{\in}@*)[{ys[0]:.0f},{ys[-1]:.0f}]）；"
              f"源位 {ts.size} 个 t(*@\ensuremath{\in}@*)[{ts[0]:.1f},{ts[-1]:.1f}] m, 权重={weight}")
    meas = expected_diameter_measured(X, Y, ts, w, eps, reach=reach)
    err = expected_diameter_error_avg(X, Y, ts, w, eps, n_e=n_e, seed=seed, reach=reach)

    t_nom = float(np.sum(w * ts))
    med = float(np.nanmedian(meas["E_det"]))
    usable_thr = 10.0 * med
    usable = np.isfinite(meas["E_det"]) & (meas["E_det"] <= usable_thr) & (meas["p_det"] > 0)
    E_use = np.where(usable, meas["E_det"], np.nan) if np.any(usable) else meas["E_det"]

    def _pick(a):
        a = np.asarray(a, float)
        if not np.any(np.isfinite(a)):
            return None
        i = int(np.nanargmin(a))
        return {"S2": [float(X[i]), float(Y[i])],
                "r_m": float(math.hypot(X[i], Y[i])),
                "phi_deg": float(math.degrees(math.atan2(Y[i], X[i])) % 360.0),
                "value_m": float(a[i]), "index": i,
                "p_det": float(meas["p_det"][i]),
                "E_mixed_m": float(meas["E"][i]),
                "E_err_det_m": (float(err["E_err_det"][i])
                                if np.isfinite(err["E_err_det"][i]) else None),
                "E_err_mixed_m": float(err["E_err"][i])}

    best_det, best_mix = _pick(E_use), _pick(meas["E"])
    n_band = int((np.isfinite(meas["E_det"]) & (meas["E_det"] > usable_thr)).sum())
    res = {
        "model": "P2 grid simulation: E[intersection-region diameter] over 2nd detection point",
        "theta1": float(theta1), "eps_deg": float(eps), "weight": weight, "domain": domain,
        "step": float(step), "reach_gate": bool(reach),
        "grid_x_range": [float(xs[0]), float(xs[-1])],
        "grid_y_range": [float(ys[0]), float(ys[-1])],
        "source_prior": sinfo, "arena_r": R_ARENA,
        "r_recv_range": [R_RECV_MIN, R_RECV_MAX], "near_r": NEAR_R,
        "S1": list(S1), "grid_points": int(X.size), "shape": list(shape),
        "xs": [float(v) for v in xs], "ys": [float(v) for v in ys],
        "usable_threshold_m": float(usable_thr),
        "X": X, "Y": Y,
        "E": meas["E"], "E_det": meas["E_det"], "cov": meas["cov"], "p_det": meas["p_det"],
        "Dmin": meas["Dmin"], "Dmax": meas["Dmax"], "n_ok": meas["n_ok"], "usable": usable,
        "E_err": err["E_err"], "E_err_det": err["E_err_det"],
        "cov_err": err["cov_err"], "n_err_pairs": err["n_pairs"],
        "summary": {
            "objective": "E[D | detectable]（|S2G| <= 1500 m 门控，按扇形内面积均匀先验加权）",
            "E_det_min_m": float(np.nanmin(E_use)),
            "E_det_median_m": med,
            "E_det_max_m": float(np.nanmax(meas["E_det"])),
            "E_mixed_min_m": float(np.nanmin(meas["E"])),
            "usable_points": int(usable.sum()),
            "usable_fraction": float(usable.mean()),
            "near_axis_band_points": n_band,
            "best_det": best_det, "best_mixed": best_mix,
            "best_S2": best_det["S2"] if best_det else None,
            "best_r_m": best_det["r_m"] if best_det else None,
            "best_phi_deg": best_det["phi_deg"] if best_det else None,
            "best_r_over_t": (best_det["r_m"] / t_nom) if best_det else None,
            "mean_source_range_m": t_nom,
            "theory_note": ("交会角 90° 时 D 取该基线下的最小值 (*@\ensuremath{\approx}@*) 2 tan(eps)|S1S2|；"
                            "一般情形 D 随基线与近简并程度增大"),
            "x_axis_degenerate": ("y(*@\ensuremath{\approx}@*)0 时两示向度近平行(|Δ(*@\ensuremath{\theta}@*)|(*@\ensuremath{\le}@*)2eps) → 交会区域无界、直径=+∞，"
                                  "被剔除；其邻域(|y| 小)直径可达 10^4-10^5 m，已用 "
                                  f"usable_threshold={usable_thr:.1f} m 标注为不可用带"),
            "reach_gate": ("|S2G| > 有效接收半径上界 1500 m 的源位样本不可检测，"
                           "不计入 E[D|可检测]；混口径 E[D·1{可检测}] 另列"),
            "grid_note": ("网格框右侧多留 0.2R（x 到 2160 m）以便右半区更密；"
                          "可行域仍为靶区圆域 R<=1800 m"),
            "runtime_s": None,
        },
    }
    if checks:
        res["checks"] = selfcheck(verbose=verbose)["checks"]
    res["summary"]["runtime_s"] = round(time.time() - t0, 3)
    if outdir:
        save_outputs(res, outdir)
        if figures:
            try:
                figs = make_figures(res, outdir)
                res["figures"] = [os.path.basename(f) for f in figs]
            except Exception as ex:                  # 图失败不影响数值结果落盘
                res["figures_error"] = f"{type(ex).__name__}: {ex}"
                if verbose:
                    print(f"[warn] 出图失败（数值结果已保存）：{ex}")
        save_outputs(res, outdir)
    if verbose:
        s = res["summary"]
        b = s["best_det"]
        print(f"\n最优第二检测点 S2* = ({b['S2'][0]:.1f}, {b['S2'][1]:.1f}) m "
              f"（r={b['r_m']:.1f} m, (*@\ensuremath{\phi}@*)={b['phi_deg']:.1f}°, r/(*@\ensuremath{\langle}@*)t(*@\ensuremath{\rangle}@*)={s['best_r_over_t']:.3f}）")
        print(f"E[D|可检测] 最小 = {s['E_det_min_m']:.3f} m；中位数 = {s['E_det_median_m']:.3f} m；"
              f"单条基线 2tan(eps)r = {2*math.tan(math.radians(eps))*b['r_m']:.3f} m")
        print(f"该点 P(可检测) = {b['p_det']:.4f}；混口径 E[D·1] = {b['E_mixed_m']:.3f} m；"
              f"误差口径 E[D|可检测] = {b['E_err_det_m']}")
        print(f"可用点 {s['usable_points']}/{res['grid_points']}"
              f"（近轴不可用带 {s['near_axis_band_points']} 个，阈值 {usable_thr:.1f} m）")
        print(f"耗时 {s['runtime_s']} s")
    return res


def main(argv=None):
    ap = argparse.ArgumentParser(
        description="B 题问题 2：第二检测点网格上的交会区域直径期望热力图")
    ap.add_argument("--theta1", type=float, default=0.0,
                    help="首次探测的示向度（度），默认 0（x 轴正向）")
    ap.add_argument("--step", type=float, default=DEF_STEP, help="网格步长（m）")
    ap.add_argument("--n-t", type=int, default=DEF_N_T, help="x 轴源位采样数")
    ap.add_argument("--n-e", type=int, default=DEF_N_E, help="误差 MC 采样数")
    ap.add_argument("--weight", choices=["area", "length"], default="area",
                    help="源位先验权重：area (*@\ensuremath{\propto}@*) t（面积均匀，默认）/ length 等权")
    ap.add_argument("--domain", choices=["arena", "reach"], default="arena",
                    help="S2 取值范围：整个靶区圆域 / 还要满足 |S2G|<=1500")
    ap.add_argument("--x-range", type=float, nargs=2, default=list(DEF_X_RANGE),
                    metavar=("XMIN", "XMAX"), help="网格 x 范围（m），默认右侧多留 0.2R")
    ap.add_argument("--y-range", type=float, nargs=2, default=list(DEF_Y_RANGE),
                    metavar=("YMIN", "YMAX"), help="网格 y 范围（m）")
    ap.add_argument("--out", default=None, help="输出目录")
    ap.add_argument("--no-reach", action="store_true", help="关闭可检测性门控")
    ap.add_argument("--no-figures", action="store_true")
    ap.add_argument("--no-checks", action="store_true", help="跳过自检（只出结果，快很多）")
    ap.add_argument("--selfcheck", action="store_true")
    a = ap.parse_args(argv)
    if a.selfcheck:
        r = selfcheck()
        if a.out:
            os.makedirs(a.out, exist_ok=True)
            with open(os.path.join(a.out, "p2_grid_selfcheck.json"), "w",
                      encoding="utf-8") as f:
                json.dump(r, f, ensure_ascii=False, indent=2, default=_json_default)
        return 0 if r["ok"] else 1
    run(theta1=a.theta1, step=a.step, n_t=a.n_t, n_e=a.n_e, weight=a.weight,
        domain=a.domain, outdir=a.out, figures=not a.no_figures,
        checks=not a.no_checks, reach=not a.no_reach,
        x_range=tuple(a.x_range), y_range=tuple(a.y_range))
    return 0


if __name__ == "__main__":
    raise SystemExit(main())

\end{lstlisting}
\subsection{\texttt{src/p3\_arena.py}}
% SHA256 790407ce57213250339f52314474d6e454f7a2a50b495bcbdafab7f093b2e8cd
\begin{lstlisting}[language=python,escapeinside={(*@}{@*)}]
"""B 题问题 3：机器狗与模拟器之间的统一接口层。

两个后端对上层暴露**完全相同的接口与返回字段**（字段名与模拟器 JSON 一致），
因此 :mod:`p3_robot` 里的策略代码在"离线演练"和"真实模拟器"上完全一致：

* :class:`MockArena` —— 完全离线复刻附件 1/附件 2 的物理与计时规则。
  用于自检、调参、批量演练（可以跑几百局，正式测试只有 3 次机会）。
* :class:`HttpArena` —— 真实模拟器的 HTTP+JSON 客户端，负责协议里最容易踩的坑：
  串行发送、``request_id`` 幂等与重试、``accepted`` 与 HTTP 状态双重校验、
  ``remaining_real_duration_s`` 动态预算、自身行为日志（论文支撑材料需要）。

计时规则（附件 1 表 1 / 附件 2 第 4 节）
----------------------------------------
* ``/measure``：移动距离/5 + （频道变化 ? 1 : 0） + 5 秒；只有 ``/measure`` 会切换频道；
* ``/clear``：移动距离/5 + （成功 5 / 未发现 3）秒；**不切换频道**；
* ``/enter``、``/exit`` 不推进虚拟时钟；
* 示向度误差在 [-1°,1°] 内，**同一地点固定**（本模块用"平滑随机场 + 局地量化项"建模）。
"""
from __future__ import annotations

import json
import math
import os
import socket
import time
import urllib.error
import urllib.request
from dataclasses import dataclass, field as dc_field

import numpy as np

# --------------------------------------------------------------------------
# 物理常数（题目附录 1/2、模拟器接口说明）
# --------------------------------------------------------------------------
R_ARENA = 1800.0            # 靶区半径（m）
V_ROBOT = 5.0               # 移动速度（m/s）
T_MEASURE = 5.0             # 一次检测耗时（s）
T_SWITCH = 1.0              # 频道切换耗时（s）
T_CLEAR_OK = 5.0            # 成功清除耗时（s）：光学 3 + 激光 2
T_CLEAR_FAIL = 3.0          # 未发现目标时的光学精确定位耗时（s）
T_CLEAR = T_CLEAR_OK
NEAR_R = 5.0                # 近场阈值（m）：更近时返回 near，无示向度
R_CLEAR = 20.0              # 清除半径（m）
R_RECV_MIN, R_RECV_MAX = 1000.0, 1500.0
BEARING_ERR = 1.0           # 示向度误差界（度）
CHANNELS = tuple(range(1, 21))
# 程序运行时间（现实时间）上限（s）：题目附录 3，/enter 响应里由 remaining_real_duration_s 给出
DEF_BUDGET_S = 1200.0
# 虚拟世界限时（s）：接口说明 1.4/4.5，默认 360000（100 小时）。
# **注意**：这是虚拟时钟的上限，与上面 1200 s 的现实时间预算是两回事。
# 离线 mock 必须用同一个值，否则等于给离线演练加了一条真实测试里不存在的约束。
DEF_VIRTUAL_LIMIT_S = 360000.0
COORD_LIMIT = 2_000_000.0   # 接口允许的坐标绝对值上限（m）


class ArenaError(RuntimeError):
    """接口层不可恢复错误（网络重试用尽、响应无法解析等）。"""


# --------------------------------------------------------------------------
# 离线模拟器
# --------------------------------------------------------------------------
@dataclass
class MockSource:
    channel: int
    x: float
    y: float
    r_recv: float
    cone_half: float = 180.0      # 覆盖半角：180 = 全向（问题 4 会用到 <180）
    dir_deg: float | None = None
    cleared: bool = False

    @property
    def xy(self):
        return (self.x, self.y)


class MockArena:
    """离线复刻的模拟器（问题 3：全向源）。"""

    def __init__(self, seed=0, n_sources=None, budget_s=DEF_VIRTUAL_LIMIT_S,
                 r_arena=R_ARENA, sources=None, error_corr_m=250.0,
                 error_local_w=0.25, verbose=False):
        """``budget_s`` = **虚拟时钟上限**（默认 360000 s，与真实模拟器一致）。"""
        self.seed = int(seed)
        self.rng = np.random.default_rng(self.seed)
        self.r_arena = float(r_arena)
        self.budget = float(budget_s)
        self.sources = list(sources) if sources is not None else self._gen_sources(n_sources)
        self._init_error_field(error_corr_m, error_local_w)
        self.pos = (0.0, 0.0)
        self.channel = 1
        self.vtime = 0.0
        self.entered = False
        self.finished = False
        self.verbose = verbose
        self.actions: list[dict] = []
        self.stats = {"n_measure": 0, "n_clear": 0, "n_clear_fail": 0,
                      "travel_m": 0.0, "n_direction": 0, "n_near": 0, "n_no_signal": 0}

    # ---------------- 场景生成 ----------------
    def _gen_sources(self, n_sources=None):
        rng = self.rng
        n = int(n_sources) if n_sources else int(rng.integers(10, 17))
        n = max(1, min(n, len(CHANNELS)))
        chans = rng.choice(np.asarray(CHANNELS), size=n, replace=False)
        out = []
        for ch in chans:
            r = self.r_arena * math.sqrt(float(rng.random()))
            a = float(rng.uniform(0, 2 * math.pi))
            out.append(MockSource(channel=int(ch), x=r * math.cos(a), y=r * math.sin(a),
                                  r_recv=float(rng.uniform(R_RECV_MIN, R_RECV_MAX))))
        return out

    def _init_error_field(self, corr_m, local_w):
        """示向度误差场：平滑随机场（模拟电磁环境）+ 局地量化项；量程严格 ±1°。"""
        rng = self.rng
        self._waves = []
        for _ in range(12):
            lam = float(rng.uniform(max(corr_m * 1.5, 120.0), 3.0 * self.r_arena))
            ang = float(rng.uniform(0, 2 * math.pi))
            self._waves.append((2 * math.pi * math.cos(ang) / lam,
                                2 * math.pi * math.sin(ang) / lam,
                                float(rng.uniform(0, 2 * math.pi)),
                                1.0 / lam))
        # 归一化：用靶区内抽样点的最大绝对值做标定，再乘 1.05 并截断到 ±1
        pts = self._sample_disc(4000, rng)
        raw = self._raw_field(pts[:, 0], pts[:, 1])
        self._scale = max(float(np.max(np.abs(raw))) * 1.05, 1e-9)
        self._local_w = float(local_w)
        self._local_seed = int(rng.integers(1, 10 ** 6))

    def _sample_disc(self, n, rng):
        r = self.r_arena * np.sqrt(rng.random(n))
        a = rng.uniform(0, 2 * math.pi, n)
        return np.stack([r * np.cos(a), r * np.sin(a)], axis=1)

    def _raw_field(self, x, y):
        x = np.asarray(x, float)
        y = np.asarray(y, float)
        v = np.zeros_like(x)
        for (kx, ky, ph, a) in self._waves:
            v = v + a * np.sin(kx * x + ky * y + ph)
        return v

    def _local_field(self, x, y):
        i = int(math.floor(x / 25.0))
        j = int(math.floor(y / 25.0))
        h = (i * 73856093) ^ (j * 19349663) ^ (self._local_seed * 83492791)
        h &= 0x7FFFFFFF
        return (h % 200001) / 100000.0 - 1.0

    def env_error(self, x, y):
        """该地点的示向度误差（度，[-1,1]，同一地点固定）。"""
        v = (float(self._raw_field(x, y)) / self._scale
             + self._local_w * self._local_field(x, y))
        return float(max(-1.0, min(1.0, v)))

    # ---------------- 状态查询 ----------------
    @property
    def virtual_time_s(self):
        return self.vtime

    def remaining_budget_s(self):
        return max(0.0, self.budget - self.vtime)

    def budget_kind(self):
        return "virtual"

    def truth(self):
        return {"n_sources": len(self.sources),
                "n_cleared": sum(1 for s in self.sources if s.cleared),
                "sources": [{"channel": s.channel, "x_m": round(s.x, 3),
                             "y_m": round(s.y, 3), "r_recv_m": round(s.r_recv, 3),
                             "cleared": bool(s.cleared)} for s in self.sources]}

    # ---------------- 内部工具 ----------------
    def _src_of(self, channel):
        for s in self.sources:
            if s.channel == int(channel):
                return s
        return None

    def _reject(self, reason, **kw):
        r = {"accepted": False, "real_timestamp_ms": int(time.time() * 1000),
             "virtual_time_s": 0.0, "reason": reason}
        r.update(kw)
        return r

    def _advance(self, x, y, dt, kind, payload):
        d = math.hypot(x - self.pos[0], y - self.pos[1])
        self.stats["travel_m"] += d
        self.vtime += dt
        self.pos = (float(x), float(y))
        rec = {"kind": kind, "x": float(x), "y": float(y), "dt": float(dt),
               "dist_m": float(d), "virtual_time_s": float(self.vtime)}
        rec.update(payload)
        self.actions.append(rec)
        if self.verbose:
            print(f"    [mock] {kind} ({x:.0f},{y:.0f}) dt={dt:.1f}s "
                  f"t={self.vtime:.1f}s {payload}")
        return rec

    # ---------------- 四条指令 ----------------
    def enter(self):
        if self.entered:
            return self._reject("already_entered")
        self.entered = True
        self.actions.append({"kind": "enter", "virtual_time_s": 0.0})
        return {"accepted": True, "real_timestamp_ms": int(time.time() * 1000),
                "virtual_time_s": 0.0, "max_virtual_duration_s": self.budget,
                "max_real_duration_s": self.budget,
                "remaining_real_duration_s": int(self.budget)}

    def measure(self, x, y, channel):
        if not self.entered:
            return self._reject("not_entered")
        if self.finished:
            return self._reject("finished")
        x, y, ch = float(x), float(y), int(channel)
        if not (math.isfinite(x) and math.isfinite(y)) \
                or abs(x) > COORD_LIMIT or abs(y) > COORD_LIMIT:
            return self._reject("bad_position")
        if ch not in CHANNELS:
            return self._reject("bad_channel")
        switch = T_SWITCH if ch != self.channel else 0.0
        d = math.hypot(x - self.pos[0], y - self.pos[1])
        dt = d / V_ROBOT + switch + T_MEASURE
        if self.vtime + dt > self.budget + 1e-9:
            self.finished = True
            return self._reject("timeout")
        src = self._src_of(ch)
        svd = None
        if src is None or src.cleared:
            result = "no_signal"
        else:
            dist = math.hypot(x - src.x, y - src.y)
            in_cone = True
            if src.cone_half < 180.0:
                a = math.degrees(math.atan2(src.y - y, src.x - x))
                ref = src.dir_deg if src.dir_deg is not None else 0.0
                in_cone = abs((a - ref + 180.0) % 360.0 - 180.0) <= src.cone_half + 1e-12
            if dist > src.r_recv or not in_cone:
                result = "no_signal"
            elif dist <= NEAR_R:
                result = "near"
            else:
                result = "direction"
                a = math.degrees(math.atan2(src.y - y, src.x - x))
                svd = round((a + self.env_error(x, y)) % 360.0, 2)
        self.channel = ch
        self.stats["n_measure"] += 1
        self.stats["n_" + result] += 1
        self._advance(x, y, dt, "measure",
                      {"channel": ch, "measure_result": result, "svd_deg": svd,
                       "switch_s": switch})
        resp = {"accepted": True, "real_timestamp_ms": int(time.time() * 1000),
                "virtual_time_s": self.vtime, "measure_result": result}
        if svd is not None:
            resp["svd_deg"] = svd
        return resp

    def clear(self, x, y, channel):
        if not self.entered:
            return self._reject("not_entered")
        if self.finished:
            return self._reject("finished")
        x, y, ch = float(x), float(y), int(channel)
        if not (math.isfinite(x) and math.isfinite(y)) \
                or abs(x) > COORD_LIMIT or abs(y) > COORD_LIMIT:
            return self._reject("bad_position")
        if ch not in CHANNELS:
            return self._reject("bad_channel")
        d = math.hypot(x - self.pos[0], y - self.pos[1])
        src = self._src_of(ch)
        ok = bool(src is not None and not src.cleared
                  and math.hypot(x - src.x, y - src.y) <= R_CLEAR + 1e-9)
        dt = d / V_ROBOT + (T_CLEAR_OK if ok else T_CLEAR_FAIL)
        if self.vtime + dt > self.budget + 1e-9:
            self.finished = True
            return self._reject("timeout")
        if ok:
            src.cleared = True
        self.stats["n_clear"] += 1
        if not ok:
            self.stats["n_clear_fail"] += 1
        self._advance(x, y, dt, "clear",
                      {"channel": ch, "clear_result": "success" if ok else "no_target_in_range"})
        return {"accepted": True, "real_timestamp_ms": int(time.time() * 1000),
                "virtual_time_s": self.vtime,
                "clear_result": "success" if ok else "no_target_in_range"}

    def exit(self):
        self.finished = True
        self.actions.append({"kind": "exit", "virtual_time_s": self.vtime})
        return {"accepted": True, "real_timestamp_ms": int(time.time() * 1000),
                "virtual_time_s": self.vtime, "exit_reason": "user_exit"}


# --------------------------------------------------------------------------
# 真实模拟器客户端
# --------------------------------------------------------------------------
class HttpArena:
    """模拟器 HTTP+JSON 客户端（默认 http://127.0.0.1:2026）。"""

    def __init__(self, base_url="http://127.0.0.1:2026", robot_id="",
                 arena_id="default", timeout=5.0, max_retries=4, retry_delay=0.4,
                 log_path=None, budget_s=DEF_BUDGET_S, safety_margin_s=3.0,
                 enter_wait_s=40.0):
        self.base_url = base_url.rstrip("/")
        self.robot_id = robot_id
        self.arena_id = arena_id
        self.timeout = float(timeout)
        self.max_retries = int(max_retries)
        self.retry_delay = float(retry_delay)
        self.log_path = log_path
        self.budget = float(budget_s)
        self.safety_margin_s = float(safety_margin_s)
        self.enter_wait_s = float(enter_wait_s)
        self._n = 0
        self.entered = False
        self.pos = (0.0, 0.0)
        self.channel = 1
        self.vtime = 0.0
        self.remaining_real0 = float(budget_s)
        self._t0 = None
        self.actions: list[dict] = []
        # 与 MockArena 同名的统计（接口不返回行程，客户端自己按动作位置累加，
        # 只统计 accepted=true 的动作，与虚拟时间的推进口径一致）
        self.stats = {"n_measure": 0, "n_clear": 0, "n_clear_fail": 0,
                      "travel_m": 0.0, "n_direction": 0, "n_near": 0, "n_no_signal": 0}
        if self.log_path:
            os.makedirs(os.path.dirname(os.path.abspath(self.log_path)), exist_ok=True)
            with open(self.log_path, "w", encoding="utf-8") as f:
                f.write(json.dumps({"kind": "session", "base_url": self.base_url,
                                    "robot_id": self.robot_id}, ensure_ascii=False) + "\n")

    # ---------------- HTTP ----------------
    def _rid(self, tag):
        self._n += 1
        return f"{tag}-{self._n}"

    def _base(self, rid):
        return {"arena_id": self.arena_id, "robot_id": self.robot_id, "request_id": rid}

    def _post(self, path, payload, tag):
        data = json.dumps(payload, ensure_ascii=False).encode("utf-8")
        last_err = None
        for attempt in range(self.max_retries):
            req = urllib.request.Request(
                self.base_url + path, data=data,
                headers={"Content-Type": "application/json"}, method="POST")
            try:
                with urllib.request.urlopen(req, timeout=self.timeout) as r:
                    body = r.read().decode("utf-8")
                resp = json.loads(body)
                self._log({"kind": "resp", "path": path, "attempt": attempt, "resp": resp})
                return resp
            except urllib.error.HTTPError as e:
                body = ""
                try:
                    body = e.read().decode("utf-8", "replace")
                    resp = json.loads(body)
                    self._log({"kind": "http_error", "path": path, "code": e.code,
                               "resp": resp})
                    return resp
                except Exception:
                    self._log({"kind": "http_error", "path": path, "code": e.code,
                               "body": body[:500]})
                    raise ArenaError(f"{path} HTTP {e.code}: {body[:200]}") from e
            except (urllib.error.URLError, socket.timeout, TimeoutError, ConnectionError) as e:
                # 网络超时/中断：按协议重试完全相同的动作（复用 request_id 与内容）
                last_err = e
                self._log({"kind": "net_error", "path": path, "attempt": attempt,
                           "error": repr(e)})
                if attempt + 1 < self.max_retries:
                    time.sleep(self.retry_delay * (attempt + 1))
        raise ArenaError(f"{path} 请求失败（已重试 {self.max_retries} 次）：{last_err!r}")

    def _log(self, rec):
        rec["wall_s"] = round(time.time() - (self._t0 or time.time()), 3)
        self.actions.append(rec)
        if self.log_path:
            with open(self.log_path, "a", encoding="utf-8") as f:
                f.write(json.dumps(rec, ensure_ascii=False) + "\n")

    def _new(self, path, extra, tag):
        rid = self._rid(tag)
        payload = self._base(rid)
        payload.update(extra)
        self._log({"kind": "req", "path": path, "payload": payload})
        return self._post(path, payload, tag)

    # ---------------- 状态 ----------------
    @property
    def virtual_time_s(self):
        return self.vtime

    def remaining_budget_s(self):
        if self._t0 is None:
            return self.remaining_real0
        used = time.monotonic() - self._t0
        return max(0.0, self.remaining_real0 - used - self.safety_margin_s)

    def budget_kind(self):
        return "real"

    def truth(self):
        return None

    # ---------------- 四条指令 ----------------
    def enter(self):
        """进入靶区。接口未开放时（倒计时未结束 / 测试未开始）会重试同一 request_id。

        协议规定"同一 request_id + 同一内容返回第一次的完整响应"，所以用固定的
        request_id 重试 /enter 是安全的（不会重复进入、也不会推进时钟）。
        """
        rid = self._rid("enter")
        payload = self._base(rid)
        self._log({"kind": "req", "path": "/enter", "payload": payload})
        t0 = time.monotonic()
        last = None
        keep_retries = self.max_retries
        self.max_retries = min(keep_retries, 2)   # 接口未开放时缩短单次重试，尽快进入下一轮
        try:
            while True:
                try:
                    resp = self._post("/enter", payload, "enter")
                except ArenaError as e:
                    resp, last = None, e
                if resp is not None and resp.get("accepted") is True:
                    self.entered = True
                    self._t0 = time.monotonic()
                    self.remaining_real0 = float(resp.get("remaining_real_duration_s",
                                                          self.budget))
                    self.vtime = float(resp.get("virtual_time_s", 0.0))
                    return resp
                if resp is not None and resp.get("accepted") is False \
                        and resp.get("reason") == "already_entered":
                    return resp
                if time.monotonic() - t0 > self.enter_wait_s:
                    if resp is not None:
                        return resp
                    raise ArenaError(
                        f"/enter 在 {self.enter_wait_s:.0f} s 内一直失败：{last!r}")
                time.sleep(2.0)
        finally:
            self.max_retries = keep_retries

    def measure(self, x, y, channel):
        resp = self._new("/measure", {"position": {"x": float(x), "y": float(y)},
                                      "channel": int(channel)}, "measure")
        if resp.get("accepted") is True:
            self.stats["travel_m"] += math.hypot(float(x) - self.pos[0],
                                                 float(y) - self.pos[1])
            self.pos = (float(x), float(y))
            self.channel = int(channel)
            self.vtime = float(resp.get("virtual_time_s", self.vtime))
            res = resp.get("measure_result")
            self.stats["n_measure"] += 1
            if res in ("direction", "near", "no_signal"):
                self.stats["n_" + res] += 1
        return resp

    def clear(self, x, y, channel):
        resp = self._new("/clear", {"position": {"x": float(x), "y": float(y)},
                                    "channel": int(channel)}, "clear")
        if resp.get("accepted") is True:
            self.stats["travel_m"] += math.hypot(float(x) - self.pos[0],
                                                 float(y) - self.pos[1])
            self.pos = (float(x), float(y))
            self.vtime = float(resp.get("virtual_time_s", self.vtime))
            self.stats["n_clear"] += 1
            if resp.get("clear_result") != "success":
                self.stats["n_clear_fail"] += 1
        return resp

    def exit(self):
        return self._new("/exit", {}, "exit")


# --------------------------------------------------------------------------
# 自检
# --------------------------------------------------------------------------
def selfcheck(verbose=True, seed=20260913):
    """核对 MockArena 是否忠实复刻附件 1/2 的规则。"""
    out = []
    ok_all = True

    def rec(name, ok, detail=""):
        nonlocal ok_all
        ok_all = ok_all and bool(ok)
        out.append({"check": name, "ok": bool(ok), "detail": detail})
        if verbose:
            print(f"[{'OK ' if ok else 'FAIL'}] {name}  {detail}")

    # A1 附件 1 表 2 的计时算例（(300,400) 频道1 -> 频道2 -> /clear 未发现 -> 频道2）
    a = MockArena(seed=1, sources=[MockSource(3, 100000.0, 0.0, 1000.0)])
    a.enter()
    r1 = a.measure(300.0, 400.0, 1)
    r2 = a.measure(300.0, 400.0, 2)
    r3 = a.clear(300.0, 0.0, 3)
    r4 = a.measure(300.0, 0.0, 2)
    ok1 = (abs(r1["virtual_time_s"] - 105.0) < 1e-9
           and abs(r2["virtual_time_s"] - 111.0) < 1e-9
           and abs(r3["virtual_time_s"] - 194.0) < 1e-9
           and abs(r4["virtual_time_s"] - 199.0) < 1e-9)
    rec("A1 计时算例与附件1表2逐条一致（105/111/194/199 s）", ok1,
        f"{r1['virtual_time_s']}, {r2['virtual_time_s']}, "
        f"{r3['virtual_time_s']}, {r4['virtual_time_s']}")

    # A2 /clear 不切换频道，也不产生切换耗时
    b = MockArena(seed=2, sources=[MockSource(1, 0.0, 0.0, 1000.0)])
    b.enter()
    b.measure(0.0, 100.0, 1)
    t0 = b.virtual_time_s
    b.clear(0.0, 50.0, 7)
    rec("A2 /clear 不切换频道、不产生切换耗时",
        abs(b.virtual_time_s - t0 - (50.0 / V_ROBOT + T_CLEAR_FAIL)) < 1e-9
        and b.channel == 1,
        f"Δt={b.virtual_time_s - t0:.1f}s, 当前频道={b.channel}")

    # A3 三种检测结果与清除半径
    src = MockSource(5, 0.0, 0.0, 1000.0)
    c = MockArena(seed=3, sources=[src])
    c.enter()
    m_far = c.measure(0.0, 0.0, 5)          # 距离 0 <= 5 m -> near
    m_dir = c.measure(0.0, 300.0, 5)        # direction（真方位 270°）
    m_no = c.measure(0.0, 1200.0, 5)        # 超出 1000 m -> no_signal
    ok3 = (m_far["measure_result"] == "near" and m_dir["measure_result"] == "direction"
           and abs(m_dir["svd_deg"] - 270.0) <= BEARING_ERR + 1e-9
           and m_no["measure_result"] == "no_signal")
    rec("A3 near / direction / no_signal 三种结果与真方位一致", ok3,
        f"near={m_far['measure_result']}, dir={m_dir.get('svd_deg')}°, "
        f"no={m_no['measure_result']}")
    cl_far = c.clear(0.0, 1000.0, 5)
    cl_ok = c.clear(0.0, 20.0, 5)
    rec("A4 清除半径 20 m：远处未发现、20 m 内成功",
        cl_far["clear_result"] == "no_target_in_range"
        and cl_ok["clear_result"] == "success",
        f"{cl_far['clear_result']} / {cl_ok['clear_result']}")
    rec("A5 同一干扰源只能清除一次",
        c.clear(0.0, 0.0, 5)["clear_result"] == "no_target_in_range",
        "二次清除返回 no_target_in_range")

    # A6 误差界与"同一地点固定"
    d = MockArena(seed=4)
    e1 = d.env_error(123.0, -456.0)
    e2 = d.env_error(123.0, -456.0)
    vals = [d.env_error(float(x), float(y)) for x, y in d._sample_disc(500, d.rng)]
    rec("A6 示向度误差：|e|<=1° 且同一地点固定",
        abs(e1 - e2) < 1e-12 and max(abs(v) for v in vals) <= 1.0 + 1e-12,
        f"e={e1:.4f}，重复一致；500 点最大 |e|={max(abs(v) for v in vals):.4f}")

    # A7 真源一定落在"测得方位 ±1°"的楔内（区间算法成立的前提）
    bad = 0
    nchk = 0
    d2 = MockArena(seed=4, budget_s=1e9)
    d2.enter()
    for s in d2.sources:
        for _ in range(3):
            aa = float(d2.rng.uniform(0, 2 * math.pi))
            dd = float(d2.rng.uniform(50.0, 0.9 * s.r_recv))
            x, y = s.x + dd * math.cos(aa), s.y + dd * math.sin(aa)
            r = d2.measure(x, y, s.channel)
            if r.get("accepted") is not True or r.get("measure_result") != "direction":
                continue
            nchk += 1
            true_a = math.degrees(math.atan2(s.y - y, s.x - x)) % 360.0
            if abs((true_a - r["svd_deg"] + 180.0) % 360.0 - 180.0) > BEARING_ERR + 1e-9:
                bad += 1
    rec("A7 真方位与示向度之差恒 <= 1°", bad == 0 and nchk > 0,
        f"{nchk} 次检测，越界 {bad} 次")

    # A8 超时拒动
    e = MockArena(seed=5, budget_s=50.0, sources=[MockSource(1, 0.0, 0.0, 1000.0)])
    e.enter()
    e.measure(0.0, 1000.0, 1)           # 200 + 5 = 205 s > 50 s
    rec("A8 超出时间预算的请求 accepted=false 且不推进时钟",
        e.virtual_time_s == 0.0 and e.finished,
        f"vtime={e.virtual_time_s}")

    if verbose:
        print(f"\n自检总结：{'全部通过' if ok_all else '存在失败项'}")
    return {"ok": ok_all, "checks": out}


if __name__ == "__main__":
    import argparse
    ap = argparse.ArgumentParser(description="问题3 接口层自检")
    ap.add_argument("--selfcheck", action="store_true")
    args = ap.parse_args()
    res = selfcheck()
    raise SystemExit(0 if res["ok"] else 1)

\end{lstlisting}
\subsection{\texttt{src/p3\_expect\_field.py}}
% SHA256 c21e97f8467b5bc77e20d977252835793a6ebc5886a0a952497dadf6001eeb35
\begin{lstlisting}[language=python,escapeinside={(*@}{@*)}]
"""B 题问题 3：机器狗"基于期望的行动"决策场。

本模块只负责问题 3 的**决策层**：把"下一步该去哪里检测"变成一个可以在全局网格上
逐格比较的标量场。它严格按下面的思路实现：

1. **全局网格**：在半径 1800 m 的靶区圆域内铺一张规则网格
   （``grid_step``，默认 100 m），每个格子记录"如果机器狗去那里做下一次检测，
   这一步有多好"。
2. **第二问的网格模拟结果**：直接复用 :mod:`p2_grid_expectation` 产出的热力图
   ``out/p2_grid/p2_grid_map.csv``。该图是在"第一检测点 :math:`S_1=(0,0)`、
   示向度 :math:`\\theta_1=0^\\circ`"的基准几何下算出的
   :math:`\\mathbb E[D\\mid\\text{可检测}]`（D = 交会定位区域直径）。
3. **变换 + 插值**：E[D] 在**刚体变换下不变**（旋转 + 平移保持距离与夹角）。
   因此真实一次检测 ``(S, (*@\ensuremath{\theta}@*))`` 对应的场，就是把基准图**绕原点旋转 (*@\ensuremath{\theta}@*)、再平移到 S**：

       b = R(-(*@\ensuremath{\theta}@*))·(g - S)     （g 是全局网格点，b 是基准图坐标）
       E_g = E_base(b)       （双线性插值）

4. **逐次叠加**：把每个"待测频道"的每一次检测都按上式投到全局网格上叠加。
5. **未覆盖的网格取最大值**：基准图里没有数据（靶区外 / 该点不可检测 / 近轴退化带）
   或数值超过上限的格子一律记 :math:`E_{\\max}=1500` m —— 也就是"最差的位置"。
   这样它们只有在附近没有更好的格子时才会被选中，天然形成探索行为。
6. **归一化**：``E_norm = E_mean / E_max (*@\ensuremath{\in}@*) [0,1]``（0 = 交会几何最好，1 = 最差）。
7. **距离惩罚层**：行动要花时间，所以在归一化场之上再叠加一层**与探测器距离线性增长**
   的代价：

       score(g) = E_norm(g) + (*@\ensuremath{\lambda}@*) · |g - p_robot| / (2·R_arena)

   (*@\ensuremath{\lambda}@*) 即"距离项系数"，可调（``dist_weight``，默认 0.35；(*@\ensuremath{\lambda}@*)=1 表示"跑到靶区对面"
   等价于把定位质量从最好拉到最差）。已访问过的扫描点邻域、以及离当前位置过近的格子
   会被直接排除（重复检测不产生新信息）。

8. **选点**：取 ``argmin score``。若整张图都被排除（例如一条方位都还没有），退化为
   "最大最小距离"探索点（离所有已访问点最远的格子）。

用法
----
    python src/p3_expect_field.py --selfcheck
    python src/p3_expect_field.py --demo --dist-weight 0.35 --out out/p3
"""
from __future__ import annotations

import argparse
import csv
import json
import math
import os
import sys
from dataclasses import dataclass, field as dc_field

import numpy as np

_HERE = os.path.dirname(os.path.abspath(__file__))
if _HERE not in sys.path:
    sys.path.insert(0, _HERE)

from p1_intersection import BEARING_ERR, ang_diff  # noqa: E402

# --------------------------------------------------------------------------
# 物理常数（题目附录 1/2 与模拟器接口说明）
# --------------------------------------------------------------------------
R_ARENA = 1800.0            # 靶区半径（m）
E_MAX = 1500.0              # 未覆盖网格的默认最大值（m）——即"最差位置"
R_RECV_MIN = 1000.0         # 有效接收半径下界（m）
R_RECV_MAX = 1500.0         # 有效接收半径上界（m）
NEAR_R = 5.0                # 近场阈值（m）
R_CLEAR = 20.0              # 清除半径（m）
DIST_REF = 2.0 * R_ARENA    # 距离惩罚的归一化尺度（靶区直径，m）

DEFAULT_BASE_CSV = os.path.normpath(
    os.path.join(_HERE, "..", "out", "p2_grid", "p2_grid_map.csv"))
DEFAULT_FAMILY_DIR = os.path.normpath(os.path.join(_HERE, "..", "out", "p3_family"))
E_COL_CANDIDATES = ("E_diam_given_detectable_m", "E_diam_mixed_mm",
                    "E_diam_mixed_m")


# --------------------------------------------------------------------------
# 基准图（第二问的网格模拟结果）
# --------------------------------------------------------------------------
@dataclass
class BaseMap:
    """问题 2 的期望热力图：``E[j, i]`` 对应 ``(xs[i], ys[j])``，nan = 无数据。"""

    xs: np.ndarray
    ys: np.ndarray
    E: np.ndarray
    source: str = "p2_grid_expectation"
    e_col: str = "E_diam_given_detectable_m"

    # ---------------- 构造 ----------------
    @property
    def step(self) -> float:
        return float(self.xs[1] - self.xs[0])

    @property
    def extent(self) -> tuple:
        return (float(self.xs[0]), float(self.xs[-1]), float(self.ys[0]), float(self.ys[-1]))

    @classmethod
    def from_points(cls, X, Y, E, source="synthetic") -> "BaseMap":
        """由散点（规则网格掩码后）构造。"""
        X = np.asarray(X, float)
        Y = np.asarray(Y, float)
        E = np.asarray(E, float)
        xs = np.unique(np.round(X, 6))
        ys = np.unique(np.round(Y, 6))
        # 步长取最小正间隔（避免浮点噪声把网格拆散）
        dx = _grid_step(xs)
        dy = _grid_step(ys)
        xs = np.round(np.arange(xs[0], xs[-1] + 1e-9, dx), 6)
        ys = np.round(np.arange(ys[0], ys[-1] + 1e-9, dy), 6)
        G = np.full((ys.size, xs.size), np.nan)
        ix = np.searchsorted(xs, np.round(X, 6))
        iy = np.searchsorted(ys, np.round(Y, 6))
        ok = (ix < xs.size) & (iy < ys.size)
        G[iy[ok], ix[ok]] = E[ok]
        return cls(xs=xs, ys=ys, E=G, source=source)

    @classmethod
    def load_csv(cls, path=DEFAULT_BASE_CSV, e_col=None) -> "BaseMap":
        if not os.path.exists(path):
            raise FileNotFoundError(
                f"找不到第二问的网格模拟结果 {path}。请先运行\n"
                f"  python src/p2_grid_expectation.py --step 40 --out out/p2_grid")
        with open(path, encoding="utf-8-sig") as f:
            rows = list(csv.DictReader(f))
        if not rows:
            raise ValueError(f"{path} 为空")
        col = e_col
        if col is None:
            for c in E_COL_CANDIDATES:
                if c in rows[0]:
                    col = c
                    break
        if col is None:
            raise ValueError(f"{path} 缺少期望值列，现有列：{list(rows[0])}")
        X = np.array([float(r["x_m"]) for r in rows])
        Y = np.array([float(r["y_m"]) for r in rows])
        E = np.array([_to_float(r.get(col)) for r in rows])
        m = cls.from_points(X, Y, E, source=os.path.basename(path))
        m.e_col = col
        return m

    @classmethod
    def synthetic(cls, step=100.0, R=R_ARENA) -> "BaseMap":
        """自检用解析基准图（形状与真实图同类：近轴退化带 + 斜前方最优翼）。"""
        xs = np.arange(-R, R + 1e-9, step)
        X, Y = np.meshgrid(xs, xs)
        r = np.hypot(X, Y)
        phi = np.abs(np.degrees(np.arctan2(Y, X)))
        E = 60.0 + 900.0 * (1.0 - np.exp(-r / 700.0)) + 700.0 * np.exp(-phi / 12.0)
        E[r > R] = np.nan
        return cls(xs=xs, ys=xs.copy(), E=E, source="synthetic")

    # ---------------- 查询 ----------------
    def lookup(self, X, Y, e_max=E_MAX):
        """双线性插值查询。返回 ``(E, covered)``；四点任一无数据即视为未覆盖。

        ``E`` 已按"未覆盖/超过上限一律取 e_max"的规则裁剪，可直接用于叠加。
        """
        X = np.asarray(X, float)
        Y = np.asarray(Y, float)
        shp = np.broadcast(X, Y).shape
        Xf = np.broadcast_to(X, shp).astype(float).ravel()
        Yf = np.broadcast_to(Y, shp).astype(float).ravel()
        x0, y0 = float(self.xs[0]), float(self.ys[0])
        dx, dy = self.step, float(self.ys[1] - self.ys[0])
        nx, ny = self.xs.size, self.ys.size
        fx = (Xf - x0) / dx
        fy = (Yf - y0) / dy
        i = np.floor(fx).astype(int)
        j = np.floor(fy).astype(int)
        inside = (i >= 0) & (i <= nx - 2) & (j >= 0) & (j <= ny - 2)
        ic = np.clip(i, 0, nx - 2)
        jc = np.clip(j, 0, ny - 2)
        tx = np.clip(fx - ic, 0.0, 1.0)
        ty = np.clip(fy - jc, 0.0, 1.0)
        G = self.E
        e00 = G[jc, ic]
        e10 = G[jc, ic + 1]
        e01 = G[jc + 1, ic]
        e11 = G[jc + 1, ic + 1]
        ok = inside & np.isfinite(e00) & np.isfinite(e10) & np.isfinite(e01) & np.isfinite(e11)
        val = (e00 * (1 - tx) * (1 - ty) + e10 * tx * (1 - ty)
               + e01 * (1 - tx) * ty + e11 * tx * ty)
        val = np.where(ok, val, np.nan)
        E_out = np.where(ok, np.minimum(val, e_max), e_max)
        return E_out.reshape(shp), ok.reshape(shp)

    # ---------------- 变换 ----------------
    def transformed_lookup(self, S, theta_deg, X, Y, e_max=E_MAX):
        """把基准图按"检测点 S + 示向度 (*@\ensuremath{\theta}@*)"旋转平移后查询全局网格点 (X, Y)。

        ``b = R(-(*@\ensuremath{\theta}@*))·(g - S)``：把全局点转回基准图坐标系（基准图的源在 +x 轴方向）。
        E[D] 在刚体变换下不变，故 ``E(g) = E_base(b)``。
        """
        S = np.asarray(S, float)
        c = math.cos(math.radians(theta_deg))
        s = math.sin(math.radians(theta_deg))
        Xf = np.asarray(X, float) - S[0]
        Yf = np.asarray(Y, float) - S[1]
        bx = Xf * c + Yf * s
        by = -Xf * s + Yf * c
        return self.lookup(bx, by, e_max=e_max)


def _grid_step(v):
    d = np.diff(np.asarray(v, float))
    d = d[d > 1e-9]
    return float(np.min(d)) if d.size else 1.0


def _to_float(v):
    if v is None or v == "" or v == "nan":
        return float("nan")
    try:
        return float(v)
    except (TypeError, ValueError):
        return float("nan")


def save_map_csv(path, X, Y, E, e_col="E_diam_given_detectable_m"):
    os.makedirs(os.path.dirname(os.path.abspath(path)), exist_ok=True)
    with open(path, "w", newline="", encoding="utf-8-sig") as f:
        w = csv.writer(f)
        w.writerow(["x_m", "y_m", e_col])
        for i in range(len(X)):
            v = E[i]
            w.writerow([round(float(X[i]), 4), round(float(Y[i]), 4),
                        "" if not np.isfinite(v) else round(float(v), 6)])
    return path


# --------------------------------------------------------------------------
# 基准图族：按"先验外径 t_hi"参数化
# --------------------------------------------------------------------------
class BaseMapFamily:
    """一族基准图 ``{t_hi: BaseMap}``。

    **为什么需要它（这一步不做对，整个"变换 + 插值"就不成立）**

    基准图里源位先验是"从检测点沿示向度射线、``t (*@\ensuremath{\in}@*) [5, t_hi]``、权重 (*@\ensuremath{\propto}@*) t"。
    真实场景中源还必须落在靶区内，所以射线上先验的**外径** ``t_hi`` 随 ``(S, (*@\ensuremath{\theta}@*))``
    变化（射线出靶区的距离可以只有几百米）。而 :math:`\\mathbb E[D]` 只依赖
    (相对位置 b, t_hi) 这一对量，因此把 t_hi 也当作族参数，旋转 + 平移 + 插值才严格成立。

    查询时按 ``t_hi`` 在族内做线性插值（先把未覆盖格子填成 e_max，再混合）。
    """

    def __init__(self, maps, t_ref=R_RECV_MAX, e_max=E_MAX):
        self.maps = {float(k): v for k, v in sorted(maps.items())}
        self.t_list = np.array(sorted(self.maps.keys()))
        self.t_ref = float(t_ref)
        self.e_max = float(e_max)
        self._cache = {}

    # ---------------- 构造 ----------------
    @classmethod
    def build(cls, outdir, t_his=(150.0, 250.0, 350.0, 450.0, 600.0, 800.0, 1000.0,
                                  1250.0, 1500.0), step=40.0, n_t=200, verbose=True):
        """按 p2 的口径重算一族基准图并落盘。"""
        from p2_grid_expectation import (  # 局部导入，避免循环依赖
            expected_diameter_measured, make_grid, source_samples)

        X, Y, xs, ys, shape = make_grid(step=step, arena_r=R_ARENA, domain="arena")
        os.makedirs(outdir, exist_ok=True)
        maps = {}
        files = []
        for T in t_his:
            ts, w, info = source_samples((0.0, 0.0), 0.0, n_t=n_t, r_arena=float(T))
            res = expected_diameter_measured(X, Y, ts, w)
            E = np.asarray(res["E_det"], float)
            m = BaseMap.from_points(X, Y, E, source=f"p3_family_thi{T:g}")
            fn = os.path.join(outdir, f"p2_grid_map_thi{int(round(T)):04d}.csv")
            save_map_csv(fn, X, Y, E)
            maps[float(T)] = m
            files.append(os.path.basename(fn))
            if verbose:
                fin = E[np.isfinite(E)]
                print(f"  t_hi={T:6.0f} m -> {os.path.basename(fn)}  "
                      f"有效格 {fin.size}/{E.size}，中位数 "
                      f"{np.median(fin) if fin.size else float('nan'):.1f} m")
        man = {"kind": "p3_base_map_family", "step": float(step), "n_t": int(n_t),
               "arena_r": R_ARENA, "t_his": [float(t) for t in t_his], "files": files,
               "note": "每个文件是 S1=(0,0)、(*@\ensuremath{\theta}@*)1=0 下，源位先验 t(*@\ensuremath{\in}@*)[5,t_hi] 的 E[D|可检测]"}
        with open(os.path.join(outdir, "p2_grid_family.json"), "w", encoding="utf-8") as f:
            json.dump(man, f, ensure_ascii=False, indent=2)
        return cls(maps)

    @classmethod
    def load_dir(cls, outdir):
        man_path = os.path.join(outdir, "p2_grid_family.json")
        if not os.path.exists(man_path):
            return None
        with open(man_path, encoding="utf-8") as f:
            man = json.load(f)
        maps = {}
        for t, fn in zip(man["t_his"], man["files"]):
            maps[float(t)] = BaseMap.load_csv(os.path.join(outdir, fn))
        return cls(maps)

    @classmethod
    def load_or_build(cls, outdir, **kw):
        fam = cls.load_dir(outdir)
        if fam is not None:
            return fam
        return cls.build(outdir, **kw)

    # ---------------- 查询 ----------------
    def select(self, t_hi):
        """返回该 t_hi 对应的（可缓存的）混合基准图。"""
        t = float(min(max(t_hi, self.t_list[0]), self.t_list[-1]))
        key = round(t / 25.0)
        if key in self._cache:
            return self._cache[key]
        tt = key * 25.0
        if tt <= self.t_list[0]:
            m = self.maps[float(self.t_list[0])]
        elif tt >= self.t_list[-1]:
            m = self.maps[float(self.t_list[-1])]
        else:
            i = int(np.searchsorted(self.t_list, tt) - 1)
            T1, T2 = float(self.t_list[i]), float(self.t_list[i + 1])
            a = (tt - T1) / max(T2 - T1, 1e-9)
            m = blend_maps(self.maps[T1], self.maps[T2], a, self.e_max)
        if len(self._cache) < 64:
            self._cache[key] = m
        return m

    def rebase(self, S, theta_deg):
        """给出 (S, (*@\ensuremath{\theta}@*)) 检测对应的基准图：先按射线出靶区距离取 t_hi，再在族内插值。"""
        S = np.asarray(S, float)
        u = np.array([math.cos(math.radians(theta_deg)), math.sin(math.radians(theta_deg))])
        b = float(S @ u)
        c = float(S @ S) - R_ARENA * R_ARENA
        disc = b * b - c
        exit_d = float(-b + math.sqrt(disc)) if disc > 0 else 0.0
        t_hi = min(R_RECV_MAX, exit_d)
        return self.select(t_hi), t_hi


def blend_maps(m1: BaseMap, m2: BaseMap, a, e_max=E_MAX) -> BaseMap:
    """把两张同网格基准图按权重 a 线性混合（未覆盖格子先填 e_max，偏保守）。"""
    a = float(min(max(a, 0.0), 1.0))
    E1 = np.where(np.isfinite(m1.E), m1.E, e_max)
    E2 = np.where(np.isfinite(m2.E), m2.E, e_max)
    covered = np.isfinite(m1.E) | np.isfinite(m2.E)
    E = np.where(covered, (1.0 - a) * E1 + a * E2, np.nan)
    return BaseMap(xs=m1.xs, ys=m1.ys, E=E, source=f"blend({m1.source},{m2.source},{a:.2f})")


# --------------------------------------------------------------------------
# 决策场
# --------------------------------------------------------------------------
@dataclass
class FieldConfig:
    grid_step: float = 100.0        # 全局网格步长（m）
    e_max: float = E_MAX            # 未覆盖网格的默认最大值（m）
    dist_weight: float = 0.35       # 距离项系数 (*@\ensuremath{\lambda}@*)（可调）
    visited_radius: float = 250.0   # 已访问扫描点的排除半径（m）
    min_move: float = 60.0          # 离当前位置太近的格子直接排除（m）
    arena_r: float = R_ARENA
    normalize: str = "mean"         # mean（期望均值归一化）| max


class ExpectField:
    """全局网格上的"期望定位质量 + 距离"标量场。"""

    def __init__(self, base, cfg: FieldConfig | None = None):
        self.base = base
        self.family = base if isinstance(base, BaseMapFamily) else None
        self.cfg = cfg or FieldConfig()
        step = float(self.cfg.grid_step)
        R = float(self.cfg.arena_r)
        xs = np.arange(-R, R + 1e-9, step)
        GX, GY = np.meshgrid(xs, xs)
        m = np.hypot(GX, GY) <= R + 1e-9
        self.gx = GX[m]
        self.gy = GY[m]
        self.terms: list[tuple[np.ndarray, float, tuple]] = []

    # ---------------- 叠加 ----------------
    def add_bearing(self, S, theta_deg, weight=1.0, tag=None):
        """叠加一次检测 ``(S, (*@\ensuremath{\theta}@*))`` 给出的期望场（已裁剪：未覆盖 = e_max）。

        有基准图族时先用 (S, (*@\ensuremath{\theta}@*)) 的"射线出靶区距离"选 t_hi，再在族内插值，
        这样"旋转 + 平移 + 插值"才是严格的刚体不变。
        """
        S = np.asarray(S, float)
        bm = self.base
        t_hi = R_RECV_MAX
        if self.family is not None:
            bm, t_hi = self.family.rebase(S, theta_deg)
        E, cov = bm.transformed_lookup(S, theta_deg, self.gx, self.gy,
                                       e_max=self.cfg.e_max)
        self.terms.append((E, float(weight),
                           (tuple(float(v) for v in S), float(theta_deg), tag, float(t_hi))))
        return cov

    def add_channel_bearings(self, bearings):
        """bearings = [(x, y, svd_deg), ...]"""
        for (x, y, th) in bearings:
            self.add_bearing((x, y), th)

    def add_raw(self, values, weight=1.0, tag=None):
        """叠加一层已经算好的、直接定义在全局网格上的值（用于"探索项"）。"""
        v = np.asarray(values, float)
        if v.shape != self.gx.shape:
            v = np.broadcast_to(v, self.gx.shape)
        self.terms.append((v, float(weight), ("raw", tag)))
        return v

    def discovery_term(self, meas_pts, min_dist_m, e_max=None):
        """探索项：某频道"在哪些格点上重新检测才有新信息"。

        对"还没听到过"的频道，若新检测点距它所有历史检测点都 < ``min_dist_m``，
        那么能新覆盖到的区域很小（有效接收半径至少 1000 m），几乎不可能听到 —— 记 e_max；
        反之记 0（无惩罚）。这样全局网格就会自动把扫描点推到"能发现新源"的地方。
        """
        e_max = float(self.cfg.e_max if e_max is None else e_max)
        if not meas_pts:
            return np.zeros(self.gx.shape)
        P = np.asarray(meas_pts, float)
        dmin = np.min(np.hypot(self.gx[:, None] - P[None, :, 0],
                               self.gy[:, None] - P[None, :, 1]), axis=1)
        return np.where(dmin >= float(min_dist_m), 0.0, e_max)

    # ---------------- 聚合 / 打分 ----------------
    def aggregate(self):
        """返回 ``(E_mean_m, coverage_count, n_terms)``。

        * 每一项先把未覆盖格子记成 ``e_max``（"最差位置"）；
        * ``normalize='mean'``：对各项取加权平均（即"总和归一化"）；
        * ``coverage_count``：至少有一项覆盖到的格子数。
        """
        n = self.gx.size
        if not self.terms:
            return (np.full(n, float(self.cfg.e_max)), np.zeros(n, dtype=int), 0)
        W = 0.0
        acc = np.zeros(n)
        cov = np.zeros(n, dtype=int)
        covered = np.zeros(n, dtype=bool)
        for (E, w, _tag) in self.terms:
            acc += w * E
            W += w
            covered |= (E < self.cfg.e_max - 1e-9)
            cov += (E < self.cfg.e_max - 1e-9).astype(int)
        E_mean = acc / max(W, 1e-12)
        if self.cfg.normalize == "max" and np.any(covered):
            ref = float(np.max(E_mean))
            E_mean = E_mean * (float(self.cfg.e_max) / max(ref, 1e-12))
            E_mean = np.minimum(E_mean, float(self.cfg.e_max))
        return E_mean, cov, len(self.terms)

    def score(self, cur_xy, visited=(), extra_penalty=None, max_dist=None):
        """计算整张网格的 score = 归一化期望 + (*@\ensuremath{\lambda}@*)·距离/靶区直径（排除项 = +inf）。"""
        E_mean, cov, n_terms = self.aggregate()
        e_norm = E_mean / max(float(self.cfg.e_max), 1e-12)
        if extra_penalty is not None:
            e_norm = e_norm + np.asarray(extra_penalty, float)
        p = np.asarray(cur_xy, float)
        d = np.hypot(self.gx - p[0], self.gy - p[1])
        score = e_norm + float(self.cfg.dist_weight) * d / DIST_REF
        mask = d < float(self.cfg.min_move)
        if max_dist is not None:
            mask |= d > float(max_dist)
        for q in visited:
            q = np.asarray(q, float)
            mask |= np.hypot(self.gx - q[0], self.gy - q[1]) < float(self.cfg.visited_radius)
        score = np.where(mask, np.inf, score)
        return score, {"E_mean_m": E_mean, "coverage": cov, "n_terms": n_terms,
                       "e_norm": e_norm, "dist_m": d, "excluded": mask}

    # ---------------- 选点 ----------------
    def argmin(self, cur_xy, visited=(), extra_penalty=None, max_dist=None):
        if not self.terms:
            return None            # 一条方位都没有：交给探索兜底
        score, info = self.score(cur_xy, visited, extra_penalty, max_dist=max_dist)
        if not np.any(np.isfinite(score)):
            return None
        i = int(np.argmin(np.where(np.isfinite(score), score, np.inf)))
        return self._describe(i, score[i], info)

    def fallback_point(self, cur_xy, visited=(), prefer_near=True, max_dist=None):
        """探索兜底：离所有已访问点"最远"的格子（最大最小距离准则）。

        一条方位都没有、或整张图都被排除时使用；也用于避免反复在同一片区域打转。
        ``max_dist`` 给定时只在"还走得到"的格子里挑（用于时间预算收尾阶段）。
        """
        p = np.asarray(cur_xy, float)
        if visited:
            V = np.asarray(visited, float)
            dmin = np.min(np.hypot(self.gx[:, None] - V[None, :, 0],
                                   self.gy[:, None] - V[None, :, 1]), axis=1)
        else:
            dmin = np.hypot(self.gx - p[0], self.gy - p[1])
        dc = np.hypot(self.gx - p[0], self.gy - p[1])
        key = dmin - (1e-6 * dc if prefer_near else 0.0)
        key = np.where(dc < float(self.cfg.min_move), -np.inf, key)
        if max_dist is not None:
            key = np.where(dc > float(max_dist), -np.inf, key)
        if not np.any(np.isfinite(key)):
            key = dmin - 1e-6 * dc          # 放宽：忽略 max_dist/min_move 再挑一次
        i = int(np.argmax(key))
        E_mean, cov, n_terms = self.aggregate()
        info = {"E_mean_m": E_mean, "coverage": cov, "n_terms": n_terms,
                "e_norm": E_mean / max(float(self.cfg.e_max), 1e-12),
                "dist_m": dc, "excluded": np.zeros(self.gx.size, bool)}
        return self._describe(i, float("nan"), info, kind="explore")

    def _describe(self, i, score, info, kind="min"):
        return {
            "kind": kind,
            "index": int(i),
            "xy": (float(self.gx[i]), float(self.gy[i])),
            "score": float(score),
            "E_mean_m": float(info["E_mean_m"][i]),
            "E_norm": float(info["e_norm"][i]),
            "dist_m": float(info["dist_m"][i]),
            "dist_penalty": float(self.cfg.dist_weight * info["dist_m"][i] / DIST_REF),
            "n_terms": int(info["n_terms"]),
            "n_cover": int(info["coverage"][i]),
            "covered": bool(info["coverage"][i] > 0),
        }

    # ---------------- 导出 ----------------
    def save_csv(self, path):
        E_mean, cov, n_terms = self.aggregate()
        score, info = self.score((0.0, 0.0))
        os.makedirs(os.path.dirname(os.path.abspath(path)), exist_ok=True)
        with open(path, "w", newline="", encoding="utf-8-sig") as f:
            w = csv.writer(f)
            w.writerow(["x_m", "y_m", "E_mean_m", "E_norm", "n_cover", "score_at_origin"])
            for k in range(self.gx.size):
                w.writerow([round(float(self.gx[k]), 3), round(float(self.gy[k]), 3),
                            round(float(E_mean[k]), 4), round(float(E_mean[k]) / self.cfg.e_max, 6),
                            int(cov[k]),
                            "" if not np.isfinite(score[k]) else round(float(score[k]), 6)])
        return path


# --------------------------------------------------------------------------
# 精确期望（独立实现）：任意 S1/(*@\ensuremath{\theta}@*)1 下，候选第二检测点的 E[D]
# --------------------------------------------------------------------------
def ray_source_samples(S, theta_deg, err=BEARING_ERR, n_t=200, r_arena=R_ARENA,
                       r_recv_max=R_RECV_MAX, near=NEAR_R):
    """源位先验：在 S 沿 (*@\ensuremath{\theta}@*) 的射线上、靶区内、可被有效接收的范围内按面积均匀采样。

    与 :func:`p2_grid_expectation.source_samples` 同一先验，但支持任意 S（不只是原点）。
    返回 ``(ts, w, t_hi)``：``ts`` 为源到 S 的距离采样。
    """
    S = np.asarray(S, float)
    u = np.array([math.cos(math.radians(theta_deg)), math.sin(math.radians(theta_deg))])
    b = float(S @ u)
    c = float(S @ S) - r_arena * r_arena
    disc = b * b - c
    exit_d = float(-b + math.sqrt(disc)) if disc > 0 else 0.0
    t_hi = min(r_recv_max, exit_d)
    if t_hi <= near:
        return np.zeros(0), np.zeros(0), 0.0
    ts = np.linspace(near, t_hi, int(n_t))
    w = ts.copy()
    return ts, w / float(w.sum()), t_hi


def exact_expected_diameter(S, theta_deg, S2, ts, w, err=BEARING_ERR,
                            r_recv_max=R_RECV_MAX, per_sample_gate=False):
    """独立实现（用于自检）：直接枚举源位样本，用精确的两楔交会算 E[D]。

    与基准图完全同构（源位先验 = 射线上的面积均匀），但**不做任何插值/变换**，
    因此可以用来校验 "旋转平移 + 双线性插值" 的正确性。

    ``per_sample_gate=False``（默认，与 :mod:`p2_grid_expectation` 一致）：
    "可检测"门控按**候选点**判定 —— 只要该点能收到先验中某个源位的信号就保留全部样本；
    置 True 则逐样本判定 ``|S2 - G| <= 有效接收半径上界``。
    """
    from p2_grid_expectation import _fast_pair_grid  # 局部导入，避免循环依赖

    S = np.asarray(S, float)
    S2 = np.atleast_2d(np.asarray(S2, float))
    ts = np.asarray(ts, float)
    w = np.asarray(w, float)
    m, n = S2.shape[0], ts.size
    if n == 0:
        return np.full(m, np.nan)
    u = np.array([math.cos(math.radians(theta_deg)), math.sin(math.radians(theta_deg))])
    G = S[None, :] + ts[:, None] * u[None, :]                    # (n,2) 源位样本
    th2 = np.degrees(np.arctan2(G[None, :, 1] - S2[:, None, 1],
                                G[None, :, 0] - S2[:, None, 0])) % 360.0
    # 注意 _fast_pair_grid 的形状约定：坐标数组要比 th 少一个尾维（(m,2) 对 (m,n)）
    S1b = np.broadcast_to(S.reshape(1, 2), (m, 2))
    th1b = np.full((m,), float(theta_deg))
    S2b = S2
    D, _st = _fast_pair_grid(S1b, th1b, S2b, th2, err)
    d2 = np.hypot(S2[:, None, 0] - G[None, :, 0], S2[:, None, 1] - G[None, :, 1])
    if per_sample_gate:
        reach = d2 <= r_recv_max + 1e-9
    else:
        reach = np.broadcast_to((d2 <= r_recv_max + 1e-9).any(axis=1)[:, None], D.shape)
    ok = np.isfinite(D) & reach
    wc = np.where(ok, w[None, :], 0.0)
    sw = wc.sum(axis=1)
    num = (np.where(ok, D, 0.0) * w[None, :]).sum(axis=1)
    with np.errstate(invalid="ignore"):
        return np.where(sw > 0, num / np.maximum(sw, 1e-12), np.nan)


def best_refine_point(rec_pts, rec_svds, cur_xy, region, cfg: FieldConfig | None = None,
                      n_t=120, n_ang=16, radii=(100.0, 200.0, 300.0, 450.0, 600.0, 800.0),
                      travel_cap_m=900.0, err=BEARING_ERR):
    """在"已有一条/多条方位"的基础上，选一个**局部**最优的下一检测点。

    与 :class:`ExpectField`（全局、用预计算基准图）互补：这里对任意 S1/(*@\ensuremath{\theta}@*)1 直接算
    精确的 E[D]，用于机器狗已经逼近目标、或刚发现一个新频道需要"再测一条方位把区域
    定下来"的场合。候选点取以当前位置为中心、半径 ``radii`` 的极角网格，
    代价 = 走到候选点 + 检测 + 走到目标区域 + 预期覆盖式清扫时间。

    ``region`` 为 ``{'center': (x, y)}``；为 ``None``（只有一条方位）时用源位先验的
    中位点作为"目标区域"，从而把"测完还要往回走多远"算进代价。

    返回 ``(xy, info)``；无可用候选时返回 ``(None, None)``。
    """
    pts = [np.asarray(p, float) for p in rec_pts]
    if not pts:
        return None, None
    S1 = pts[-1]
    th1 = float(rec_svds[-1])
    ts, w, t_hi = ray_source_samples(S1, th1, n_t=n_t)
    if ts.size == 0:
        return None, None
    cur = np.asarray(cur_xy, float)
    ang0 = math.degrees(math.atan2(cur[1] - S1[1], cur[0] - S1[0]))
    cand = [cur]
    for r in radii:
        for k in range(n_ang):
            a = math.radians(ang0 + 360.0 * k / n_ang)
            p = cur + r * np.array([math.cos(a), math.sin(a)])
            if np.hypot(*p) <= R_ARENA - 1.0:
                cand.append(p)
    C = np.array(cand, float)
    keep = np.hypot(C[:, 0] - cur[0], C[:, 1] - cur[1]) <= travel_cap_m
    C = C[keep]
    if C.size == 0:
        return None, None
    E = exact_expected_diameter(S1, th1, C, ts, w, err=err)
    dcur = np.hypot(C[:, 0] - cur[0], C[:, 1] - cur[1])
    if region is not None:
        c = np.asarray(region["center"], float)
        back_target = "region"
    else:
        u = np.array([math.cos(math.radians(th1)), math.sin(math.radians(th1))])
        c = S1 + float(ts[ts.size // 2]) * u        # 先验中位源位
        back_target = "prior_median"
    dreg = np.hypot(C[:, 0] - c[0], C[:, 1] - c[1])
    # 代价（秒）= 走到候选点 + 检测 + 走到目标 + 按 E[D] 估计的覆盖式清扫
    #   清扫路程 (*@\ensuremath{\approx}@*) 圆域面积 / 覆盖行距 = (*@\ensuremath{\pi}@*)·(r)²/(2·20)，r 取 E[D]/2 + 20 m
    v = 5.0
    sweep_m = np.pi * np.maximum(E / 2.0 + 20.0, 20.0) ** 2 / (2.0 * R_CLEAR)
    cost = dcur / v + 6.0 + dreg / v + np.where(np.isfinite(E), sweep_m / v, 1e6)
    i = int(np.argmin(cost))
    info = {"S1": tuple(float(v) for v in S1), "th1": th1,
            "xy": (float(C[i, 0]), float(C[i, 1])),
            "E_diam_m": float(E[i]), "cost_s": float(cost[i]),
            "move_s": float(dcur[i] / v), "back_s": float(dreg[i] / v),
            "n_cand": int(C.shape[0]), "t_hi": float(t_hi), "back_target": back_target}
    return (float(C[i, 0]), float(C[i, 1])), info


def _ring_candidates(S, theta_deg, radii=(400.0, 800.0, 1200.0), n_ang=8,
                     dphi0=30.0, arena_r=R_ARENA):
    """绕检测点 S 的环形候选点（相对 (*@\ensuremath{\theta}@*) 偏开 dphi0，避免正落在退化射线上）。"""
    S = np.asarray(S, float)
    C = []
    for r in radii:
        for k in range(n_ang):
            a = math.radians(theta_deg + dphi0 + (360.0 - 2 * dphi0) * k / max(n_ang - 1, 1))
            p = S + r * np.array([math.cos(a), math.sin(a)])
            if np.hypot(*p) <= arena_r - 1.0:
                C.append(p)
    return np.array(C) if C else np.zeros((0, 2))


# --------------------------------------------------------------------------
# 自检
# --------------------------------------------------------------------------
def selfcheck(base_csv=DEFAULT_BASE_CSV, verbose=True, seed=20260913):
    out = []
    ok_all = True

    def rec(name, ok, detail=""):
        nonlocal ok_all
        ok_all = ok_all and bool(ok)
        out.append({"check": name, "ok": bool(ok), "detail": detail})
        if verbose:
            print(f"[{'OK ' if ok else 'FAIL'}] {name}  {detail}")

    rng = np.random.default_rng(seed)
    base = BaseMap.synthetic(step=50.0)
    E, cov = base.lookup(0.0, 0.0)
    rec("T1 基准图基本可查（原点有值）", bool(cov) and np.isfinite(float(E)),
        f"E(0,0)={float(E):.3f}, covered={bool(cov)}")

    # T2 旋转平移不变性：把基准网格点做刚体变换后查回来，应完全一致
    S = np.array([137.0, -421.0])
    th = 37.0
    c, s = math.cos(math.radians(th)), math.sin(math.radians(th))
    worst = 0.0
    for (x, y) in ((-500.0, 300.0), (0.0, 0.0), (1000.0, -800.0), (-1200.0, -1200.0)):
        e0, c0 = base.lookup(x, y)
        if not bool(c0):
            continue
        g = S + np.array([x * c - y * s, x * s + y * c])         # 正向变换到全局
        e1, c1 = base.transformed_lookup(S, th, g[0], g[1])
        worst = max(worst, abs(float(e1) - float(e0)))
    rec("T2 旋转+平移不变性（刚体变换下 E 不变）", worst < 1e-9,
        f"最大绝对差 {worst:.3e}")

    # T3 双线性插值精度：中心值应落在四邻点区间内，且与真值相差小
    bad = 0
    worst_rel = 0.0
    for _ in range(400):
        i = int(rng.integers(1, base.xs.size - 2))
        j = int(rng.integers(1, base.ys.size - 2))
        if not np.all(np.isfinite(base.E[j - 1:j + 2, i - 1:i + 2])):
            continue
        x = float(base.xs[i]) + float(rng.uniform(-0.49, 0.49)) * base.step
        y = float(base.ys[j]) + float(rng.uniform(-0.49, 0.49)) * base.step
        v, cv = base.lookup(x, y)
        if not bool(cv):
            continue
        nb = base.E[j - 1:j + 2, i - 1:i + 2]
        if not (nb.min() - 1e-9 <= float(v) <= nb.max() + 1e-9):
            bad += 1
        # 与最近邻的差应不超过 1 个格距内的局部变化
        worst_rel = max(worst_rel, abs(float(v) - base.E[j, i]) / max(abs(base.E[j, i]), 1e-9))
    rec("T3 双线性插值落在邻域值域内", bad == 0,
        f"越界 {bad} 次；与中心格最大相对差 {worst_rel:.3f}")

    # T4 未覆盖 / 超上限 -> e_max
    v_out, c_out = base.lookup(5000.0, 0.0)
    rec("T4 靶区外/无数据 -> 默认最大值 e_max",
        (not bool(c_out)) and abs(float(v_out) - E_MAX) < 1e-9,
        f"lookup(5000,0)={float(v_out)}, covered={bool(c_out)}")

    # T5 叠加 + 归一化 + 距离惩罚的结构性质
    f = ExpectField(base, FieldConfig(grid_step=100.0, dist_weight=0.35))
    d0 = f.argmin((0.0, 0.0))
    fb = f.fallback_point((0.0, 0.0), visited=[(0.0, 0.0)])
    rec("T5a 无任何方位时退化为探索点（离已访问点最远）",
        d0 is None and fb["kind"] == "explore"
        and math.hypot(fb["xy"][0], fb["xy"][1]) > 1200.0,
        f"argmin=None, 探索点 {np.round(fb['xy'], 1).tolist()}（半径 "
        f"{math.hypot(*fb['xy']):.0f} m）")
    for (x, y, t) in ((0.0, 0.0, 0.0), (0.0, 0.0, 90.0)):
        f.add_bearing((x, y), t)
    sel = f.argmin((0.0, 0.0), visited=[(0.0, 0.0)])
    sc, info = f.score((0.0, 0.0), visited=[(0.0, 0.0)])
    i = sel["index"]
    rec("T5b argmin 与全网格暴力最小一致",
        abs(float(sc[i]) - float(np.min(sc))) < 1e-9,
        f"argmin score={sel['score']:.6f}")
    rec("T5c 归一化性质：E_norm(*@\ensuremath{\in}@*)[0,1] 且未覆盖格子=1",
        bool(np.all(info["e_norm"] <= 1.0 + 1e-9)) and
        float(np.max(info["e_norm"][info["coverage"] == 0])) == 1.0,
        f"max E_norm={float(np.max(info['e_norm'])):.4f}")
    # 距离项的定量性质：score (*@\ensuremath{\equiv}@*) E_norm + (*@\ensuremath{\lambda}@*)·d/(2R_arena)，且 (*@\ensuremath{\lambda}@*)=0 时与距离无关
    c0 = FieldConfig(grid_step=100.0, dist_weight=0.35, min_move=0.0)
    f1 = ExpectField(base, c0)
    for (x, y, t) in ((0.0, 0.0, 0.0), (0.0, 0.0, 90.0)):
        f1.add_bearing((x, y), t)
    sc1, inf1 = f1.score((120.0, -340.0))
    manual = inf1["e_norm"] + 0.35 * inf1["dist_m"] / DIST_REF
    f0 = ExpectField(base, FieldConfig(grid_step=100.0, dist_weight=0.0, min_move=0.0))
    for (x, y, t) in ((0.0, 0.0, 0.0), (0.0, 0.0, 90.0)):
        f0.add_bearing((x, y), t)
    sc0, _ = f0.score((120.0, -340.0))
    rec("T5d 距离层：score = E_norm + (*@\ensuremath{\lambda}@*)·d/(2R_arena)（(*@\ensuremath{\lambda}@*)=0 时与距离无关）",
        bool(np.all(np.isfinite(sc1)))
        and float(np.max(np.abs(sc1 - manual))) < 1e-12
        and float(np.max(np.abs(sc0 - inf1["e_norm"]))) < 1e-12,
        f"最大偏差 {float(np.max(np.abs(sc1 - manual))):.3e}")

    # T6 与精确实现交叉校验（用基准图族；含"射线被靶区截断"的病态情形）
    fam = BaseMapFamily.load_dir(DEFAULT_FAMILY_DIR) if os.path.isdir(DEFAULT_FAMILY_DIR) else None
    rec("T6a 基准图族可用（按先验外径 t_hi 分族）", fam is not None,
        "族目录 " + DEFAULT_FAMILY_DIR if fam is not None
        else f"未找到 {DEFAULT_FAMILY_DIR}（可运行 --build-family 生成）；"
             f"此时只能对 t_hi=1500 的几何严格成立")

    def _cmp(rb_or_fam, cases, label):
        rels, n_local = [], 0
        worst = None
        for (Sd, thd, C) in cases:
            ts, w, _ = ray_source_samples(Sd, thd, n_t=200)
            if ts.size == 0:
                continue
            E_ex = exact_expected_diameter(Sd, thd, C, ts, w)
            if rb_or_fam is None:
                continue
            if isinstance(rb_or_fam, BaseMapFamily):
                rb, _t = rb_or_fam.rebase(Sd, thd)
            else:
                rb = rb_or_fam
            E_bm, _cov = rb.transformed_lookup(Sd, thd, C[:, 0], C[:, 1])
            u = np.array([math.cos(math.radians(thd)), math.sin(math.radians(thd))])
            G = np.asarray(Sd)[None, :] + ts[:, None] * u[None, :]
            for k in range(C.shape[0]):
                if not np.isfinite(E_ex[k]) or not np.isfinite(E_bm[k]):
                    continue
                u1 = G - np.asarray(Sd)[None, :]
                u2 = G - C[k][None, :]
                cosang = (u1 * u2).sum(1) / np.maximum(
                    np.linalg.norm(u1, axis=1) * np.linalg.norm(u2, axis=1), 1e-9)
                ang = np.degrees(np.arccos(np.clip(cosang, -1.0, 1.0)))
                gmin = min(float(np.min(ang)), float(np.min(180.0 - ang)))
                if gmin < 15.0:      # 近简并：E 是尖峰，插值无意义（决策场里会截断成 e_max）
                    continue
                rel = abs(float(E_bm[k]) - float(E_ex[k])) / max(float(E_ex[k]), 1e-9)
                rels.append(rel)
                n_local += 1
                if worst is None or rel > worst[0]:
                    worst = (rel, float(E_bm[k]), float(E_ex[k]),
                             float(np.hypot(*(C[k] - Sd))))
        if not rels:
            return None
        return {"n": n_local, "med": float(np.median(rels)),
                "p90": float(np.percentile(rels, 90)), "max": max(rels),
                "worst": worst, "label": label}

    # 生成两组几何：S 在靶区中部（先验基本不被截断）与 S 贴近靶区边界且朝外测（被截断）
    cases_mid, cases_edge = [], []
    for _ in range(6):
        Sd = np.round(rng.uniform(-700, 700, 2), 1)
        thd = float(rng.uniform(0, 360))
        cases_mid.append((Sd, thd, _ring_candidates(Sd, thd)))
    for _ in range(6):
        a = float(rng.uniform(0, 360))
        Sd = 1560.0 * np.array([math.cos(math.radians(a)), math.sin(math.radians(a))])
        thd = float((a + rng.uniform(-25, 25)) % 360.0)      # 朝靶区外测
        cases_edge.append((Sd, thd, _ring_candidates(Sd, thd)))
    cases = [(S, t, C) for (S, t, C) in cases_mid + cases_edge if len(C)]

    r_fam = _cmp(fam, cases, "family") if fam is not None else None
    r_plain = _cmp(BaseMap.load_csv(base_csv) if os.path.exists(base_csv) else None,
                   cases, "plain1500") if os.path.exists(base_csv) else None
    if r_fam is not None:
        ok6 = (r_fam["med"] < 0.01) and (r_fam["p90"] < 0.05)
        detail6 = (f"{r_fam['n']} 个良态候选点（含边界截断几何）：相对差 中位数 "
                   f"{r_fam['med']:.4f}、P90 {r_fam['p90']:.4f}、最大 {r_fam['max']:.4f}")
        if r_plain is not None:
            detail6 += (f"；不用族（只按 t_hi=1500 的图）时中位数 {r_plain['med']:.4f}、"
                        f"最大 {r_plain['max']:.4f}")
    elif r_plain is not None:
        ok6 = (r_plain["med"] < 0.01) and (r_plain["p90"] < 0.05)
        detail6 = (f"{r_plain['n']} 个良态候选点：中位数 {r_plain['med']:.4f}、"
                   f"P90 {r_plain['p90']:.4f}")
    else:
        ok6, detail6 = True, "跳过（既无基准图族也无单张基准图）"
    rec("T6b 变换+插值 (*@\ensuremath{\approx}@*) 精确期望（良态几何）", ok6, detail6)

    # T7 反例：忽略 t_hi（只用 t_hi=1500 的图）在边界朝外几何上会明显失真
    if fam is not None and r_plain is not None and r_fam is not None:
        edge_rel = [x for x in
                    [_cmp(fam, [(S, t, C)], "f") for (S, t, C) in cases_edge
                     if len(C)] if x]
        edge_plain = [x for x in
                      [_cmp(BaseMap.load_csv(base_csv), [(S, t, C)], "p")
                       for (S, t, C) in cases_edge if len(C)] if x]
        if edge_rel and edge_plain:
            med_f = float(np.median([e["med"] for e in edge_rel]))
            med_p = float(np.median([e["med"] for e in edge_plain]))
            rec("T7 边界朝外几何：按 t_hi 分族显著优于固定 t_hi=1500",
                med_f < 0.05 and med_f < med_p / 5.0,
                f"族：中位相对差 {med_f:.4f}；固定 1500：{med_p:.4f}（"
                f"差 {med_p / max(med_f, 1e-6):.0f} 倍）")

    if verbose:
        print(f"\n自检总结：{'全部通过' if ok_all else '存在失败项'}")
    return {"ok": ok_all, "checks": out}


# --------------------------------------------------------------------------
# 演示 / 命令行
# --------------------------------------------------------------------------
def field_figure(f, cur_xy, sel, path, title=None):
    """把决策场画出来：叠加后的 E 场、距离惩罚、合成 score、选点与探测器位置。

    图注用英文（本机 matplotlib 无中文字体，避免缺字告警）。
    """
    import matplotlib
    matplotlib.use("Agg")
    import matplotlib.pyplot as plt

    E_mean, cov, n_terms = f.aggregate()
    score, info = f.score(cur_xy)
    fig, axes = plt.subplots(1, 3, figsize=(16.5, 5.6))
    panels = [(E_mean, "E[intersection diameter] (m)\n(uncovered cells = 1500)",
               "viridis_r"),
              (np.where(np.isfinite(info["dist_m"]),
                        info["dist_m"] / DIST_REF, np.nan),
               f"distance term  (*@\ensuremath{\lambda}@*)·d/(2R),  (*@\ensuremath{\lambda}@*)={f.cfg.dist_weight}", "magma"),
              (np.where(np.isfinite(score), score, np.nan),
               "score = normalised E + distance term\n((*@\ensuremath{\blacktriangle}@*) = argmin)", "cividis_r")]
    for ax, (V, ttl, cmap) in zip(axes, panels):
        sc = ax.scatter(f.gx, f.gy, c=V, s=26, cmap=cmap, marker="s")
        fig.colorbar(sc, ax=ax, shrink=0.8)
        ax.set_aspect("equal")
        ax.set_title(ttl, fontsize=10)
        ax.set_xlim(-1950, 1950)
        ax.set_ylim(-1950, 1950)
        ax.set_xlabel("x (m)")
        ax.set_ylabel("y (m)")
        th = np.linspace(0, 2 * math.pi, 400)
        ax.plot(f.cfg.arena_r * np.cos(th), f.cfg.arena_r * np.sin(th), "k--", lw=0.8)
        ax.plot([cur_xy[0]], [cur_xy[1]], "r*", ms=14)
        for (S, thd, _tag, _thi) in [t[2] for t in f.terms]:
            a = math.radians(thd)
            ax.plot([S[0], S[0] + 300 * math.cos(a)],
                    [S[1], S[1] + 300 * math.sin(a)], "w-", lw=1.2, alpha=0.8)
        if sel is not None:
            ax.plot([sel["xy"][0]], [sel["xy"][1]], "w^", ms=11, mec="k")
        ax.grid(alpha=0.2)
    fig.suptitle(title or (f"Q3 decision field: {n_terms} bearing/discovery terms "
                           f"superimposed on the global grid"), fontsize=11)
    fig.tight_layout()
    os.makedirs(os.path.dirname(os.path.abspath(path)), exist_ok=True)
    fig.savefig(path, dpi=150)
    plt.close(fig)
    return path


def demo(base_csv=DEFAULT_BASE_CSV, outdir=None, dist_weight=1.5, grid_step=100.0,
         bearings=((0.0, 0.0, 0.0),), verbose=True, family_dir=DEFAULT_FAMILY_DIR):
    fam = BaseMapFamily.load_dir(family_dir) if os.path.isdir(family_dir) else None
    base = fam if fam is not None else BaseMap.load_csv(base_csv)
    cfg = FieldConfig(grid_step=grid_step, dist_weight=dist_weight)
    f = ExpectField(base, cfg)
    for (x, y, t) in bearings:
        f.add_bearing((x, y), t)
    sel = f.argmin((0.0, 0.0), visited=[(0.0, 0.0)])
    if verbose:
        kind = "基准图族" if fam is not None else f"单张基准图 {base.source}"
        print("基准：", kind)
        print(f"叠加 {len(f.terms)} 条方位；选点 {np.round(sel['xy'], 1).tolist()}，"
              f"E_mean={sel['E_mean_m']:.1f} m，E_norm={sel['E_norm']:.3f}，"
              f"距离 {sel['dist_m']:.0f} m，score={sel['score']:.4f}")
    if outdir:
        os.makedirs(outdir, exist_ok=True)
        p = f.save_csv(os.path.join(outdir, "p3_field.csv"))
        with open(os.path.join(outdir, "p3_field_report.json"), "w", encoding="utf-8") as fh:
            json.dump({"cfg": cfg.__dict__, "terms": len(f.terms), "pick": sel},
                      fh, ensure_ascii=False, indent=2)
        try:
            fig = field_figure(f, (0.0, 0.0), sel,
                               os.path.join(outdir, "p3_field.png"))
        except Exception as e:      # 画图失败不影响主流程
            fig = f"figure_error: {e!r}"
        if verbose:
            print("写出", p, "与", fig)
        return f, sel
    return f, sel


def main():
    ap = argparse.ArgumentParser(description="问题3 期望决策场")
    ap.add_argument("--selfcheck", action="store_true")
    ap.add_argument("--demo", action="store_true")
    ap.add_argument("--build-family", action="store_true",
                    help="按第二问口径重算一族基准图（不同先验外径 t_hi）")
    ap.add_argument("--base-map", default=DEFAULT_BASE_CSV)
    ap.add_argument("--family-dir", default=DEFAULT_FAMILY_DIR)
    ap.add_argument("--family-step", type=float, default=40.0)
    ap.add_argument("--family-nt", type=int, default=200)
    ap.add_argument("--out", default=None)
    ap.add_argument("--dist-weight", type=float, default=1.5,
                    help="距离项系数 (*@\ensuremath{\lambda}@*)（问题 3 策略默认 1.5）")
    ap.add_argument("--grid-step", type=float, default=100.0)
    args = ap.parse_args()
    if args.selfcheck:
        res = selfcheck(args.base_map)
        raise SystemExit(0 if res["ok"] else 1)
    if args.build_family:
        print(f"计算基准图族 -> {args.family_dir}")
        BaseMapFamily.build(args.family_dir, step=args.family_step, n_t=args.family_nt)
    if args.demo or not args.build_family:
        demo(args.base_map, args.out, args.dist_weight, args.grid_step,
             family_dir=args.family_dir)


if __name__ == "__main__":
    main()

\end{lstlisting}
\subsection{\texttt{src/p3\_robot.py}}
% SHA256 90852f65461abe477f872fe6d12f8b02ebbea60010072e4db25575156e688496
\begin{lstlisting}[language=python,escapeinside={(*@}{@*)}]
"""B 题问题 3：机器狗"基于期望的行动"在线搜索—定位—清除策略。

策略流程（与设计说明一一对应）
------------------------------
1. **第一轮扫描**：机器狗从原点出发，测向机初始频道 1，按频道 1..20 依次 ``/measure``，
   得到若干条"检测点 + 示向度"（部分频道无信号属正常）。
2. **期望场选点**：对每个"只测到一次方位"的频道，用第二问的网格模拟结果
   （``out/p3_family`` 或 ``out/p2_grid``）做旋转 + 平移 + 双线性插值，把
   :math:`\\mathbb E[D]` 叠加到全局网格上；没有数据/超阈值的格子取最大值 1500 m；
   归一化后再叠加一层**与探测器距离线性增长**的代价（系数 (*@\ensuremath{\lambda}@*) 可调），取 ``argmin``
   作为下一个检测点，随后开始新一轮 20 频道扫描。
3. **清理**：某些频道已经测到 (*@\ensuremath{\ge}@*)2 条方位 —— 用问题 1 的交会算法算出定位区域
   （最小包围圆 = 中心 + 半径 r*）。对所有可清理的频道按最近邻顺序规划航路，
   依次前往；到达区域后继续检测把区域压小，直到进入清除范围，再用
   **覆盖式试探**（六边形格点，间距 (*@\ensuremath{\le}@*) 20√3 m，保证圆域内任一点到某个试探点 (*@\ensuremath{\le}@*) 20 m）
   逐个 ``/clear``，必然有一点成功。
4. **清除后继续扫描**；已经清除的频道、已经定位够的频道、以及"在同一个有效接收半径内
   重复测过的频道"都会被跳过，从而节省时间。
5. 反复执行 2–4，直到时间预算耗尽或没有新的可用信息。

关键设计
--------
* **为什么"未覆盖 = 1500 m"是对的**：E[D] 越小定位越好；把没有信息的格子记成最差，
  只有在附近确实没有更好格子时才会选它 —— 这天然形成了"顺路探索"的行为。
* **为什么不能沿示向度方向逼近**：那会退化成问题 1 的"无界区域"几何。所以清理时的
  侧移精定位点取在**与上一条视线近似正交**的方向上。
* **清理阶段的覆盖式试探是"保证成功"的**：真源必落在 ±1° 锥的交集（定位区域）内，
  区域内任意一点到某个试探点的距离 (*@\ensuremath{\le}@*) 20 m = 清除半径。
"""
from __future__ import annotations

import argparse
import json
import math
import os
import sys
import time
from dataclasses import dataclass, field as dc_field

import numpy as np

_HERE = os.path.dirname(os.path.abspath(__file__))
if _HERE not in sys.path:
    sys.path.insert(0, _HERE)

from p1_intersection import analyse  # noqa: E402
from p3_arena import (  # noqa: E402
    ArenaError, BEARING_ERR, CHANNELS, MockArena, NEAR_R, R_ARENA, R_CLEAR,
    T_CLEAR_FAIL, T_CLEAR_OK, T_MEASURE, T_SWITCH, V_ROBOT,
)
from p3_expect_field import (  # noqa: E402
    DEFAULT_BASE_CSV, DEFAULT_FAMILY_DIR, BaseMap, BaseMapFamily, ExpectField,
    FieldConfig, best_refine_point,
)


# --------------------------------------------------------------------------
# 覆盖式试探的三种结局
# --------------------------------------------------------------------------
# 必须区分"预算不够、一个试探点都没来得及发"与"试探了但没打中"：
# 前者不是策略失败（不该累加 clear_fail、不该写进论文的失败归因），
# 后者才是真正的"区域内没找到目标"。
SWEEP_OK = "success"        # 某一点 /clear 成功
SWEEP_MISS = "miss"         # 试探了全部候选点（或中途超时）都没打中
SWEEP_BUDGET = "no_budget"  # 剩余预算不足，**一个 /clear 都没发出**，主动放弃


# --------------------------------------------------------------------------
# 配置
# --------------------------------------------------------------------------
@dataclass
class P3Config:
    # ---- 期望场（全局网格 + 距离惩罚）----
    grid_step: float = 100.0        # 全局网格步长（m）
    e_max: float = 1500.0           # 未覆盖网格的默认最大值（m）
    dist_weight: float = 1.5        # 距离项系数 (*@\ensuremath{\lambda}@*)（离线演练调参得到，见 review/P3-design.md）
    visited_radius: float = 250.0   # 已扫描点邻域排除半径（m）
    min_move: float = 60.0          # 距当前位置过近的格子排除（m）
    field_terms: str = "pending"    # pending：只叠加"还缺方位"的频道
    field_discovery: bool = False   # 始终叠加"探索项"（把扫描点推向能发现新源的地方）
    auto_discovery_when_empty: bool = True  # 场上一条可用方位都没有时自动启用探索项

    # ---- 扫描 ----
    channels: tuple = CHANNELS
    scan_1bearing_first: bool = True
    repeat_gap_m: float = 60.0          # 同一频道在附近重复测量的最小间隔（m）
    nosignal_gap_m: float = 1000.0      # "无信号"频道要换足够远的地方重测
                                        # （= 有效接收半径下界：在 1000 m 外重测才有
                                        #   真正的新信息；实测 450→1000 使检测次数
                                        #   289→186、平均定位清除时间 398→315 s，
                                        #   清除率 0.992→0.995，见 review/P3-design.md §5.1）（m）
    enroute_scan: bool = True            # 清理途中顺路补测
    enroute_scan_max: int = 2
    second_bearing_refine: bool = False  # 新发现频道立刻补第二条方位（实测收益不稳，默认关）
    second_bearing_max: int = 2          # 每个停点最多补几条
    second_bearing_budget_s: float = 320.0  # 剩余时间低于它就不再补测
    max_rounds: int = 30
    tour: str = "nn2opt"                 # 清理航路：nn | nn2opt
    dynamic_replan: bool = True          # 清理阶段每步重规划（滚动时域）；False=阶段内固定顺序
    clear_policy: str = "tour"           # tour（NN+2-opt 全局结构）| rollout（一步前瞻+SPT）
    clear_budget_truncation: bool = False  # 是否按剩余预算裁剪航路（实测关掉更好）
    optimistic_probe_steps: int = 3      # 打开裁剪时的试探点数封顶
    min_scan_rounds_before_clear: int = 1  # 先扫描几轮再开始清理

    # ---- 定位与清除 ----
    clear_min_bearings: int = 2
    clear_max_radius_m: float = 600.0   # 定位区域半径超过它就先不清理
                                        # （400 → 600：扫圈在 r(*@\ensuremath{\approx}@*)1100 m 处交会、区域半径
                                        #   常在 400–650 m，400 会让这些源被"跳过"；
                                        #   600 实测把扫圈的清除率从 0.958 提到 0.989–1.000）
    arrive_radius_m: float = 30.0
    direct_probe_radius_m: float = 60.0  # r* 不超过它就无需侧移精定位，直接试探
    probe_spacing_m: float = 34.0        # (*@\ensuremath{\le}@*) 20√3 (*@\ensuremath{\approx}@*) 34.6 m 才能保证覆盖
    probe_max_points: int = 60           # 候选点数上限（仅"可覆盖"档位，见下）
    probe_unlimited: bool = True         # True=不按点数截断（保证覆盖；截断会破坏 20 m 保证）
    probe_max_radius_m: float = 200.0
    refine_leg_m: float = 320.0
    refine_travel_cap_m: float = 700.0
    refine_mode: str = "perp"            # perp（正交侧移）| exact（精确期望最优）
    max_approach_steps: int = 3

    # ---- 时间预算 ----
    time_reserve_s: float = 3.0
    min_action_budget_s: float = 10.0
    min_scan_channels: int = 2       # 单轮扫描至少要能测这么多频道才值得跑一趟

    # ---- 预算感知的扫描闸门（"剩余时间不够就不要再扫"）----
    # 诊断依据：30 局里扫描阶段吃掉约 1016 s / 1200 s，而末尾有 23% 的未清除源
    # **已经定位好却没时间走过去**。所以关键不是"怎么选扫描点"，而是"还要不要扫"。
    scan_budget_gate: bool = False       # 开：本轮扫描后必须还够清最近的一个目标，否则不扫
    scan_gate_margin: float = 0.0        # 额外预留秒数（越大越倾向"别扫了，去清除"）
    scan_gate_factor: float = 1.0        # 对"清最近目标"耗时估计的放大系数
    scan_only_toward_targets: bool = False  # 扫描点若背着已知目标走，则该轮不扫
    # 扫描时间的硬性上限（占总预算比例）。诊断显示扫描占掉约 1016 s / 1200 s，
    # 而 29/49 个"已定位未清除"的目标在变成可清目标时**剩余时间完全够用**却没去清。
    scan_budget_frac: float = 1.0        # 1.0 = 不限制；0.5 = 扫描累计不超过一半预算

    # ---- 清理航路：可达近点子集（"优先清路程近的点"）----
    clear_shortlist: bool = False        # 开：只保留"预算内可达"的近点子集，再排序
    clear_shortlist_reserve_s: float = 0.0  # 为后续扫描预留的秒数


# --------------------------------------------------------------------------
# 频道台账
# --------------------------------------------------------------------------
@dataclass
class ChannelRec:
    channel: int
    pts: list = dc_field(default_factory=list)     # 测到方位的检测点
    svds: list = dc_field(default_factory=list)    # 对应的示向度
    meas_pts: list = dc_field(default_factory=list)  # 所有检测过的位置（含无信号）
    meas_log: list = dc_field(default_factory=list)  # (x, y, virtual_time, result)
    meas_stations: list = dc_field(default_factory=list)  # 测过该频道的**站位编号**
                                                     # （问题 4 的覆盖记账用：只有"这一站
                                                     #   本来能听到它"才抵扣覆盖）
    cleared: bool = False
    cleared_at: float | None = None
    near_hits: int = 0
    no_signal: int = 0
    clear_fail: int = 0                # 自最近一次"新增方位"以来清理失败的次数
    last_bounded: tuple | None = None  # (center, radius, n_bearings) 最近一次有界区域
    tried_probes: list = dc_field(default_factory=list)
    history: list = dc_field(default_factory=list)

    @property
    def n_bearings(self) -> int:
        return len(self.pts)


# --------------------------------------------------------------------------
# 机器狗
# --------------------------------------------------------------------------
def load_base(family_dir=DEFAULT_FAMILY_DIR, base_csv=DEFAULT_BASE_CSV):
    """优先用"基准图族"（按先验外径 t_hi 分族，变换才严格成立），否则退回单张基准图。"""
    if family_dir and os.path.isdir(family_dir):
        fam = BaseMapFamily.load_dir(family_dir)
        if fam is not None:
            return fam
    if os.path.exists(base_csv):
        return BaseMap.load_csv(base_csv)
    return None


def _dist_point_segment(p, a, b):
    ab = b - a
    L2 = float(ab @ ab)
    if L2 < 1e-12:
        return float(np.linalg.norm(p - a))
    t = float(np.clip((p - a) @ ab / L2, 0.0, 1.0))
    return float(np.linalg.norm(p - (a + t * ab)))


def _point_in_polygon(p, P):
    """射线法（P 为凸多边形顶点，逆时针或顺时针均可）。"""
    n = len(P)
    inside = False
    for i in range(n):
        a, b = P[i], P[(i + 1) % n]
        if (a[1] > p[1]) != (b[1] > p[1]):
            x = a[0] + (p[1] - a[1]) * (b[0] - a[0]) / (b[1] - a[1])
            if x > p[0]:
                inside = not inside
    return inside


def _dist_point_polygon(p, P):
    p = np.asarray(p, float)
    P = np.asarray(P, float)
    if _point_in_polygon(p, P):
        return 0.0
    n = len(P)
    return min(_dist_point_segment(p, P[i], P[(i + 1) % n]) for i in range(n))


class P3Robot:
    def __init__(self, arena, cfg: P3Config | None = None, base=None, log=None,
                 log_prefix=""):
        self.arena = arena
        self.cfg = cfg or P3Config()
        self.base = base if base is not None else load_base()
        self._log_fn = log
        self.log_prefix = log_prefix
        self.recs = {int(ch): ChannelRec(int(ch)) for ch in self.cfg.channels}
        self.scan_points: list[tuple] = []
        self.meas_pts: list[tuple] = []      # **真实做过的检测**位置（覆盖度评估口径：
                                             # "到过的点"(*@\ensuremath{\ne}@*)"测过的点"，覆盖度只能按后者算）
        self.events: list[dict] = []
        self.aborted = False
        # 覆盖式试探因"预算不足"放弃时置位：让 clear_phase 立刻收敛，
        # 不要对同一个目标反复重规划（实测会把重规划次数抬高一到两个数量级）。
        self.abort_clear_phase = False
        self.stats = {"n_measure": 0, "n_clear": 0, "n_clear_fail": 0,
                      "n_scan_rounds": 0, "n_bearings": 0, "n_rejected": 0,
                      "n_refine": 0, "n_second_bearing": 0, "n_replans": 0,
                      "n_midphase_targets": 0, "n_sweep_aborted": 0,
                      "n_sweep_miss": 0, "n_probe_issued": 0,
                      "n_scan_gated": 0, "n_shortlist_kept": 0, "n_shortlist_in": 0,
                      "t_scan_s": 0.0, "t_clear_s": 0.0, "t_clear_ends_s": 0.0,
                      "vtime_at_last_clear": 0.0}

    # ---------------- 基础设施 ----------------
    def log(self, msg):
        self.events.append({"t": float(self.arena.virtual_time_s), "msg": msg})
        if self._log_fn:
            # 日志里避免使用 GBK 控制台无法编码的符号（(*@\ensuremath{\Rightarrow}@*) (*@\ensuremath{\star}@*) √ × 等），
            # 万一还有，退化为替换字符也不要让整局崩掉。
            try:
                self._log_fn(f"{self.log_prefix}{msg}")
            except UnicodeEncodeError:
                self._log_fn(f"{self.log_prefix}{msg}".encode(
                    "gbk", "replace").decode("gbk"))

    def remaining(self):
        return float(self.arena.remaining_budget_s())

    def can_afford(self, seconds):
        return self.remaining() > float(seconds) + self.cfg.time_reserve_s

    @property
    def pos(self):
        return tuple(self.arena.pos)

    # ---------------- 检测 / 清除 ----------------
    def measure(self, x, y, ch):
        resp = self.arena.measure(x, y, ch)
        self.stats["n_measure"] += 1
        if resp.get("accepted") is not True:
            self.stats["n_rejected"] += 1
            self.aborted = True
            self.log(f"  ! /measure 频道{ch} 未执行（{resp.get('reason', '?')}）")
            return resp
        res = resp.get("measure_result")
        rec = self.recs[int(ch)]
        rec.meas_pts.append((float(x), float(y)))
        rec.meas_log.append((float(x), float(y), float(self.arena.virtual_time_s), res))
        if res == "direction":
            self.add_bearing(int(ch), x, y, float(resp["svd_deg"]))
        elif res == "near":
            rec.near_hits += 1
        else:
            rec.no_signal += 1
        return resp

    def add_bearing(self, ch, x, y, svd):
        rec = self.recs[int(ch)]
        rec.pts.append((float(x), float(y)))
        rec.svds.append(float(svd))
        rec.clear_fail = 0
        rec.tried_probes = []
        self.stats["n_bearings"] += 1
        reg = self.region(rec)
        if reg.get("status") == "bounded":
            rec.last_bounded = (tuple(reg["min_enclosing_center"]),
                                float(reg["min_enclosing_radius_m"]), len(rec.pts))
        rec.history.append({"n": len(rec.pts), "x": float(x), "y": float(y),
                            "svd": float(svd), "status": reg.get("status"),
                            "r_star_m": (rec.last_bounded[1] if rec.last_bounded else None)})
        return reg

    def do_clear(self, x, y, ch):
        resp = self.arena.clear(x, y, ch)
        self.stats["n_probe_issued"] = self.stats.get("n_probe_issued", 0) + 1
        if resp.get("accepted") is not True:
            self.stats["n_rejected"] += 1
            self.aborted = True
            self.log(f"  ! /clear 频道{ch} 未执行（{resp.get('reason', '?')}）")
            return False
        if resp.get("clear_result") == "success":
            rec = self.recs[int(ch)]
            rec.cleared = True
            rec.cleared_at = float(self.arena.virtual_time_s)
            self.stats["n_clear"] += 1
            self.stats["vtime_at_last_clear"] = float(self.arena.virtual_time_s)
            self.log(f"  [OK] 频道{ch} 已清除（t={self.arena.virtual_time_s:.1f}s）")
            return True
        self.stats["n_clear_fail"] += 1
        return False

    # ---------------- 定位区域 ----------------
    def region(self, rec: ChannelRec):
        if not rec.pts:
            return {"status": "none", "n_bearings": 0}
        return analyse(rec.pts, rec.svds, BEARING_ERR)

    def bounded_region(self, rec: ChannelRec):
        """当前（或最近一次）有界区域；返回 ``{'center','radius','poly','status'}``。"""
        if len(rec.pts) >= 2:
            reg = self.region(rec)
            if reg.get("status") == "bounded":
                return {"center": tuple(reg["min_enclosing_center"]),
                        "radius": float(reg["min_enclosing_radius_m"]),
                        "poly": [tuple(v) for v in reg.get("vertices", [])],
                        "status": "bounded", "n_bearings": len(rec.pts)}
        if rec.last_bounded is not None:
            c, r, n = rec.last_bounded
            return {"center": tuple(c), "radius": float(r), "poly": None,
                    "status": "fallback", "n_bearings": n}
        return None

    # ---------------- 扫描 ----------------
    def needed_channels(self):
        """还需要（且值得）拿到新方位的频道：已清除、已够定位两次的跳过。"""
        need, done = [], []
        for ch in self.cfg.channels:
            rec = self.recs[int(ch)]
            if rec.cleared or rec.n_bearings >= self.cfg.clear_min_bearings:
                continue
            (done if rec.pts else need).append(int(ch))
        order = (done + need) if self.cfg.scan_1bearing_first else sorted(done + need)
        return order

    def channels_to_measure(self, limit=None, pos=None, subset=None, apply_gap=True):
        """在 ``pos`` 处值得检测的频道。

        **跳过规则**（"节省时间"的关键）：已清除 / 已定位两次的不测；
        对"无信号"频道，新检测点必须距它所有历史检测点 (*@\ensuremath{\ge}@*) ``nosignal_gap_m``
        （有效接收半径至少 1000 m，在同一个圈里重复扫几乎拿不到新信息）；
        对已有一条方位的频道，新检测点距历史点 (*@\ensuremath{\ge}@*) ``repeat_gap_m``（太小交会角没意义）。
        """
        p = np.asarray(pos if pos is not None else self.arena.pos, float)
        out = []
        for ch in (subset if subset is not None else self.needed_channels()):
            rec = self.recs[int(ch)]
            if apply_gap and rec.meas_pts:
                gap = self.cfg.repeat_gap_m if rec.pts else self.cfg.nosignal_gap_m
                d = min(math.hypot(p[0] - q[0], p[1] - q[1]) for q in rec.meas_pts)
                if d < gap:
                    continue
            out.append(int(ch))
        return out[:limit] if limit else out

    def build_field(self, grid_step=None, dist_weight=None):
        cfg = FieldConfig(grid_step=grid_step or self.cfg.grid_step,
                          e_max=self.cfg.e_max,
                          dist_weight=(self.cfg.dist_weight if dist_weight is None
                                       else dist_weight),
                          visited_radius=self.cfg.visited_radius,
                          min_move=self.cfg.min_move)
        f = ExpectField(self.base, cfg)
        if self.base is None:
            return f, 0
        n = 0
        for ch, rec in self.recs.items():
            if rec.cleared:
                continue
            if self.cfg.field_terms == "pending" and rec.n_bearings >= self.cfg.clear_min_bearings:
                continue
            if rec.n_bearings == 0:
                continue
            for ((x, y), th) in zip(rec.pts, rec.svds):
                f.add_bearing((x, y), th, tag=ch)
                n += 1
        if n == 0:
            # 一条"缺第二方位"的频道都没有（例如后半程只剩没听到过的频道）：
            # 自动退化为**探索项**，否则场上没有任何信息、选点无从谈起。
            for ch, rec in self.recs.items():
                if rec.cleared or rec.n_bearings > 0 or not rec.meas_pts:
                    continue
                if not self.cfg.field_discovery and not self.cfg.auto_discovery_when_empty:
                    continue
                f.add_raw(f.discovery_term(rec.meas_pts, self.cfg.nosignal_gap_m),
                          tag=("disc", ch))
                n += 1
        return f, n

    def choose_scan_point(self, dist_weight=None, max_dist=None):
        f, n = self.build_field(dist_weight=dist_weight)
        visited = list(self.scan_points)
        sel = f.argmin(self.pos, visited=visited, max_dist=max_dist)
        if sel is None:
            sel = f.fallback_point(self.pos, visited=visited, max_dist=max_dist)
        sel["n_bearing_terms"] = n
        return sel, f

    def plan_scan(self, dist_weight=None):
        """选下一个检测点，并确定该点上值得测哪些频道。

        顺序很重要：**先**用期望场选点，**再**判断该点上哪些频道还没被"测透"，
        否则会出现"因为要换地方才有信息、但没地方可去"的死锁。
        """
        base = self.needed_channels()
        if not base:
            return None
        # 时间不够时把可选点限制在"还能走到"的范围内（给检测与安全余量留出空间）
        slack = 2.0 * (T_MEASURE + T_SWITCH) + 4.0
        dmax = max(0.0, (self.remaining() - self.cfg.time_reserve_s - slack) * V_ROBOT)
        sel, f = self.choose_scan_point(dist_weight=dist_weight, max_dist=dmax)
        chans = self.channels_to_measure(pos=sel["xy"], subset=base)
        if not chans:
            sel2 = f.fallback_point(self.pos, visited=self.scan_points)
            chans2 = self.channels_to_measure(pos=sel2["xy"], subset=base)
            if chans2:
                sel, chans = sel2, chans2
        if not chans:                     # 兜底：宁可多测几次也不原地空转
            chans = base[:3]
        # 时间不够时自动缩减本轮的频道数（至少测 min_scan_channels 个），把边角时间也用上
        travel = math.dist(self.pos, sel["xy"]) / V_ROBOT
        avail = self.remaining() - self.cfg.time_reserve_s - travel - 4.0
        kmax = max(1, int(avail // (T_MEASURE + T_SWITCH)))
        if len(chans) > kmax:
            if kmax < self.cfg.min_scan_channels:
                return None        # 只剩时间测一两个频道：不值得为它专门跑一趟
            self.log(f"  时间只剩 {self.remaining():.0f} s -> 本轮只测 {kmax}/{len(chans)} 个频道")
            chans = chans[:kmax]
        cost = travel + len(chans) * (T_MEASURE + T_SWITCH) + 2.0
        plan = {"chans": chans, "sel": sel, "field": f, "cost_s": cost}
        if self.cfg.scan_budget_gate and not self._scan_worth_it(plan):
            return None
        return plan

    # ---------------- 预算感知的扫描闸门 ----------------
    def nearest_clear_cost_after(self, at_xy):
        """估计"到达 ``at_xy`` 之后，再清掉最近一个已知目标"所需秒数；没有已知目标则 None。"""
        cand = self.clear_targets()
        if not cand:
            return None
        nxt = min(cand, key=lambda t: math.dist(at_xy, t[2]["center"]))
        return self.estimate_clear_cost(nxt[2], from_pos=at_xy,
                                        probe_cap=self.cfg.optimistic_probe_steps)

    def _scan_budget_left(self):
        """允许继续扫描的剩余秒数（= frac × 总预算 − 已花在扫描上的时间）。"""
        if self.cfg.scan_budget_frac >= 1.0:
            return float("inf")
        return self.cfg.scan_budget_frac * float(self.arena.budget) \
            - self.stats.get("t_scan_s", 0.0)

    def _scan_worth_it(self, plan):
        """本轮扫描值不值得做：**扫完还必须够清最近的一个已知目标**，否则不扫。

        动机（诊断数据）：30 局里扫描阶段平均吃掉约 1016 s / 1200 s，末尾有 23% 的
        未清除源**已经定位好却走不到**。也就是说"再多发现一个源"的边际价值，在
        剩余时间已经不足以走过去清掉它时等于 0 —— 这种扫描纯属消耗预算。

        注意这里只做"取舍"不做"折中"：要么整轮扫，要么整轮不扫；
        因为扫描点本身的移动无法中途变成对目标的逼近（两者方向不同）。
        """
        rem = self.remaining()
        left = self._scan_budget_left()
        if plan["cost_s"] > left:
            self.stats["n_scan_gated"] = self.stats.get("n_scan_gated", 0) + 1
            self.log(f"  (*@\ensuremath{\blacksquare}@*) 放弃本轮扫描：扫描时间额度已用尽（本轮需 {plan['cost_s']:.0f} s，"
                     f"额度剩 {left:.0f} s = {self.cfg.scan_budget_frac:.2f}×总预算）"
                     f"-> 预算留给已定位目标")
            return False
        cand = self.clear_targets()
        need = None
        if cand:
            nxt = min(cand, key=lambda t: math.dist(plan["sel"]["xy"], t[2]["center"]))
            need = self.estimate_clear_cost(nxt[2], from_pos=plan["sel"]["xy"],
                                            probe_cap=self.cfg.optimistic_probe_steps)
            # 可选：扫描点不能"背着"已知目标走（否则这轮的移动无法回收为进度）
            if self.cfg.scan_only_toward_targets:
                u = np.asarray(plan["sel"]["xy"], float) - np.asarray(self.pos, float)
                w = np.asarray(nxt[2]["center"], float) - np.asarray(self.pos, float)
                nu, nw = float(np.hypot(*u)), float(np.hypot(*w))
                if nu > 1e-9 and nw > 1e-9 and float(u @ w) / (nu * nw) < 0.0:
                    self.stats["n_scan_gated"] = self.stats.get("n_scan_gated", 0) + 1
                    self.log(f"  (*@\ensuremath{\blacksquare}@*) 放弃本轮扫描：扫描点 ({plan['sel']['xy'][0]:.0f},"
                             f"{plan['sel']['xy'][1]:.0f}) 与已知目标方向相反 -> 预算留给清除")
                    return False
        if need is None:
            return True                      # 手上没有可清目标：扫描是唯一的进展方式
        required = (plan["cost_s"] + self.cfg.scan_gate_margin
                    + self.cfg.scan_gate_factor * need)
        if rem >= required:
            return True
        self.stats["n_scan_gated"] = self.stats.get("n_scan_gated", 0) + 1
        self.log(f"  (*@\ensuremath{\blacksquare}@*) 放弃本轮扫描：剩余 {rem:.0f} s < 本轮扫描 {plan['cost_s']:.0f} s"
                 f" + 清最近目标 {need:.0f} s（阈值 {required:.0f} s）"
                 f"-> 预算留给已定位目标")
        return False

    def execute_scan(self, plan):
        sel, chans = plan["sel"], plan["chans"]
        x, y = sel["xy"]
        if sel["kind"] == "explore":
            self.log(f"扫描轮：期望场无可用方位（{sel.get('n_bearing_terms', 0)} 项）"
                     f"-> 探索点 ({x:.0f},{y:.0f})，距 {sel['dist_m']:.0f} m")
        else:
            self.log(f"扫描轮：选点 ({x:.0f},{y:.0f})｜E={sel['E_mean_m']:.1f} m  "
                     f"E_norm={sel['E_norm']:.3f}  距离惩罚={sel['dist_penalty']:.3f}  "
                     f"score={sel['score']:.4f}  移动 {sel['dist_m']:.0f} m｜"
                     f"待测 {len(chans)} 个频道 {chans}")
        self.scan_points.append((float(x), float(y)))
        n = self.scan_at(x, y, chans)
        self.stats["n_scan_rounds"] += 1
        return n

    def scan_round(self, dist_weight=None):
        plan = self.plan_scan(dist_weight=dist_weight)
        if plan is None:
            return 0
        if not self.can_afford(plan["cost_s"]):
            return 0
        return self.execute_scan(plan)

    def scan_at(self, x, y, chans):
        """在 (x, y) 依次检测这些频道（第一次调用会带来移动耗时）。"""
        n = 0
        for i, ch in enumerate(chans):
            if self.aborted or not self.can_afford(T_MEASURE + T_SWITCH + 3.0):
                break
            resp = self.measure(x, y, ch)
            n += 1
            if resp.get("measure_result") == "near":
                # 已在 5 m 内：直接清除
                self.do_clear(x, y, ch)
        return n

    # ---------------- 清理 ----------------
    def clear_targets(self):
        out = []
        for ch, rec in self.recs.items():
            if rec.cleared or rec.n_bearings < self.cfg.clear_min_bearings:
                continue
            if rec.clear_fail > 0:
                continue                       # 上次清理失败，等下一次新方位再试
            br = self.bounded_region(rec)
            if br is None or br["radius"] > self.cfg.clear_max_radius_m:
                continue
            out.append((ch, rec, br))
        out.sort(key=lambda t: math.dist(self.pos, t[2]["center"]))
        return out

    # ---------------- 清理：滚动时域动态规划 ----------------
    def _probe_count(self, br):
        """该定位区域需要多少个覆盖式试探点（带缓存：区域不变就不重算）。

        这里**故意**对点数封顶：它只用于"值不值得跑一趟"的代价预估，
        真实覆盖由 :meth:`probe_points` 保证。不封顶会让大 r* 上的枚举变得很贵
        （r*=200 m 时 151 个点 × 每步重规划），实测把程序运行时间抬高两个数量级。
        """
        key = (round(float(br["center"][0]), 1), round(float(br["center"][1]), 1),
               round(float(br["radius"]), 1))
        cache = getattr(self, "_probe_cache", None)
        if cache is None:
            cache = self._probe_cache = {}
        if key not in cache:
            r = min(float(br["radius"]), self.cfg.probe_max_radius_m)
            n = len(self.probe_points(np.asarray(br["center"], float), r,
                                      poly=br.get("poly")))
            cache[key] = min(int(n), int(self.cfg.probe_max_points))
        return cache[key]

    def estimate_clear_cost(self, br, from_pos=None, n_probe=None, probe_cap=None):
        """估计"从 ``from_pos`` 出发清掉某个目标"的耗时（秒）。

        组成：走到区域中心 + 一次逼近检测 + 覆盖式试探（试探点数 × 平均步长/速度 + 3 s/次）
        + 成功清除 5 s。

        ``probe_cap`` 用于**乐观估计**：实际执行时试探是逐点进行的、随时可以中止
        （多数目标前几个试探点就成功），所以判断"值不值得跑一趟"时把试探数封顶，
        否则会过于保守、把还有 400 s 的尾巴白白浪费掉。
        """
        p = np.asarray(from_pos if from_pos is not None else self.pos, float)
        c = np.asarray(br["center"], float)
        if n_probe is None:
            n_probe = self._probe_count(br)
        if probe_cap is not None:
            n_probe = min(int(n_probe), int(probe_cap))
        t_probe = n_probe * (self.cfg.probe_spacing_m / V_ROBOT + T_CLEAR_FAIL)
        return (math.dist(p, c) / V_ROBOT + T_MEASURE + T_SWITCH
                + t_probe + T_CLEAR_OK)

    def budget_feasible_shortlist(self, targets, start=None, budget=None):
        """从 ``targets`` 里挑出"剩余预算内走得完"的**近点子集**（贪心插入）。

        对"近点优先"的回答：滚动时域下每步只执行第一个目标，所以**只保留一个近点**
        在语义上已经等价于"优先清最近的点"；本函数真正的作用是判断
        "这一轮清理是否还排得下任何一个目标"，以及给排序一个**可达性**约束。

        排序键是"距当前位置"（近点优先），逐个累计
        ``走到区域中心 + 一次逼近检测 + 覆盖式试探 + 清除`` 的乐观代价，
        超预算就停止插入（后面更远的点更进不来）。
        """
        if not targets:
            return []
        p = np.asarray(start if start is not None else self.pos, float)
        rem = self.remaining() if budget is None else float(budget)
        room = rem - self.cfg.time_reserve_s - self.cfg.min_action_budget_s \
            - self.cfg.clear_shortlist_reserve_s
        ordered = sorted(targets, key=lambda t: math.dist(p, t[2]["center"]))
        out, used = [], 0.0
        for t in ordered:
            src = p if not out else np.asarray(out[-1][2]["center"], float)
            # estimate_clear_cost 已含"从 from_pos 走到区域中心"的行程
            cost = self.estimate_clear_cost(
                t[2], from_pos=src, probe_cap=self.cfg.optimistic_probe_steps)
            if used + cost > room:
                continue          # 这个点排不进，试下一个更远的（更近的已进来）
            out.append(t)
            used += cost
        return out

    def plan_clear_tour(self, targets, remaining_s=None):
        """**动态规划（滚动时域）**清理航路：每执行一步就重新规划一次。

        与旧版（清理阶段开始时规划一次、阶段内顺序冻结）的区别：

        * 每次调用都重读台账 (*@\ensuremath{\Rightarrow}@*) **途中新定位出来的频道立刻进入航路**；
        * 各频道用**最新**的定位区域中心（逼近过程中区域会变小、中心会移动）；
        * 以**当前位置**为起点重新做最近邻 + 2-opt（可用 ``tour`` 配置切换）；
        * 本步只执行规划结果的第一个，执行完立刻重规划（rolling horizon）。

        **为什么默认不按剩余预算裁剪目标表**：目标是"边走边判"执行的 ——
        逼近、一次检测、逐点试探各自都检查剩余时间、随时可以中止，多数目标在
        前几个试探点就成功。用统一的预估代价去裁剪，会把还剩 300–400 s 的尾巴
        整段浪费掉（实测：裁剪 39.8% vs 不裁剪 41.2%）。故默认只排序、不裁剪，
        由 :meth:`clear_target` / :meth:`sweep_clear` 在动作粒度上兜底。
        ``clear_budget_truncation=True`` 可打开裁剪做对照。

        ``clear_policy`` 两种策略：
        * ``tour``（默认）：NN + 2-opt 的**全局航路结构**，按顺序执行；
        * ``rollout``：一步前瞻 + SPT 推演，直接最大化"预算内可清除个数"
          （理论更贴合目标函数，但实测更差 36.3% vs 41.2%：贪心总挑最近的，
          最后把最远的目标拖成"永远清不掉"）。
        """
        if not targets:
            return []
        self.stats["n_replans"] = self.stats.get("n_replans", 0) + 1
        remaining = self.remaining() if remaining_s is None else float(remaining_s)
        budget = remaining - self.cfg.time_reserve_s - self.cfg.min_action_budget_s
        if self.cfg.clear_policy == "rollout":
            return self._plan_rollout(targets, budget)
        if self.cfg.clear_policy == "tour_greedy_cost":
            # 贪心"每次挑预计剩余预算花得完、且最近的"目标，滚动推进。
            # 与 'tour' 的差别：这里显式用**剩余预算**逐目标判定可达性，
            # 而不是先排一条全局航路、再边走边判。
            return self._plan_greedy_cost(targets, remaining)
        # ---- 可选：先用"可达近点子集"过滤，再做航路排序 ----
        pool = list(targets)
        if self.cfg.clear_shortlist:
            pool = self.budget_feasible_shortlist(pool, start=self.pos, budget=remaining)
            self.stats["n_shortlist_kept"] = self.stats.get("n_shortlist_kept", 0) \
                + len(pool)
            self.stats["n_shortlist_in"] = self.stats.get("n_shortlist_in", 0) \
                + len(targets)
            if not pool:
                return []
        # ---- tour 策略：全局航路结构（可选预算截断）----
        order = self.order_targets(list(pool), from_pos=self.pos)
        by_ch = {t[0]: t for t in pool}
        ordered = [by_ch[ch] for ch in order if ch in by_ch]
        if not self.cfg.clear_budget_truncation:
            return ordered
        out, used, pos = [], 0.0, tuple(self.pos)
        for t in ordered:
            c = self.estimate_clear_cost(t[2], from_pos=pos,
                                         probe_cap=self.cfg.optimistic_probe_steps)
            if used + c > budget:
                continue
            out.append(t)
            used += c
            pos = tuple(t[2]["center"])
        return out

    def _plan_greedy_cost(self, targets, remaining):
        """贪心 + 剩余预算可达性（``clear_policy='tour_greedy_cost'``）。

        诊断动机：末尾 26/49 个"已定位未清除"的目标在变成可清目标时剩余时间完全够用，
        却因为全局航路把它们排在后面而没轮到（清完前面的目标后预算就没了）。
        本策略每一步只在"当前剩余预算真的能走完"的目标里挑最近的，
        从而避免"先承诺一个远点、结果后面全部落空"。

        与 ``clear_budget_truncation`` 的区别：那个是对已排好的全局航路做事后裁剪，
        本策略是**在挑选时就**把可达性作为约束。
        """
        pos = np.asarray(self.pos, float)
        left = list(targets)
        out, used = [], 0.0
        room = remaining - self.cfg.time_reserve_s - self.cfg.min_action_budget_s
        while left:
            best, best_cost = None, None
            for t in left:
                c = self.estimate_clear_cost(t[2], from_pos=pos,
                                             probe_cap=self.cfg.optimistic_probe_steps)
                if best_cost is None or c < best_cost:
                    best, best_cost = t, c
            if best is None or used + best_cost > room:
                break
            out.append(best)
            used += best_cost
            pos = np.asarray(best[2]["center"], float)
            left.remove(best)
        return out

    def _plan_rollout(self, targets, budget):
        """一步前瞻 + SPT 推演（``clear_policy='rollout'``，见 :meth:`plan_clear_tour`）。"""
        k = len(targets)
        pos0 = tuple(self.pos)
        centers = [tuple(t[2]["center"]) for t in targets]
        nprobe = []
        for t in targets:
            nprobe.append(self._probe_count(t[2]))
        t_probe = [n * (self.cfg.probe_spacing_m / V_ROBOT + T_CLEAR_FAIL) for n in nprobe]

        def leg(a, b):
            return (math.dist(a, centers[b]) / V_ROBOT + T_MEASURE + T_SWITCH
                    + t_probe[b] + T_CLEAR_OK)

        cost0 = [leg(pos0, j) for j in range(k)]

        def rollout(cur, left, tb):
            left = set(left)
            n = 0
            while True:
                best_c, best_j = None, None
                for j in left:
                    c = leg(cur, j)
                    if c <= tb and (best_c is None or c < best_c):
                        best_c, best_j = c, j
                if best_j is None:
                    return n
                tb -= best_c
                n += 1
                cur = centers[best_j]
                left.discard(best_j)

        best_j, best_key = None, None
        for j in range(k):
            if cost0[j] > budget:
                continue
            n = 1 + rollout(centers[j], [i for i in range(k) if i != j], budget - cost0[j])
            key = (n, -cost0[j])
            if best_key is None or key > best_key:
                best_key, best_j = key, j
        if best_j is None:
            return []
        out = [targets[best_j]]
        used, pos = cost0[best_j], centers[best_j]
        for j in sorted((i for i in range(k) if i != best_j), key=lambda i: leg(pos, i)):
            c = leg(pos, j)
            if used + c > budget:
                continue
            out.append(targets[j])
            used += c
            pos = centers[j]
        return out

    def order_targets(self, targets, from_pos=None):
        """航路排序：最近邻 + 2-opt（以 ``from_pos`` 为固定起点）。

        目标点用各自的（最新）定位区域中心；k (*@\ensuremath{\le}@*) 20，暴力 2-opt 代价可忽略。

        ``tour`` 三种模式（对照用）：
        * ``nn``      ：纯最近邻 —— 第一步必然最近，但总路程偏长；
        * ``nn2opt``  ：最近邻 + 2-opt（**默认**）—— 总路程最短，但 2-opt 的
          区间反转会改写**第一步**（60 局实测有 14.1% 的排序其第一步不是最近的）；
        * ``nn2opt_keepfirst``：先钉住最近的点作第一步，只对**其后**的子序列做 2-opt。
          动机：滚动时域下**只执行第一步**，后面重排一次；所以"总路程最短"并不是
          正确的目标函数，"第一步最近的 + 剩余结构合理"才是。
        """
        start = tuple(self.pos if from_pos is None else from_pos)
        if len(targets) <= 1 or self.cfg.tour == "nn":
            targets.sort(key=lambda t: math.dist(start, t[2]["center"]))
            return [t[0] for t in targets]
        pts = [tuple(t[2]["center"]) for t in targets]
        k = len(pts)
        cur = start
        # 最近邻构造
        left = list(range(k))
        seq = []
        p = cur
        while left:
            j = min(left, key=lambda i: math.dist(p, pts[i]))
            seq.append(j)
            left.remove(j)
            p = pts[j]

        def total(order):
            d = math.dist(cur, pts[order[0]])
            for a, b in zip(order, order[1:]):
                d += math.dist(pts[a], pts[b])
            return d

        # 钉住第一步时，只允许反转 index >= 1 的区间
        lo = 1 if (self.cfg.tour == "nn2opt_keepfirst" and k >= 2) else 0
        best = total(seq)
        improved = True
        while improved:
            improved = False
            for i in range(lo, k - 1):
                for j in range(i + 1, k):
                    cand = seq[:i] + seq[i:j + 1][::-1] + seq[j + 1:]
                    d = total(cand)
                    if d < best - 1e-9:
                        seq, best, improved = cand, d, True
        return [targets[i][0] for i in seq]

    def clear_phase(self):
        """清理阶段：**每清一个目标就重新规划一次**（滚动时域）。

        循环体：读最新台账 → 动态规划 → 只执行第一步 → 顺路补测（可能新增可清理频道）
        → 回到循环顶部重新规划。这样"清除之后扫描过程中新定位的位置"会立刻进入航路，
        而且航路始终以当前位置为起点、用最新的区域中心优化。
        """
        first = self.clear_targets()
        if not first:
            return 0
        self.abort_clear_phase = False
        tour0 = self.plan_clear_tour(first) if self.cfg.dynamic_replan else None
        if self.cfg.dynamic_replan and not tour0:
            return 0          # 预算内一个都排不下：静默返回，让主循环把时间用于扫描
        if not self.can_afford(self.cfg.min_action_budget_s):
            return 0
        self.log(f"清理阶段（动态重规划）：{len(first)} 个频道已定位两次 -> "
                 f"{[(t[0], round(t[2]['radius'])) for t in first]}")
        n = 0
        guard = 0
        static_order = None
        tour = tour0
        initial_chs = {t[0] for t in (tour0 or first)}
        midphase = 0
        while not self.aborted and guard < 4 * len(self.recs) + 8:
            guard += 1
            if self.cfg.dynamic_replan:
                tour = self.plan_clear_tour(self.clear_targets())   # (*@\textcircled{1}@*) 每步重新读台账并规划
                if not tour:
                    self.log("剩余时间不足以完成任何一条清理航段 -> 结束清理阶段")
                    break
                ch, rec, _br = tour[0]                 # (*@\textcircled{2}@*) 只执行第一步
                if ch not in initial_chs:
                    midphase += 1                      # 本阶段中途才定位出来的目标
                    self.log(f"  动态航路：频道{ch} 是本阶段中途新定位的，立即插队执行")
                elif len(tour) > 1:
                    self.log(f"  动态航路：本步 频道{ch}，随后 "
                             f"{[t[0] for t in tour[1:4]]}{'…' if len(tour) > 4 else ''}")
            else:                                      # 静态对照（消融用）
                targets = self.clear_targets()
                if not targets:
                    break
                if not self.can_afford(self.cfg.min_action_budget_s):
                    break
                if static_order is None:
                    static_order = self.order_targets(list(targets))
                if not static_order:
                    break
                ch = static_order.pop(0)
                t = next((x for x in targets if x[0] == ch), None)
                if t is None:
                    continue
                rec = t[1]
            if self.clear_target(ch, rec):
                n += 1
            if self.abort_clear_phase:
                self.log("预算不足 -> 结束清理阶段（剩余时间留给扫描）")
                break
            self.enroute_scan()                        # (*@\textcircled{3}@*) 顺路补测（可能新增可清理频道）
        self.stats["n_midphase_targets"] = self.stats.get("n_midphase_targets", 0) + midphase
        return n

    def enroute_scan(self):
        """清理途中在当前位置顺手补测：既补充方位，也给新发现的频道补第二条方位。

        这一步很关键：发现一个新频道只要 6 s，但只有拿到**第二条**方位才能进入清理队列。
        所以一旦在本停点发现新频道，就立刻用第二问的"局部最优第二检测点"再测一条，
        让它本轮就能被清除，而不是等到下一次扫描（往往等不到）。
        """
        if not self.cfg.enroute_scan or self.aborted:
            return 0
        chans = self.channels_to_measure(limit=self.cfg.enroute_scan_max)
        if not chans:
            return 0
        cost = len(chans) * (T_MEASURE + T_SWITCH) + 2.0
        if not self.can_afford(cost + self.cfg.min_action_budget_s * 0.5):
            return 0
        self.log(f"  顺路补测 {len(chans)} 个频道 {chans}")
        n = self.scan_at(self.pos[0], self.pos[1], chans)
        if self.cfg.second_bearing_refine:
            fresh = [ch for ch in chans
                     if not self.recs[ch].cleared and self.recs[ch].n_bearings == 1]
            fresh.sort(key=lambda ch: -self.recs[ch].near_hits)
            for ch in fresh[:self.cfg.second_bearing_max]:
                if self.remaining() < self.cfg.second_bearing_budget_s + self.cfg.time_reserve_s:
                    break
                self.acquire_second_bearing(ch)
        return n

    def acquire_second_bearing(self, ch):
        """给"只有一条方位"的频道补第二检测点（局部最优，含行驶代价）。"""
        rec = self.recs[int(ch)]
        if rec.n_bearings != 1 or rec.cleared:
            return False
        p, info = best_refine_point(rec.pts, rec.svds, self.pos, None, n_t=100,
                                    travel_cap_m=self.cfg.refine_travel_cap_m * 1.4)
        if p is None:
            return False
        trav = math.dist(self.pos, (float(p[0]), float(p[1])))
        if not self.can_afford(trav / V_ROBOT + T_MEASURE + T_SWITCH + 20.0):
            return False
        self.log(f"  频道{ch} 新方位 -> 第二检测点 ({p[0]:.0f},{p[1]:.0f})，"
                 f"移动 {trav:.0f} m（E[D](*@\ensuremath{\approx}@*){info['E_diam_m']:.0f} m，"
                 f"预计代价 {info['cost_s']:.0f} s）")
        resp = self.measure(float(p[0]), float(p[1]), ch)
        if resp.get("measure_result") == "near":
            return self.do_clear(float(p[0]), float(p[1]), ch)
        self.stats["n_second_bearing"] = self.stats.get("n_second_bearing", 0) + 1
        return rec.n_bearings >= 2

    def choose_refine_point(self, ch, rec, br):
        """侧移精定位点：与上一条视线近似正交，把交会角拉到 60–90°。"""
        c = np.asarray(br["center"], float)
        if self.cfg.refine_mode == "exact":
            p, info = best_refine_point(rec.pts, rec.svds, self.pos,
                                        {"center": c}, n_t=120,
                                        travel_cap_m=self.cfg.refine_travel_cap_m)
            if p is not None:
                return np.asarray(p, float), info
        if not rec.pts:
            return None, None
        S_last = np.asarray(rec.pts[-1], float)
        u = c - S_last
        nu = float(np.linalg.norm(u))
        if nu < 1e-6:
            return None, None
        u = u / nu
        perp = np.array([-u[1], u[0]])
        rho = min(max(1.5 * br["radius"], 80.0), self.cfg.refine_leg_m)
        best = None
        for sgn in (1.0, -1.0):
            p = c + sgn * rho * perp
            if np.hypot(*p) > R_ARENA - 1.0:
                continue
            trav = float(np.hypot(*(p - np.asarray(self.pos, float))))
            if trav > self.cfg.refine_travel_cap_m:
                continue
            if best is None or trav < best[0]:
                best = (trav, p)
        if best is None:
            return None, None
        return best[1], {"mode": "perp", "rho": rho, "travel_m": best[0]}

    def probe_points(self, c, r, poly=None, tried=()):
        """覆盖式试探点。

        * 有定位区域多边形 ``poly`` 时：只在"到多边形距离 (*@\ensuremath{\le}@*) 20 m"的格点上试
          —— 交会区域常常是细长四边形，用外接圆（面积 (*@\ensuremath{\pi}@*)r*²）试探会浪费一倍以上的点；
        * 没有多边形时退化为覆盖半径 r 的圆域。
        格点间距 (*@\ensuremath{\le}@*) 20√3 m，保证区域内任一点到某个试探点 (*@\ensuremath{\le}@*) 20 m = 清除半径，**必然成功**。

        注意"必然成功"只在**候选点不被截断**时成立：``probe_max_points`` 会按"离中心
        由近及远"截断，等于丢掉覆盖区域的**外圈**，保证随之失效（实测 r* (*@\ensuremath{\ge}@*) 120 m 即开始
        截断，60 个点对应 r* (*@\ensuremath{\approx}@*) 120 m，而截断前需要 61 个）。因此默认
        ``probe_unlimited=True`` 取全部候选点；``probe_max_points`` 只在显式关掉
        ``probe_unlimited``（消融对照）时才起作用。
        """
        c = np.asarray(c, float)
        s = float(self.cfg.probe_spacing_m)
        rr = min(float(r), self.cfg.probe_max_radius_m) + R_CLEAR
        u = np.array([s, 0.0])
        v = np.array([s / 2.0, s * math.sqrt(3.0) / 2.0])
        N = int(math.ceil(rr / s)) + 1
        cand = []
        for i in range(-N, N + 1):
            for j in range(-N, N + 1):
                p = c + i * u + j * v
                if float(np.hypot(*(p - c))) <= rr + 1e-9:
                    cand.append(p)
        if poly is not None and len(poly) >= 3:
            P = np.asarray(poly, float)
            keep = []
            for p in cand:
                if _dist_point_polygon(p, P) <= R_CLEAR + 1e-9:
                    keep.append(p)
            cand = keep if keep else cand
        cand.sort(key=lambda p: float(np.hypot(*(p - c))))
        out = []
        for p in cand:
            if any(float(np.hypot(*(p - np.asarray(q, float)))) < 2.0 for q in tried):
                continue
            out.append(p)
        if self.cfg.probe_unlimited:
            return out
        return out[:self.cfg.probe_max_points]

    def sweep_clear(self, ch, rec, c, r, poly=None):
        """在区域 ``(c, r, poly)`` 上做覆盖式试探。

        返回 :data:`SWEEP_OK` / :data:`SWEEP_MISS` / :data:`SWEEP_BUDGET` 三者之一。
        **预算不足时立刻返回 ``SWEEP_BUDGET``，一个 ``/clear`` 都不发** ——
        否则会污染 ``clear_fail``、并在论文里被误读成"区域覆盖了却没找到目标"。
        """
        probes = self.probe_points(c, r, poly=poly, tried=rec.tried_probes)
        full = self.probe_points(c, r, poly=poly)
        n_try = 0
        for p in probes:
            if self.aborted:
                self.log(f"  覆盖式试探中断：程序已中止（已试 {n_try} 点）")
                return SWEEP_MISS if n_try else SWEEP_BUDGET
            if not self.can_afford(T_CLEAR_FAIL + 5.0):
                self.log(f"  覆盖式试探放弃：剩余 {self.remaining():.0f} s 不足以支付一次 "
                         f"/clear（候选 {len(full)} 个，实际尝试 {n_try} 个）")
                self.stats["n_sweep_aborted"] += 1
                return SWEEP_BUDGET if n_try == 0 else SWEEP_MISS
            trav = math.dist(self.pos, (float(p[0]), float(p[1])))
            if not self.can_afford(trav / V_ROBOT + T_CLEAR_FAIL + 3.0):
                self.log(f"  覆盖式试探放弃：剩余 {self.remaining():.0f} s 不足以走到下一个"
                         f"试探点（{trav:.0f} m；候选 {len(full)} 个，实际尝试 {n_try} 个）")
                self.stats["n_sweep_aborted"] += 1
                return SWEEP_BUDGET if n_try == 0 else SWEEP_MISS
            n_try += 1
            rec.tried_probes.append((float(p[0]), float(p[1])))
            if self.do_clear(float(p[0]), float(p[1]), ch):
                self.log(f"  覆盖式试探命中：第 {n_try}/{len(full)} 个候选点"
                         f"（间距 {self.cfg.probe_spacing_m:.0f} m）")
                return SWEEP_OK
        # 候选点全部试完仍未命中：这才是真正的"区域内没找到目标"
        self.stats["n_sweep_miss"] += 1
        self.log(f"  覆盖式试探未命中：{n_try}/{len(full)} 个候选点全部试过"
                 f"（{'多边形区域' if poly else '外接圆'}，r*={r:.0f} m）")
        return SWEEP_MISS

    def clear_target(self, ch, rec):
        """定位→逼近→（必要时）侧移精定位→覆盖式试探清除。

        失败归因**必须三分**（见 :data:`SWEEP_BUDGET`）：

        * ``SWEEP_OK``          —— 清除成功；
        * ``SWEEP_MISS``        —— 候选点全试过仍没打中 (*@\ensuremath{\Rightarrow}@*) 累加 ``clear_fail``（真的失败了）；
        * ``SWEEP_BUDGET``      —— 预算不足、**一个 /clear 都没发** (*@\ensuremath{\Rightarrow}@*) **不**累加 ``clear_fail``，
          也不报"未清除成功"，而是收敛回清理阶段由 ``plan_clear_tour`` 决定收工。
          这正是"预算耗尽后还调用 sweep_clear 并记一笔失败"那个统计/日志缺陷的修复点。
        """
        self.log(f"清理频道{ch}：当前 {rec.n_bearings} 条方位，"
                 f"起点 ({self.pos[0]:.0f},{self.pos[1]:.0f})")
        budget_stop = False
        # (1) 逼近定位区域中心（每次逼近顺带测一条方位，移动本来就要花）
        for _ in range(self.cfg.max_approach_steps):
            br = self.bounded_region(rec)
            if br is None:
                rec.clear_fail += 1
                return False
            c = np.asarray(br["center"], float)
            d = math.dist(self.pos, tuple(c))
            if d <= max(br["radius"], self.cfg.arrive_radius_m):
                break
            cost = d / V_ROBOT + T_MEASURE + T_SWITCH + 4.0
            if not self.can_afford(cost):
                # 走不过去了：进入第 (3) 步，由覆盖式试探统一判定"预算不足"
                budget_stop = True
                break
            resp = self.measure(float(c[0]), float(c[1]), ch)
            if resp.get("measure_result") == "near":
                return self.do_clear(float(c[0]), float(c[1]), ch)
            if self.aborted:
                return False
        # (2) 区域还大：侧移一条正交视线，把 r* 压小
        if not budget_stop:
            br = self.bounded_region(rec)
            if br is not None and br["radius"] > self.cfg.direct_probe_radius_m:
                p, info = self.choose_refine_point(ch, rec, br)
                if p is not None:
                    trav = math.dist(self.pos, (float(p[0]), float(p[1])))
                    cost = trav / V_ROBOT + T_MEASURE + T_SWITCH + 4.0
                    if self.can_afford(cost):
                        self.stats["n_refine"] += 1
                        self.log(f"  侧移精定位 -> ({p[0]:.0f},{p[1]:.0f})，"
                                 f"移动 {trav:.0f} m（{info}）")
                        resp = self.measure(float(p[0]), float(p[1]), ch)
                        if resp.get("measure_result") == "near":
                            return self.do_clear(float(p[0]), float(p[1]), ch)
        # (3) 覆盖式试探清除
        br = self.bounded_region(rec)
        if br is None:
            rec.clear_fail += 1
            return False
        c = np.asarray(br["center"], float)
        r = min(br["radius"], self.cfg.probe_max_radius_m)
        poly = br.get("poly")
        st = self.sweep_clear(ch, rec, c, r, poly=poly)
        if st == SWEEP_OK:
            return True
        if st == SWEEP_BUDGET:
            # 预算不足、"一个 /clear 都没发出"：不是清除失败，不计 clear_fail、
            # 也不上升为空/随机搜索；交由调用方收敛本阶段。
            self.abort_clear_phase = True
            self.log(f"  (*@\ensuremath{\curvearrowright}@*) 频道{ch} 未尝试清除（预算不足），本阶段收敛")
            return False
        # (4) 区域内没找到：说明区域被算小了（或误差超界），扩大半径再试一轮
        if br["radius"] <= self.cfg.probe_max_radius_m:
            r2 = min(br["radius"] * 1.6 + 20.0, self.cfg.probe_max_radius_m)
            self.log(f"  区域内未发现目标 -> 扩大到 r={r2:.0f} m 再试")
            st2 = self.sweep_clear(ch, rec, c, r2, poly=None)
            if st2 == SWEEP_OK:
                return True
            if st2 == SWEEP_BUDGET:
                self.abort_clear_phase = True
                self.log(f"  (*@\ensuremath{\curvearrowright}@*) 频道{ch} 扩大半径后未尝试（预算不足），本阶段收敛")
                return False
        rec.clear_fail += 1
        self.log(f"  [x] 频道{ch} 本轮未清除成功（候选点已试遍仍无目标，留待后续扫描后再清）")
        return False

    # ---------------- 主循环 ----------------
    def run(self):
        t0 = time.time()
        resp = self.arena.enter()
        if resp.get("accepted") is not True:
            raise ArenaError(f"/enter 未被接受：{resp}")
        self.log(f"进入靶区：可用预算 {self.remaining():.0f} s"
                 f"（{self.arena.budget_kind()}）")
        try:
            # 第 0 轮：机器狗本来就在当前位置（原点），直接扫描全部频道
            chans0 = self.channels_to_measure()
            self.log(f"第 0 轮扫描（原地）：待测 {len(chans0)} 个频道 {chans0}")
            self.scan_points.append(self.pos)
            self.scan_at(self.pos[0], self.pos[1], chans0)
            rounds = 1
            stall = 0
            while (not self.aborted and rounds < self.cfg.max_rounds
                   and self.remaining() > self.cfg.time_reserve_s):
                did = 0
                t_before = float(self.arena.virtual_time_s)
                if self.clear_targets() and (rounds >= self.cfg.min_scan_rounds_before_clear
                                             or not self.plan_scan()):
                    did += self.clear_phase()
                    self.stats["t_clear_s"] += (float(self.arena.virtual_time_s)
                                                - t_before)
                    self.stats["t_clear_ends_s"] = float(self.arena.virtual_time_s)
                if self.aborted or self.remaining() <= self.cfg.time_reserve_s:
                    break
                plan = self.plan_scan()
                if plan is None:
                    # 扫描被预算闸门否掉时，plan_scan() 也可能只是"暂无值得测的频道"；
                    # 此时若还有可清目标，就再走一次清理阶段，而不是直接判为 stall。
                    if self.clear_targets():
                        did += self.clear_phase()
                        self.stats["t_clear_s"] += (float(self.arena.virtual_time_s)
                                                    - t_before)
                        self.stats["t_clear_ends_s"] = float(self.arena.virtual_time_s)
                    if self.aborted:
                        break
                elif self.can_afford(plan["cost_s"]) \
                        and (did == 0 or not self.clear_targets()):
                    did += self.execute_scan(plan)
                    self.stats["t_scan_s"] += (float(self.arena.virtual_time_s)
                                               - t_before)
                    rounds += 1
                if did == 0:
                    stall += 1
                    if stall >= 2:
                        self.log("没有可清理的频道、也没有值得检测的频道 -> 提前结束")
                        break
                else:
                    stall = 0
        finally:
            try:
                self.arena.exit()
            except ArenaError as e:  # 接口已关闭等
                self.log(f"/exit 异常：{e}")
        self.stats["program_runtime_s"] = time.time() - t0
        return self.report()

    # ---------------- 汇报 ----------------
    def report(self):
        truth = self.arena.truth()
        n_cleared = sum(1 for r in self.recs.values() if r.cleared)
        total = (truth or {}).get("n_sources")
        vt = float(self.arena.virtual_time_s)
        last = float(self.stats.get("vtime_at_last_clear", vt))
        st = getattr(self.arena, "stats", {})
        rep = {
            "cleared": int(n_cleared),
            "n_sources": total,
            "clear_ratio": (n_cleared / total) if total else None,
            "virtual_time_s": vt,
            "avg_clear_time_s": (last / n_cleared) if n_cleared else None,
            "program_runtime_s": float(self.stats.get("program_runtime_s", 0.0)),
            "budget_kind": self.arena.budget_kind(),
            "remaining_budget_s": self.remaining(),
            "n_measure": self.stats["n_measure"],
            "n_bearings": self.stats["n_bearings"],
            "n_clear": self.stats["n_clear"],
            "n_clear_fail": self.stats["n_clear_fail"],
            "n_scan_rounds": self.stats["n_scan_rounds"],
            "n_refine": self.stats["n_refine"],
            "n_second_bearing": self.stats.get("n_second_bearing", 0),
            "n_replans": self.stats.get("n_replans", 0),
            "n_midphase_targets": self.stats.get("n_midphase_targets", 0),
            "n_rejected": self.stats["n_rejected"],
            "travel_m": float(st.get("travel_m", 0.0)),
            "mock_stats": dict(st),
            "n_sweep_aborted": self.stats.get("n_sweep_aborted", 0),
            "n_sweep_miss": self.stats.get("n_sweep_miss", 0),
            "n_probe_issued": self.stats.get("n_probe_issued", 0),
            "n_scan_gated": self.stats.get("n_scan_gated", 0),
            "t_scan_s": self.stats.get("t_scan_s", 0.0),
            "t_clear_s": self.stats.get("t_clear_s", 0.0),
            "t_clear_ends_s": self.stats.get("t_clear_ends_s", 0.0),
            "n_leftover_located": sum(
                1 for r in self.recs.values()
                if not r.cleared and r.n_bearings >= self.cfg.clear_min_bearings),
            "scan_points": [[round(float(x), 2), round(float(y), 2)]
                            for (x, y) in self.scan_points],
            "channels": {str(ch): {"n_bearings": rec.n_bearings,
                                   "cleared": bool(rec.cleared),
                                   "no_signal": rec.no_signal,
                                   "near": rec.near_hits,
                                   "status": self.region(rec).get("status"),
                                   "r_star_m": (self.bounded_region(rec) or {}).get("radius")}
                         for ch, rec in self.recs.items() if rec.n_bearings or rec.cleared},
            "truth": truth,
        }
        return rep


# --------------------------------------------------------------------------
# 自检：策略层的结构性断言（离线 mock 上跑）
# --------------------------------------------------------------------------
def selfcheck_mock(verbose=True, seed=20260913, cases=3):
    out = []
    ok_all = True

    def rec(name, ok, detail=""):
        nonlocal ok_all
        ok_all = ok_all and bool(ok)
        out.append({"check": name, "ok": bool(ok), "detail": detail})
        if verbose:
            print(f"[{'OK ' if ok else 'FAIL'}] {name}  {detail}")

    # S1 覆盖式试探点确实能覆盖圆域（任意点 20 m 内有试探点）
    rb = P3Robot(MockArena(seed=1), P3Config())
    rng = np.random.default_rng(7)
    worst = 0.0
    for (cx, cy, r) in ((0.0, 0.0, 40.0), (500.0, -300.0, 80.0), (-900.0, 900.0, 25.0)):
        pts = np.array(rb.probe_points((cx, cy), r))
        for _ in range(400):
            a = rng.uniform(0, 2 * math.pi)
            rad = r * math.sqrt(rng.random())
            p = np.array([cx + rad * math.cos(a), cy + rad * math.sin(a)])
            d = float(np.min(np.hypot(pts[:, 0] - p[0], pts[:, 1] - p[1])))
            worst = max(worst, d)
    rec("S1 覆盖式试探：区域内任意点到最近试探点 <= 20 m（清除半径）", worst <= R_CLEAR + 1e-6,
        f"最大覆盖距离 {worst:.2f} m（试探点间距 {rb.cfg.probe_spacing_m} m）")

    # S2 端到端：mock 上跑通、清除率、时间预算
    ratios, times, cleared = [], [], []
    for k in range(cases):
        arena = MockArena(seed=seed + 100 * k)
        rb = P3Robot(arena, P3Config(), log=None)
        rep = rb.run()
        ratios.append(rep["clear_ratio"])
        times.append(rep["avg_clear_time_s"])
        cleared.append(rep["cleared"])
        if verbose:
            print(f"   案例{seed + 100 * k}: 清除 {rep['cleared']}/{rep['n_sources']}，"
                  f"虚拟时间 {rep['virtual_time_s']:.0f}s，"
                  f"平均 {rep['avg_clear_time_s']:.1f}s，"
                  f"扫描 {rep['n_scan_rounds']} 轮，检测 {rep['n_measure']} 次")
    rec("S2 mock 演练：清除率>0 且每个被清除源平均时间有限",
        all(r is not None and r > 0 for r in ratios) and all(t and t > 0 for t in times),
        f"{cases} 局：清除率 {[round(r, 3) for r in ratios]}，平均时间 "
        f"{[round(t, 1) for t in times]} s")

    # S3 策略不会超预算（mock 会计时；超时会 accepted=false）
    arena = MockArena(seed=seed + 7, budget_s=400.0)
    rb = P3Robot(arena, P3Config(), log=None)
    rep = rb.run()
    rec("S3 策略在时间预算内收工（不触发超时拒动）",
        rep["n_rejected"] == 0 and rep["virtual_time_s"] <= 400.0 + 1e-9,
        f"虚拟时间 {rep['virtual_time_s']:.1f}s / 400s，拒动 {rep['n_rejected']} 次")

    # S4 已清除的频道不再被检测（"跳过已清理过的点，节省时间"）
    arena = MockArena(seed=seed + 11)
    rb = P3Robot(arena, P3Config(), log=None)
    rep = rb.run()
    after = {}
    for ch, r in rb.recs.items():
        if r.cleared and r.cleared_at is not None:
            k = sum(1 for (_x, _y, t, _res) in r.meas_log if t > r.cleared_at + 1e-9)
            if k:
                after[ch] = k
    rec("S4 已清除频道不再被检测", not after,
        f"已清除 {sum(1 for r in rb.recs.values() if r.cleared)} 个频道；"
        f"清除后仍被检测的频道 {after or '无'}")

    # S5 定位区域一定包含真源（问题 1 的交会算法 + 模拟器误差界的一致性）
    bad5 = []
    n5 = 0
    src_map = {s["channel"]: (s["x_m"], s["y_m"]) for s in (rep["truth"] or {}).get("sources", [])}
    for ch, r in rb.recs.items():
        if r.n_bearings < 2:
            continue
        br = rb.bounded_region(r)
        if not br or not br.get("poly"):
            continue
        n5 += 1
        xy = src_map.get(ch)
        if xy is None:
            continue
        P = np.asarray(br["poly"], float)
        if not _point_in_polygon(np.asarray(xy, float), P):
            bad5.append(ch)
    rec("S5 定位区域（多边形）包含真源", not bad5,
        f"检查 {n5} 个已定位频道，越界 {bad5 or '无'}")

    # S6 探索兜底：没有任何方位时选"离已访问点最远"的点
    rb6 = P3Robot(MockArena(seed=seed + 3), P3Config(), log=None)
    sel6, _f = rb6.choose_scan_point()
    rec("S6 无方位时退化为探索点（离已访问点最远）",
        sel6["kind"] == "explore" and math.hypot(*sel6["xy"]) > 1000.0,
        f"探索点 {tuple(round(v) for v in sel6['xy'])}，距原点 "
        f"{math.hypot(*sel6['xy']):.0f} m")

    # S7 动态重规划：每清一个目标都重规划一次；静态对照为 0
    arena7 = MockArena(seed=seed + 13)
    rb7 = P3Robot(arena7, P3Config(), log=None)
    rep7 = rb7.run()
    arena7s = MockArena(seed=seed + 13)
    rb7s = P3Robot(arena7s, P3Config(dynamic_replan=False), log=None)
    rep7s = rb7s.run()
    rec("S7 动态重规划：清除途中每步重规划，静态对照 0 次",
        rep7["n_replans"] >= rep7["cleared"] >= 1 and rep7s["n_replans"] == 0,
        f"动态：清除 {rep7['cleared']} 个、重规划 {rep7['n_replans']} 次、"
        f"中途新定位插队 {rep7['n_midphase_targets']} 个；"
        f"静态：重规划 {rep7s['n_replans']} 次、中途插队 {rep7s['n_midphase_targets']} 个")

    # S8 动态航路：默认不裁剪（目标齐全）、顺序与航路排序一致
    rb8 = P3Robot(MockArena(seed=seed + 21), P3Config(), log=None)
    rb8.arena.enter()
    for ch, (gx, gy) in ((1, (300.0, 200.0)), (2, (-500.0, 700.0)), (3, (900.0, -400.0))):
        r8 = ChannelRec(ch)
        # 在真源周围取两条正交视线（交会角 90°）(*@\ensuremath{\Rightarrow}@*) 定位区域小且有界
        for ang, dist in ((0.0, 500.0), (90.0, 500.0)):
            sx = gx + dist * math.cos(math.radians(ang))
            sy = gy + dist * math.sin(math.radians(ang))
            a = math.degrees(math.atan2(gy - sy, gx - sx)) % 360.0
            r8.pts.append((sx, sy))
            r8.svds.append(a)
        rb8.recs[ch] = r8
    tg8 = rb8.clear_targets()
    tour8 = rb8.plan_clear_tour(tg8)
    rec("S8 动态航路：默认不裁剪、顺序与航路排序一致",
        len(tg8) == 3 and len(tour8) == 3
        and [t[0] for t in tour8] == rb8.order_targets(list(tg8)),
        f"{len(tg8)} 个已定位频道 -> 航路 {[t[0] for t in tour8]}，"
        f"各区域半径 {[round(t[2]['radius']) for t in tour8]} m")

    # S9 覆盖式试探**不得被点数上限截断**（截断会丢掉外圈、破坏 20 m 覆盖保证）
    rb9 = P3Robot(MockArena(seed=seed + 31), P3Config(), log=None)
    n9_uncapped = len(rb9.probe_points(np.zeros(2), 195.0))
    rb9c = P3Robot(MockArena(seed=seed + 31),
                   P3Config(probe_unlimited=False, probe_max_points=60), log=None)
    n9_capped = len(rb9c.probe_points(np.zeros(2), 195.0))
    # 覆盖保证：r*=195 m 的外接圆内任一点到最近候选点必须 <= 20 m
    pts9 = np.asarray(rb9.probe_points(np.zeros(2), 195.0), float)
    rng9 = np.random.default_rng(11)
    worst9 = 0.0
    for _ in range(400):
        a = rng9.uniform(0, 2 * math.pi)
        rad = 195.0 * math.sqrt(rng9.random())
        p = np.array([rad * math.cos(a), rad * math.sin(a)])
        worst9 = max(worst9, float(np.min(np.hypot(pts9[:, 0] - p[0],
                                                   pts9[:, 1] - p[1]))))
    rec("S9 覆盖式试探不被截断（大 r* 时仍满足 20 m 覆盖保证）",
        worst9 <= R_CLEAR + 1e-6 and n9_uncapped > n9_capped,
        f"r*=195 m：{n9_uncapped} 个候选点（封顶对照 {n9_capped} 个），"
        f"最大覆盖距离 {worst9:.2f} m")

    # S10 预算不足时**不得**把"未尝试"记成"清除失败"（失败归因三分）
    arena10 = MockArena(seed=seed + 41)
    rb10 = P3Robot(arena10, P3Config(), log=None)
    rb10.arena.enter()
    rec10 = ChannelRec(4)
    for (sx, sy) in ((0.0, 0.0), (0.0, 300.0)):     # 两条正交视线 -> 小而紧的区域
        a = math.degrees(math.atan2(150.0 - sy, 100.0 - sx)) % 360.0
        rec10.pts.append((sx, sy))
        rec10.svds.append(a)
    rb10.recs[4] = rec10
    br10 = rb10.bounded_region(rec10)
    arena10.vtime = arena10.budget - 0.5          # 只剩 0.5 s：连一次 /clear 都付不起
    f10 = rb10.remaining()
    st10 = rb10.sweep_clear(4, rec10, np.asarray(br10["center"], float),
                            float(br10["radius"]), poly=br10.get("poly"))
    rec("S10 预算不足时不误记 clear_fail / 不发 /clear",
        st10 == SWEEP_BUDGET and rec10.clear_fail == 0
        and rb10.stats["n_probe_issued"] == 0
        and rb10.stats["n_sweep_aborted"] == 1,
        f"剩余预算 {f10:.1f} s（付不起一次 /clear）(*@\ensuremath{\Rightarrow}@*) 返回 {st10!r}，"
        f"clear_fail={rec10.clear_fail}，发出 /clear {rb10.stats['n_probe_issued']} 次，"
        f"n_sweep_aborted={rb10.stats['n_sweep_aborted']}")

    # S11 区域中心在"排序决策 → 执行"之间**不发生漂移**（否则排序依据会过期）
    arena11 = MockArena(seed=seed + 51)
    rb11 = P3Robot(arena11, P3Config(), log=None)
    orig_plan11 = rb11.plan_clear_tour
    orig_clear11 = rb11.clear_target
    snap11 = {}
    drift11 = []

    def plan11(targets, remaining_s=None):
        for t in (targets or []):
            snap11[t[0]] = tuple(t[2]["center"])
        return orig_plan11(targets, remaining_s=remaining_s)

    def clear11(ch, rec):
        br = rb11.bounded_region(rec)
        if br is not None and ch in snap11:
            drift11.append(math.dist(snap11[ch], tuple(br["center"])))
        return orig_clear11(ch, rec)

    rb11.plan_clear_tour = plan11
    rb11.clear_target = clear11
    rb11.run()
    rec("S11 区域中心在'排序决策→执行'之间不漂移",
        bool(drift11) and max(drift11) == 0.0,
        f"检查 {len(drift11)} 次执行，最大漂移 {max(drift11) if drift11 else float('nan'):.1f} m"
        f"（新方位只在 /measure 后产生，逼近途中不产生）")

    # S12 航路排序的三种模式：第一步必为最近邻（nn / nn2opt_keepfirst），
    #     或允许 2-opt 改写第一步但总路程不更差（nn2opt）
    rb12 = P3Robot(MockArena(seed=seed + 61), P3Config(), log=None)
    rb12.arena.enter()
    pts12 = [(300.0, 200.0), (-500.0, 700.0), (900.0, -400.0),
             (-1200.0, -300.0), (200.0, -1500.0)]
    for i, (gx, gy) in enumerate(pts12, start=1):
        r12 = ChannelRec(i)
        for ang in (0.0, 90.0):
            sx = gx + 500.0 * math.cos(math.radians(ang))
            sy = gy + 500.0 * math.sin(math.radians(ang))
            r12.pts.append((sx, sy))
            r12.svds.append(math.degrees(math.atan2(gy - sy, gx - sx)) % 360.0)
        rb12.recs[i] = r12
    tg12 = rb12.clear_targets()
    start12 = tuple(rb12.pos)
    nearest12 = min(tg12, key=lambda t: math.dist(start12, t[2]["center"]))[0]

    def plen12(order):
        seq = [next(t for t in tg12 if t[0] == ch) for ch in order]
        d = math.dist(start12, seq[0][2]["center"])
        for a, b in zip(seq, seq[1:]):
            d += math.dist(a[2]["center"], b[2]["center"])
        return d

    rb12.cfg.tour = "nn"
    o_nn = rb12.order_targets(list(tg12))
    rb12.cfg.tour = "nn2opt_keepfirst"
    o_kf = rb12.order_targets(list(tg12))
    rb12.cfg.tour = "nn2opt"
    o_2o = rb12.order_targets(list(tg12))
    rec("S12 航路排序：nn 与 keepfirst 第一步为最近邻；2-opt 不劣于 nn",
        o_nn[0] == nearest12 and o_kf[0] == nearest12
        and plen12(o_2o) <= plen12(o_nn) + 1e-9,
        f"最近目标=频道{nearest12}；nn 第一步={o_nn[0]}，keepfirst={o_kf[0]}，"
        f"2-opt={o_2o[0]}；路程 nn={plen12(o_nn):.0f} m，2-opt={plen12(o_2o):.0f} m")

    if verbose:
        print(f"\n自检总结：{'全部通过' if ok_all else '存在失败项'}")
    return {"ok": ok_all, "checks": out}


def main():
    ap = argparse.ArgumentParser(description="问题3 机器狗策略（离线自检 / 单局演练）")
    ap.add_argument("--selfcheck", action="store_true")
    ap.add_argument("--mock", action="store_true", help="跑一局离线演练并打印明细")
    ap.add_argument("--seed", type=int, default=20260913)
    ap.add_argument("--dist-weight", type=float, default=1.5)
    ap.add_argument("--grid-step", type=float, default=100.0)
    ap.add_argument("--quiet", action="store_true")
    args = ap.parse_args()
    if args.selfcheck:
        res = selfcheck_mock(verbose=True)
        raise SystemExit(0 if res["ok"] else 1)
    cfg = P3Config(dist_weight=args.dist_weight, grid_step=args.grid_step)
    arena = MockArena(seed=args.seed)
    rb = P3Robot(arena, cfg, log=None if args.quiet else print)
    rep = rb.run()
    print(json.dumps({k: v for k, v in rep.items()
                      if k not in ("truth", "channels", "mock_stats")},
                     ensure_ascii=False, indent=2))


if __name__ == "__main__":
    main()

\end{lstlisting}
\subsection{\texttt{src/p3\_sweep.py}}
% SHA256 f4534a86ca66a9ffecc0dc1934dd00df79701b3286ff4d923b8725f1eb7ea807
\begin{lstlisting}[language=python,escapeinside={(*@}{@*)}]
"""问题 3 的"覆盖扫圈"策略（与 `p3_robot.P3Robot` 的"期望场贪心"并列可选）。

设计动机（来自 2026-09-11 两局在线演练的实测）
--------------------------------------------
1. **不必"发现即清除"**：实测第 2 局的 14 个源在 3353 s 就全部清完，程序又空转
   1932 s（37%）——因为它不知道"已经清完"。反过来，第 1 局有 3114 s（45%）花在
   "场上一条方位都没有"的阶段，决策场 `E_norm (*@\ensuremath{\equiv}@*) 0`、score 只剩距离项，
   机器人就在圆心附近 500 m 步长乱蹭。
2. **"1000 m 已扫过"是可以判定的**：干扰源的有效接收半径 (*@\ensuremath{\ge}@*) `R_RECV_MIN = 1000 m`。
   因此"某点距所有已扫检测点都 > 1000 m"(*@\ensuremath{\Rightarrow}@*)"该点不可能存在未清除的源"是**严格**的。
   把靶区按 1000 m 的覆盖半径铺开，就能得到一个**可证明完备**的扫描航线：
   走完即"靶区内不存在未清除源"，可以立刻 `/exit`，不需要盲探、不需要走回头路。
3. **原点那一次检测白送了内圈**：原点已扫 (*@\ensuremath{\Rightarrow}@*) 半径 1000 m 以内不必再去。
   需要覆盖的只剩 r (*@\ensuremath{\in}@*) [1000, 1800] 的环带，因此停靠点不必贴边，
   沿 r (*@\ensuremath{\approx}@*) 1100 / 1500 / 1800 三圈铺开即可（见 :func:`lattice_candidates`）。

策略本体
--------
    score(g) = w_cover · (g 的 1000 m 圆内"未覆盖"面积占比)
             − w_dir   · cos(方位 − 期望前进方位)      # 保持顺时针推进，不回头
             − w_near  · (1 − d/最大可能距离)          # 就近取点，别横穿靶区

    pick = argmax score     （只从格点候选里选；已停靠点 300 m 邻域排除）

每到一个停靠点：先 `clear_targets()`（顺路清除已经定位好的源），再测该点值得测的频道。
覆盖率达到 100%（或时间不够）就收工——**因为此时可以确定没有残留源**。
"""
from __future__ import annotations

import math
import time
from dataclasses import dataclass

import numpy as np

from p3_arena import (  # noqa: E402
    ArenaError, R_ARENA, R_CLEAR, R_RECV_MIN, R_RECV_MAX, T_CLEAR_FAIL, T_CLEAR_OK,
    T_MEASURE, T_SWITCH, V_ROBOT,
)
from p3_robot import P3Config, P3Robot  # noqa: E402
from p3_expect_field import (  # noqa: E402
    best_refine_point, exact_expected_diameter, ray_source_samples,
)


# --------------------------------------------------------------------------
# 候选停靠点：同心环格点（"扫圈"航线就是在这些点之间走）
# --------------------------------------------------------------------------
DEFAULT_RINGS = ((1100.0, 8),)
"""``(半径 m, 该圈上的点数)``。默认只用一圈；如有需要可在后面追加"补漏圈"。

为什么一圈就够 —— 这是把题目的物理条件用足之后的结论：

* 干扰源有效接收半径 (*@\ensuremath{\ge}@*) ``R_RECV_MIN = 1000 m`` (*@\ensuremath{\Rightarrow}@*) 某点距**所有**检测点都 > 1000 m 时，
  该处不可能存在未清除源（**严格**，不是启发式）；
* 原点已经测过一次（20 个频道全扫），所以 **r (*@\ensuremath{\le}@*) 1000 m 的整个内圈不必再去**；
* 于是只需要覆盖环带 r (*@\ensuremath{\in}@*) [1000, 1800]。数值枚举（`python src/p3_sweep.py --coverage`）：

  | 单圈 | 环带内最大未覆盖距离 | 30 局虚拟时间 | 30 局清除率 | 30 局全清 |
  |---|---|---|---|---|
  | **r=1100, n=8（默认）** | **853 m** | **4243 s** | 0.978 | 23/30 |
  | r=1150, n=8 | — | 4237 s | 0.970 | 21/30 |
  | r=1250, n=8 | 788 m | 4516 s | 0.973 | 22/30 |
  | r=1400, n=8 | 721 m | 4621 s | 0.981 | **25/30** |

  都完备（< 1000 m）。**r=1100 最快、全清数居中**；r=1400 最稳（全清 25/30）但慢 9%。
  若把"全清"看得比时间更重，可一行改成 `DEFAULT_RINGS = ((1400.0, 8),)`。
"""


def lattice_candidates(rings=DEFAULT_RINGS, stagger=True):
    """同心环格点：``[(x, y), ...]``。相邻环之间做半格错开，避免放射状缝隙。"""
    out = []
    for k, (r, n) in enumerate(rings):
        off = (math.pi / n) if (stagger and k % 2 == 1) else 0.0
        for i in range(int(n)):
            a = off + 2.0 * math.pi * i / n
            out.append((float(r) * math.cos(a), float(r) * math.sin(a)))
    return out


def rotate_lattice(cand, level, frac, ring_n):
    """把某个圈的候选点整体旋转 ``frac`` 个站位（第 2 圈错开半格用）。"""
    out = []
    step = 2.0 * math.pi / max(int(ring_n), 1)
    for k, (x, y) in enumerate(cand):
        a = math.atan2(float(y), float(x)) + frac * step
        r = math.hypot(float(x), float(y))
        out.append((r * math.cos(a), r * math.sin(a)))
    return out


def coverage_report(pts, cover_r=R_RECV_MIN, step=25.0, r_arena=R_ARENA,
                    r_inner=0.0):
    """给定停靠点集合，返回 ``(最大未覆盖距离, 未覆盖面积占比, 覆盖率)``。

    "未覆盖距离" = 评估区内某点到最近停靠点的距离；它 > ``cover_r`` 即说明该处
    若有源则必然已被听到（源的接收半径 (*@\ensuremath{\ge}@*) cover_r）—— 所以这是**完备性度量**。

    ``r_inner``：只评估 ``r_inner (*@\ensuremath{\le}@*) |p| (*@\ensuremath{\le}@*) r_arena`` 的环带。原点已扫过时，
    ``r (*@\ensuremath{\le}@*) R_RECV_MIN`` 的内圈无需再覆盖，评估环带才是正确的口径。
    """
    R = float(r_arena)
    xs = np.arange(-R, R + 1e-9, step)
    GX, GY = np.meshgrid(xs, xs)
    rr = np.hypot(GX, GY)
    m = (rr <= R + 1e-9) & (rr >= float(r_inner) - 1e-9)
    gx, gy = GX[m], GY[m]
    if not pts:
        return float("inf"), 1.0, 0.0
    P = np.asarray(pts, float)
    d2 = ((gx[:, None] - P[None, :, 0]) ** 2 + (gy[:, None] - P[None, :, 1]) ** 2)
    dmin = np.sqrt(np.min(d2, axis=1))
    unc = dmin > float(cover_r)
    return float(np.max(dmin)), float(unc.mean()), float(1.0 - unc.mean())


# --------------------------------------------------------------------------
# 配置
# --------------------------------------------------------------------------
@dataclass
class SweepConfig(P3Config):
    """在 `P3Config` 上追加"扫圈"参数（继承 `P3Config` 的全部字段 + 下面这些）。"""
    sweep_rings: tuple = DEFAULT_RINGS
    cover_radius_m: float = R_RECV_MIN   # 覆盖判定半径（有效接收半径下界）
    cover_w: float = 1.0                 # 未覆盖收益的权重
    near_w: float = 0.25                 # 就近权重（避免横穿靶区）
    jump_w: float = 0.00035              # 跳跃惩罚（每米），抑制"绕回已走过的扇区"
    min_gain: float = 0.25               # 低于此边际收益就不值得专门跑一趟
    stop_visited_radius_m: float = 300.0  # 已停靠点邻域排除半径
    # ---- 清除时机 ----
    clear_mode: str = "detour"           # detour（只清顺路的）| station（每站都清）| none
    clear_detour_max_m: float = 350.0    # detour 模式：绝对距离上限（m）
    onlap_clear: bool = False            # **边绕边清**：一圈走完就清完，不再有"圈后统一清除"
    onlap_cos: float = 0.10              # 只在"前方"这个夹角内（cos）的候选才顺路处理
    onlap_max_m: float = 1500.0          # 候选(或单方位源的推算位置)离本站的最远距离
    onlap_min_m: float = 120.0           # 比这更近的目标不单独绕（到站测量时顺手清）
    onlap_single: bool = True            # 单方位频道也用"推算位置"参与顺路判定
    onlap_after_ring: bool = False       # 走完圈后再清一次剩余的已定位目标
    onlap_step_m: float = 320.0          # (*@\ensuremath{\star}@*)"取第二方位"的侧移步长上限（m）：小步就够，
                                         #   不要跑到推算位置去（实测原来跑 1400-1600 m）
    onlap_step_frac: float = 0.25        # 侧移步长 = min(onlap_step_m, 该比例 × 到目标的距离)
    onlap_min_step_m: float = 140.0      # 步长下限（太短交会角没改善）
    onlap_angle_deg: float = 50.0        # 侧移方向与旧视线的夹角（交会角 (*@\ensuremath{\approx}@*) 这个值）
    onlap_max_steps: int = 3             # 最多走几步去把区域压到可清除
    skip_station_near_target_m: float = 260.0  # 站位离"顺路目标"这么近时可考虑省掉该站
    skip_station_max_uncovered: float = 1150.0  # 省站后"最大未覆盖距离"仍须低于它
    start_toward_sources: bool = False   # 第一站朝"信息方向"出发（实测有 bug 且无收益，见 review §6.2）
    debug_onlap: bool = False            # 打印"边绕边清"的候选判定过程
    start_min_sources: int = 1           # 至少这么多条单方位才据此定起始方向
    # 下面两个"顺路"判据默认**关闭**（实测在 30 局上都没有正向收益，见 review/p3-sweep-design.md §6）：
    enroute_min_cos: float = 1.1         # >1 表示关闭"按方向判定顺路"（仅用 350 m 兜底）
    enroute_max_m: float = 1400.0        # 按方向判定时的绕行距离上限
    initial_locate: bool = True          # **开场先按期望场行动一次**
    initial_locate_max_m: float = 1400.0  # 只对"首条方位距离 (*@\ensuremath{\le}@*) 此值"的频道做开场定点
    batch_locate: bool = False           # 多个"只有一条方位"的频道**合并规划**第二测点
                                         # （实测每局触发约 2 次、每次确实更省，但会把
                                         #   机器狗拉离环线，整局行程反而 +12%，默认关）
    batch_max_m: float = 700.0           # 聚类阈值：彼此最优测点相距 (*@\ensuremath{\le}@*) 此值才算"可以合并"
    batch_min_channels: int = 2          # 至少这么多条单方位才值得合并规划
    batch_max_rounds: int = 6            # 一局里最多做几次"合并定位"
    batch_every: int = 2                 # 每走这么多站就考虑一次合并定位
    batch_min_gain_s: float = 60.0       # 联合方案至少要省这么多秒才值得做
    bear_cone_deg: float = 0.0           # >0 才启用"单方位频道按行进方向就地补测"
    sweep_all_rounds: int = 1            # 完整走圈的遍数（第 2 遍只补缺口）
    lap_offsets: tuple = (0.0, 0.5)      # 每一圈的角向错开量（单位：站位格）
    measured_revisit_gap_m: float = 400.0  # 已测过该频道但没听到 → 换这么远再测


# --------------------------------------------------------------------------
# 机器狗：覆盖扫圈
# --------------------------------------------------------------------------
class SweepRobot(P3Robot):
    """覆盖扫圈。注意：`SweepConfig` 是 dataclass，其额外参数直接传给 `P3Config` 的
    ``__init__``（会把未知键忽略），所以覆盖参数通过 ``cfg`` 对象设置，见 :meth:`_apply_cfg`。
    """

    def _apply_cfg(self, cfg):
        self.cfg = cfg

    def __init__(self, arena, cfg: SweepConfig | None = None, base=None, log=None):
        super().__init__(arena, cfg or SweepConfig(), base=base, log=log)
        self.cand = np.asarray(lattice_candidates(self.cfg.sweep_rings), float)
        # 每个候选点属于第几圈（0 = 主圈；>0 = 补漏圈，只在主圈走完还有缺口时启用）
        self.cand_level = []
        for k, (_r, n) in enumerate(self.cfg.sweep_rings):
            self.cand_level += [k] * int(n)
        self.cand_level = np.asarray(self.cand_level, int)
        self.ring_n = int(self.cfg.sweep_rings[0][1]) if self.cfg.sweep_rings else 0
        self.ring_r = float(self.cfg.sweep_rings[0][0]) if self.cfg.sweep_rings else 0.0
        self.ring_dir = 0.0        # 0 = 还没出发；±1 = 已定的绕圈方向
        self.ring_last_a = None    # 上一站的极角
        self.ring_laps = 0
        self.lap = 0               # 第几圈（0 起）：每圈把站位错开 lap_offsets[lap] 格
        self._last_dir = None      # 上一段的行进方向（用于"单方位频道是否顺路补测"）
        self.visited = []          # **真实**到过的位置（含清除绕行点），用于排除已停靠点
        self.meas_pts = []         # **真实做过的检测**位置，用于可证明的覆盖度评估
        self._cg = None            # 覆盖度评估用细网格（缓存）
        self.coverage = 0.0
        self.max_uncov_dist = float("inf")
        self.phase = "scan"

    # ---------------- 覆盖度 ----------------
    def _grid(self):
        if self._cg is None:
            step = 25.0
            R = R_ARENA
            xs = np.arange(-R, R + 1e-9, step)
            GX, GY = np.meshgrid(xs, xs)
            m = np.hypot(GX, GY) <= R + 1e-9
            self._cg = (GX[m], GY[m])
        return self._cg

    def note_pos(self, ch=None):
        """记录一次**真实做过的检测**位置，用于覆盖度评估。

        为什么要记"检测位置"而不是"到过的位置"：清除阶段会把机器人带到源附近的
        试探格点，那些点**确实测过**（我们在那儿测到方位、也听过无信号），所以算有效；
        但"到过"并不等于"测过每个频道"，所以覆盖度必须按**检测记录**算才是可证明的：
        "每个靶区点到某个检测点的距离 (*@\ensuremath{\le}@*) 1000 m" (*@\ensuremath{\Rightarrow}@*) 该点若有源且当时在测它，必已被听到。
        """
        p = (float(self.arena.pos[0]), float(self.arena.pos[1]))
        if not self.meas_pts or math.dist(self.meas_pts[-1], p) > 1.0:
            self.meas_pts.append(p)

    def coverage_now(self):
        """覆盖率 + 最大未覆盖点距最近**检测点**的距离（可证明完备的判据）。"""
        gx, gy = self._grid()
        pts = list(self.meas_pts) or list(self.scan_points)
        if not pts:
            return 0.0, float("inf")
        P = np.asarray(pts, float)
        d2 = ((gx[:, None] - P[None, :, 0]) ** 2 + (gy[:, None] - P[None, :, 1]) ** 2)
        dmin = np.sqrt(np.min(d2, axis=1))
        unc = dmin > float(self.cfg.cover_radius_m)
        self.max_uncov_dist = float(np.max(dmin))
        self.coverage = float(1.0 - unc.mean())
        return self.coverage, self.max_uncov_dist

    # ---------------- 选点（沿圈单调推进，不回头） ----------------
    def _cand_xy(self, idx):
        """第 ``lap`` 圈、第 ``idx`` 个候选点的坐标（按 ``lap_offsets`` 旋转）。"""
        offs = self.cfg.lap_offsets or (0.0,)
        frac = float(offs[min(self.lap, len(offs) - 1)])
        c = self.cand[idx]
        if abs(frac) < 1e-12:
            return float(c[0]), float(c[1])
        step = 2.0 * math.pi / max(self.ring_n, 1)
        a = math.atan2(float(c[1]), float(c[0])) + frac * step
        r = math.hypot(float(c[0]), float(c[1]))
        return float(r * math.cos(a)), float(r * math.sin(a))

    def pick_stop(self, level=0):
        """选下一站：**沿当前绕圈方向单调推进**，只在"还没覆盖"的站里挑。

        打分 ``score = w_cover·边际覆盖 − w_near·(d/2R) − w_jump·前向角增量(米)``，
        并对"反向站"直接剔除（``adv < min_step``）。这样航线就自然是
        "走一段 → 沿一个方向绕圈 → 走满一圈收工"，不会出现来回横跳或走回头路。

        ``level``：只用第 ``level`` 圈（以及更低层）的候选点。主圈走完仍有缺口时，
        调用方会升到补漏圈。
        """
        gx, gy = self._grid()
        P = np.asarray(self.visited, float) if self.visited else None
        if P is None:
            covered = np.zeros(gx.size, bool)
        else:
            d2 = ((gx[:, None] - P[None, :, 0]) ** 2 + (gy[:, None] - P[None, :, 1]) ** 2)
            covered = np.min(d2, axis=1) <= float(self.cfg.cover_radius_m) ** 2
        p = np.asarray(self.arena.pos, float)
        step = 2.0 * math.pi / max(self.ring_n, 1)
        min_step = 0.35 * step          # 至少要往前走这么多角度
        best, best_score, best_info = None, -1e18, {}
        for idx, c0 in enumerate(self.cand):
            if int(self.cand_level[idx]) > int(level):
                continue
            c = self._cand_xy(idx)
            dc = float(np.hypot(c[0] - p[0], c[1] - p[1]))
            if dc < self.cfg.min_move:
                continue
            if any(float(np.hypot(c[0] - q[0], c[1] - q[1]))
                   < self.cfg.stop_visited_radius_m for q in self.scan_points):
                continue
            in_disk = ((gx - c[0]) ** 2 + (gy - c[1]) ** 2) <= float(self.cfg.cover_radius_m) ** 2
            cnt = int(np.count_nonzero(in_disk))
            if cnt == 0:
                continue
            gain = float(np.count_nonzero(in_disk & ~covered)) / cnt     # [0,1]
            if gain < self.cfg.min_gain:
                continue
            a = math.atan2(float(c[1]), float(c[0]))
            if self.ring_last_a is None:
                adv = 0.0                       # 第一站：方向未定，就近去
            else:
                raw = a - self.ring_last_a
                if self.ring_dir >= 0:          # 未定方向时按"逆时针为正"试探
                    adv = (raw + 2.0 * math.pi) % (2.0 * math.pi)
                else:
                    adv = -(((-raw) + 2.0 * math.pi) % (2.0 * math.pi))
                if adv < min_step:
                    continue                    # 反向或原地：不回头
            near_term = 1.0 - dc / (2.0 * R_ARENA)
            jump = max(0.0, adv) * self.ring_r
            score = (self.cfg.cover_w * gain + self.cfg.near_w * near_term
                     - self.cfg.jump_w * jump)
            if score > best_score:
                best, best_score = c, score
                best_info = {"gain": gain, "dist": dc, "adv_deg": math.degrees(adv),
                             "jump_m": jump, "level": int(self.cand_level[idx]),
                             "lap": self.lap,
                             "n_new": int(np.count_nonzero(in_disk & ~covered)),
                             "cnt": cnt}
        return best, best_info

    def commit_stop(self, c):
        """记录这一站，并据此确定/保持绕圈方向。"""
        a = math.atan2(float(c[1]), float(c[0]))
        step = 2.0 * math.pi / max(self.ring_n, 1)
        if self.ring_last_a is None:
            self.ring_last_a = a
            return
        raw = (a - self.ring_last_a + math.pi) % (2.0 * math.pi) - math.pi
        if self.ring_dir == 0.0:
            # 第一段的方向即定为绕圈方向；角度增量接近 0（小于半格）时沿用正向
            self.ring_dir = 1.0 if raw >= 0 else -1.0
        if (self.ring_dir > 0 and raw < -step / 2.0) or \
           (self.ring_dir < 0 and raw > step / 2.0):
            self.ring_laps += 1        # 说明已经绕回这一圈的起点附近
        self.ring_last_a = a

    # ---------------- 到站动作 ----------------
    # ---------------- 何时可以确定"已经清完了" ----------------
    def all_heard(self):
        """是否每个频道都"给出过明确结论"：已清除，或已测到过方位。

        这是比"覆盖 100%"更早、也更容易达到的**完工判据**：一圈走下来，
        每个频道要么被测到过方位（说明确有源、随后会被清掉），要么一次都没测到。
        配合"整圈已走完"这一事实（整圈 (*@\ensuremath{\Rightarrow}@*) 每个靶区点都距某个检测点 (*@\ensuremath{\le}@*) 1000 m
        (*@\ensuremath{\Rightarrow}@*) 若有源必被听到过），就能断定**没有再多的源了**。
        """
        for rec in self.recs.values():
            if rec.cleared:
                continue
            if rec.n_bearings == 0 and not rec.meas_pts:
                return False
        return True

    def channels_here(self, pos=None, incoming_dir=None):
        """本停点值得测的频道。

        * 还缺方位的频道（未清除且方位数 < 2）一律要测；
        * **只有一条方位、且那条第方位大致指向"我们的前进方向"** 的频道也测
          —— 从当前位置看它、和第一次看它的方向差了接近 90°，交会角好，
          拿到第二条方位就能当场定位并顺路清掉（对应"不要事后走回头路"）；
        * "以前测过但没听到"的频道，若本点到它所有历史检测点都 (*@\ensuremath{\ge}@*)
          ``measured_revisit_gap_m`` / ``repeat_gap_m``，也再测一次。

        注意用的是**真实位置** ``pos``（清除绕行之后的位置），否则会把"计划坐标"
        当成已测过，从而在新站点漏测大量频道。
        """
        p = np.asarray(self.arena.pos if pos is None else pos, float)
        cos_need = math.cos(math.radians(self.cfg.bear_cone_deg))
        out = []
        for ch in self.needed_channels():
            rec = self.recs[int(ch)]
            if not rec.meas_pts:
                out.append(int(ch))
                continue
            gap = (self.cfg.repeat_gap_m if rec.pts
                   else self.cfg.measured_revisit_gap_m)
            d = min(math.hypot(p[0] - q[0], p[1] - q[1]) for q in rec.meas_pts)
            if d >= gap:
                out.append(int(ch))
                continue
            # 已有 1 条方位但"离历史检测点太近"被 gap 挡住的频道：
            # 若那条方位与行进方向夹角 < bear_cone_deg，就地补第二方位。
            if rec.n_bearings == 1 and incoming_dir is not None:
                (sx, sy), a = rec.pts[-1], rec.svds[-1]
                u = np.array([math.cos(math.radians(a)), math.sin(math.radians(a))])
                if float(np.asarray(incoming_dir, float) @ u) >= cos_need:
                    out.append(int(ch))
        return out

    # ---------------- 边绕边清：候选目标（含"只有一条方位"的推算位置）----------------
    def _predicted_source(self, rec):
        """给"只有一条方位"的频道一个**位置估计**，好让它在顺路判定里参与竞争。

        做法：把那条视线与"以原点为心、半径 ``R_RECV_MAX`` 的圆"求交，取较远的交点。
        这个估计误差最大约 ±R_RECV_MAX·tan1° (*@\ensuremath{\approx}@*) 26 m 的横向、以及沿视线的未知距离，
        但用来判断"这个源是不是在我前方、值不值得顺路去清"已经足够。
        """
        (sx, sy), a = rec.pts[-1], rec.svds[-1]
        u = np.array([math.cos(math.radians(a)), math.sin(math.radians(a))])
        S = np.array([float(sx), float(sy)])
        b = float(S @ u)
        c = float(S @ S) - (R_ARENA * R_ARENA)
        disc = b * b - c
        t = (-b + math.sqrt(disc)) if disc > 0 else R_RECV_MAX
        t = min(t, R_RECV_MAX * 1.0)
        return S + t * u

    def _onlap_candidates(self, x, y, travel_dir):
        """本段"顺路就该处理"的候选，按"越正前方、越近"排序。

        返回 ``[(kind, ch, point, cos, dist), ...]``，``kind`` (*@\ensuremath{\in}@*) {``'clear'``, ``'locate'``}：

        * ``'clear'``：已经定位（(*@\ensuremath{\ge}@*)2 条方位且有界）——直接去清；
        * ``'locate'``：**只有一条方位**——先按推算位置过去，在那儿补一条方位把它定住再清。
          这正是用户指出的缺口：旧实现只挑"已定位"的目标，于是前方那些"只测到一次"的
          点全被跳过，最后只能绕完圈再回头找。
        """
        out = []
        p = np.asarray((x, y), float)
        for (ch, rec, br) in self.clear_targets():
            c = np.asarray(br["center"], float)
            v, d = c - p, float(np.linalg.norm(c - p))
            if d > self.cfg.onlap_max_m or d < self.cfg.onlap_min_m:
                continue
            cs = 1.0 if d < 1e-9 else float((v / d) @ travel_dir)
            if cs >= self.cfg.onlap_cos:
                out.append(("clear", ch, c, cs, d))
        if self.cfg.onlap_single:
            for ch, rec in self.recs.items():
                if rec.cleared or rec.n_bearings != 1:
                    continue
                if rec.clear_fail > 0:      # 刚失败过，先换别的
                    continue
                c = self._predicted_source(rec)
                v, d = c - p, float(np.linalg.norm(c - p))
                if d > self.cfg.onlap_max_m or d < self.cfg.onlap_min_m:
                    continue
                cs = 1.0 if d < 1e-9 else float((v / d) @ travel_dir)
                if cs >= self.cfg.onlap_cos:
                    out.append(("locate", ch, c, cs, d))
        out.sort(key=lambda t: (-t[3], t[4]))
        return out

    def _bearing_step(self, rec, p):
        """**小步定位点**：从当前位置沿"与旧视线成 (*@\ensuremath{\approx}@*) onlap_angle_deg"的方向走一小步。

        * 交会角就是这个夹角（几何上最优是 90°，但 50–60° 已经把横向误差压到 1/sin 倍，
          再大就要多走路）；
        * 步长取 ``min(onlap_step_m, frac × 到目标推算位置的距离)`` 并有下限与上限。
        旧实现直接跑到"推算位置"（距离上限 R_RECV_MAX=1500 m！），一次侧移就 1400–1600 m，
        这正是路径图里那些长支线的来源。
        """
        (sx, sy), a = rec.pts[-1], rec.svds[-1]
        u = np.array([math.cos(math.radians(a)), math.sin(math.radians(a))])
        p = np.asarray(p, float)
        # 目标大致在视线上、距离未知：用"当前位置沿视线方向的投影"当作尺度
        t_hat = float((p - np.array([sx, sy])) @ u)
        t_hat = min(max(t_hat, 200.0), R_RECV_MAX)
        ang = math.radians(self.cfg.onlap_angle_deg)
        best = None
        for sgn in (+1.0, -1.0):
            perm = np.array([u[0] * math.cos(sgn * ang) - u[1] * math.sin(sgn * ang),
                             u[0] * math.sin(sgn * ang) + u[1] * math.cos(sgn * ang)])
            step = min(self.cfg.onlap_step_m,
                       max(self.cfg.onlap_min_step_m, self.cfg.onlap_step_frac * t_hat))
            q = p + perm * step
            if np.hypot(*q) > R_ARENA - 1.0:
                continue
            d = float(np.linalg.norm(q - p))
            if best is None or d < best[0]:
                best = (d, q)
        return (best[1] if best else None), (best[0] if best else None)

    def _locate_and_clear(self, ch, predicted):
        """顺路把一个"只有一条方位"的频道定位并清掉：**小步侧移 → 量 → 再清**。

        与旧实现的区别：旧版走到"推算位置"（可能 1500 m 外）再量，量完还常常不在
        清除半径内，于是又要侧移/试探，来回好几趟；新版每次只走一小步、
        量一次就把区域压小一半以上，通常 1 步就能进清除流程。
        """
        rec = self.recs[int(ch)]
        n = 0
        for _step in range(max(1, self.cfg.onlap_max_steps)):
            if rec.cleared or self.aborted:
                break
            br = self.bounded_region(rec)
            if br is not None and br["radius"] <= self.cfg.clear_max_radius_m:
                self.phase = "clear"
                n += 1 if self.clear_target(int(ch), rec) else 0
                self.phase = "scan"
                if rec.cleared:
                    return n
            q, trav = self._bearing_step(rec, self.arena.pos)
            if q is None or trav is None:
                break
            if not self.can_afford(trav / V_ROBOT + T_MEASURE + T_SWITCH + 10.0):
                break
            self.log(f"  顺路定位频道{ch}：侧移 {trav:.0f} m 到 "
                     f"({q[0]:.0f},{q[1]:.0f}) 取方位"
                     f"（步长上限 {self.cfg.onlap_step_m:.0f} m，交会角 "
                     f"(*@\ensuremath{\approx}@*){self.cfg.onlap_angle_deg:.0f}°）")
            resp = self.measure(float(q[0]), float(q[1]), int(ch))
            self.note_pos()
            if resp.get("measure_result") == "near":
                return n + (1 if self.do_clear(float(q[0]), float(q[1]), int(ch)) else 0)
            if resp.get("measure_result") != "direction":
                rec.no_signal += 1
                break
        br = self.bounded_region(rec)
        if br is not None and not rec.cleared:
            self.phase = "clear"
            n += 1 if self.clear_target(int(ch), rec) else 0
            self.phase = "scan"
        return n

    def _onlap_handle(self, kind, ch, pt):
        """顺路处理一个候选：``locate`` 走小步补方位再清，``clear`` 直接清。"""
        rec = self.recs[int(ch)]
        if kind == "locate":
            return self._locate_and_clear(int(ch), pt)
        br = self.bounded_region(rec)
        if br is None:
            return 0
        self.phase = "clear"
        n = 1 if self.clear_target(int(ch), rec) else 0
        self.phase = "scan"
        return n

    # ---------------- 起始站：朝"信息最多的方向"出发 ----------------
    def pick_start_stop(self):
        """把第一站定在"单方位源聚集的那个方向"，并据此确定绕圈方向。

        为什么重要：旧实现的第一站恒为格点里角度最小的那个（正东），于是当所有方位
        都指向西/南时，机器人**背着源**出发 —— 一整圈里都要绕回来（用户看到的
        "绕了两圈"里就有这一段）。原点的 20 次检测已经把每个源的方向给出来了，
        用它来定起始方向是"免费的"。
        """
        if not self.cfg.start_toward_sources:
            return None, None
        vecs = []
        for rec in self.recs.values():
            if rec.cleared or rec.n_bearings != 1:
                continue
            p = self._predicted_source(rec)
            d = float(np.hypot(*p))
            if d < 1e-6:
                continue
            vecs.append(p / d)
        if len(vecs) < self.cfg.start_min_sources:
            return None, None
        mean = np.sum(np.asarray(vecs, float), axis=0)
        if float(np.linalg.norm(mean)) < 1e-6:
            return None, None
        a_mean = math.atan2(float(mean[1]), float(mean[0]))
        # 在格点里挑与"信息方向"最接近的一站
        offs = self.cfg.lap_offsets or (0.0,)
        frac = float(offs[0])
        step = 2.0 * math.pi / max(self.ring_n, 1)
        best, best_d = None, None
        for idx in range(len(self.cand)):
            ang = math.atan2(float(self.cand[idx][1]), float(self.cand[idx][0])) + frac * step
            diff = abs((ang - a_mean + math.pi) % (2 * math.pi) - math.pi)
            if best_d is None or diff < best_d:
                best, best_d = idx, diff
        c = self._cand_xy(best)
        # 绕圈方向：选"下一站"更靠近信息方向的那个方向
        nxt = []
        for sgn in (+1, -1):
            j = (best + sgn) % len(self.cand)
            cj = self._cand_xy(j)
            nxt.append((abs((math.atan2(cj[1], cj[0]) - a_mean + math.pi)
                            % (2 * math.pi) - math.pi), sgn))
        self.ring_dir = 1.0 if nxt[0][0] <= nxt[1][0] else -1.0
        # 注意：**不要**在这里设置 ring_last_a —— 第一站的角度要由 commit_stop 记录，
        # 否则"单调前进"的约束会从"信息方向"起算，导致真正的第一站被当成"往回走"而全部淘汰。
        self.log(f"起始方向：{len(vecs)} 个频道只有单条方位，其平均方向 "
                 f"{math.degrees(a_mean):.0f}° -> 第一站定在 "
                 f"({c[0]:.0f},{c[1]:.0f})，绕圈方向 "
                 f"{'逆时针' if self.ring_dir > 0 else '顺时针'}")
        return c, best

    def _target_near_station(self, c):
        """如果有"已定位/可定位目标"的位置就在计划站位 ``c`` 附近，返回 ``(ch, center)``。

        用途：**用目标位置顶替环上站位**。清一个离站位 200 m 的源时，机器狗已经站在
        那一带并做过检测 —— 让这次检测顶替"专门跑到站位再测一次"，只要覆盖度仍然够，
        就省下一个路径点（用户指出的"两个点走一遍"）。
        """
        best = None
        cands = []
        for (ch, rec, br) in self.clear_targets():
            cands.append((ch, np.asarray(br["center"], float)))
        for ch, rec in self.recs.items():
            if rec.cleared or rec.n_bearings != 1 or rec.clear_fail > 0:
                continue
            cands.append((ch, self._predicted_source(rec)))
        for ch, ctr in cands:
            d = float(np.linalg.norm(ctr - np.asarray(c, float)))
            if d <= self.cfg.skip_station_near_target_m:
                if best is None or d < best[2]:
                    best = (ch, ctr, d)
        return best

    def station_skippable(self, c, alt_pts):
        """若用 ``alt_pts``（顺路已经会去的点）顶替站位 ``c`` 后，覆盖度仍然够，则返回 True。

        判据：去掉 ``c``、加上 ``alt_pts`` 后，"靶区内最大未覆盖距离"仍 <
        ``skip_station_max_uncovered``（比严格的 1000 m 留一点余地，因为这个判据在
        过程中会反复评估，留余量可以避免把后面的补测机会也一起砍掉）。
        """
        if not self.cfg.onlap_clear:
            return False
        try:
            cov_rep, _ = self.coverage_now()
        except Exception:
            return False
        keep = list(self.scan_points) + [(float(p[0]), float(p[1])) for p in alt_pts]
        keep = [q for q in keep if float(np.hypot(q[0] - c[0], q[1] - c[1]))
                > self.cfg.stop_visited_radius_m]
        if not keep:
            return False
        gx, gy = self._grid()
        P = np.asarray(keep, float)
        d2 = ((gx[:, None] - P[None, :, 0]) ** 2 + (gy[:, None] - P[None, :, 1]) ** 2)
        dmin = np.sqrt(np.min(d2, axis=1))
        if float(np.max(dmin)) > self.cfg.skip_station_max_uncovered:
            return False
        return True

    def _enroute_target(self, x, y, travel_dir):
        """挑一个"顺路"的目标：位于本段行进方向的前方（夹角 cos (*@\ensuremath{\ge}@*) enroute_min_cos）。

        为什么不能只看绝对距离：站与站之间相距 842–2000 m，目标常常"就在前方 600–900 m"
        却离本站超过 350 m —— 旧判据直接跳过，最后只能绕完一圈再回头（就是"几乎转了两圈"
        的来源）。改成按**方向**判定后，这类目标会当场清掉，回程自然消失。
        """
        best, best_key = None, None
        for (ch, rec, br) in self.clear_targets():
            c = np.asarray(br["center"], float)
            v = c - np.asarray((x, y), float)
            d = float(np.linalg.norm(v))
            if d > self.cfg.enroute_max_m:
                continue
            cos = 1.0 if d < 1e-9 else float((v / d) @ travel_dir)
            if cos < self.cfg.enroute_min_cos:
                continue
            key = (-cos, d)       # 越"正前方"越好；同方向取近的
            if best_key is None or key < best_key:
                best, best_key = (ch, rec, br), key
        return best

    def visit_stop(self, x, y, travel_dir=None, incoming_dir=None, next_dir=None):
        """移动到 (x,y) 并测量。

        ``travel_dir``：**本段**（上一站→本站）的行进方向；
        ``next_dir``：**下一段**（本站→下一站）的计划方向 —— 顺路判定要用它，
        因为到达本站后"朝本站的方向"已经是零向量（机器狗就站在站上）。

        **清除时机很关键**：第一版每站都调 `clear_phase()`，结果为了清一个源跑到
        靶区另一头，回来时"绕圈"已被打断、剩下的站点全被跳过（实测只清 6/10）。
        现在的 `detour` 模式按**方向**挑顺路目标（见 :meth:`_enroute_target`）：
        在前方就先清掉，背对行进方向的才留到一圈之后统一处理。
        """
        n = 0
        do_detour = (self.cfg.clear_mode == "detour" and not self.aborted)
        if self.cfg.clear_mode == "station" and not self.aborted:
            targets = self.clear_targets()
            if targets and self.can_afford(self.cfg.min_action_budget_s):
                self.phase = "clear"
                n += self.clear_phase()
                self.phase = "scan"
                self.note_pos()          # 清除绕行点也算一次有效检测点
        # ---- 先测本点（保证"这一站真的被扫过"），再顺路清除 ----
        if not self.aborted:
            p = np.asarray(self.arena.pos, float)
            if not self.meas_pts or math.dist(self.meas_pts[-1], tuple(p)) > 1.0:
                self.meas_pts.append((float(p[0]), float(p[1])))
            chans = self.channels_here(pos=p, incoming_dir=incoming_dir)
            if not chans:
                chans = [int(c) for c in self.needed_channels()]
            if chans:
                travel_cost = math.hypot(x - self.arena.pos[0],
                                         y - self.arena.pos[1]) / V_ROBOT
                avail = self.remaining() - self.cfg.time_reserve_s - travel_cost - 4.0
                kmax = max(1, int(avail // (T_MEASURE + T_SWITCH)))
                if len(chans) > kmax:
                    self.log(f"  剩余时间只够测 {kmax}/{len(chans)} 个频道")
                    chans = chans[:kmax]
                n += self.scan_at(float(x), float(y), chans)
                self.note_pos()

        # ---- 边绕边清：本点测完后，把"前方该清的"清掉（含只有一条方位的）----
        if self.cfg.debug_onlap:
            self.log(f"    [onlap] 本站 ({x:.0f},{y:.0f}) 当前 pos="
                     f"({self.arena.pos[0]:.0f},{self.arena.pos[1]:.0f}) "
                     f"onlap_clear={self.cfg.onlap_clear} aborted={self.aborted}")
        if self.cfg.onlap_clear and not self.aborted:
            guard = 0
            while guard < 4 and not self.aborted:
                guard += 1
                d_now = None
                if next_dir is not None:
                    d_now = np.asarray(next_dir, float)      # 用"计划前往的下一站"方向
                else:
                    c0 = np.asarray(self.arena.pos, float)
                    v = np.asarray((x, y), float) - c0
                    dv = float(np.linalg.norm(v))
                    if dv > 1e-6:
                        d_now = v / dv
                if d_now is None:
                    break
                cands = self._onlap_candidates(float(x), float(y), d_now)
                if self.cfg.debug_onlap:
                    self.log(f"    [onlap] 本站 ({x:.0f},{y:.0f}) 行进方向 "
                             f"{math.degrees(math.atan2(d_now[1], d_now[0])):.0f}°，"
                             f"候选 {[(k, ch, round(cs, 2), round(d)) for k, ch, _p, cs, d in cands[:3]]}")
                if not cands:
                    break
                kind, ch, pt, cs, d = cands[0]
                if not self.can_afford(self.cfg.min_action_budget_s + 20.0):
                    break
                n += self._onlap_handle(kind, ch, pt)
                self.note_pos()

        # ---- 兼容旧行为：按方向/距离的"顺路清除" ----
        if do_detour and not self.aborted:
            tgt = self._enroute_target(x, y, travel_dir) if travel_dir is not None else None
            if tgt is None:               # 方向判据没命中时，保留"很近就清"的兜底
                for t in self.clear_targets():
                    if math.dist((x, y), t[2]["center"]) <= self.cfg.clear_detour_max_m:
                        tgt = t
                        break
            if tgt is not None and self.can_afford(self.cfg.min_action_budget_s):
                self.phase = "clear"
                self.log(f"  顺路清除频道{tgt[0]}（中心距本站 "
                         f"{math.dist((x, y), tgt[2]['center']):.0f} m）")
                n += self.clear_target(int(tgt[0]), tgt[1])
                self.phase = "scan"
                self.note_pos()
        return n

    # ---------------- 开场 / 批量：按期望场行动 ----------------
    def _singletons(self, max_r=None):
        """"只有一条方位"且未清除的频道（可选限制首条方位距离）。"""
        out = []
        for ch, rec in self.recs.items():
            if rec.cleared or rec.n_bearings != 1:
                continue
            if max_r is not None and math.hypot(rec.pts[0][0], rec.pts[0][1]) > max_r:
                continue
            out.append(int(ch))
        return out

    def _field_batch_locate(self, reason=""):
        """**按期望场行动**：把当前所有"只有一条方位"的频道**合并**规划第二测点。

        关键认识：**一个测点可以同时给多个源提供好的交会几何**。问题 2 的结论是
        "最优第二检测点在斜前方约 1.2–1.3 km、与源方向成 ±30° 的两翼"，而这些"两翼"
        对不同方向的源常常落在同一片区域。所以：

        1. 对每个待定频道求**局部最优第二测点**（`best_refine_point`，代价里含
           "测完还要往回走多远"）；
        2. **聚类**：彼此最优测点相距 (*@\ensuremath{\le}@*) ``batch_max_m`` 的频道算一批；
        3. 每批用一个**共同测点**（各成员最优点的中点/本身，取联合代价最小者）；
        4. 去那儿按批测一遍，然后立刻清除已经定下来的目标。

        这样一次移动能定住好几个源，而不是"为一个源跑一趟、再为一个源跑一趟"。
        """
        if not self.cfg.batch_locate or self.base is None or self.aborted:
            return 0
        if self.stats.get("n_batch", 0) >= self.cfg.batch_max_rounds:
            return 0
        pend = (self._singletons(max_r=self.cfg.initial_locate_max_m)
                if reason == "开场" else self._singletons())
        if len(pend) < self.cfg.batch_min_channels:
            return 0

        per = {}
        for ch in pend:
            rec = self.recs[ch]
            p, info = best_refine_point(rec.pts, rec.svds, self.pos, None, n_t=60,
                                        radii=(150.0, 300.0, 500.0, 800.0),
                                        travel_cap_m=self.cfg.refine_travel_cap_m)
            if p is None:
                continue
            per[ch] = {"p": np.asarray(p, float), "info": info, "E": None}
        if len(per) < self.cfg.batch_min_channels:
            return 0

        # (*@\textcircled{2}@*) 聚类：**完全连接**（与组内每个成员都要够近），避免"链式"把不相干的点串成一批
        groups = []
        for ch, d in per.items():
            for g in groups:
                if all(float(np.linalg.norm(d["p"] - per[c]["p"])) <= self.cfg.batch_max_m
                       for c in g):
                    g.append(ch)
                    break
            else:
                groups.append([ch])

        best = None
        for g in groups:
            if len(g) < self.cfg.batch_min_channels:
                continue
            pts = [per[c]["p"] for c in g]
            cand = [np.mean(np.asarray(pts, float), axis=0)] + pts
            bt, bc = None, None
            for c0 in cand:
                tot = float(np.linalg.norm(np.asarray(c0) - np.asarray(self.pos))) / V_ROBOT
                tot += len(g) * (T_MEASURE + T_SWITCH)
                bad = False
                for ch in g:
                    rec = self.recs[ch]
                    S, th = rec.pts[-1], rec.svds[-1]
                    ts, w, _hi = ray_source_samples(S, th, n_t=60)
                    if ts.size == 0:
                        bad = True
                        break
                    E = float(exact_expected_diameter(S, th, np.asarray(c0, float),
                                                      ts, w)[0])
                    if not math.isfinite(E):
                        bad = True
                        break
                    tot += math.pi * max(E / 2.0 + 20.0, 20.0) ** 2 / (2.0 * R_CLEAR) / V_ROBOT
                if bad:
                    continue
                if bt is None or tot < bt:
                    bt, bc = tot, np.asarray(c0, float)
            if bc is None:
                continue
            gain = sum(per[c]["info"]["cost_s"] for c in g) - bt
            if best is None or gain > best[0]:
                best = (gain, bc, g, bt)
        if best is None:
            return 0
        gain, c, group, cost = best
        if gain <= self.cfg.batch_min_gain_s:
            return 0
        trav = float(np.linalg.norm(c - np.asarray(self.pos)))
        if not self.can_afford(trav / V_ROBOT + len(group) * (T_MEASURE + T_SWITCH) + 10.0):
            return 0
        self.stats["n_batch"] = self.stats.get("n_batch", 0) + 1
        self.log(f"定位批次{'（' + reason + '）' if reason else ''}：{len(group)} 个频道只有一条方位 "
                 f"{group} -> 合并测点 ({c[0]:.0f},{c[1]:.0f})，移动 {trav:.0f} m，"
                 f"预计比各测各的省 {gain:.0f} s")
        self.scan_points.append((float(c[0]), float(c[1])))
        n = 0
        for ch in group:
            if self.aborted or not self.can_afford(T_MEASURE + T_SWITCH + 3.0):
                break
            resp = self.measure(float(c[0]), float(c[1]), int(ch))
            n += 1
            if resp.get("measure_result") == "near":
                self.do_clear(float(c[0]), float(c[1]), int(ch))
        self.note_pos()
        if self.clear_targets() and self.can_afford(self.cfg.min_action_budget_s):
            self.phase = "clear"
            self.clear_phase()
            self.phase = "scan"
            self.note_pos()
        return n

    def _initial_locate(self):
        """开场定点：委托给 :meth:`_field_batch_locate`（它自带合并规划）。"""
        if not self.cfg.initial_locate or self.base is None or self.aborted:
            return 0
        return self._field_batch_locate(reason="开场")

    # ---------------- 主循环 ----------------
    def run(self):
        t0 = time.time()
        resp = self.arena.enter()
        if resp.get("accepted") is not True:
            raise ArenaError(f"/enter 未被接受：{resp}")
        self.log(f"进入靶区：可用预算 {self.remaining():.0f} s"
                 f"（{self.arena.budget_kind()}）；策略=覆盖扫圈")
        try:
            # 第 0 轮：原点（已经覆盖了 r<=1000 m 的整个内圈）
            chans0 = self.channels_to_measure()
            self.log(f"第 0 轮扫描（原地）：待测 {len(chans0)} 个频道 {chans0}")
            self.scan_points.append(self.pos)
            self.scan_at(self.pos[0], self.pos[1], chans0)
            self.note_pos()
            cov, mdist = self.coverage_now()
            self.log(f"  初始覆盖率 {cov:.3f}（原点一次检测即覆盖 r<=1000 m 的内圈）")

            # ---- 开场定点（方案 1）：先用决策场把近处那一簇源定位并清掉 ----
            self._initial_locate()
            cov, mdist = self.coverage_now()
            self.log(f"  开场定点结束：覆盖率 {cov:.3f}，"
                     f"已清除 {sum(1 for r in self.recs.values() if r.cleared)} 个源")

            rounds = 1
            level = 0
            laps_at_level = 0
            done = False
            first, _first_idx = self.pick_start_stop()   # 朝"信息最多的方向"出发
            while (not self.aborted and rounds < self.cfg.max_rounds
                   and self.remaining() > self.cfg.time_reserve_s):
                is_first = first is not None
                target_hit = None
                if is_first:
                    c = first
                    first = None
                    info = {"dist": float(np.hypot(float(c[0]) - self.arena.pos[0],
                                                   float(c[1]) - self.arena.pos[1])),
                            "adv_deg": 0.0, "gain": 1.0, "n_new": 0, "cnt": 0}
                else:
                    cov, mdist = self.coverage_now()
                    if mdist <= self.cfg.cover_radius_m:
                        self.log(f"* 覆盖率 100%（最大未覆盖点距最近检测点 {mdist:.0f} m "
                                 f"<= {self.cfg.cover_radius_m:.0f} m）"
                                 f"-> 靶区内不可能存在未清除源，提前收工")
                        break
                    c, info = self.pick_stop(level=level)
                if c is None:
                    # 这一圈走完了（该圈所有站位要么已到过、要么没有新覆盖）
                    if level + 1 < len(self.cfg.sweep_rings):
                        level += 1
                        laps_at_level = 0
                        self.ring_last_a = None      # 换圈后重新确定绕行方向
                        self.ring_dir = -self.ring_dir if self.ring_dir else 1.0
                        self.log(f"主圈走完仍有缺口（覆盖率 {cov:.3f}）"
                                 f"-> 启用补漏圈 level={level}"
                                 f"（r={self.cfg.sweep_rings[level][0]:.0f} m）")
                        continue
                    # 所有圈都走过了：换下一圈（角向错开 lap_offsets），再走一遍
                    if (self.lap + 1) < max(2, len(self.cfg.lap_offsets)):
                        self.lap += 1
                        level = 0
                        laps_at_level = 0
                        self.ring_last_a = None
                        self.log(f"所有环的站位都已到过，但仍有频道从未测到信号 "
                                 f"-> 第 {self.lap + 1} 圈（角向错开 "
                                 f"{self.cfg.lap_offsets[min(self.lap, len(self.cfg.lap_offsets) - 1)]} 格）")
                        continue
                    targets = self.clear_targets()
                    if targets and self.can_afford(self.cfg.min_action_budget_s):
                        self.log(f"扫圈完成但仍剩 {len(targets)} 个已定位未清除的源 -> 清除")
                        self.clear_phase()
                        continue
                    self.log(f"没有可选的停靠点、也没有待清除的源 -> 结束"
                             f"（覆盖率 {cov:.3f}，最大未覆盖 {mdist:.0f} m）")
                    break
                cost = (info["dist"] / V_ROBOT + T_MEASURE + T_SWITCH
                        + self.cfg.min_action_budget_s)
                if not self.can_afford(cost):
                    self.log(f"剩余时间不足以再走一站（还需约 {cost:.0f} s）-> 收工")
                    break
                cov, mdist = self.coverage_now()
                self.log(f"扫圈第 {rounds} 站(lv{level})：选点 ({c[0]:.0f},{c[1]:.0f})｜"
                         f"移动 {info['dist']:.0f} m｜前进 {info['adv_deg']:.0f}°｜"
                         f"本点新覆盖 {info['n_new']}/{info['cnt']} 格（收益 {info['gain']:.2f}）｜"
                         f"当前覆盖率 {cov:.3f}（最大未覆盖 {mdist:.0f} m）")
                # ---- 用"目标位置"顶替环上站位（省一个路径点）----
                if self.cfg.onlap_clear and info.get("gain", 1.0) < 0.5:
                    hit = self._target_near_station(c)
                    if hit is not None and self.station_skippable(c, [hit[1]]):
                        self.scan_points.append((float(c[0]), float(c[1])))
                        self.log(f"  * 站位 ({c[0]:.0f},{c[1]:.0f}) 与频道{hit[0]} 的目标位置只差 "
                                 f"{hit[2]:.0f} m，且顶替后覆盖仍够（最大未覆盖 < "
                                 f"{self.cfg.skip_station_max_uncovered:.0f} m）-> 省掉这一站，"
                                 f"直接去清频道{hit[0]}")
                        self._onlap_handle("locate" if self.recs[int(hit[0])].n_bearings < 2
                                           else "clear", int(hit[0]), hit[1])
                        self.note_pos()
                        rounds += 1
                        continue
                p0 = np.asarray(self.arena.pos, float)
                v = np.asarray(c, float) - p0
                dv = float(np.linalg.norm(v))
                travel_dir = (v / dv) if dv > 1e-9 else None      # 本段计划行进方向
                incoming_dir = self._last_dir                       # 上一段的行进方向
                # 下一段的方向：先"试推"一次 pick_stop（只读），用它来判定"前方该清的点"
                next_dir = None
                try:
                    _saved_last, _saved_dir = self.ring_last_a, self.ring_dir
                    _saved_pts = list(self.scan_points)
                    if not is_first:
                        self.scan_points.append((float(c[0]), float(c[1])))
                    _c2, _ = self.pick_stop(level=0)
                    if _c2 is not None:
                        v2 = np.asarray(_c2, float) - np.asarray(c, float)
                        d2 = float(np.linalg.norm(v2))
                        if d2 > 1e-9:
                            next_dir = v2 / d2
                except Exception:
                    next_dir = None
                finally:
                    self.ring_last_a, self.ring_dir = _saved_last, _saved_dir
                    self.scan_points = _saved_pts
                if not is_first:
                    self.commit_stop(c)
                    self.scan_points.append((float(c[0]), float(c[1])))
                self.visit_stop(float(c[0]), float(c[1]), travel_dir=travel_dir,
                                incoming_dir=incoming_dir, next_dir=next_dir)
                self._last_dir = travel_dir
                # 起始站在访问之后才登记：`pick_stop` 用 `scan_points` 做"已到过"的排除，
                # 若提前登记，第一站的邻域会立刻生效，下一站就选不出来了（零收益）
                if is_first:
                    self.commit_stop(c)
                    self.scan_points.append((float(c[0]), float(c[1])))
                # ---- 按期望场行动：每走几站就把"只有一条方位"的频道合并定一次 ----
                if (self.cfg.batch_locate and not self.aborted
                        and rounds % max(1, self.cfg.batch_every) == 0):
                    self._field_batch_locate(reason=f"第 {rounds} 站后")
                rounds += 1
                laps_at_level += 1
                if laps_at_level > 4 * max(self.ring_n, 8) + 8:
                    break

            # ---- 走完圈之后：把所有已定位的源统一清掉（边绕边清模式下通常已清完）----
            if (not self.cfg.onlap_clear or self.cfg.onlap_after_ring) \
                    and not self.aborted and self.clear_targets() \
                    and self.can_afford(self.cfg.min_action_budget_s):
                self.log(f"扫圈结束（覆盖率 {self.coverage:.3f}）：统一清除 "
                         f"{len(self.clear_targets())} 个已定位的源")
                self.phase = "clear"
                self.clear_phase()
                self.phase = "scan"
                self.note_pos()

            # ---- 完工判据：整圈已走完 + 每个频道都有明确结论 -> 直接收工 ----
            if not self.aborted and self.all_heard():
                self.log(f"整圈已走完，且每个频道都有明确结论"
                         f"（已清除 {sum(1 for r in self.recs.values() if r.cleared)} 个 / "
                         f"其余从未测到信号）-> 靶区内无残留源")
                self.stats["n_scan_rounds"] = rounds
                done = True

            # ---- 额外圈：整圈走完但还有频道从未测到（角向缝隙）时再走 -----------------
            while (not done and not self.aborted and rounds < self.cfg.max_rounds
                   and self.remaining() > self.cfg.time_reserve_s):
                cov, mdist = self.coverage_now()
                if self.lap + 1 >= max(2, len(self.cfg.lap_offsets)):
                    break
                c, info = self.pick_stop(level=0)
                if c is None:
                    self.lap += 1
                    self.level = 0
                    self.ring_last_a = None
                    self.log(f"第 {self.lap} 圈走完仍有频道从未测到 -> "
                             f"第 {self.lap + 1} 圈（角向错开 "
                             f"{self.cfg.lap_offsets[min(self.lap, len(self.cfg.lap_offsets) - 1)]} 格）")
                    continue
                cost = (info["dist"] / V_ROBOT + T_MEASURE + T_SWITCH
                        + self.cfg.min_action_budget_s)
                if not self.can_afford(cost):
                    break
                self.log(f"补测圈第 {rounds} 站：({c[0]:.0f},{c[1]:.0f})｜"
                         f"覆盖率 {cov:.3f}（最大未覆盖 {mdist:.0f} m）")
                self.commit_stop(c)
                self.visit_stop(float(c[0]), float(c[1]))
                rounds += 1
                if self.all_heard() and not self.clear_targets():
                    self.log("补测后每个频道都有明确结论、且没有待清除目标 -> 结束")
                    break
            self.stats["n_scan_rounds"] = rounds

            # ---- 收尾：把"听到了但还没清掉"的少量频道用自适应方式解决 ----
            guards = 0
            while (not self.aborted and rounds < self.cfg.max_rounds
                   and self.remaining() > self.cfg.time_reserve_s and guards < 8):
                guards += 1
                leftover = [ch for ch, rec in self.recs.items()
                            if not rec.cleared and rec.n_bearings >= 1]
                if not leftover:
                    break
                # 只关心"有界区域但还没清"的频道；其余先补一条方位
                todo = [t for t in self.clear_targets()]
                if todo:
                    self.log(f"收尾：清除剩余 {len(todo)} 个已定位目标 "
                             f"{[t[0] for t in todo]}")
                    self.phase = "clear"
                    n = self.clear_phase()
                    self.phase = "scan"
                    self.note_pos()
                    if n == 0:
                        break
                    continue
                plan = self.plan_scan()
                if plan is None or not self.can_afford(plan["cost_s"]):
                    self.log(f"收尾：剩余频道 {leftover} 无法在预算内继续定位 -> 结束")
                    break
                self.log(f"收尾：为剩余频道 {leftover} 补测方位 -> "
                         f"选点 {tuple(round(v) for v in plan['sel']['xy'])}")
                self.execute_scan(plan)
                rounds += 1
            self.stats["n_scan_rounds"] = rounds
        finally:
            try:
                self.arena.exit()
            except ArenaError as e:
                self.log(f"/exit 异常：{e}")
        self.stats["program_runtime_s"] = time.time() - t0
        rep = self.report()
        rep["coverage"] = self.coverage
        rep["max_uncovered_dist_m"] = self.max_uncov_dist
        return rep


class HybridRobot(SweepRobot):
    """混合策略：**先用扫圈把该听的都听到，再交给"期望场贪心"去清除**。

    动机（离线实测）：
    * 纯扫圈的**发现**效率明显更好：同样两局在线真源，虚拟时间 3.7–4.0k vs 4.0–4.4k，
      检测次数 117–145 vs 143–198，行程少 3–5 km，而且不会乱蹭、有完备性判据；
    * 但纯扫圈的**清除**偏弱（外圈环带场景 0.972），因为它把清除也绑在"绕圈"结构上。
    * 所以让两者各干自己最擅长的：扫圈负责"把 20 个频道都听到一遍"（发现），
      之后切到自适应循环负责"把听到的源清干净"（决策场选点 + 覆盖式试探 + 顺路补测）。
    """

    def run(self):
        t0 = time.time()
        resp = self.arena.enter()
        if resp.get("accepted") is not True:
            raise ArenaError(f"/enter 未被接受：{resp}")
        self.log(f"进入靶区：可用预算 {self.remaining():.0f} s"
                 f"（{self.arena.budget_kind()}）；策略=扫圈发现 + 自适应清除")
        try:
            chans0 = self.channels_to_measure()
            self.log(f"第 0 轮扫描（原地）：待测 {len(chans0)} 个频道 {chans0}")
            self.scan_points.append(self.pos)
            self.scan_at(self.pos[0], self.pos[1], chans0)
            self.note_pos()

            # ---- 阶段 1：绕一圈，把该听的频道都听一遍（不清除） ----
            self.phase = "scan"
            rounds = 1
            while (not self.aborted and rounds < self.cfg.max_rounds
                   and self.remaining() > self.cfg.time_reserve_s):
                cov, mdist = self.coverage_now()
                c, info = self.pick_stop(level=0)
                if c is None:
                    break
                cost = (info["dist"] / V_ROBOT + T_MEASURE + T_SWITCH
                        + self.cfg.min_action_budget_s)
                if not self.can_afford(cost):
                    break
                self.log(f"扫圈第 {rounds} 站：({c[0]:.0f},{c[1]:.0f})｜"
                         f"移动 {info['dist']:.0f} m｜收益 {info['gain']:.2f}｜"
                         f"覆盖率 {cov:.3f}（最大未覆盖 {mdist:.0f} m）")
                self.commit_stop(c)
                self.scan_points.append((float(c[0]), float(c[1])))
                self.visit_stop(float(c[0]), float(c[1]))
                rounds += 1
                if self.all_heard():
                    self.log("整圈走完，每个频道都有明确结论 -> 结束扫描阶段，"
                             "转交自适应清除")
                    break
            self.stats["n_scan_rounds"] = rounds

            # ---- 阶段 2：切回自适应主循环（决策场选点 + 清理 + 顺路补测） ----
            self.phase = "adaptive"
            stall = 0
            while (not self.aborted and rounds < self.cfg.max_rounds
                   and self.remaining() > self.cfg.time_reserve_s):
                did = 0
                if self.clear_targets() and (rounds >= self.cfg.min_scan_rounds_before_clear
                                             or not self.plan_scan()):
                    did += self.clear_phase()
                if self.aborted or self.remaining() <= self.cfg.time_reserve_s:
                    break
                plan = self.plan_scan()
                if plan is not None and self.can_afford(plan["cost_s"]) \
                        and (did == 0 or not self.clear_targets()):
                    did += self.execute_scan(plan)
                    rounds += 1
                if did == 0:
                    stall += 1
                    if stall >= 2:
                        self.log("自适应阶段：没有可清理的频道、也没有值得检测的频道 "
                                 "-> 提前结束")
                        break
                else:
                    stall = 0
            self.stats["n_scan_rounds"] = rounds
        finally:
            try:
                self.arena.exit()
            except ArenaError as e:
                self.log(f"/exit 异常：{e}")
        self.stats["program_runtime_s"] = time.time() - t0
        rep = self.report()
        rep["coverage"] = self.coverage
        rep["max_uncovered_dist_m"] = self.max_uncov_dist
        return rep


def selfcheck(verbose=True):
    out, ok_all = [], True

    def rec(name, ok, detail=""):
        nonlocal ok_all
        ok_all = ok_all and bool(ok)
        out.append({"check": name, "ok": bool(ok), "detail": detail})
        if verbose:
            print(f"[{'OK ' if ok else 'FAIL'}] {name}  {detail}")

    cand = lattice_candidates()
    mdist, unc, cov = coverage_report(cand, r_inner=R_RECV_MIN)
    rec("C1 一圈格点即可覆盖环带 r(*@\ensuremath{\in}@*)[1000,1800]（最大未覆盖距离 <= 1000 m）",
        mdist <= R_RECV_MIN,
        f"{len(cand)} 个候选点；环带内最大未覆盖距离 {mdist:.0f} m；"
        f"未覆盖面积 {unc*100:.2f}%")

    mdist1, unc1, cov1 = coverage_report([(0.0, 0.0)] + cand, r_inner=0.0)
    rec("C2 原点 + 一圈覆盖整个靶区（这才是实际航线的完备性口径）",
        mdist1 <= R_RECV_MIN,
        f"最大未覆盖 {mdist1:.0f} m；覆盖率 {cov1*100:.2f}%")

    from p3_arena import MockArena
    a = MockArena(seed=20260913)
    rb = SweepRobot(a, SweepConfig(), base=None, log=None)
    sel, info = rb.pick_stop()
    rec("C3 第一站只从还没覆盖的格点里选，且不选已停靠点邻域",
        sel is not None and info["gain"] >= rb.cfg.min_gain,
        f"第一个停靠点 ({sel[0]:.0f},{sel[1]:.0f})，距原点 {info['dist']:.0f} m，"
        f"边际覆盖 {info['gain']:.2f}（新覆盖 {info['n_new']}/{info['cnt']} 格）")

    return {"ok": ok_all, "checks": out}


if __name__ == "__main__":
    import sys
    if "--selfcheck" in sys.argv:
        res = selfcheck()
        raise SystemExit(0 if res["ok"] else 1)
    if "--coverage" in sys.argv:
        for rings in (DEFAULT_RINGS,):
            cand = lattice_candidates(rings)
            mdist, unc, cov = coverage_report(cand)
            print(f"候选点 {len(cand)} 个：最大未覆盖 {mdist:.0f} m，未覆盖 {unc*100:.2f}%")
        print("\n只用原点一个点：", coverage_report([(0.0, 0.0)])[:2])
        print("原点 + 一圈 r=1500 (8 点)：",
              coverage_report([(0.0, 0.0)] + [(1500 * math.cos(2 * math.pi * i / 8),
                                              1500 * math.sin(2 * math.pi * i / 8))
                                              for i in range(8)])[:2])
        print("原点 + 一圈 r=1800 (12 点)：",
              coverage_report([(0.0, 0.0)] + [(1800 * math.cos(2 * math.pi * i / 12),
                                               1800 * math.sin(2 * math.pi * i / 12))
                                              for i in range(12)])[:2])

\end{lstlisting}
\subsection{\texttt{src/p3\_tour.py}}
% SHA256 8ff880cce160ecd96f7c64c88c95e145dff864a0d5b1cdb1d9b2fd947c86e8e5
\begin{lstlisting}[language=python,escapeinside={(*@}{@*)}]
"""B 题问题 3：**统一滚动航路**策略（Tour）——"绕圈"与"清除"合成一条动态航路。

为什么需要它（旧两套策略的实测缺陷，详见 `review/p3-tour-design.md`）
------------------------------------------------------------------
* `P3Robot`（期望场贪心）：虚拟时间 5270 s，其中 **29% 是清完源之后的空转**，
  检测 186 次、行程 20.5 km。
* `SweepRobot`（覆盖扫圈）：虚拟时间降到 4243 s，但结构是**"先绕完一圈、再统一清除"**：
  绕圈的 6.4 km 和清除的 5.2 km 是两段互不重叠的路，实测行程 16.8 km，
  是"已知全部源位后的最优航路"（8.8 km）的 **1.92 倍**，其中折返（转角 >120°）
  1.5 km/局。按题目的"平均定位清除时间 = 定位清除总时间 ÷ 清除个数"口径，
  这一半行程全是白走的。

本模块把扫描与清除合成**一条滚动时域的统一航路**：

```
每个决策点重算一次航路，只执行第一个停点，然后重规划
  航路 = (*@\textcircled{1}@*) 环带骨架停点（保证"没有残留源"的可证明完备性，按角向单调推进）
         ↑ 目标（清除 / 定位）用"最便宜插入"插进骨架的空隙里
  目标 = 已定位频道的最新定位区域中心（clear）
         只有一条方位 / 区域无界 / 区域过大的频道（locate，用第二问的 E[D] 代价模型）
```

于是航线天然就是用户描述的**锯齿状一圈**：从原点出发 →（期望场）选第二个检测点、
顺手清掉靠内圈的目标 → 就近切入环带 → 沿一个方向绕一圈、边走边清 →
走完一圈即收工。不会出现"绕完再回头"的长距离折返。

相对旧实现的六个关键修正
------------------------
1. **不再分"绕圈/清除"两个阶段**：骨架停点与目标一起排序，清完一个就地重规划。
2. **环带骨架按角向单调推进**（锁定方向 + 角向游标）：这是"不折返"的结构性保证；
   旧实现每一步都重新贪心选点，导致航路在角向上来回横跳（实测折返 2.5 km/局）。
3. **覆盖判据按频道算**：某点是否值得测，取决于"这个频道在那里是否还可能有
   没被听到的源"，用该频道**真实测过的位置**（`rec.meas_pts`）判定 ——
   清除绕行时顺带做的检测会自动抵扣掉后面的骨架停点，不重复走。
4. **修掉 `visit_stop` 的 gap 判定 bug**：旧实现用"上一个动作的位置"而不是
   "本停点坐标"判断"这个频道在这儿测过没有"，导致整站漏测（实测有源因为
   连续漏测而**永远没被听到**）。
5. **"只有一条方位"的频道不再无脑重测**：只有当"本停点相对上一条视线的张角"
   落在 25°–155° 时才测（否则交会角接近 0，纯属浪费 6 s）。
6. **清除子过程可迭代**：区域无界 / 区域半径过大不再被永久跳过（旧实现在
   `clear_targets()` 里把 r*>600 m 的目标过滤掉，是清除率只有 0.978 的主因），
   而是"逼近中心补方位 → 区域压小 → 覆盖式试探"循环，直到成功或用尽预算。
"""
from __future__ import annotations

import math
import time
from dataclasses import dataclass

import numpy as np

from p3_arena import (  # noqa: E402
    ArenaError, R_ARENA, R_CLEAR, R_RECV_MIN, T_CLEAR_FAIL, T_MEASURE, T_SWITCH,
    V_ROBOT,
)
from p3_expect_field import best_refine_point, ray_source_samples  # noqa: E402
from p3_sweep import SweepConfig, SweepRobot  # noqa: E402


# --------------------------------------------------------------------------
# 配置
# --------------------------------------------------------------------------
@dataclass
class TourConfig(SweepConfig):
    """在 `SweepConfig`（含 `P3Config` 全部字段）之上追加"统一航路"参数。"""

    # ---- 环带骨架 ----
    cover_rings: tuple = ((1000.0, 8),)
    """骨架候选停点的同心环 ``(半径 m, 点数)``。

    为什么是"小半径 + 稀疏停点"：原点那一次全频道检测已经把 r<=1000 的内圈覆盖了，
    只剩环带 r(*@\ensuremath{\in}@*)[1000,1800] 要覆盖；半径 (*@\ensuremath{\rho}@*) 的停点最多覆盖到
    ``1800² + (*@\ensuremath{\rho}@*)² - 2·1800·(*@\ensuremath{\rho}@*)·cosΔ <= 1000²`` 的角向半宽 Δ（Δ 越大，环上可以越稀），
    于是

    | 配置 | 环带最大未覆盖 | 覆盖余量 | 单圈弦长 | 含进入段 | 每圈要测的频道数 |
    |---|---|---|---|---|---|
    | 7 点 r=1050 | 968 m | 32 m | 6378 m | 7428 m | 7×N |
    | **8 点 r=1000（默认）** | **943 m** | **57 m** | **6123 m** | **7123 m** | 8×N |
    | 9 点 r=950 | 964 m | 36 m | 5849 m | 6799 m | 9×N |
    | 13 点 r=880 | 969 m | 31 m | 5476 m | 6356 m | 13×N |

    半径越小弦长越短，但停点越多、每站都要把"从没听到过"的频道整批测一遍
    （6 s/频道）。四批共 150 局实测（`python src/p3_tour_sweep.py --rings ...`）：

    | 配置 | 全清局数 | T/N 均值 |
    |---|---|---|
    | 7 点 r=1050 | 146/150 | 333.2 s |
    | **8 点 r=1000（默认）** | **149/150** | **332.5 s** |
    | 9 点 r=950 | 149/150 | 333.3 s |

    r=1000/8 点同时拿到最好的清除率、最低的 T/N，以及最大的覆盖余量（57 m），故为默认。
    """
    cover_grid_step_m: float = 25.0     # 覆盖度评估网格步长
    cover_slack_m: float = 0.0          # 覆盖判据余量（>0 更保守）
    ring_min_advance_rad: float = 0.06  # 角向单调推进的最小步长（约 50 m）
    entry_tie_m: float = 0.0
    """入口选择的"并列"阈值：航路长度相差不超过它就算并列，此时才按信息方向挑。

    取 0 的含义是"**优先就近进入**；只有长度**完全相等**（机器狗正停在原点、
    8 个入口一样远）时，才用信息方向打破僵局" —— 这正好对应用户指出的现象：
    进场位置之所以"固定"，就是因为原点扫描之后机器狗还在原点，
    `atan2(0,0)` 恒为 0。阈值放大到 250/400 m 会把"信息方向"的权重放大成
    主导项，实测反而变差（常规随机 327.8 → 327.9 → 332.2）。"""
    entry_select: bool = True           # 消融对照：False = 旧行为（用当前位置极角当游标）
    cover_duty_frac: float = 0.95       # 目标停点半径 >= 骨架半径×它时，代劳覆盖检测

    # ---- 路点合并与排序 ----
    merge_radius_m: float = 250.0       # 目标这么近的骨架停点直接省掉
    route_mode: str = "polish"          # insert（只做最便宜插入）| polish（+2-opt，骨架不倒序）
                                        # | tsp（完全自由 NN+2-opt，消融对照）
    locate_max_detour_m: float = 260.0  # 定位路点的绕行超过它就本轮不排（等航路自己给方位）
    locate_force_after: int = 2         # 连续跳过这么多轮之后强制排入（避免永远排不进）
    clear_max_detour_m: float = 1e18    # **清除路点一律排进航路**（这个上限只是记录用）
    """清除目标绝不允许"因为绕行太远就不排"。

    实测（`p3_paths` 场景 C）：ch6 在环带 (0,-1000) 处就拿到了第二条方位、
    $r^*=42$ m，完全可以就地清掉；但当时"最便宜插入"的额外代价是 2382 m，
    超过了旧的 1500 m 硬上限，于是它被踢出航路，一直留到收尾阶段的
    `clear_leftovers` —— 那时机器人已经跑到场地另一头（1051,1295），
    再横穿 3.4 km 回去清它。**"绕行大"只应该影响排序，绝不能变成"不清"**。
    """

    # ---- 期望场（第二问 E[D] 模型）作为可选项 ----
    field_min_pending: int = 1
    """只有"待定（只有一条方位的）频道数 >= 它"时才动用期望场/批量合并定位。

    期望场的好处是"一次移动定住好几个源"，也**顺带把"走多远拉基线"与
    "定位精度差要多扫多少面积"放在同一个代价里权衡**（`best_refine_point`
    的 cost 里含 $\\pi(E[D]/2+20)^2/(2R_{clear})/v$）。所以它不是"为了精度可以
    无限跑远"—— 只是因为"基线短 (*@\ensuremath{\Rightarrow}@*) $r^*$ 大 (*@\ensuremath{\Rightarrow}@*) 后面搜索面积按平方增长"，
    权衡下来它的最优点常常在几百米到 1 km 之间。

    待定频道很少时改用 `_cheap_locate_point`（候选半径压到
    ``short_locate_cap_m``）这一档**实测是亏的**（四类情景 T/N 全部更差：
    常规随机 330.0→344.4、外圈环带 359.6→371.7、在线真源 394.6→426.3），
    所以默认取 1（= 始终允许期望场）。真正解决"为单个目标跑很远"的是
    `locate_patience_stops`（环形航路本来会给的方位就先等着）与
    `locate_max_detour_m`（绕行太大的定位点本轮不排）。
    """

    # ---- 主动"防过期"补测（默认关闭：实测在 4 类情景上都略亏，见 §"清除时机"）----
    risk_locate: bool = False
    risk_lookahead: int = 3             # 往后看几个骨架停点
    risk_min_lost: float = 0.5          # 先验里有这么大比例的源位"往后几个站都听不到"就该补测
    risk_step_m: float = 300.0          # 补测只走这么远（小步）
    risk_max_per_stop: int = 1
    sector_clear_m: float = 0.0
    """"刚走过的扇区里已经能清的目标"的就地清除半径（0 = 关闭）。

    这两个机制（`risk_locate` / `sector_clear`）是照"每个角度只经过一次、
    所以覆盖到的目标要尽快清"的思路写的，但四类情景实测都是**略亏**
    （常规随机 T/N 334.6 → 327.5、外圈环带 4575 s → 4085 s、行程 16.1 km → 13.4 km
    都是关掉更好）。真正解决"该清没清"的是 `offroute_clear`（见下）。"""

    # ---- "不在行进方向上就地补测并清除" ----
    offroute_clear: bool = True
    wait_angle_min_deg: float = 30.0
    """夹角粗筛：第一条方位与"本站->下一站"方向的夹角小于它就直接等下一站。

    用户的直觉（"夹角小就等、夹角大就立刻走一小步定位清除"）方向完全正确，
    但**只按夹角判太粗**：夹角大其实是常态，而"拐回来清"到底多花多少，
    取决于源离哪一站更近。所以夹角在这里只当粗筛，真正的判据是
    `offroute_clear` 里用射线先验直接算的"两种做法的额外绕行"
    （实测：只用夹角判 → 常规随机 330 → 350、外圈环带 4085 → 5238 s；
    换成绕行比较之后才真正有收益）。
    """
    offroute_step_m: float = 320.0      # 就地补测只走这么远
    offroute_margin_m: float = 60.0     # "现在清"要比"等下一站"至少省这么多米才做
    offroute_max_thi_m: float = 900.0
    """射线先验外径超过它就放弃判断。

    判断依赖"源大致在哪"，而 ``t_hi``（射线在靶区内的长度）越长，先验越摊得开、
    估计越不可靠 —— 实测不加这条闸门时，常规随机会从 330 掉到 350 s
    （对着一堆"其实在 1.4 km 外"的源白跑 320 m 的补测步）。
    ``t_hi`` 小（源被靶区边界卡住、不可能远）时判断才可信。
    """
    offroute_max_per_stop: int = 1
    offroute_max_dist_m: float = 450.0
    """源位先验的加权平均距离超过它就放弃"现在清"。

    单条方位对源距几乎没有分辨力（只有射线先验），源越远估计越不可靠；
    而且"现在清"的价值来自"顺手"，源太远时顺手也顺不了。
    """

    # ---- 清除 ----
    probe_enter_r_m: float = 110.0      # r* 小到它就进入覆盖式试探
    probe_max_radius_m: float = 300.0   # 试探圆域半径上限（区域更大就先补方位）
    probe_max_points: int = 150
    clear_max_iters: int = 6
    approach_eps_m: float = 25.0        # 与既有测量点的最小间隔
    fail_cooldown: int = 2              # 清除失败后冷却几个停点
    big_r_m: float = 1200.0             # r* 超过它就不"去区域中心"，改用横向定位点
                                        # （两条近似平行的方位会把最小包围圆撑到几公里，
                                        #   中心甚至落在靶区外 —— 旧实现会真的走过去，
                                        #   实测有一局为此多走 13.5 km、多花 2700 s）

    # ---- 定位路点（"只有一条方位"的频道）----
    locate_mode: str = "stale"          # always | stale | never
    locate_stale_stops: int = 3
    locate_patience_stops: int = 6
    """只有一条方位的频道，最多"等"环形航路几个停点。

    环形航路每个角度只经过一次，所以一条只有单方位的频道**通常会在它所在扇区的
    下一个骨架停点上拿到第二条方位** —— 白等就行，专门跑一趟是纯浪费
    （实测 `p3_paths` 场景 D 里那些"偏离预定路线的尖锐拐角"就是这么来的：
    为了给 ch9 补方位，从 (−1716,58) 往靶区里折了 1020 m，而环形航路本来
    就会在 (1000,0) 给它第二条方位）。
    因此：先按 `_lost_fraction` 判断"环形航路还会不会再听到它"，
    只有确实听不到了、或者已经等了 `locate_patience_stops` 站，才专门去定位。
    """
    locate_travel_cap_m: float = 1100.0
    locate_bearing_cap_m: float = 700.0  # "回到上一条方位的检测点附近横向补测"的半径上限
    locate_radii: tuple = (150.0, 300.0, 500.0, 750.0, 1000.0)
    short_locate_cap_m: float = 420.0
    """`_cheap_locate_point` 的候选半径上限（"单点定位只走一小步"）。"""
    cheap_locate_angles: tuple = (55.0, 80.0, 105.0)  # 小步横向补测的候选夹角（度）
    min_cross_deg: float = 25.0         # 单方位频道：张角小于它就不值得重测

    # ---- 开场定点（期望场第二检测点）----
    initial_rounds: int = 3             # 原点附近最多做几轮"合并定位 + 顺手清除"
    initial_move_growth: float = 2.0    # 每轮把移动上限放宽这么多倍
    initial_locate_max_move_m: float = 350.0   # 单轮定位最多走多远
    """开场定点单轮的最远移动距离。

    第二问的 E[D] 代价模型用的是"源位沿射线、按面积均匀"的先验，其**中位源距约 1 km**，
    于是对"其实就在原点附近"的那些源，它会给出一个 1 km 开外的第二检测点。
    实测（30 局 × 3 种情景）把它从 1300 m 收到 350 m：

    | 情景 | 1300 m | **350 m** |
    |---|---|---|
    | 常规随机（新种子批） | T/N 324.4 | **312.5** |
    | 中心密集 12 源 | T/N 282.1（还漏 1 局） | **224.9（30/30 全清）** |
    | 外圈环带 12 源 | T/N 364.2 | 365.5（打平） |

    原因：近处的源只要在 300–400 m 外测一条方位就能定得很准；跑 1 km 去测既费时，
    又把它后面的清除路线拉乱（"先清较靠近原点的目标"靠的是**顺路**，不是专门跑一趟）。
    """
    initial_clear_radius_m: float = 900.0      # 开场顺手清除的半径（"先清近处的源"）

    # ---- 结束 ----
    max_stops: int = 80
    min_stop_budget_s: float = 15.0


# --------------------------------------------------------------------------
# 机器狗
# --------------------------------------------------------------------------
class TourRobot(SweepRobot):
    """统一滚动航路（环带骨架 + 目标最便宜插入 + 逐站重规划）。"""

    def __init__(self, arena, cfg: TourConfig | None = None, base=None, log=None):
        super().__init__(arena, cfg or TourConfig(), base=base, log=log)
        self._ag = None                 # 环带评估网格缓存
        self._unsafe_cache = {}         # (ch, len(meas_pts)) -> 未排除掩码
        self._ring = None               # 骨架（(角度, 半径) 按角度排序）
        self.ring_dir = 0.0             # 锁定的绕行方向 ±1
        self.ring_cursor_a = None       # 角向游标（上一个骨架停点的极角）
        self.ring_served = set()        # 已经被"顶替并完成覆盖职责"的骨架停点（角度键）
        self.stop_seq = 0
        self.last_stop_xy = None
        self.n_cover_stops = 0
        self.n_locate_stops = 0
        self.n_clear_iter = 0
        self._hold_until = {}
        self._locate_skips = {}
        self._last_bearing_stop = {}
        self._prev_stop_a = None
        self.tour_log = []

    # ------------------------------------------------------------------
    # 台账
    # ------------------------------------------------------------------
    def pending(self):
        return [int(ch) for ch, rec in self.recs.items() if not rec.cleared]

    def unheard(self):
        return [int(ch) for ch, rec in self.recs.items()
                if not rec.cleared and rec.n_bearings == 0]

    def heard_pending(self):
        return [int(ch) for ch, rec in self.recs.items()
                if not rec.cleared and rec.n_bearings > 0]

    def _n_cleared(self):
        return sum(1 for r in self.recs.values() if r.cleared)

    def on_hold(self, ch):
        return self.stop_seq < self._hold_until.get(int(ch), -1)

    def hold(self, ch, stops=None):
        self._hold_until[int(ch)] = self.stop_seq + int(
            self.cfg.fail_cooldown if stops is None else stops)

    def add_bearing(self, ch, x, y, svd):
        reg = super().add_bearing(ch, x, y, svd)
        self._last_bearing_stop[int(ch)] = self.stop_seq
        self._hold_until.pop(int(ch), None)
        return reg

    def region_info(self, rec):
        """当前（或最近一次）有界区域 ``{'center','radius','poly','status'}``。"""
        if rec.n_bearings >= 2:
            reg = self.region(rec)
            if reg.get("status") == "bounded":
                return {"center": tuple(reg["min_enclosing_center"]),
                        "radius": float(reg["min_enclosing_radius_m"]),
                        "poly": [tuple(v) for v in reg.get("vertices", [])],
                        "status": "bounded", "n_bearings": rec.n_bearings}
        if rec.last_bounded is not None:
            c, r, n = rec.last_bounded
            return {"center": tuple(c), "radius": float(r), "poly": None,
                    "status": "fallback", "n_bearings": n}
        return None

    # ------------------------------------------------------------------
    # 覆盖：哪些环带格点"还可能藏着从没听到过的源"
    # ------------------------------------------------------------------
    def annulus_grid(self):
        if self._ag is None:
            step = float(self.cfg.cover_grid_step_m)
            R = R_ARENA
            xs = np.arange(-R, R + 1e-9, step)
            GX, GY = np.meshgrid(xs, xs)
            rr = np.hypot(GX, GY)
            m = (rr <= R + 1e-9) & (rr >= R_RECV_MIN - 1e-9)
            self._ag = (GX[m].copy(), GY[m].copy())
        return self._ag

    def unsafe_of(self, ch):
        """频道 ``ch`` 仍然"没被排除"的环带格点掩码。

        判据：该格点到频道 ``ch`` **真实测过的任一位置**的距离 <= 有效接收半径下界
        （1000 m）(*@\ensuremath{\Rightarrow}@*) 若那里真有源，当时就会被听到 (*@\ensuremath{\Rightarrow}@*) 可以排除。
        """
        rec = self.recs[int(ch)]
        key = (int(ch), len(rec.meas_pts))
        got = self._unsafe_cache.get(key)
        if got is not None:
            return got
        gx, gy = self.annulus_grid()
        radius = float(self.cfg.cover_radius_m) + float(self.cfg.cover_slack_m)
        if not rec.meas_pts:
            m = np.ones(gx.size, bool)
        else:
            P = np.asarray(rec.meas_pts, float)
            d2 = ((gx[:, None] - P[None, :, 0]) ** 2
                  + (gy[:, None] - P[None, :, 1]) ** 2)
            m = np.sqrt(np.min(d2, axis=1)) > radius
        if len(self._unsafe_cache) > 600:
            self._unsafe_cache.clear()
        self._unsafe_cache[key] = m
        return m

    def uncovered_mask(self):
        """所有"从没听到过"的频道的联合未排除掩码。"""
        gx, gy = self.annulus_grid()
        unsafe = np.zeros(gx.size, bool)
        for ch in self.unheard():
            unsafe |= self.unsafe_of(ch)
        return gx, gy, unsafe

    def adds_coverage(self, ch, x, y):
        """在 ``(x, y)`` 测频道 ``ch`` 是否还能排除掉新的环带格点。"""
        m = self.unsafe_of(ch)
        if not m.any():
            return False
        gx, gy = self.annulus_grid()
        r2 = float(self.cfg.cover_radius_m) ** 2
        return bool(np.any(m & (((gx - x) ** 2 + (gy - y) ** 2) <= r2)))

    def is_complete(self):
        """能否断定"靶区内已无未清除源"。"""
        if self.heard_pending():
            return False
        return not bool(self.uncovered_mask()[2].any())

    # ------------------------------------------------------------------
    # 环带骨架
    # ------------------------------------------------------------------
    def ring_skeleton(self):
        """骨架停点 ``[(角度, 半径), ...]``，按角度排序（只构造一次）。"""
        if self._ring is None:
            pts = []
            for k, (r, n) in enumerate(self.cfg.cover_rings):
                off = (math.pi / float(n)) if (k % 2 == 1) else 0.0
                for i in range(int(n)):
                    a = (off + 2.0 * math.pi * i / float(n)) % (2.0 * math.pi)
                    pts.append((float(a), float(r)))
            pts.sort()
            self._ring = pts
        return self._ring

    def ring_xy(self, a, r):
        return (float(r) * math.cos(a), float(r) * math.sin(a))

    def forward_ring(self, dir_sign, cursor=None):
        """角向单调推进的骨架停点（已过去的不再选，方向锁死 (*@\ensuremath{\Rightarrow}@*) 不折返）。"""
        sk = self.ring_skeleton()
        a0 = self.ring_cursor_a if cursor is None else cursor
        if a0 is None:
            a0 = math.atan2(self.pos[1], self.pos[0]) % (2.0 * math.pi)
        out = []
        for (a, r) in sk:
            d = ((a - a0) * dir_sign) % (2.0 * math.pi)
            if d < float(self.cfg.ring_min_advance_rad):
                continue
            out.append((d, a, r))
        out.sort()
        res = []
        for _d, a, r in out:
            if self._akey(a) in self.ring_served:
                continue            # 这个角向站点的覆盖职责已经由某个目标停点代劳
            x, y = self.ring_xy(a, r)
            if not self.adds_any_coverage(x, y):
                continue
            res.append({"kind": "cover", "ch": None, "xy": (x, y), "r": None})
        return res

    @staticmethod
    def _akey(a):
        return round(float(a) % (2.0 * math.pi), 6)

    def mark_served(self, x, y):
        """某个停点（半径不小于骨架半径）已经代劳覆盖 (*@\ensuremath{\Rightarrow}@*) 附近的骨架点不必再走。

        不做这一步会留下一个隐蔽的浪费：目标停点测完之后目标被清除、不再出现在
        路点里，于是它旁边的骨架点又冒出来 —— 实测会多走一站、多测十几个频道
        （同一片区域在 70 m 内被整批测了两遍）。
        """
        if math.hypot(x, y) < min((r for (_a, r) in self.ring_skeleton()),
                                  default=0.0) * float(self.cfg.cover_duty_frac):
            return
        for (a, r) in self.ring_skeleton():
            px, py = self.ring_xy(a, r)
            if math.hypot(px - x, py - y) <= self.cfg.merge_radius_m:
                self.ring_served.add(self._akey(a))

    def adds_any_coverage(self, x, y):
        """该点是否还能为"某些从没听到过的频道"排除新的环带格点。"""
        need = self.unheard()
        if not need:
            return False
        gx, gy = self.annulus_grid()
        r2 = float(self.cfg.cover_radius_m) ** 2
        near = ((gx - x) ** 2 + (gy - y) ** 2) <= r2
        if not near.any():
            return False
        for ch in need:
            if bool(np.any(self.unsafe_of(ch) & near)):
                return True
        return False

    # ------------------------------------------------------------------
    # 目标路点
    # ------------------------------------------------------------------
    def is_long_sliver(self, br):
        """区域是不是"长条形"（最小包围圆半径巨大，或圆心落到靶区外）。

        两条近似平行的方位会把最小包围圆撑到几公里，圆心甚至落在靶区外十几公里。
        旧实现会**真的朝那个圆心走过去** —— 实测有一局为此多走 13.5 km、多花 2700 s。
        注意：判据只看 r*，不能看"圆心离原点近不近"—— 贴着靶区边界、r* 只有 1 m
        的正常区域圆心也可能在 1770 m 处，那是完全正常的可清除目标。
        """
        return bool(br is not None and float(br["radius"]) > float(self.cfg.big_r_m))

    def pending_singletons(self):
        """还只有一条方位的频道数（决定要不要动用期望场/批量定位）。"""
        return sum(1 for ch in self.heard_pending()
                   if self.recs[ch].n_bearings == 1)

    def _cheap_locate_point(self, rec, step=None):
        """**单点定位的最省做法**：只在当前位置周围一小圈里找最优第二检测点。

        和"期望场"用的是同一个代价模型（`best_refine_point` 的 E[D]），
        区别只在**候选半径**：这里把候选点限制在 ``short_locate_cap_m`` 以内
        （默认 420 m），不再允许"跑到 1 km 外去把基线拉长"。
        交会角小一点、$r^*$ 大一点没关系 —— 到了目标附近还会再定向
        （`clear_target` 的逼近—测量迭代会把误差迅速压下去），
        而**路上多走的时间是实打实丢掉的**。

        实测（`p3_paths` 场景 D）：旧做法为了给一个只知方向的源补方位，
        从 (−1716,58) 往靶区里折了 1020 m；用小步之后行程 16.1 km → 13.4 km。

        返回 ``(xy, info)``；无可用候选时 ``(None, None)``。
        """
        if not rec.pts:
            return None, None
        cap = float(self.cfg.short_locate_cap_m if step is None else step)
        p, info = best_refine_point(rec.pts, rec.svds, self.pos, None, n_t=60,
                                    radii=tuple(r for r in self.cfg.locate_radii
                                                if r <= cap) or (cap,),
                                    travel_cap_m=cap)
        if p is not None:
            if "mode" not in info:
                info = dict(info, mode="short")
            return p, info
        # 代价模型给不出候选（例如射线先验为空）时退化到"几何小步"
        return self._basic_transverse_point(rec, cap)

    def _basic_transverse_point(self, rec, cap):
        """纯几何兜底：从 $S_1$ 出发、与上一条视线成 55–105° 走一小步。"""
        S1 = np.asarray(rec.pts[-1], float)
        th1 = float(rec.svds[-1])
        ts, _w, _t_hi = ray_source_samples(S1, th1, n_t=40)
        if ts.size == 0:
            return None, None
        rho = min(max(0.5 * float(self.cfg.locate_radii[0]), 0.4 * float(np.median(ts))),
                  cap)
        cur = np.asarray(self.pos, float)
        best = None
        for ang in self.cfg.cheap_locate_angles:
            for sgn in (1.0, -1.0):
                a = math.radians(th1 + sgn * ang)
                p = S1 + rho * np.array([math.cos(a), math.sin(a)])
                if float(np.hypot(*p)) > R_ARENA - 1.0:
                    continue
                dd = float(np.linalg.norm(p - cur))
                if best is None or dd < best[0]:
                    best = (dd, p, ang)
        if best is None:
            return None, None
        _d, p, ang = best
        return ((float(p[0]), float(p[1])),
                {"mode": "basic_transverse", "rho": rho, "angle_deg": ang,
                 "move_m": _d})

    def _locate_point(self, ch, rec, br):
        """给"还需要更多方位"的频道选下一个检测点。

        两个候选来源，取总代价小者：

        * **以当前位置为中心**（`best_refine_point` 的默认做法）—— 走到候选点的
          代价已经算在 `cost_s` 里，适合"顺手就能测"的场合；
        * **以"上一条方位的检测点" $S_1$ 为中心**（半径 150–700 m）——
          这对"贴靶区边界、有效接收半径又小"的源是关键：环带骨架在 r=1000，
          而 r(*@\ensuremath{\approx}@*)1800 处与骨架站位同方向的源，相邻两个站位离它有 1.3 km，
          早就超出它的有效接收半径了 —— 骨架只能给它**一条**方位，
          第二条必须专门回到 $S_1$ 附近横向补测。

        **待定频道很少时（< field_min_pending）不走期望场**，直接用
        `_cheap_locate_point` 的小步横向补测。
        """
        if self.pending_singletons() < int(self.cfg.field_min_pending):
            p, info = self._cheap_locate_point(rec)
            if p is not None:
                return p, info
        region = {"center": br["center"]} if br is not None else None
        p, info = best_refine_point(rec.pts, rec.svds, self.pos, region, n_t=90,
                                    radii=tuple(self.cfg.locate_radii),
                                    travel_cap_m=self.cfg.locate_travel_cap_m)
        if p is None:
            return self._cheap_locate_point(rec)
        best = (float(info["cost_s"]) if "cost_s" in info else float("inf"), p, info)
        S1 = tuple(rec.pts[-1])
        if math.dist(S1, self.pos) > 1.0:
            p2, info2 = best_refine_point(
                rec.pts, rec.svds, S1, region, n_t=90,
                radii=tuple(self.cfg.locate_radii),
                travel_cap_m=self.cfg.locate_bearing_cap_m)
            if p2 is not None:
                total = (float(info2.get("cost_s", 1e9))
                         + math.dist(self.pos, p2) / V_ROBOT)
                if total < best[0]:
                    best = (total, p2, info2)
        p, info = best[1], best[2]
        if p is None and br is not None:
            p = tuple(self._safe_point(br["center"]))
            info = {"mode": "center_fallback"}
        return p, info

    @staticmethod
    def _safe_point(p):
        """把点夹到靶区圆盘内（两条近平行方位会把区域中心撑到靶区外十几公里）。"""
        x, y = float(p[0]), float(p[1])
        d = math.hypot(x, y)
        lim = R_ARENA - 1.0
        if d > lim:
            return (x * lim / d, y * lim / d)
        return (x, y)

    def target_waypoints(self):
        """当前所有"听到了但没清掉"的频道 → 下一批目标路点。

        每个频道给的是 ``clear``（有界且不是长条形 (*@\ensuremath{\Rightarrow}@*) 直接去它的定位区域中心）
        或 ``locate``（只有一条方位 / 区域无界 / 长条形 (*@\ensuremath{\Rightarrow}@*) 先去横向补一条方位）。
        **长条形绝不能去最小包围圆的圆心** —— 圆心可能在靶区外十几公里。

        ``locate`` 的门槛（"该不该专门跑一趟去补方位"）分两种：

        * **只有一条方位**：先问"环形航路还会不会再听到它"（`_lost_fraction`），
          还会听到就**耐心等**（环形航路每个角度只经过一次，下一个骨架停点
          通常就在它那个方向）；只有确实听不到了、或已经等了
          ``locate_patience_stops`` 站，才专门去补测；
        * **有 >=2 条方位但区域无界/长条形**：环形航路再给方位也压不下去
          （问题在于交会角太差），必须专门换一个横向检测点。
        """
        cfg = self.cfg
        fut_xy = None
        wps = []
        for ch in self.heard_pending():
            if self.on_hold(ch):
                continue
            rec = self.recs[ch]
            br = self.region_info(rec)
            if br is not None and not self.is_long_sliver(br):
                wps.append({"kind": "clear", "ch": ch, "xy": self._safe_point(br["center"]),
                            "r": float(br["radius"]), "n_bearings": rec.n_bearings})
                continue
            # 区域无界/退化/长条形：必须换一个横向检测点
            if cfg.locate_mode == "never" or self.stop_seq - \
                    self._last_bearing_stop.get(int(ch), -99) < int(cfg.locate_stale_stops):
                continue
            if rec.n_bearings < 2:
                # 还会被环形航路听到 (*@\ensuremath{\Rightarrow}@*) 白等就行，别专门跑一趟
                waited = self.stop_seq - self._last_bearing_stop.get(int(ch), -99)
                if waited < int(cfg.locate_patience_stops):
                    if fut_xy is None:
                        fut_xy = [w["xy"] for w in
                                  self.forward_ring(self.ring_dir if self.ring_dir else 1.0)]
                    if self._lost_fraction(rec, fut_xy) < float(cfg.risk_min_lost):
                        continue
            p, info = self._locate_point(ch, rec, br)
            if p is None:
                continue
            p = self._safe_point(p)
            if math.hypot(p[0] - self.pos[0], p[1] - self.pos[1]) \
                    > cfg.locate_travel_cap_m:
                continue
            wps.append({"kind": "locate", "ch": ch, "xy": p,
                        "r": (float(br["radius"]) if br is not None else None),
                        "info": info, "n_bearings": rec.n_bearings})
        return wps

    # ------------------------------------------------------------------
    # 航路组装：骨架 + 目标"最便宜插入"
    # ------------------------------------------------------------------
    def assemble_route(self, dir_sign, targets=None, cursor=None):
        """返回 ``(route, length)``：骨架按角向单调排列，目标插在最便宜的空隙里。

        ``cover_duty``：被目标"顶替"掉的骨架停点，其覆盖职责转交给该目标 ——
        到那一站时必须把"从没听到过"的频道一并测掉，否则会留下覆盖空洞。
        """
        cfg = self.cfg
        saved = self.ring_cursor_a
        self.ring_cursor_a = cursor
        ring_all = self.forward_ring(dir_sign, cursor=cursor)
        self.ring_cursor_a = saved
        tgs = [dict(t) for t in (self.target_waypoints() if targets is None else targets)]
        for t in tgs:
            t["cover_duty"] = any(math.dist(w["xy"], t["xy"]) <= cfg.merge_radius_m
                                  for w in ring_all)
        ring = [w for w in ring_all
                if all(math.dist(w["xy"], t["xy"]) > cfg.merge_radius_m for t in tgs)]
        seq = list(ring)
        start = np.asarray(self.pos, float)

        def L(a, b):
            return math.dist(tuple(a), tuple(b))

        for t in sorted(tgs, key=lambda w: math.dist(self.pos, w["xy"])):
            best_i, best_extra = 0, None
            for i in range(len(seq) + 1):
                a = tuple(start) if i == 0 else seq[i - 1]["xy"]
                b = seq[i]["xy"] if i < len(seq) else None
                extra = L(a, t["xy"]) + (L(t["xy"], b) - L(a, b) if b else 0.0)
                if best_extra is None or extra < best_extra:
                    best_i, best_extra = i, extra
            cap = (cfg.locate_max_detour_m if t["kind"] == "locate"
                   else cfg.clear_max_detour_m)
            forced = (t["kind"] == "locate"
                      and self._locate_skips.get(t["ch"], 0) >= int(cfg.locate_force_after))
            if best_extra is not None and best_extra > float(cap) and not forced:
                # **定位**路点绕行太狠：本轮先不排它（航路绕下去往往自然会给它第二条方位）。
                # 但连续跳过 locate_force_after 轮之后就强制排进去。
                # **清除**路点不在这个规则里 —— 听到了的源必须清，绕行大只影响排序，
                # 不能变成"不清"（`clear_max_detour_m` 默认取一个极大的值）。
                self._locate_skips[t["ch"]] = self._locate_skips.get(t["ch"], 0) + 1
                continue
            seq.insert(best_i, t)
        mode = str(self.cfg.route_mode)
        if mode == "tsp":
            seq = self.nn_2opt(seq)
        elif mode == "polish":
            seq = self.polish_route(seq)
        if not seq:
            return [], 0.0
        total = L(start, seq[0]["xy"])
        for a, b in zip(seq, seq[1:]):
            total += L(a["xy"], b["xy"])
        return seq, total

    # ---- 航路后处理：允许目标在骨架空隙间自由换位，但不许骨架倒序 ----
    def _order_total(self, order, pts, s):
        d = float(np.linalg.norm(pts[order[0]] - s))
        for a, b in zip(order, order[1:]):
            d += float(np.linalg.norm(pts[b] - pts[a]))
        return d

    def polish_route(self, seq):
        """2-opt 后处理：**骨架停点的相对顺序不变**，目标可以在空隙间换位。

        实测（12 局）"骨架顺序 + 最便宜插入"比"同样这些路点的最优航路"多走
        3.0 km/局（13841 vs 10806 m），其中大部分是"某个目标本该排在后面一站
        之后"这种局部次序问题；允许目标换位、但禁止骨架倒序，既拿回这部分行程，
        又保留"绕圈不折返"的结构性保证。
        """
        if len(seq) < 4:
            return seq
        s = np.asarray(self.pos, float)
        pts = [np.asarray(w["xy"], float) for w in seq]
        is_cover = [w["kind"] == "cover" for w in seq]
        cid, c = {}, 0
        for i, f in enumerate(is_cover):
            if f:
                cid[i] = c
                c += 1

        def ring_ok(order):
            last = -1
            for i in order:
                if is_cover[i]:
                    if cid[i] < last:
                        return False
                    last = cid[i]
            return True

        order = list(range(len(seq)))
        best = self._order_total(order, pts, s)
        improved = True
        while improved:
            improved = False
            for i in range(len(order) - 1):
                for j in range(i + 1, len(order)):
                    cand = order[:i] + order[i:j + 1][::-1] + order[j + 1:]
                    if not ring_ok(cand):
                        continue
                    d = self._order_total(cand, pts, s)
                    if d < best - 1e-9:
                        order, best, improved = cand, d, True
        return [seq[i] for i in order]

    def nn_2opt(self, seq):
        """完全自由的最近邻 + 2-opt（消融对照用；结构约束全部放开）。"""
        if len(seq) < 4:
            return seq
        s = np.asarray(self.pos, float)
        pts = [np.asarray(w["xy"], float) for w in seq]
        k = len(seq)
        left, cur, order = list(range(k)), s, []
        while left:
            j = min(left, key=lambda i: float(np.linalg.norm(pts[i] - cur)))
            order.append(j)
            left.remove(j)
            cur = pts[j]
        best = self._order_total(order, pts, s)
        improved = True
        while improved:
            improved = False
            for i in range(k - 1):
                for j in range(i + 1, k):
                    cand = order[:i] + order[i:j + 1][::-1] + order[j + 1:]
                    d = self._order_total(cand, pts, s)
                    if d < best - 1e-9:
                        order, best, improved = cand, d, True
        return [seq[i] for i in order]

    def pick_direction(self):
        """选定**绕行方向**与**进入环带的位置**。

        旧实现直接用"当前位置的极角"当游标：原点扫描之后机器人还在原点附近，
        `atan2(0,0)` 恒为 0，于是每次都是"从正东开始绕" —— 这就是路径图上
        "进入环的位置是固定的"的原因。

        现在枚举 **(每个骨架停点 × 两个方向)** 共 2n 个入口，各组装一次航路，
        取总长最短的；长度相差在 ``entry_tie_m`` 以内的，再按"进入方向是否
        指向信息方向"（所有单方位频道候选源位的平均方向）来选 ——
        从源多的那一侧切进环带，源会更早被找到、更早被清掉，
        这也正是"平均定位清除时间"这个指标最敏感的地方。
        """
        cfg = self.cfg
        step = 2.0 * math.pi / max(len(self.ring_skeleton()), 1)
        if not cfg.entry_select:
            # 消融对照：只用"当前位置极角"当游标（旧行为）
            best = None
            for s in (1.0, -1.0):
                _seq, L = self.assemble_route(s)
                if best is None or L < best[0]:
                    best = (L, s)
            self.ring_dir = best[1]
            return self.ring_dir
        cands = []
        for s in (1.0, -1.0):
            for i, (a, _r) in enumerate(self.ring_skeleton()):
                cursor = (a - s * 0.5 * step) % (2.0 * math.pi)
                seq, L = self.assemble_route(s, cursor=cursor)
                cands.append((L, s, a))
        if not cands:
            self.ring_dir = 1.0
            return self.ring_dir
        best_len = min(c[0] for c in cands)
        a_info = self.info_direction_deg()
        pool = [c for c in cands if c[0] <= best_len + float(cfg.entry_tie_m)]
        if a_info is None:
            chosen = min(pool, key=lambda c: c[0])
        else:
            def key(c):
                diff = abs((math.degrees(c[2]) - a_info + 180.0) % 360.0 - 180.0)
                return (round(diff / 15.0), c[0])   # 先看是否朝信息方向，再看长度
            chosen = min(pool, key=key)
        _L, s, a = chosen
        self.ring_dir = s
        self.ring_cursor_a = (a - s * 0.5 * step) % (2.0 * math.pi)
        self.log(f"进入环带：选 ({self.cfg.cover_rings[0][0] * math.cos(a):.0f},"
                 f"{self.cfg.cover_rings[0][0] * math.sin(a):.0f})"
                 f"（角 {math.degrees(a):.0f}°），绕行方向 "
                 f"{'逆时针' if s > 0 else '顺时针'}；"
                 f"{len(cands)} 个入口里航路 {best_len:.0f} m 最短"
                 + (f"，信息方向 {a_info:.0f}°" if a_info is not None else ""))
        return self.ring_dir

    def info_direction_deg(self):
        """信息方向：所有"只有一条方位"的频道按先验中位源位加权求平均方向。"""
        acc = np.zeros(2)
        used = 0
        for ch in self.heard_pending():
            rec = self.recs[ch]
            if rec.n_bearings != 1:
                continue
            S1 = np.asarray(rec.pts[-1], float)
            ts, _w, _t_hi = ray_source_samples(S1, float(rec.svds[-1]), n_t=30)
            if ts.size == 0:
                continue
            t = float(np.median(ts))
            u = np.array([math.cos(math.radians(float(rec.svds[-1]))),
                          math.sin(math.radians(float(rec.svds[-1])))])
            g = S1 + t * u
            n = float(np.linalg.norm(g))
            if n > 1e-6:
                acc += g / n
                used += 1
        if used == 0 or float(np.linalg.norm(acc)) < 1e-6:
            return None
        return math.degrees(math.atan2(float(acc[1]), float(acc[0]))) % 360.0

    # ------------------------------------------------------------------
    # 停点测量
    # ------------------------------------------------------------------
    def _cross_ok(self, rec, x, y):
        """从 ``(x, y)`` 再测这条频道，与上一条视线的张角是否够大（否则白测）。"""
        if not rec.pts:
            return True
        (sx, sy), a = rec.pts[-1], rec.svds[-1]
        dx, dy = x - sx, y - sy
        n = math.hypot(dx, dy)
        if n < 1e-6:
            return False
        u = (math.cos(math.radians(a)), math.sin(math.radians(a)))
        cosang = abs((dx * u[0] + dy * u[1]) / n)
        cosang = max(-1.0, min(1.0, cosang))
        return math.degrees(math.acos(cosang)) >= float(self.cfg.min_cross_deg)

    def chans_at(self, x, y, ch=None, cover_duty=False):
        """本停点该测哪些频道。

        * 目标频道一定测；
        * **0 方位**频道（这些频道是"没有源"的频道，测它们纯粹为了证明没有残留源）：
          只有在本停点**负有覆盖职责**、且本点还能排除新的环带格点时才测。
          "负有覆盖职责"= 骨架停点 / 顶替了骨架停点的目标 / **半径不小于骨架半径的
          外圈目标停点** —— 最后一条很关键：半径越大的检测点，对 r=1800 外圈 rim 的
          角向覆盖半宽越大（r=880 时 ±18°，r=1500 时 ±34°），所以**顺路清外圈源时
          顺手一测，就能顶掉两个骨架停点**，比"先跑完骨架再清"省一倍检测；
        * **1 方位**频道：张角够大（>= ``min_cross_deg``）才测。
        """
        out = []
        duty = bool(cover_duty) or ch is None
        if not duty:
            rr = min((r for (_a, r) in self.ring_skeleton()), default=0.0)
            duty = math.hypot(x, y) >= rr * float(self.cfg.cover_duty_frac)
        for c in self.pending():
            rec = self.recs[c]
            if rec.n_bearings == 0:
                if duty and self.adds_coverage(c, x, y):
                    out.append(c)
            elif rec.n_bearings == 1:
                if self._cross_ok(rec, x, y):
                    out.append(c)
            if ch is not None and c == int(ch) and c not in out:
                out.append(c)
        if ch is not None and int(ch) in out:
            out.remove(int(ch))
            out.insert(0, int(ch))
        return out

    def visit(self, wp):
        x, y = wp["xy"]
        chans = self.chans_at(x, y, wp.get("ch"),
                              cover_duty=bool(wp.get("cover_duty")))
        if not chans:
            return 0
        n = self.scan_at(float(x), float(y), chans)
        self.note_pos()
        return n

    # ------------------------------------------------------------------
    # 清除：逼近—测量—试探 的迭代
    # ------------------------------------------------------------------
    def probe_path(self, c, r, poly=None, tried=()):
        """覆盖式试探点，按"里圈起、环内按极角、贪心最近邻"排成一条短折线。"""
        from p3_robot import _dist_point_polygon
        cfg = self.cfg
        c = np.asarray(c, float)
        s = float(cfg.probe_spacing_m)
        rr = min(float(r), float(cfg.probe_max_radius_m)) + R_CLEAR
        u = np.array([s, 0.0])
        v = np.array([s / 2.0, s * math.sqrt(3.0) / 2.0])
        N = int(math.ceil(rr / s)) + 1
        cand = []
        for i in range(-N, N + 1):
            for j in range(-N, N + 1):
                p = c + i * u + j * v
                if float(np.hypot(*(p - c))) <= rr + 1e-9:
                    cand.append(p)
        if poly is not None and len(poly) >= 3:
            P = np.asarray(poly, float)
            keep = [p for p in cand
                    if _dist_point_polygon(p, P) <= R_CLEAR + 1e-9]
            cand = keep if keep else cand
        if tried:
            T = np.asarray(tried, float)
            cand = [p for p in cand
                    if float(np.min(np.hypot(T[:, 0] - p[0], T[:, 1] - p[1]))) >= 2.0]
        if not cand:
            return []
        cand.sort(key=lambda p: (round(float(np.hypot(*(p - c))) / (2 * s)),
                                 math.atan2(float(p[1] - c[1]), float(p[0] - c[0]))))
        cand = cand[:int(cfg.probe_max_points)]
        cur = np.asarray(self.pos, float)
        rest, out = list(cand), []
        while rest:
            j = min(range(len(rest)),
                    key=lambda i: float(np.hypot(*(rest[i] - cur))))
            p = rest.pop(j)
            out.append(p)
            cur = p
        return out

    # 兼容父类 `_probe_count` / `estimate_clear_cost` 的调用名
    def probe_points(self, c, r, poly=None, tried=()):
        return self.probe_path(c, r, poly=poly, tried=tried)

    def sweep_clear(self, ch, rec, c, r, poly=None):
        probes = self.probe_path(c, r, poly=poly, tried=rec.tried_probes)
        self.log(f"  覆盖式试探：中心 ({c[0]:.0f},{c[1]:.0f})，r={r:.0f} m，"
                 f"{'多边形区域' if poly else '外接圆'}，试探点 {len(probes)}"
                 f"（间距 {self.cfg.probe_spacing_m:.0f} m）")
        for p in probes:
            if self.aborted or not self.can_afford(T_CLEAR_FAIL + 5.0):
                return False
            trav = math.dist(self.pos, (float(p[0]), float(p[1])))
            if not self.can_afford(trav / V_ROBOT + T_CLEAR_FAIL + 3.0):
                return False
            rec.tried_probes.append((float(p[0]), float(p[1])))
            self.stats["n_probe"] = self.stats.get("n_probe", 0) + 1
            self.stats["probe_m"] = self.stats.get("probe_m", 0.0) + trav
            if self.do_clear(float(p[0]), float(p[1]), ch):
                return True
        return False

    def _measure_far(self, ch, x, y, min_gap=None):
        """在 ``(x, y)`` 给频道补一条方位（与既有测量点太近就跳过）。"""
        rec = self.recs[int(ch)]
        gap = float(self.cfg.approach_eps_m if min_gap is None else min_gap)
        for q in rec.meas_pts:
            if math.hypot(x - q[0], y - q[1]) < gap:
                return None
        if not self.can_afford(
                math.dist(self.pos, (x, y)) / V_ROBOT + T_MEASURE + T_SWITCH + 4.0):
            return None
        return self.measure(float(x), float(y), int(ch))

    def clear_target(self, ch, rec):
        """定位 → 逼近 → 覆盖式试探（可迭代，直到成功或用尽预算/迭代数）。

        旧实现在 `clear_targets()` 里把 ``r* > 600 m`` 的目标永久过滤掉，
        碰到"区域无界"就直接 `clear_fail += 1` 走人 —— 这正是清除率上不去的主因。
        这里改成迭代：区域大就往中心逼近并补方位（每补一条方位区域都会明显缩小），
        区域小到 ``probe_enter_r_m`` 再走覆盖式试探。
        """
        ch = int(ch)
        cfg = self.cfg
        self.log(f"清除频道{ch}：{rec.n_bearings} 条方位，"
                 f"起点 ({self.pos[0]:.0f},{self.pos[1]:.0f})")
        for it in range(int(cfg.clear_max_iters)):
            if rec.cleared or self.aborted:
                break
            self.n_clear_iter += 1
            br = self.region_info(rec)
            # (A) 没有可用区域：换横向点补方位
            if br is None:
                p, info = self._locate_point(ch, rec, None)
                if p is None:
                    break
                resp = self._measure_far(ch, p[0], p[1])
                if resp is None:
                    break
                self.stats["n_refine"] += 1
                self.log(f"  [{it+1}] 区域无界/退化 -> 横向补测 "
                         f"({p[0]:.0f},{p[1]:.0f})")
                if resp.get("measure_result") == "near":
                    return self.do_clear(float(p[0]), float(p[1]), ch)
                continue
            c = np.asarray(br["center"], float)
            r = float(br["radius"])
            d = math.dist(self.pos, tuple(c))
            long_sliver = self.is_long_sliver(br)
            # (A0) 区域还大时先靠近区域中心（顺便测一条方位把区域压小）。
            #      **区域已经够小时不要走这一步**：覆盖式试探的第一发就是"移动到
            #      区域内的试探点"，本身就是逼近；而"逼近测量"要求新测点离既有
            #      测点至少 approach_eps_m，对"源就在原点附近"这类目标会永远返回
            #      None，于是 clear_target 空转、clear_leftovers 反复重试
            #      （实测 seed 20262713：源在 (−4,15)、r* 只有 2.2 m，却跑了 37.6 km）。
            if (not long_sliver) and r > cfg.probe_enter_r_m \
                    and d > max(r, cfg.arrive_radius_m) + cfg.approach_eps_m:
                if not self.can_afford(d / V_ROBOT + T_MEASURE + T_SWITCH + 6.0):
                    break
                resp = self._measure_far(ch, float(c[0]), float(c[1]))
                if resp is None:
                    p, _info = self._locate_point(ch, rec, br)
                    if p is not None:
                        resp = self._measure_far(ch, p[0], p[1])
                if resp is not None:
                    self.log(f"  [{it+1}] 逼近区域中心 ({c[0]:.0f},{c[1]:.0f})："
                             f"移动 {d:.0f} m（r*={r:.0f} m）")
                    if resp.get("measure_result") == "near":
                        return self.do_clear(float(self.pos[0]), float(self.pos[1]), ch)
                    continue
                # 逼近测量做不了（附近已测过）(*@\ensuremath{\Rightarrow}@*) 直接进入覆盖式试探
            # (B) 区域够小：覆盖式试探
            if r <= cfg.probe_enter_r_m and not long_sliver:
                if self.sweep_clear(ch, rec, c, r, poly=br.get("poly")):
                    return True
                r2 = min(max(r * 1.8, r + 40.0), cfg.probe_max_radius_m)
                if r2 > r + 1.0:
                    self.log(f"  区域内未发现目标 -> 扩大到 r={r2:.0f} m 再试")
                    if self.sweep_clear(ch, rec, c, r2, poly=None):
                        return True
                p, _info = self._locate_point(ch, rec, br)
                if p is None:
                    break
                resp = self._measure_far(ch, p[0], p[1])
                if resp is None:
                    break
                if resp.get("measure_result") == "near":
                    return self.do_clear(float(p[0]), float(p[1]), ch)
                continue
            # (C) 区域还大 / 长条形：补一条横向方位（对长条形**不去圆心**，那可能在靶区外）
            if long_sliver:
                p, _info = self._locate_point(ch, rec, None)
                if p is None:
                    break
                target = self._safe_point(p)
            else:
                target = tuple(c)
                if d < cfg.approach_eps_m:
                    p, _info = self._locate_point(ch, rec, br)
                    if p is None:
                        break
                    target = self._safe_point(p)
            resp = self._measure_far(ch, target[0], target[1])
            if resp is None:
                p, _info = self._locate_point(ch, rec, None if long_sliver else br)
                if p is None:
                    break
                target = self._safe_point(p)
                resp = self._measure_far(ch, target[0], target[1])
                if resp is None:
                    break
            self.stats["n_refine"] += 1
            self.log(f"  [{it+1}] 逼近测量 ({target[0]:.0f},{target[1]:.0f})："
                     f"移动 {math.dist(self.pos, target):.0f} m，r* {r:.0f} -> ?")
            if resp.get("measure_result") == "near":
                return self.do_clear(float(target[0]), float(target[1]), ch)
        if not rec.cleared:
            rec.clear_fail += 1
            self.hold(ch)
            self.log(f"  [x] 频道{ch} 本轮未清除成功（冷却 {cfg.fail_cooldown} 个停点）")
        return rec.cleared

    # ------------------------------------------------------------------
    # 开场：期望场定点 + 先清近处目标
    # ------------------------------------------------------------------
    def initial_stage(self):
        """原点扫描之后：**用期望场选第二个检测点**，顺手清掉靠内圈的目标。

        为什么必须做：原点那一测只听得见一定范围内的源；这些"内圈源"若等到沿环带
        走的时候再处理，就得从环带上专门往回跑一趟。实测（seed 20261113）一次
        "专门跑一趟定位"的绕行高达 2.1 km，等于整局行程的 11%。趁还在原点附近，
        用第二问的 E[D] 代价模型在"两翼"上选第二检测点（对多个单方位频道的候选点
        做完全连接聚类，一次移动定住好几个源），定位完立刻清掉，再往环带走。

        做成**多轮**而不是一轮：一轮只能处理一个聚类，剩下的单方位频道会被留到
        环带上，正是上面那种绕行的来源。
        """
        if self.aborted or self.cfg.locate_mode == "never":
            return 0
        n_total = 0
        for rnd in range(int(self.cfg.initial_rounds)):
            # 每轮的移动上限递增：第一轮只走 350 m（近处的源这样最省），
            # 还有没解决的频道才逐轮放宽（远处的源需要更长的基线）
            cap = min(float(self.cfg.locate_travel_cap_m),
                      float(self.cfg.initial_locate_max_move_m)
                      * (float(self.cfg.initial_move_growth) ** rnd))
            pend = []
            for ch in self.heard_pending():
                rec = self.recs[ch]
                if rec.n_bearings != 1 or self.on_hold(ch):
                    continue
                if math.hypot(rec.pts[0][0], rec.pts[0][1]) \
                        > self.cfg.initial_locate_max_m:
                    continue
                pend.append(ch)
            if not pend:
                break
            per = {}
            for ch in pend:
                p, info = self._locate_point(ch, self.recs[ch], None)
                if p is not None:
                    per[ch] = np.asarray(p, float)
            if not per:
                break
            groups = []
            for ch, p in per.items():
                for g in groups:
                    if all(float(np.linalg.norm(p - per[c2])) <= self.cfg.batch_max_m
                           for c2 in g):
                        g.append(ch)
                        break
                else:
                    groups.append([ch])
            best = None
            for g in groups:
                if len(g) < int(self.cfg.batch_min_channels) \
                        and len(groups) > 1:
                    continue
                P = np.asarray([per[c2] for c2 in g], float)
                for c0 in [P.mean(axis=0)] + [P[i] for i in range(P.shape[0])]:
                    trav = float(np.linalg.norm(c0 - np.asarray(self.pos, float)))
                    if trav > cap:
                        continue
                    cost = trav / V_ROBOT + len(g) * (T_MEASURE + T_SWITCH)
                    key = (-len(g), cost)          # 先看能一次定住几个频道
                    if best is None or key < best[0]:
                        best = (key, c0, g, trav)
            if best is None:
                break
            _key, c0, group, trav = best
            if not self.can_afford(
                    trav / V_ROBOT + len(group) * (T_MEASURE + T_SWITCH) + 20.0):
                break
            self.log(f"开场定点（期望场，第 {rnd+1} 轮）：{len(group)} 个单方位频道 "
                     f"{group} -> 合并第二检测点 ({c0[0]:.0f},{c0[1]:.0f})，"
                     f"移动 {trav:.0f} m")
            self.n_locate_stops += 1
            self.scan_points.append((float(c0[0]), float(c0[1])))
            before = self._state_key()
            for ch in group:
                if self.aborted or not self.can_afford(T_MEASURE + T_SWITCH + 3.0):
                    break
                resp = self.measure(float(c0[0]), float(c0[1]), int(ch))
                n_total += 1
                if resp.get("measure_result") == "near":
                    self.do_clear(float(c0[0]), float(c0[1]), int(ch))
            self.note_pos()
            self.stop_seq += 1
            self.last_stop_xy = (float(c0[0]), float(c0[1]))
            # 顺手清掉已经定下来、又在附近的目标（"先清较靠近原点的目标"）
            self.opportunistic_clear(radius=self.cfg.initial_clear_radius_m, max_n=4)
            if self._state_key() == before:
                break
        self.opportunistic_clear(radius=self.cfg.initial_clear_radius_m, max_n=4)
        return n_total

    def opportunistic_clear(self, radius=250.0, max_n=3):
        """脚下已经定位好、区域也够小的目标立刻清掉（"边走边清"的落实）。"""
        done = 0
        for _ in range(int(max_n)):
            if self.aborted or not self.can_afford(self.cfg.min_action_budget_s):
                break
            cands = []
            for ch in self.heard_pending():
                if self.on_hold(ch):
                    continue
                br = self.region_info(self.recs[ch])
                if br is None or br["radius"] > self.cfg.probe_enter_r_m * 2.0:
                    continue
                d = math.dist(self.pos, br["center"])
                if d <= max(radius, br["radius"]):
                    cands.append((d, ch))
            if not cands:
                break
            cands.sort()
            ch = cands[0][1]
            if not self.clear_target(ch, self.recs[ch]):
                break
            done += 1
        return done

    # ------------------------------------------------------------------
    # 角向扇区：**走过的扇区里"该清没清"的目标不许留到下一圈**
    # ------------------------------------------------------------------
    def offroute_clear(self, here, nxt, nxt2=None, max_n=None):
        """**"现在清"还是"等下一站再清"** —— 用"哪一个的额外绕行更小"来判。

        用户给的直觉是"看第一条方位与行进方向的夹角"：夹角小（源基本在行进方向上）
        就等下一站，夹角大就**立刻走一小步、定位并清除**。方向完全正确 ——
        环形航路每个角度只经过一次，等下一站 = 让下一站替我们定位、
        但之后要**拐回来清**，那个回头距离由"源离环带多远"决定。

        但**只按夹角判太粗**（实测：常规随机 330 → 350、外圈环带 4085 → 5238 s），
        因为夹角大是常态，而"拐回来"到底多花多少，取决于源离哪一站更近。
        所以这里把夹角当**粗筛**，真正的判据是拿射线先验直接算两种做法的
        **额外绕行**：

        * 现在清：``|A→G| + |G→B| − |A→B|``
        * 等下一站：``|B→G| + |G→C| − |B→C|``（没有 C 时按"从 B 出再回 B"算）

        按面积均匀的射线先验对 ``G`` 加权平均，前者明显更小才就地清。
        这样"源就在正前方"（B 更近）自然选等待，"源在侧后方"（A 更近）自然选就地清，
        夹角判据只是它的一个粗糙近似。
        """
        cfg = self.cfg
        if not cfg.offroute_clear or self.aborted or nxt is None:
            return 0
        ax, ay = here["xy"]
        bx, by = nxt["xy"]
        A = np.array([ax, ay], float)
        B = np.array([bx, by], float)
        C = (np.array(nxt2["xy"], float) if nxt2 is not None
             else B + (B - A))
        phi = math.atan2(by - ay, bx - ax)
        cands = []
        for ch in self.heard_pending():
            rec = self.recs[ch]
            if rec.n_bearings != 1 or self.on_hold(ch):
                continue
            # 只处理"就是在这一站听到的"（上一条方位的检测点离本站很近）
            if math.dist(rec.pts[-1], (ax, ay)) > 80.0:
                continue
            th = float(rec.svds[-1])
            sep = abs(math.degrees((th - math.degrees(phi) + 180.0) % 360.0 - 180.0))
            if sep < float(cfg.wait_angle_min_deg):
                continue            # 源几乎在行进方向上 (*@\ensuremath{\Rightarrow}@*) 等下一站肯定更省
            ts, w, _t_hi = ray_source_samples(np.asarray(rec.pts[-1], float), th,
                                              n_t=50)
            if ts.size == 0:
                continue
            if float(_t_hi) > float(cfg.offroute_max_thi_m):
                continue        # 射线太长 (*@\ensuremath{\Rightarrow}@*) 源位估计不可靠，别拿它做判断
            u = np.array([math.cos(math.radians(th)), math.sin(math.radians(th))])
            S1 = np.asarray(rec.pts[-1], float)
            G = S1[None, :] + ts[:, None] * u[None, :]

            def d(X):
                return np.hypot(G[:, 0] - X[0], G[:, 1] - X[1])
            # 源位先验的加权平均距离太远时，估计不可靠、"现在清"多半是白跑
            if float(np.dot(w, d(A))) > float(cfg.offroute_max_dist_m):
                continue
            extra_now = d(A) + d(B) - float(np.linalg.norm(B - A))
            extra_wait = d(B) + d(C) - float(np.linalg.norm(C - B))
            e_now = float(np.dot(w, extra_now))
            e_wait = float(np.dot(w, extra_wait))
            if e_now + float(cfg.offroute_margin_m) >= e_wait:
                continue            # 等下一站更省 (*@\ensuremath{\Rightarrow}@*) 耐心等
            cands.append((e_wait - e_now, ch, sep))
        cands.sort(reverse=True)
        n_max = int(max_n or cfg.offroute_max_per_stop)
        done = 0
        for gain, ch, _sep in cands[:max(0, n_max)]:
            rec = self.recs[ch]
            p, _info = self._cheap_locate_point(rec, step=cfg.offroute_step_m)
            if p is None:
                continue
            dd = math.dist(self.pos, p)
            if dd > float(cfg.offroute_step_m) * 1.3:
                continue
            if not self.can_afford(dd / V_ROBOT + T_MEASURE + T_SWITCH + 8.0):
                continue
            resp = self.measure(float(p[0]), float(p[1]), int(ch))
            self.note_pos()
            self.log(f"  就地定位 频道{ch}：预计省 {gain:.0f} m（现在清 vs 等下一站）"
                     f"-> 走 {dd:.0f} m 到 ({p[0]:.0f},{p[1]:.0f}) 补第二条方位")
            if resp.get("measure_result") == "near":
                self.do_clear(float(p[0]), float(p[1]), int(ch))
                done += 1
                continue
            br = self.region_info(rec)
            if br is None:
                continue
            if self.clear_target(int(ch), rec):
                done += 1
        return done

    def risk_locate(self, max_n=None):
        """**防过期补测**：对"马上就要听不到"的单方位频道，就地走一小步补第二条方位。

        为什么必须做：这条航路每个角度**大致只经过一次**。某个频道如果在它所在
        扇区只拿到一条方位、之后环形航路越走越远，就只能等下一圈（或收尾阶段
        横穿全场）—— 实测 `p3_paths` 场景 C 里 ch6 就是这样横穿了 3.4 km。

        判据（按射线先验加权）：把源位按面积均匀撒在射线上，统计"往后
        ``risk_lookahead`` 个骨架停点**都听不到**（距离 > 1000 m，而任何源的
        有效接收半径都 >= 1000 m）"的比例；比例超过 ``risk_min_lost`` 就说明
        这个源有相当大概率要"过期"。此刻我们还在它的扇区里，走一小步补一条方位
        最便宜；步长只有 ``risk_step_m``（默认 300 m），基线短一点无所谓 ——
        等清到目标附近还会再定向（`clear_target` 的逼近—测量迭代）。
        """
        cfg = self.cfg
        if not cfg.risk_locate or self.aborted:
            return 0
        fut = self.forward_ring(self.ring_dir if self.ring_dir else 1.0)
        fut_xy = [w["xy"] for w in fut[:int(cfg.risk_lookahead)]]
        if not fut_xy:
            return 0
        todo = []
        for ch in self.heard_pending():
            rec = self.recs[ch]
            if rec.n_bearings != 1 or self.on_hold(ch):
                continue
            if self.stop_seq - self._last_bearing_stop.get(int(ch), -99) <= 0:
                continue            # 刚在本站拿到方位，先给下一站一次机会
            if self._lost_fraction(rec, fut_xy) < float(cfg.risk_min_lost):
                continue
            p, info = self._cheap_locate_point(rec)
            if p is None:
                continue
            d = math.dist(self.pos, p)
            if d > float(cfg.risk_step_m) * 1.5:
                continue
            todo.append((d, ch, p, info))
        todo.sort()
        n = 0
        for _d, ch, p, _info in todo[:int(max_n or cfg.risk_max_per_stop)]:
            if not self.can_afford(math.dist(self.pos, p) / V_ROBOT
                                   + T_MEASURE + T_SWITCH + 8.0):
                break
            resp = self.measure(float(p[0]), float(p[1]), int(ch))
            n += 1
            self.log(f"  防过期补测 频道{ch}：走到 ({p[0]:.0f},{p[1]:.0f})"
                     f"（{_d:.0f} m）取第二条方位")
            if resp.get("measure_result") == "near":
                self.do_clear(float(p[0]), float(p[1]), int(ch))
                break
        if n:
            self.note_pos()
            self.opportunistic_clear(radius=self.cfg.sector_clear_m, max_n=2)
        return n

    def _lost_fraction(self, rec, fut_xy):
        """先验中"往后这些骨架停点都听不到该源"的概率（面积均匀的射线先验）。"""
        S1 = np.asarray(rec.pts[-1], float)
        th1 = float(rec.svds[-1])
        ts, w, _t_hi = ray_source_samples(S1, th1, n_t=60)
        if ts.size == 0:
            return 0.0
        u = np.array([math.cos(math.radians(th1)), math.sin(math.radians(th1))])
        G = S1[None, :] + ts[:, None] * u[None, :]
        heard = np.zeros(ts.size, bool)
        for (x, y) in fut_xy:
            heard |= (np.hypot(G[:, 0] - x, G[:, 1] - y) <= R_RECV_MIN)
        return float(w[~heard].sum())

    def sector_clear(self, a_from, a_to):
        """把"刚刚走过的角向扇区"里已经定位好的目标就地清掉。

        扇区定义：从上一站到本站的方位角区间（沿绕行方向）。这就是用户说的
        "每个角度大致只经过一次，覆盖的点要尽快清理"—— 目标已经能清了、
        又正好在我们刚扫过的扇区里，就地清掉比留给下一圈便宜得多。
        """
        if a_from is None or a_to is None or self.aborted:
            return 0
        cap = float(self.cfg.sector_clear_m)
        todo = []
        for ch in self.heard_pending():
            if self.on_hold(ch):
                continue
            br = self.region_info(self.recs[ch])
            if br is None:
                continue
            xy = br["center"]
            d = math.dist(self.pos, xy)
            if d > max(cap, br["radius"]):
                continue
            a = math.atan2(xy[1], xy[0]) % (2.0 * math.pi)
            if self._in_sector(a_from, a_to, a):
                todo.append((d, ch))
        todo.sort()
        done = 0
        for _d, ch in todo:
            if self.aborted or not self.can_afford(self.cfg.min_action_budget_s):
                break
            if self.clear_target(ch, self.recs[ch]):
                done += 1
        return done

    def _in_sector(self, a0, a1, a):
        """``a`` 是否落在"从 ``a0`` 走到 ``a1``"的短弧区间里。"""
        span = (a1 - a0) % (2.0 * math.pi)
        if span > math.pi:
            span = (a0 - a1) % (2.0 * math.pi)
            rel = (a0 - a) % (2.0 * math.pi)
        else:
            rel = (a - a0) % (2.0 * math.pi)
        return rel <= span + 1e-9

    # ------------------------------------------------------------------
    # 主循环
    # ------------------------------------------------------------------
    def run(self):
        t0 = time.time()
        resp = self.arena.enter()
        if resp.get("accepted") is not True:
            raise ArenaError(f"/enter 未被接受：{resp}")
        self.log(f"进入靶区：可用预算 {self.remaining():.0f} s"
                 f"（{self.arena.budget_kind()}）；策略=统一滚动航路")
        try:
            # ---- 第 0 轮：原点全频道扫描 ----
            chans0 = self.chans_at(0.0, 0.0)
            self.log(f"第 0 轮扫描（原地）：{len(chans0)} 个频道 {chans0}")
            self.scan_points.append(self.pos)
            if chans0:
                self.scan_at(self.pos[0], self.pos[1], chans0)
            self.note_pos()
            self.stop_seq += 1

            # ---- 开场：期望场定点，先清近处目标 ----
            self.initial_stage()

            # ---- 统一滚动航路 ----
            stall = 0
            while (not self.aborted and self.stop_seq < int(self.cfg.max_stops)
                   and self.remaining() > self.cfg.min_stop_budget_s):
                if self.is_complete():
                    self.log("* 所有频道要么已清除、要么已被 1000 m 覆盖排除 -> 收工")
                    break
                targets = self.target_waypoints()
                if self.ring_dir == 0.0:
                    self.pick_direction()
                route, _L = self.assemble_route(self.ring_dir, targets=targets,
                                                cursor=self.ring_cursor_a)
                if not route:
                    if self.opportunistic_clear(radius=1e9, max_n=4):
                        continue
                    if self.clear_leftovers():
                        continue
                    self.log("没有可用路点、也没有可处理的目标 -> 结束")
                    break
                wp = route[0]
                # 防死循环：与上一个停点几乎重合时退到下一站
                if (self.last_stop_xy is not None
                        and math.dist(wp["xy"], self.last_stop_xy) < self.cfg.min_move):
                    alt = next((w for w in route[1:]
                                if math.dist(w["xy"], self.last_stop_xy)
                                >= self.cfg.min_move), None)
                    if alt is None:
                        self.log("航路剩余点都与当前停点重合 -> 结束")
                        break
                    wp = alt
                cost = (math.dist(self.pos, wp["xy"]) / V_ROBOT
                        + T_MEASURE + T_SWITCH + 4.0)
                if not self.can_afford(cost):
                    self.log(f"剩余时间不足以再走一站（约 {cost:.0f} s）-> 收工")
                    break
                before = self._state_key()
                nxt = route[1] if len(route) > 1 else None
                nxt2 = route[2] if len(route) > 2 else None
                n = self.visit(wp)
                self.stop_seq += 1
                self.last_stop_xy = tuple(wp["xy"])
                self.mark_served(float(wp["xy"][0]), float(wp["xy"][1]))
                if wp["kind"] == "cover":
                    self.ring_cursor_a = math.atan2(wp["xy"][1], wp["xy"][0]) \
                        % (2.0 * math.pi)
                self.tour_log.append({"stop": self.stop_seq, "kind": wp["kind"],
                                      "ch": wp.get("ch"),
                                      "xy": [round(v, 1) for v in wp["xy"]],
                                      "n_chans": n,
                                      "t": round(float(self.arena.virtual_time_s), 1)})
                self.log(f"航路第 {self.stop_seq} 站｜{self._wp_desc(wp)}｜"
                         f"测 {n} 个频道｜已排除 "
                         f"{1.0 - float(self.uncovered_mask()[2].mean()):.3f}｜"
                         f"已清除 {self._n_cleared()}")
                if n == 0:
                    break
                # 到站之后：本站目标立刻处理，附近可清目标顺手清
                a_prev = self._prev_stop_a
                a_now = math.atan2(wp["xy"][1], wp["xy"][0]) % (2.0 * math.pi)
                self._prev_stop_a = a_now
                ch = wp.get("ch")
                if ch is not None and not self.recs[int(ch)].cleared \
                        and self.region_info(self.recs[int(ch)]) is not None:
                    self.clear_target(int(ch), self.recs[int(ch)])
                self.opportunistic_clear(radius=self.cfg.merge_radius_m, max_n=2)
                # (*@\textcircled{1}@*) 不在行进方向上的新目标：就地补测并清除（见 wait_angle_min_deg）
                self.offroute_clear(wp, nxt, nxt2)
                # (*@\textcircled{2}@*) 刚走过的角向扇区里"已经能清"的目标就地清掉（默认关闭）
                self.sector_clear(a_prev, a_now)
                # (*@\textcircled{3}@*) 马上就要"过期"的单方位频道（默认关闭）
                self.risk_locate()
                if self._state_key() == before:
                    stall += 1
                    if stall >= 3:
                        self.log("连续 3 站没有产生任何新信息 -> 结束")
                        break
                else:
                    stall = 0

            # ---- 收尾 ----
            self.clear_leftovers()
        finally:
            try:
                self.arena.exit()
            except ArenaError as e:
                self.log(f"/exit 异常：{e}")
        self.stats["program_runtime_s"] = time.time() - t0
        self.stats["n_scan_rounds"] = self.stop_seq
        rep = self.report()
        cov, mdist = self.coverage_now()
        rep["coverage"] = cov
        rep["max_uncovered_dist_m"] = mdist
        rep["complete_proof"] = bool(self.is_complete())
        rep["n_cover_stops"] = self.n_cover_stops
        rep["n_locate_stops"] = self.n_locate_stops
        rep["n_clear_iter"] = self.n_clear_iter
        return rep

    def clear_leftovers(self, max_steps=14):
        """收尾：剩下的"听到过但没清掉"的频道，逐个自适应定位清除。"""
        done = 0
        for _ in range(int(max_steps)):
            if self.aborted or self.remaining() <= self.cfg.min_stop_budget_s:
                break
            left = [ch for ch in self.heard_pending() if not self.on_hold(ch)]
            if not left:
                for ch in self.heard_pending():
                    self._hold_until.pop(int(ch), None)
                left = self.heard_pending()
            if not left:
                break
            left.sort(key=lambda c: math.dist(
                self.pos, (self.region_info(self.recs[c]) or {"center": (0.0, 0.0)})["center"]))
            ch = left[0]
            self.log(f"收尾：处理频道{ch}")
            if self.clear_target(ch, self.recs[ch]):
                done += 1
            else:
                self.hold(ch, 1)
        return done

    # ------------------------------------------------------------------
    def _state_key(self):
        return (self._n_cleared(), sum(len(r.pts) for r in self.recs.values()),
                sum(len(r.meas_pts) for r in self.recs.values()))

    def _wp_desc(self, wp):
        x, y = wp["xy"]
        if wp["kind"] == "cover":
            self.n_cover_stops += 1
            return f"骨架覆盖 ({x:.0f},{y:.0f})"
        if wp["kind"] == "clear":
            return (f"清除 ch{wp['ch']}（{wp['n_bearings']} 方位，"
                    f"r*={wp['r']:.0f} m）({x:.0f},{y:.0f})")
        self.n_locate_stops += 1
        return f"定位 ch{wp['ch']}（{wp['n_bearings']} 方位）({x:.0f},{y:.0f})"


# --------------------------------------------------------------------------
# 自检
# --------------------------------------------------------------------------
def _nn_2opt_len(pts, start=(0.0, 0.0)):
    """从 ``start`` 出发、遍历 ``pts`` 的最近邻 + 2-opt 航路长度（下界参考）。"""
    P = [np.asarray(p, float) for p in pts]
    S = np.asarray(start, float)
    k = len(P)
    if k == 0:
        return 0.0
    left, cur, seq = list(range(k)), S, []
    while left:
        j = min(left, key=lambda i: float(np.linalg.norm(P[i] - cur)))
        seq.append(j)
        left.remove(j)
        cur = P[j]

    def tot(order):
        d = float(np.linalg.norm(P[order[0]] - S))
        for a, b in zip(order, order[1:]):
            d += float(np.linalg.norm(P[b] - P[a]))
        return d
    best = tot(seq)
    improved = True
    while improved:
        improved = False
        for i in range(k - 1):
            for j in range(i + 1, k):
                cand = seq[:i] + seq[i:j + 1][::-1] + seq[j + 1:]
                d = tot(cand)
                if d < best - 1e-9:
                    seq, best, improved = cand, d, True
    return best


def selfcheck(verbose=True, seed=20260913, cases=3):
    """T1–T6：统一航路策略的结构性断言。"""
    from p3_arena import MockArena

    out, ok_all = [], True

    def rec(name, ok, detail=""):
        nonlocal ok_all
        ok_all = ok_all and bool(ok)
        out.append({"check": name, "ok": bool(ok), "detail": detail})
        if verbose:
            print(f"[{'OK ' if ok else 'FAIL'}] {name}  {detail}")

    rb = TourRobot(MockArena(seed=1), TourConfig(), base=None, log=None)
    gx, gy, unsafe = rb.uncovered_mask()
    rec("T1 只有原点一次检测时，环带全部未经排除", bool(unsafe.all()),
        f"{gx.size} 个环带格点，未排除 {int(unsafe.sum())} 个")

    sk = rb.ring_skeleton()
    P = np.asarray([rb.ring_xy(a, r) for (a, r) in sk] + [(0.0, 0.0)], float)
    d = np.sqrt(np.min((gx[:, None] - P[None, :, 0]) ** 2
                       + (gy[:, None] - P[None, :, 1]) ** 2, axis=1))
    rec("T2 原点 + 环带骨架即可覆盖环带（最大未覆盖 <= 1000 m）",
        float(d.max()) <= 1000.0,
        f"骨架 {len(sk)} 点，最大未覆盖 {float(d.max()):.1f} m")

    # T3 角向单调：锁定方向后，forward_ring 的极角严格递增
    rb3 = TourRobot(MockArena(seed=2), TourConfig(), base=None, log=None)
    rb3.arena.enter()
    rb3.ring_dir = 1.0
    rb3.ring_cursor_a = 0.3
    seq = rb3.forward_ring(1.0)
    angs = [math.atan2(w["xy"][1], w["xy"][0]) % (2 * math.pi) for w in seq]
    ok3 = len(seq) > 0 and all((b - a) % (2 * math.pi) > 0 for a, b in zip(angs, angs[1:]))
    rec("T3 骨架停点按角向单调推进（不折返的结构性保证）", ok3,
        f"{len(seq)} 个停点，角向增量 "
        f"{[round(math.degrees((b - a) % (2 * math.pi)), 1) for a, b in zip(angs, angs[1:])][:6]}…")

    # T4 端到端
    res = []
    for k in range(cases):
        a = MockArena(seed=seed + 100 * k)
        r4 = TourRobot(a, TourConfig(), base=None, log=None)
        res.append(r4.run())
    rec("T4 mock 端到端跑通且清除率 > 0.8",
        all(r["clear_ratio"] is not None and r["clear_ratio"] > 0.8 for r in res),
        "; ".join(f"seed{seed+100*k}: {r['cleared']}/{r['n_sources']} "
                  f"T={r['virtual_time_s']:.0f}s" for k, r in enumerate(res)))

    # T5 完工判据与真值一致
    bad = [seed + 100 * k for k, r in enumerate(res)
           if r["complete_proof"] and r["cleared"] != r["n_sources"]]
    rec("T5 判定『靶区无残留源』时确实全部清除", not bad, f"反例 {bad or '无'}")

    # T6 行程相对"结构下界"的倍数
    #    下界 = TSP(原点 + 骨架停点 + 全部真源位置)（预知全部源位、且骨架点免费给的理想航路）
    #    实测 30 局该比值中位数约 1.2–1.5；最坏出现在"源数取到上限 16 个"的个别案例上
    #    （路点最多、组合最乱，实测 seed 20261113 约 1.9）。这是"滚动时域在线决策"
    #    相对"离线最优"的真实差距，也是尾部风险，故留到 2.1 作为报警线。
    worst, ratios = 1.0, []
    for k, r in enumerate(res):
        a = MockArena(seed=seed + 100 * k)
        pts = [(rr * math.cos(2 * math.pi * i / n), rr * math.sin(2 * math.pi * i / n))
               for (rr, n) in TourConfig().cover_rings for i in range(int(n))]
        pts += [(x.x, x.y) for x in a.sources]
        lb = _nn_2opt_len(pts)
        ratios.append(r["travel_m"] / max(lb, 1.0))
        worst = max(worst, ratios[-1])
    rec("T6 行程不超过『骨架 + 源位最优航路』的 2.1 倍（30 局中位数 1.2–1.5）",
        worst <= 2.1,
        f"最坏 {worst:.2f}，逐局 {[round(x, 2) for x in ratios]}")

    if verbose:
        print(f"\n自检总结：{'全部通过' if ok_all else '存在失败项'}")
    return {"ok": ok_all, "checks": out}


if __name__ == "__main__":
    import argparse
    ap = argparse.ArgumentParser(description="问题3 统一滚动航路策略自检")
    ap.add_argument("--selfcheck", action="store_true")
    ap.add_argument("--cases", type=int, default=3)
    args = ap.parse_args()
    r = selfcheck(cases=args.cases)
    raise SystemExit(0 if r["ok"] else 1)

\end{lstlisting}
\subsection{\texttt{src/p3\_homing.py}}
% SHA256 e52ddc45a9db120325a07d13af10e24048024d2ffb4431640f5d3b95cdcbb9ec
\begin{lstlisting}[language=python,escapeinside={(*@}{@*)}]
"""Experimental Q3 homing: plan through source estimates, localize on the way.

All decisions use measurement replies only. Ground truth is read by inherited
reporting after /exit, as in TourRobot.
"""
from dataclasses import dataclass
import math
import numpy as np

from p1_intersection import clip_halfplane, wedge_halfplanes, min_enclosing_circle
from p3_arena import R_ARENA, R_RECV_MAX, R_CLEAR
from p3_tour import TourConfig, TourRobot


@dataclass
class HomingConfig(TourConfig):
    initial_rounds: int = 0
    offroute_clear: bool = False
    route_mode: str = 'tsp'
    homing_lateral_m: float = 100.0
    homing_fraction: float = 0.55
    homing_all_single: bool = True
    homing_negatives: bool = True
    homing_max_steps: int = 12


class HomingRobot(TourRobot):
    def __init__(self, arena, cfg=None, base=None, log=None):
        super().__init__(arena, cfg or HomingConfig(), base=base, log=log)
        self._home_regions = {}

    def region(self, rec):
        br = self.region_info(rec)
        if br is None:
            return {'status':'none', 'n_bearings':rec.n_bearings}
        return dict(status='bounded', min_enclosing_center=br['center'],
                    min_enclosing_radius_m=br['radius'], vertices=br['poly'] or [])

    def region_info(self, rec):
        if rec.near_hits:
            x,y,_t,_res = next(row for row in reversed(rec.meas_log) if row[3]=='near')
            return dict(center=(x,y),radius=5.0,poly=None,status='bounded',n_bearings=rec.n_bearings)
        if not rec.pts:
            return None
        key = (rec.channel, len(rec.pts), len(rec.meas_pts))
        if key in self._home_regions:
            return self._home_regions[key]
        P = [np.array(p, float) for p in [(-1800,-1800),(1800,-1800),(1800,1800),(-1800,1800)]]
        # Tangent half planes circumscribe disks, so rounding never excludes truth.
        for S, radius in [(np.zeros(2), R_ARENA)] + [(np.array(s), R_RECV_MAX) for s in rec.pts]:
            for a in np.arange(0, 2*math.pi, math.pi/16):
                u = np.array([math.cos(a), math.sin(a)])
                P = clip_halfplane(P, u, -float(u @ S) - radius)
        for S, theta in zip(rec.pts, rec.svds):
            for a, b in wedge_halfplanes(S, theta, 1.005 + 1e-9):
                P = clip_halfplane(P, a, b)
        if not P:
            return None
        arr = np.array(P)
        origin = arr.mean(axis=0)
        circle = min_enclosing_circle(arr-origin)
        center, radius = circle if circle is not None else (np.zeros(2), np.inf)
        center = center + origin
        c = np.asarray(center, float)
        # Negative replies exclude the 1000m disk. Use them only to improve an
        # estimate, never to certify completion or shrink the outer polygon.
        if rec.n_bearings == 1 and self.cfg.homing_negatives:
            S = np.array(rec.pts[0]); a = math.radians(rec.svds[0])
            u = np.array([math.cos(a), math.sin(a)])
            t = np.linspace(5, 1500, 301)
            G = S + t[:,None]*u
            keep = np.linalg.norm(G,axis=1) <= 1800
            for x,y,_time,res in rec.meas_log:
                if res == 'no_signal':
                    keep &= np.linalg.norm(G - (x,y),axis=1) >= 1000
            if keep.any():
                c = np.average(G[keep], axis=0, weights=t[keep])
        radius = float(np.max(np.linalg.norm(arr-c,axis=1)))
        out = dict(center=tuple(c), radius=radius, poly=arr.tolist(),
                   status='bounded', n_bearings=rec.n_bearings)
        self._home_regions[key] = out
        return out

    def target_waypoints(self):
        if not self.cfg.homing_all_single:
            return super().target_waypoints()
        return [dict(kind='clear', ch=ch, xy=self._safe_point(br['center']),
                     r=br['radius'], n_bearings=self.recs[ch].n_bearings)
                for ch in self.heard_pending() if not self.on_hold(ch)
                if (br := self.region_info(self.recs[ch])) is not None]

    def _home_step(self, ch):
        rec = self.recs[ch]
        br = self.region_info(rec)
        if br is None:
            return self.pos
        c = np.array(br['center']); cur = np.array(self.pos)
        delta = c-cur; d = float(np.linalg.norm(delta))
        if rec.n_bearings == 1:
            a = math.radians(rec.svds[-1]); u = np.array([math.cos(a),math.sin(a)])
            normal = np.array([-u[1],u[0]])
            ahead = cur + self.cfg.homing_fraction*delta
            lat = min(self.cfg.homing_lateral_m, max(35., 0.25*d))
            opts = [self._safe_point(ahead + sign*lat*normal) for sign in [-1,1]]
            return min(opts,key=lambda p:math.dist(p,c)+math.dist(p,cur))
        if d > 25:
            return self._safe_point(c)
        a = math.radians(rec.svds[-1]); normal = np.array([-math.sin(a),math.cos(a)])
        lat = min(80., max(25., br['radius']*0.3))
        return self._safe_point(cur + lat*normal)

    def visit(self, wp):
        ch = wp.get('ch')
        if ch is not None and self.recs[ch].n_bearings == 1:
            p = self._home_step(ch)
            wp['xy'] = p
            chans = [ch] + [k for k in self.heard_pending() if k != ch
                           and self.recs[k].n_bearings == 1 and self._cross_ok(self.recs[k],*p)]
            n = self.scan_at(*p, chans)
            self.note_pos()
            return n
        return super().visit(wp)

    def mark_served(self, x, y):
        # The actual channel observations, not proximity, certify coverage.
        pass

    def clear_target(self, ch, rec):
        for _ in range(self.cfg.homing_max_steps):
            if rec.cleared or self.aborted or not self.can_afford(15):
                break
            br = self.region_info(rec)
            if br is None:
                break
            c, radius = br['center'], br['radius']
            if radius <= R_CLEAR or (rec.n_bearings >= 2 and radius < 65 and math.dist(self.pos,c)<100):
                if self.do_clear(*c, ch):
                    return True
            p = self._home_step(ch)
            if not self.can_afford(math.dist(self.pos,p)/5+12):
                break
            resp = self.measure(*p,ch)
            if resp.get('measure_result') == 'near':
                return self.do_clear(*p,ch)
            if resp.get('measure_result') == 'no_signal':
                # Retry toward the last positive position; never infer absence.
                S=np.array(rec.pts[-1]); p=(np.array(p)+S)/2
                resp = self.measure(*p,ch)
                if resp.get('measure_result') == 'near':
                    return self.do_clear(*p,ch)
        if not rec.cleared:
            return super().clear_target(ch,rec)
        return True

\end{lstlisting}
\subsection{\texttt{src/p3\_frontier.py}}
% SHA256 a0ebe7e00ea5859ac9174cb4a8698b229d92915d445c08b7f2e787f8fad9cca9
\begin{lstlisting}[language=python,escapeinside={(*@}{@*)}]
"""Experimental receding-horizon covering-tour planner for B question 3.

Unlike TourRobot, this planner has no fixed ring stations or angular cursor.
At each decision it builds a route through known targets, then solves a greedy
covering-tour problem over the *remaining per-channel coverage gap*. Target
stops compete with temporary measurement stations, including their scan cost.
Only the first stop is executed; actual observations determine the next plan.

This module deliberately inherits the existing localization and clearance
procedure so routing changes can be compared separately. No truth is consulted.
The inherited completion predicate evaluates a 25 m grid; its report field
``complete_proof`` is not a continuous geometric coverage certificate.
"""
from __future__ import annotations

import argparse
import json
import math
from dataclasses import dataclass, fields, replace
from pathlib import Path

import numpy as np

from p3_arena import MockArena, T_MEASURE, T_SWITCH, V_ROBOT
from p3_bench import (LIVE1, LIVE2, HDR, arena_from, center_cluster, fmt_row,
                      outer_annulus, run_one, summarize)
from p3_robot import load_base
from p3_tour import TourConfig, TourRobot


@dataclass
class FrontierConfig(TourConfig):
    candidate_radii: tuple = (900.0, 1100.0, 1300.0, 1500.0)
    candidate_angles: int = 72
    scan_cost_weight: float = 0.3
    travel_cost_floor_m: float = 100.0
    gain_power: float = 2.0
    target_first: bool = False
    target_cover_duty: bool = False
    route_mode: str = "tsp"
    locate_mode: str = "never"
    offroute_clear: bool = False
    enroute_spacing_m: float = 0.0


class FrontierRobot(TourRobot):
    """Dynamic set cover with cheapest insertion and unconstrained 2-opt."""

    def __init__(self, arena, cfg=None, base=None, log=None):
        super().__init__(arena, cfg or FrontierConfig(), base=base, log=log)
        self._frontier_candidates = None
        self._frontier_cover = None

    def pick_direction(self):
        # Retained only to satisfy the inherited loop; no angular restriction.
        self.ring_dir = 1.0

    def mark_served(self, x, y):
        # Coverage is recomputed from real observations, never by proximity.
        return None

    def forward_ring(self, dir_sign, cursor=None):
        # Stale-localization mode is intentionally disabled in this experiment.
        return []

    def chans_at(self, x, y, ch=None, cover_duty=False):
        out = []
        for c in self.pending():
            rec = self.recs[c]
            if rec.n_bearings == 0:
                if (cover_duty or ch is None) and self.adds_coverage(c, x, y):
                    out.append(c)
            elif rec.n_bearings == 1 and self._cross_ok(rec, x, y):
                out.append(c)
            if ch is not None and c == int(ch) and c not in out:
                out.append(c)
        if ch is not None and int(ch) in out:
            out.remove(int(ch))
            out.insert(0, int(ch))
        return out

    def _candidate_cover(self):
        if self._frontier_candidates is None:
            angles = np.arange(self.cfg.candidate_angles) * (2 * math.pi /
                      self.cfg.candidate_angles)
            self._frontier_candidates = np.asarray([
                (r * math.cos(a), r * math.sin(a))
                for r in self.cfg.candidate_radii for a in angles])
            gx, gy = self.annulus_grid()
            self._frontier_cover = ((self._frontier_candidates[:, 0, None] - gx) ** 2
                    + (self._frontier_candidates[:, 1, None] - gy) ** 2
                    <= self.cfg.cover_radius_m ** 2)
        return self._frontier_candidates, self._frontier_cover

    def _insert_cost(self, pts, seq):
        """Cost and index for inserting every candidate into the open route."""
        prev = np.asarray([self.pos] + [w["xy"] for w in seq], float)
        d = np.linalg.norm(pts[:, None, :] - prev[None, :, :], axis=2)
        cost = d.copy()
        if seq:
            cost[:, :-1] += d[:, 1:] - np.linalg.norm(np.diff(prev, axis=0), axis=1)
        ind = np.argmin(cost, axis=1)
        return cost[np.arange(len(pts)), ind], ind

    def assemble_route(self, dir_sign, targets=None, cursor=None):
        seq = [dict(w, cover_duty=False) for w in
               (self.target_waypoints() if targets is None else targets)]
        seq = self.nn_2opt(seq)
        lattice, lattice_cover = self._candidate_cover()
        gx, gy, unsafe = self.uncovered_mask()
        if not unsafe.any():
            return seq, self._route_length(seq)

        if self.cfg.enroute_spacing_m > 0.0 and seq:
            path = [self.pos] + [w["xy"] for w in seq]
            extras = []
            for a, b in zip(path, path[1:]):
                parts = int(math.ceil(math.dist(a, b) / self.cfg.enroute_spacing_m))
                extras.extend(((1-t)*a[0] + t*b[0], (1-t)*a[1] + t*b[1])
                              for t in np.arange(1, parts) / parts)
            if extras:
                extra_pts = np.asarray(extras, float)
                extra_cover = ((extra_pts[:, 0, None] - gx) ** 2
                         + (extra_pts[:, 1, None] - gy) ** 2
                         <= self.cfg.cover_radius_m ** 2)
                lattice = np.concatenate((lattice, extra_pts), axis=0)
                lattice_cover = np.concatenate((lattice_cover, extra_cover), axis=0)

        # Known targets are mandatory route stops but their coverage scan is
        # optional. The set-cover objective decides which should perform it.
        targets_copy = list(seq)
        if targets_copy:
            target_pts = np.asarray([w["xy"] for w in targets_copy], float)
            target_cover = ((target_pts[:, 0, None] - gx) ** 2
                    + (target_pts[:, 1, None] - gy) ** 2
                    <= self.cfg.cover_radius_m ** 2)
            pts = np.concatenate((lattice, target_pts), axis=0)
            covers = np.concatenate((lattice_cover, target_cover), axis=0)
        else:
            pts, covers = lattice, lattice_cover
        available = np.ones(len(pts), bool)
        n_lattice = len(lattice)
        scan_cost = (T_MEASURE + T_SWITCH) * max(1, len(self.unheard())) * V_ROBOT
        scan_cost *= self.cfg.scan_cost_weight

        # Retain the inherited grid coverage requirement, without early stopping.
        # The score decides route shape only.
        for _ in range(30):
            if not unsafe.any():
                break
            gains = np.count_nonzero(covers[:, unsafe], axis=1).astype(float)
            costs, inserts = self._insert_cost(pts, seq)
            costs[n_lattice:] = 0.0
            score = gains ** self.cfg.gain_power / (scan_cost +
                    self.cfg.travel_cost_floor_m + np.maximum(0.0, costs))
            score[~available] = -1.0
            best = int(np.argmax(score))
            if gains[best] == 0:
                break
            available[best] = False
            if best >= n_lattice:
                targets_copy[best - n_lattice]["cover_duty"] = True
            else:
                seq.insert(int(inserts[best]), {"kind": "cover", "ch": None,
                    "xy": tuple(pts[best]), "r": None, "cover_duty": True})
            unsafe = unsafe & ~covers[best]

        # Reorder the joint route without the original angular constraints.
        seq = self.nn_2opt(seq)
        if self.cfg.target_first and targets_copy:
            nearest = min(targets_copy, key=lambda w: math.dist(self.pos, w["xy"]))
            seq.remove(nearest)
            seq.insert(0, nearest)
        return seq, self._route_length(seq)

    def _route_length(self, seq):
        pts = [self.pos] + [w["xy"] for w in seq]
        return sum(math.dist(a, b) for a, b in zip(pts, pts[1:]))


class OriginalTourRobot(TourRobot):
    """Freeze the user's original singleton measurement behavior for controls."""

    def _singleton_reachable(self, rec, x, y):
        return True


def original_config():
    vals = dict(cover_rings=((1000.0, 8),), route_mode="polish",
                locate_mode="stale", offroute_clear=True, target_cover_duty=True)
    return TourConfig(**{k: v for k, v in vals.items() if k in
                         {f.name for f in fields(TourConfig)}})


def main():
    ap = argparse.ArgumentParser(description=__doc__)
    ap.add_argument("--cases", type=int, default=5)
    ap.add_argument("--seed", type=int, default=20260913)
    ap.add_argument("--step", type=int, default=100)
    ap.add_argument("--scenarios", default="random,annulus,center,live1,live2")
    ap.add_argument("--variants", default="original,tour,frontier")
    ap.add_argument("--cfg", default="{}")
    ap.add_argument("--json-out", default="")
    args = ap.parse_args()
    makers = {"random": lambda s: MockArena(seed=s), "annulus": outer_annulus,
              "center": center_cluster, "live1": lambda s: arena_from(LIVE1, seed=s),
              "live2": lambda s: arena_from(LIVE2, seed=s)}
    cfg = replace(FrontierConfig(), **json.loads(args.cfg))
    variants = {"original": (OriginalTourRobot, original_config()),
                "tour": (TourRobot, TourConfig(cover_rings=((1050.0, 7),),
                         route_mode="tsp", locate_mode="never", offroute_clear=False)),
                "frontier": (FrontierRobot, cfg)}
    base = load_base()
    output = {"arguments": vars(args), "results": {}}
    for scenario in args.scenarios.split(","):
        print(scenario, flush=True)
        print(HDR, flush=True)
        for name in args.variants.split(","):
            cls, config = variants[name]
            reps = [run_one(cls, config, makers[scenario](args.seed + i * args.step),
                            base, want_path=True) for i in range(args.cases)]
            summary = summarize(reps)
            summary["complete_proof_cases"] = sum(bool(r["complete_proof"]) for r in reps)
            print(fmt_row(name, summary) + " proof=" + str(summary["complete_proof_cases"]),
                  flush=True)
            output["results"][scenario + "/" + name] = {**summary,
                "per_case": [{k: v for k, v in r.items() if not k.startswith("_")
                    and k not in ("truth", "channels", "mock_stats")} for r in reps]}
            if args.json_out:
                p = Path(args.json_out)
                p.parent.mkdir(parents=True, exist_ok=True)
                p.write_text(json.dumps(output, ensure_ascii=False, indent=2, default=str),
                             encoding="utf-8")


if __name__ == "__main__":
    main()

\end{lstlisting}
\subsection{\texttt{src/p3\_coverage.py}}
% SHA256 6a4969451a99e66dafea2aeab74c15384b3142f387b736f34139b4321bae551b
\begin{lstlisting}[language=python,escapeinside={(*@}{@*)}]
"""Conservative continuous-disk coverage from actual no-signal replies.

Closed square cells tile [-1800,1800]^2, including every cell intersecting
the arena disk. A center covered at radius 1000 - step/sqrt(2) - epsilon
certifies the entire cell by the triangle inequality. No positive/near
reply is ever used to certify absence. This may over-cover boundary cells.
"""
import math
import numpy as np
from p3_arena import R_ARENA, R_RECV_MIN


class CertifiedCoverageMixin:
    def annulus_grid(self):
        if self._ag is None:
            requested = float(self.cfg.cover_grid_step_m)
            if requested <= 0:
                raise ValueError('cover_grid_step_m must be positive')
            n = int(math.ceil(2 * R_ARENA / requested))
            step = 2 * R_ARENA / n
            xs = -R_ARENA + (np.arange(n)+0.5)*step
            gx, gy = np.meshgrid(xs,xs)
            intersects = (np.maximum(np.abs(gx)-step/2,0)**2 +
                          np.maximum(np.abs(gy)-step/2,0)**2 <= R_ARENA**2+1e-6)
            self._ag = (gx[intersects],gy[intersects])
            self._cert_radius = R_RECV_MIN-step/math.sqrt(2)-1e-6
            if self._cert_radius <= 0:
                raise ValueError('coverage cells are too large')
        return self._ag

    def certified_radius(self):
        self.annulus_grid()
        return self._cert_radius

    def unsafe_of(self,ch):
        rec = self.recs[int(ch)]
        key = (int(ch),len(rec.meas_log))
        if key not in self._unsafe_cache:
            gx,gy = self.annulus_grid()
            out = np.ones(len(gx),dtype=bool)
            r2 = self.certified_radius()**2
            for x,y,_t,result in rec.meas_log:
                if result == 'no_signal':
                    out &= (gx-x)**2+(gy-y)**2 > r2
            self._unsafe_cache[key] = out
        return self._unsafe_cache[key]

    def heard_pending(self):
        return [int(ch) for ch,r in self.recs.items()
                if not r.cleared and (r.n_bearings or r.near_hits)]

    def unheard(self):
        return [int(ch) for ch,r in self.recs.items()
                if not r.cleared and not r.n_bearings and not r.near_hits]

    def adds_coverage(self,ch,x,y):
        gx,gy = self.annulus_grid()
        return bool(np.any(self.unsafe_of(ch) &
                    ((gx-x)**2+(gy-y)**2 <= self.certified_radius()**2)))

    def adds_any_coverage(self,x,y):
        return any(self.adds_coverage(ch,x,y) for ch in self.unheard())

    def is_complete(self):
        return (not self.heard_pending() and
                all(not self.unsafe_of(ch).any() for ch in self.unheard()))

\end{lstlisting}
\subsection{\texttt{src/p3\_adaptive.py}}
% SHA256 f1f9d33e31695aac558b8aab390a5016231c645c8c1339afe1f459693c1c01ae
\begin{lstlisting}[language=python,escapeinside={(*@}{@*)}]
"""Experimental combination of homing and dynamic covering-tour planning."""
from dataclasses import dataclass
import math
import numpy as np
from p3_homing import HomingConfig, HomingRobot
from p3_frontier import FrontierConfig, FrontierRobot
from p3_coverage import CertifiedCoverageMixin


@dataclass
class AdaptiveConfig(HomingConfig, FrontierConfig):
    locate_mode: str = 'never'
    quick_baseline_m: float = 400.0
    cover_grid_step_m: float = 50.0


class AdaptiveRobot(CertifiedCoverageMixin, HomingRobot, FrontierRobot):
    def __init__(self, arena, cfg=None, base=None, log=None):
        super().__init__(arena, cfg or AdaptiveConfig(), base=base, log=log)
        # The route's candidate masks must use the same conservative radius.
        from dataclasses import replace
        self.cfg = replace(self.cfg,cover_radius_m=self.certified_radius(),cover_slack_m=0.0)

    def initial_stage(self):
        if self.cfg.quick_baseline_m <= 0:
            return
        pending = self.heard_pending()
        if not pending:
            return
        angles = np.arange(72)*2*math.pi/72
        bearings = np.radians([self.recs[c].svds[0] for c in pending])
        # One short shared transverse baseline, then approach while measuring.
        score = np.sum(np.sin(angles[:,None]-bearings)**2,axis=1)
        a = angles[int(np.argmax(score))]
        p = self.cfg.quick_baseline_m*np.array([math.cos(a),math.sin(a)])
        self.scan_at(*p,pending)
        self.note_pos()
        self.stop_seq += 1

\end{lstlisting}
\subsection{\texttt{src/p3\_joint.py}}
% SHA256 6bd22c9ed32e41bbea634193c66c4fa68476367eaf855e94313bef3968732efe
\begin{lstlisting}[language=python,escapeinside={(*@}{@*)}]
"""Q3 joint routing and scan scheduling with continuous coverage certificates.

Only actual replies enter decisions. Probability gates affect optional bearing
measurements, never mandatory clearance or the deterministic completion proof.
"""
from dataclasses import dataclass
import math
import numpy as np
from p3_adaptive import AdaptiveRobot,AdaptiveConfig


@dataclass
class JointConfig(AdaptiveConfig):
    after_service_scan: str = 'duty'
    reachable_filter: bool = True
    polish_small_routes: bool = True
    candidate_radii: tuple = (1050.,1200.,1350.,1500.)
    speculative_radius_m: float = 400.0
    cover_polish: bool = True
    planner_multistart: bool = True
    opportunistic_receive_probability: float = 0.5
    continuous_cover: bool = True
    probe_fail_measure: bool = True
    min_move: float = 0.0


class JointRobot(AdaptiveRobot):
    def __init__(self,arena,cfg=None,base=None,log=None):
        super().__init__(arena,cfg or JointConfig(),base=base,log=log)


    def clear_target(self,ch,rec):
        if self.cfg.speculative_radius_m and rec.n_bearings>=2 and not rec.cleared:
            br=self.region_info(rec)
            if br and br['radius'] <= self.cfg.speculative_radius_m:
                cost=math.dist(self.pos,br['center'])/5+15
                if self.can_afford(cost):
                    if self.do_clear(*br['center'],ch):
                        return True
                    if self.cfg.probe_fail_measure and self.can_afford(8):
                        resp=self.measure(*self.pos,ch)
                        if resp.get('measure_result')=='near':
                            return self.do_clear(*self.pos,ch)
        return super().clear_target(ch,rec)


    def nn_2opt(self,seq):
        if self.cfg.polish_small_routes and len(seq)<4:
            from itertools import permutations
            return list(min(permutations(seq),key=self._route_length)) if seq else []
        return super().nn_2opt(seq)


    def assemble_route(self,dir_sign,targets=None,cursor=None):
        if self.cfg.planner_multistart:
            from dataclasses import replace
            cfg=self.cfg
            proposals=[]
            try:
                for power,weight in ((1.,.3),(1.,1.),(2.,.3),(2.,1.),(3.,.3)):
                    self.cfg=replace(cfg,planner_multistart=False,gain_power=power,scan_cost_weight=weight)
                    seq,length=self.assemble_route(dir_sign,targets=targets,cursor=cursor)
                    proposals.append((self._plan_cost(seq),seq,length))
            finally:
                self.cfg=cfg
            _cost,seq,length=min(proposals,key=lambda q:q[0])
            return seq,length
        seq,length=super().assemble_route(dir_sign,targets=targets,cursor=cursor)
        if not self.cfg.cover_polish or not seq:
            return seq,length
        pts,covers=self._candidate_cover()
        gx,gy,needed=self.uncovered_mask()
        if not needed.any():
            return seq,length
        r2=self.cfg.cover_radius_m**2
        for _ in range(3):
            changed=False
            for w in list(seq):
                if not w.get('cover_duty'):
                    continue
                others=np.zeros(len(gx),bool)
                for v in seq:
                    if v is not w and v.get('cover_duty'):
                        x,y=v['xy'];others |= (gx-x)**2+(gy-y)**2<=r2
                must=needed&~others
                if not must.any():
                    if w['kind']=='cover':
                        seq.remove(w)
                    else:
                        w['cover_duty']=False
                    changed=True
                    continue
                if w['kind']!='cover':
                    continue
                allowed=np.all(covers[:,must],axis=1)
                if not allowed.any():
                    continue
                reduced=[v for v in seq if v is not w]
                costs,places=self._insert_cost(pts,reduced)
                costs[~allowed]=np.inf
                j=int(np.argmin(costs))
                old_cost=self._route_length(seq)-self._route_length(reduced)
                if costs[j] < old_cost-1e-5:
                    w['xy']=tuple(pts[j]);reduced.insert(int(places[j]),w)
                    seq=reduced;changed=True
                if self.cfg.continuous_cover:
                    reduced=[v for v in seq if v is not w]
                    start=np.asarray(w['xy'],float)
                    Q=np.column_stack([gx[must],gy[must]])
                    delta=start-Q
                    c=np.einsum('ij,ij->i',delta,delta)-r2
                    if c.max(initial=-np.inf)>1e-5:
                        continue
                    chain=np.array([self.pos]+[v['xy'] for v in reduced])
                    destinations=[chain[-1]]
                    for a,b in zip(chain,chain[1:]):
                        v=b-a;den=v@v
                        t=float(np.clip((start-a)@v/max(den,1e-12),0,1))
                        destinations.append(a+t*v)
                    candidates=[]
                    for z in destinations:
                        v=z-start;aa=float(v@v)
                        if aa<1e-10:continue
                        bb=2*delta@v
                        roots=(-bb+np.sqrt(np.maximum(0,bb*bb-4*aa*c)))/(2*aa)
                        t=max(0.,min(1.,float(roots.min()))-1e-8)
                        candidates.append(start+t*v)
                    if candidates:
                        P=np.asarray(candidates)
                        costs,places=self._insert_cost(P,reduced)
                        j=int(np.argmin(costs))
                        old_cost=self._route_length(seq)-self._route_length(reduced)
                        if costs[j]<old_cost-1e-5:
                            w['xy']=tuple(P[j]);reduced.insert(int(places[j]),w)
                            seq=reduced;changed=True
            seq=self.nn_2opt(seq)
            if not changed:
                break
        return seq,self._route_length(seq)

    def _plan_cost(self,seq):
        gx,gy=self.annulus_grid();r2=self.cfg.cover_radius_m**2
        pending={ch:self.unsafe_of(ch).copy() for ch in self.unheard()}
        scans=0
        for w in seq:
            if not w.get('cover_duty'):
                continue
            x,y=w['xy'];mask=(gx-x)**2+(gy-y)**2<=r2
            for ch,m in pending.items():
                if np.any(m&mask):
                    scans+=1;pending[ch]=m&~mask
        return self._route_length(seq)/5+6*scans

    def _cross_ok(self,rec,x,y):
        if self.cfg.opportunistic_receive_probability>0 and rec.n_bearings==1:
            s=np.array(rec.pts[0]);a=math.radians(rec.svds[0])
            t=np.linspace(5.,1500.,151);u=np.array([math.cos(a),math.sin(a)])
            p=s+t[:,None]*u
            lo=np.maximum(1000.,t);hi=np.full(len(t),1500.)
            for qx,qy,_t,res in rec.meas_log:
                if res=='no_signal':hi=np.minimum(hi,np.linalg.norm(p-(qx,qy),axis=1))
            density=t*(np.linalg.norm(p,axis=1)<=1800.)
            total=np.sum(density*np.maximum(0.,hi-lo))
            heard=np.sum(density*np.maximum(0.,hi-np.maximum(lo,np.linalg.norm(p-(x,y),axis=1))))
            if total>0 and heard/total < self.cfg.opportunistic_receive_probability:
                return False
        if self.cfg.reachable_filter and rec.n_bearings==1:
            br=self.region_info(rec)
            if br:
                poly=br['poly']
                from p3_robot import _dist_point_polygon
                if poly and _dist_point_polygon(np.array([x,y]),np.array(poly))>1500.:
                    return False
        return super()._cross_ok(rec,x,y)

    def visit(self,wp):
        ch=wp.get('ch')
        if ch is None or self.cfg.after_service_scan=='off':
            return super().visit(wp)
        rec=self.recs[ch]
        before=self.stats['n_measure']
        self.clear_target(ch,rec)
        # The route's target was an estimate; perform its scan at the actual
        # endpoint after localization, then replan using that measured position.
        wp['xy']=tuple(self.pos)
        chans=[]
        duty=wp.get('cover_duty') or self.cfg.after_service_scan=='all'
        if self.cfg.after_service_scan=='replan':
            here=dict(kind='scan_here',ch=None,xy=tuple(self.pos),r=None)
            route,_=self.assemble_route(1.,targets=[here]+self.target_waypoints())
            duty=any(w['kind']=='scan_here' and w.get('cover_duty') for w in route)
        for k in self.pending():
            r=self.recs[k]
            if not r.n_bearings and duty and self.adds_coverage(k,*self.pos):
                chans.append(k)
            elif r.n_bearings==1 and self._cross_ok(r,*self.pos):
                chans.append(k)
        if chans:
            self.scan_at(*self.pos,chans)
        self.note_pos()
        # The inherited loop also accepts a successful clear as progress.
        return max(1,self.stats['n_measure']-before)

\end{lstlisting}
\subsection{\texttt{src/p3\_bench.py}}
% SHA256 6e9d7f665b79f7d970a9c6e3add9c0aa0f4af73628411cd39d0b510ccd142e7a
\begin{lstlisting}[language=python,escapeinside={(*@}{@*)}]
"""问题 3 改进用的批量对照台（不进入正式交付，只服务于调参与回归）。

用法示例::

    python src/p3_bench.py --cases 30                       # 默认对照集
    python src/p3_bench.py --cases 30 --set onlap           # 指定对照集
    python src/p3_bench.py --cases 10 --only A,onlap        # 只跑其中几项

输出：一行一配置的汇总表（清除率 / 平均定位清除时间 / 虚拟时间 / 检测 / 行程 /
折返指标），并可选写出 JSON 供后续比对。
"""
from __future__ import annotations

import argparse
import json
import math
import os
import statistics as st
import sys
import time
from dataclasses import replace

import numpy as np

_HERE = os.path.dirname(os.path.abspath(__file__))
sys.path.insert(0, _HERE)
_ROOT = os.path.normpath(os.path.join(_HERE, ".."))

# Windows 控制台默认 GBK：中文与表格符号混排时容易抛 UnicodeEncodeError
for _s in (sys.stdout, sys.stderr):
    try:
        _s.reconfigure(encoding="utf-8", errors="replace")
    except Exception:
        pass

from p3_arena import MockArena, MockSource, R_ARENA, V_ROBOT  # noqa: E402
from p3_robot import P3Config, P3Robot, load_base  # noqa: E402
from p3_sweep import SweepConfig, SweepRobot  # noqa: E402

LIVE1 = [(2, 937, -782), (5, 363, 418), (6, -939, -1491), (9, 1652, -112),
         (10, 1606, -590), (11, -591, -82), (13, 18, -197), (14, 470, -1292),
         (18, 1024, 1242), (20, -316, 241)]
LIVE2 = [(1, 693, -1157), (2, 722, 1401), (4, 288, 1723), (6, 622, 86),
         (8, 1066, 1008), (10, 1060, -550), (11, -386, 628), (14, 1378, 265),
         (15, 56, 734), (16, -722, -719), (17, 279, 945), (18, 1746, -247),
         (19, -509, 1384), (20, 498, 769)]


def arena_from(sources, r_recv=1400.0, seed=0, budget=360000.0):
    return MockArena(seed=seed, budget_s=budget,
                     sources=[MockSource(c, x, y, r_recv) for (c, x, y) in sources])


def outer_annulus(seed=0, n=12, r_lo=1300.0, r_hi=1800.0, budget=360000.0):
    rng = np.random.default_rng(seed)
    chans = rng.choice(np.arange(1, 21), size=n, replace=False)
    out = []
    for ch in chans:
        r = math.sqrt(rng.uniform(r_lo ** 2, r_hi ** 2))
        a = rng.uniform(0, 2 * math.pi)
        out.append(MockSource(int(ch), r * math.cos(a), r * math.sin(a),
                              float(rng.uniform(1000.0, 1500.0))))
    return MockArena(seed=seed, sources=out, budget_s=budget)


def center_cluster(seed=0, n=12, budget=360000.0):
    out = []
    rng = np.random.default_rng(seed)
    for c in range(1, n + 1):
        r = rng.uniform(0, 260.0)
        a = rng.uniform(0, 2 * math.pi)
        out.append(MockSource(c, r * math.cos(a), r * math.sin(a),
                              float(rng.uniform(1000.0, 1500.0))))
    return MockArena(seed=seed, sources=out, budget_s=budget)


# --------------------------------------------------------------------------
def path_metrics(acts, robot=None):
    """路径形态指标：折返（相邻两段夹角 > 120°）次数与折返累计距离。

    用户关心的"大量折返 / 绕远路"需要一个可量化的代理指标：把动作序列看成折线，
    统计"转弯角 > 120° 且折返段长度 > 150 m"的段数（记为 ``u_turns``）以及这些段
    的总长度（``u_m``）。纯绕圈 + 顺路清除的理想路径此值应接近 0。
    """
    pts = [(0.0, 0.0)] + [(a["x"], a["y"]) for a in acts
                          if a.get("kind") in ("measure", "clear")
                          and "x" in a and "y" in a]
    # 去掉原地重复点
    P = []
    for p in pts:
        if not P or math.dist(P[-1], p) > 1.0:
            P.append(p)
    u_turns, u_m, max_u = 0, 0.0, 0.0
    for i in range(1, len(P) - 1):
        a = np.asarray(P[i - 1], float)
        b = np.asarray(P[i], float)
        c = np.asarray(P[i + 1], float)
        v1, v2 = b - a, c - b
        n1, n2 = float(np.linalg.norm(v1)), float(np.linalg.norm(v2))
        if n1 < 30.0 or n2 < 30.0:
            continue
        cosang = float((v1 / n1) @ (v2 / n2))
        ang = math.degrees(math.acos(max(-1.0, min(1.0, cosang))))
        if ang > 120.0:
            u_turns += 1
            u_m += n2
            max_u = max(max_u, n2)
    return {"u_turns": u_turns, "u_m": u_m, "u_max_m": max_u, "n_seg": len(P) - 1}


def run_one(cls, cfg, arena, base, want_path=False):
    rb = cls(arena, cfg, base=base, log=None)
    t0 = time.time()
    rep = rb.run()
    rep["_wall"] = time.time() - t0
    st_ = getattr(arena, "stats", {}) or {}
    rep["travel_m"] = float(st_.get("travel_m", 0.0))
    if want_path:
        rep.update(path_metrics(getattr(arena, "actions", [])))
        rep["_actions"] = [(a["x"], a["y"], a.get("kind"), a.get("virtual_time_s"))
                           for a in getattr(arena, "actions", [])
                           if "x" in a and "y" in a]
    return rep


def summarize(reps):
    cl = sum(r["cleared"] for r in reps)
    tot = sum(r["n_sources"] for r in reps)
    tm = [r["avg_clear_time_s"] for r in reps if r["avg_clear_time_s"]]
    idle = []
    for r in reps:
        last = (r["avg_clear_time_s"] or 0) * r["cleared"]
        idle.append((r["virtual_time_s"] - last) / max(r["virtual_time_s"], 1e-9))
    return {
        "clear_ratio": cl / tot if tot else 0.0,
        "cleared_mean": st.mean(r["cleared"] for r in reps),
        "full_clear_cases": sum(1 for r in reps if r["cleared"] == r["n_sources"]),
        "cases": len(reps),
        "avg_clear_time_s": st.mean(tm) if tm else 0.0,
        "vtime_s": st.mean(r["virtual_time_s"] for r in reps),
        "n_measure": st.mean(r["n_measure"] for r in reps),
        "travel_m": st.mean(r["travel_m"] for r in reps),
        "idle": st.mean(idle),
        "u_turns": st.mean(r.get("u_turns", 0) for r in reps),
        "u_m": st.mean(r.get("u_m", 0.0) for r in reps),
        "wall_s": st.mean(r["_wall"] for r in reps),
        "n_rejected": sum(r["n_rejected"] for r in reps),
    }


def fmt_row(name, s):
    return (f"{name:<26}{s['clear_ratio']:>8.3f}{s['cleared_mean']:>8.2f}"
            f"{s['full_clear_cases']:>6}/{s['cases']:<3}"
            f"{s['avg_clear_time_s']:>9.1f}{s['vtime_s']:>9.0f}"
            f"{s['n_measure']:>7.0f}{s['travel_m']:>9.0f}"
            f"{s['idle']:>8.1%}{s['u_turns']:>7.1f}{s['u_m']:>8.0f}"
            f"{s['wall_s']:>7.2f}")


HDR = (f"{'配置':<26}{'清除率':>8}{'平均清除':>8}{'全清':>9}"
       f"{'定位清除/s':>9}{'虚拟/s':>9}{'检测':>7}{'行程/m':>9}"
       f"{'空转':>8}{'折返':>7}{'折返/m':>8}{'实跑':>7}")


# --------------------------------------------------------------------------
def variant_sets(names=None, cfg_over=None):
    """候选配置。键为名字，值为 ``(robot_cls, cfg, extra_kwargs)``。"""
    from dataclasses import replace as _rep
    from p3_tour import TourConfig, TourRobot
    over = cfg_over or {}
    tc = _rep(TourConfig(), **over) if over else TourConfig()
    tag = "E_统一航路" + ("" if not over else "(" + ",".join(
        f"{k}={v}" for k, v in over.items() if k in ("cover_rings",)) + ")")
    A = ("A_基线(扫圈)", SweepRobot, SweepConfig(), {})
    B = ("B_边绕边清", SweepRobot,
         SweepConfig(onlap_clear=True), {})
    C = ("C_边绕边清+开场定点", SweepRobot,
         SweepConfig(onlap_clear=True, initial_locate=True), {})
    D = ("D_期望场贪心", P3Robot, P3Config(), {})
    E = (tag, TourRobot, tc, {})
    sets = {"base": [A], "onlap": [A, B, C, D], "all": [A, D, E],
            "tour": [E], "cmp": [A, D, E]}
    if names:
        keys = [n.strip() for n in names.split(",") if n.strip()]
        out = []
        for k in keys:
            for item in sets.get(k, []):
                out.append(item)
        return out
    return sets["all"]


def main():
    ap = argparse.ArgumentParser()
    ap.add_argument("--cases", type=int, default=30)
    ap.add_argument("--seed", type=int, default=20260913)
    ap.add_argument("--step", type=int, default=100)
    ap.add_argument("--scenario", default="random",
                    choices=["random", "annulus", "live1", "live2", "center"])
    ap.add_argument("--set", default="all")
    ap.add_argument("--only", default="")
    ap.add_argument("--json-out", default="")
    ap.add_argument("--cfg", default="", help='TourConfig 覆盖（JSON），如 {"cover_rings":[[950,9]]}')
    args = ap.parse_args()

    base = load_base()
    seeds = [args.seed + args.step * i for i in range(args.cases)]

    def mk(seed):
        if args.scenario == "random":
            return MockArena(seed=seed)
        if args.scenario == "annulus":
            return outer_annulus(seed=seed)
        if args.scenario == "center":
            return center_cluster(seed=seed)
        if args.scenario == "live1":
            return arena_from(LIVE1, seed=seed)
        return arena_from(LIVE2, seed=seed)

    rows = variant_sets(args.set, cfg_over=(json.loads(args.cfg) if args.cfg else None))
    if args.only:
        keep = {k.strip() for k in args.only.split(",") if k.strip()}
        rows = [r for r in rows if r[0].split("_")[0] in keep or r[0] in keep]

    print(f"场景 {args.scenario}，{args.cases} 局（seed {seeds[0]}..{seeds[-1]}）")
    print(HDR)
    print("-" * len(HDR))
    store = {}
    for name, cls, cfg, kw in rows:
        reps = [run_one(cls, cfg, mk(s), base, want_path=True) for s in seeds]
        s = summarize(reps)
        store[name] = {**s, "per_case": [
            {k: v for k, v in r.items() if not k.startswith("_") and k != "truth"
             and k != "channels" and k != "mock_stats"} for r in reps]}
        print(fmt_row(name, s), flush=True)
    if args.json_out:
        os.makedirs(os.path.dirname(os.path.abspath(args.json_out)), exist_ok=True)
        with open(args.json_out, "w", encoding="utf-8") as f:
            json.dump(store, f, ensure_ascii=False, indent=2, default=str)
        print(f"\n-> {args.json_out}")


if __name__ == "__main__":
    main()

\end{lstlisting}
\subsection{\texttt{src/p3\_adaptive\_bench.py}}
% SHA256 fa6c8b3d564caf86e52848dc6da802ba287b41746c6f65886e26a5d65515d86e
\begin{lstlisting}[language=python,escapeinside={(*@}{@*)}]
"""Frozen paired benchmark; all policies receive fresh copies of each scene."""
from __future__ import annotations
import argparse
import csv
from dataclasses import asdict
import hashlib
import json
from pathlib import Path
import numpy as np
from p3_arena import MockArena
from p3_adaptive import AdaptiveConfig, AdaptiveRobot
from p3_tour import TourConfig, TourRobot
from p3_bench import LIVE1,LIVE2,arena_from,outer_annulus,center_cluster,run_one,load_base


class HashErrorArena(MockArena):
    """Different bounded error model: deterministic per location, not smooth."""
    def env_error(self,x,y):
        key=f'{self._local_seed}:{float(x).hex()}:{float(y).hex()}'.encode()
        value=int.from_bytes(hashlib.sha256(key).digest()[:8],'little')/(2**64-1)
        return (2*value-1)*0.995


def make_scene(kind,seed):
    if kind=='annulus':
        return outer_annulus(seed)
    if kind=='center':
        return center_cluster(seed)
    if kind in ('live1','live2'):
        return arena_from(LIVE1 if kind=='live1' else LIVE2,seed=seed)
    if kind=='hash':
        return HashErrorArena(seed=seed)
    arena=MockArena(seed=seed)
    if kind=='worstrecv':
        for source in arena.sources:
            source.r_recv=1000.0
    elif kind!='random':
        raise ValueError(kind)
    return arena


def statistics(rows):
    good=[r for r in rows if 'error' not in r]
    out={'cases':len(rows),'paired_completed':len(good)}
    if not good:
        return out
    for name in ('tour','adaptive'):
        vals=[r[name] for r in good]
        times=[r['virtual_time_s'] for r in vals]
        out[name]={
            'mean_s':float(np.mean(times)),
            'p90_s':float(np.percentile(times,90)),
            'max_s':float(np.max(times)),
            'last_clear_mean_s':float(np.mean([r['last_clear_s'] for r in vals])),
            'last_clear_per_source_mean_s':float(np.mean([r['avg_clear_time_s'] or 0 for r in vals])),
            'exit_per_source_mean_s':float(np.mean([r['virtual_time_s']/max(r['cleared'],1) for r in vals])),
            'travel_mean_m':float(np.mean([r['travel_m'] for r in vals])),
            'measure_mean':float(np.mean([r['n_measure'] for r in vals])),
            'wall_mean_s':float(np.mean([r['_wall'] for r in vals])),
            'full_clear_cases':sum(r['cleared']==r['n_sources'] for r in vals),
            'cleared':sum(r['cleared'] for r in vals),
            'sources':sum(r['n_sources'] for r in vals),
            'complete_predicate_cases':sum(bool(r['complete_proof']) for r in vals),
            'rejections':sum(r['n_rejected'] for r in vals),
        }
    pairs=np.asarray([[r['tour']['virtual_time_s'],r['adaptive']['virtual_time_s']] for r in good])
    out['improvement_fraction']=float(1-pairs[:,1].mean()/pairs[:,0].mean())
    out['adaptive_faster_cases']=int(np.count_nonzero(pairs[:,1]<pairs[:,0]))
    sampled=np.random.default_rng(121212).integers(0,len(pairs),size=(4000,len(pairs)))
    means=pairs[sampled].mean(axis=1)
    out['paired_bootstrap_95pct']=np.percentile(1-means[:,1]/means[:,0],[2.5,97.5]).tolist()
    return out


def main():
    ap=argparse.ArgumentParser(description=__doc__)
    ap.add_argument('--counts',default='random:50,annulus:20,center:20,live1:15,live2:15,hash:20,worstrecv:20')
    ap.add_argument('--seed',type=int,default=777000)
    ap.add_argument('--step',type=int,default=37)
    ap.add_argument('--out',default=str(Path(__file__).resolve().parents[1] / 'out/p3_adaptive_validation'))
    args=ap.parse_args()
    outdir=Path(args.out);outdir.mkdir(parents=True,exist_ok=True)
    configs={'tour':TourConfig(),'adaptive':AdaptiveConfig()}
    classes={'tour':TourRobot,'adaptive':AdaptiveRobot}
    source_dir=Path(__file__).resolve().parent
    files=['p3_tour.py','p3_homing.py','p3_frontier.py','p3_coverage.py','p3_adaptive.py','p3_arena.py']
    result={'arguments':vars(args),'configs':{k:asdict(v) for k,v in configs.items()},
            'source_sha256':{f:hashlib.sha256((source_dir/f).read_bytes()).hexdigest() for f in files},
            'note':'live1/live2 are reconstructed coordinates with assumed reception radius 1400 m, not live runs; tour complete_proof is only its original grid predicate.',
            'scenarios':{}}
    base=load_base()
    for part in args.counts.split(','):
        kind,count=part.split(':');rows=[]
        for i in range(int(count)):
            seed=args.seed+args.step*i
            row={'seed':seed}
            for name in ('tour','adaptive'):
                arena=make_scene(kind,seed)
                try:
                    r=run_one(classes[name],configs[name],arena,base,want_path=True)
                    r['last_clear_s']=(r['avg_clear_time_s'] or 0)*r['cleared']
                    row[name]={k:v for k,v in r.items() if k not in ('truth','channels','mock_stats','_actions')}
                    if i<3:
                        trace={'scenario':kind,'seed':seed,'policy':name,'actions':arena.actions,'truth':arena.truth()}
                        (outdir/f'{kind}_{seed}_{name}_trace.json').write_text(json.dumps(trace,ensure_ascii=False,indent=2,default=str),encoding='utf-8')
                except Exception as exc:
                    row['error']=f'{name}: {type(exc).__name__}: {exc}'
                    print(kind,seed,row['error'],flush=True)
                    break
            rows.append(row)
            if (i+1)%10==0:
                print(f'{kind}: {i+1}/{count}',flush=True)
        summary=statistics(rows)
        result['scenarios'][kind]={'summary':summary,'rows':rows}
        (outdir/'paired_results.json').write_text(json.dumps(result,ensure_ascii=False,indent=2,default=str),encoding='utf-8')
        print(kind,json.dumps(summary,ensure_ascii=False),flush=True)
    with (outdir/'paired_results.csv').open('w',newline='',encoding='utf-8-sig') as f:
        w=csv.writer(f)
        w.writerow(['scenario','seed','tour_s','adaptive_s','saved_s','tour_cleared','adaptive_cleared','n_sources','adaptive_certified','error'])
        for kind,data in result['scenarios'].items():
            for row in data['rows']:
                if 'error' in row:
                    w.writerow([kind,row['seed'],'','','','','','','',row['error']]);continue
                a,b=row['tour'],row['adaptive']
                w.writerow([kind,row['seed'],a['virtual_time_s'],b['virtual_time_s'],a['virtual_time_s']-b['virtual_time_s'],a['cleared'],b['cleared'],a['n_sources'],b['complete_proof'],''])


if __name__=='__main__':
    main()

\end{lstlisting}
\subsection{\texttt{scripts/run\_p3\_paper.py}}
% SHA256 175b03467b90f7b8735a17845374d432a0e80f2f3b5bd074197c3c08d89e8749
\begin{lstlisting}[language=python,escapeinside={(*@}{@*)}]
#!/usr/bin/env python3
"""Frozen Q3 paired experiment and independent, replay-only audits.

Default execution resumes missing runs; --audit-only never constructs a robot.
Changing this reporting script does not change the experiment definition or
invalidate completed runs. --smoke always writes to a separate smoke directory.
Only NumPy and the standard library are required. See out/paper_q3/README.md.
"""
from __future__ import annotations
import argparse
import ast
import csv
import hashlib
import json
import math
import platform
import shutil
import signal
import sys
import time
from collections import Counter, defaultdict
from dataclasses import asdict, replace
from pathlib import Path

sys.dont_write_bytecode = True
import numpy as np

HERE = Path(__file__).resolve()
Q3 = HERE.parents[1]
SRC = Q3 / 'src'
OUT_DEFAULT = Q3 / 'out/paper_q3'
sys.path.insert(0, str(SRC))
from p3_robot import P3Config, P3Robot, load_base
from p3_sweep import SweepConfig, SweepRobot
from p3_tour import TourConfig, TourRobot
from p3_adaptive import AdaptiveConfig, AdaptiveRobot
from p3_joint import JointConfig, JointRobot
from p3_adaptive_bench import make_scene

SCENARIOS = (('random', 50), ('annulus', 20), ('center', 10), ('hash', 20), ('worstrecv', 20))
POLICIES = (('field', P3Robot, P3Config), ('sweep', SweepRobot, SweepConfig),
            ('tour', TourRobot, TourConfig), ('adaptive', AdaptiveRobot, AdaptiveConfig),
            ('joint', JointRobot, JointConfig))
ABLATIONS = (('joint_no_after_service', {'after_service_scan': 'off'}),
             ('joint_single_plan', {'planner_multistart': False}),
             ('joint_no_cover_polish', {'cover_polish': False}),
             ('joint_no_probability_gate', {'opportunistic_receive_probability': 0.0}))
SCHEMA_VERSION = 'q3-paper-v1'
AUDIT_VERSION = 'q3-independent-audit-v2'
SEED_BASE = 2026091200
BOOTSTRAP_N = 4000
BOOTSTRAP_SEED = 20260912
TOL = 1e-7
R_ARENA, R_RECV_MIN = 1800.0, 1000.0
# Literal physical rules for independent verification; never import simulator
# timing constants, reported switch_s, or reported clocks as calculation inputs.
SPEED, MEASURE_S, SWITCH_S, CLEAR_OK_S, CLEAR_FAIL_S = 5.0, 5.0, 1.0, 5.0, 3.0
POLICY_CLASSES = {'P3Robot', 'SweepRobot', 'TourRobot', 'AdaptiveRobot', 'JointRobot',
                  'HomingRobot', 'FrontierRobot', 'CertifiedCoverageMixin'}
POLICY_FILES = ('p3_robot.py', 'p3_sweep.py', 'p3_tour.py', 'p3_homing.py',
                'p3_frontier.py', 'p3_coverage.py', 'p3_adaptive.py', 'p3_joint.py')


def clean(v):
    if isinstance(v, np.ndarray): return clean(v.tolist())
    if isinstance(v, np.generic): return clean(v.item())
    if isinstance(v, dict): return {str(k): clean(x) for k, x in v.items()}
    if isinstance(v, (list, tuple)): return [clean(x) for x in v]
    if isinstance(v, float) and not math.isfinite(v): return None
    return v


def digest(obj):
    return hashlib.sha256(json.dumps(clean(obj), ensure_ascii=False, sort_keys=True,
                                     separators=(',', ':'), allow_nan=False).encode()).hexdigest()


def sha256(path): return hashlib.sha256(Path(path).read_bytes()).hexdigest()


def atomic_json(path, obj):
    path = Path(path); path.parent.mkdir(parents=True, exist_ok=True)
    tmp = path.with_suffix(path.suffix + '.tmp')
    tmp.write_text(json.dumps(clean(obj), ensure_ascii=False, indent=2, allow_nan=False), encoding='utf-8')
    tmp.replace(path)


def read_json(path): return json.loads(Path(path).read_text(encoding='utf-8'))


def frozen_inputs():
    source = {str(p.relative_to(Q3)): sha256(p) for p in sorted(SRC.glob('*.py'))}
    # Obsolete scratch entry point was never imported by the paper runner.
    source.pop('src/run_p3_paper.py', None)
    files = sorted((Q3/'out/p3_family').glob('*.csv'))
    files += [Q3/'out/p3_family/p2_grid_family.json', Q3/'out/p2_grid/p2_grid_map.csv']
    fields = {str(p.relative_to(Q3)): sha256(p) for p in files if p.exists()}
    configs = {name: asdict(cfg()) for name, _, cfg in POLICIES}
    configs.update({name: asdict(replace(JointConfig(), **override)) for name, override in ABLATIONS})
    return source, fields, clean(configs)


def make_manifest(prior=None):
    source, fields, configs = frozen_inputs()
    definition = {'source_sha256': source, 'field_sha256': fields, 'configs': configs,
                  'scenarios': SCENARIOS, 'seed_base': SEED_BASE, 'seed_step': 97,
                  'scenario_seed_stride': 100000, 'virtual_limit_s': 360000.0,
                  'ablation_random_cases': 20}
    fingerprint = digest(definition)
    if prior:
        # Compare all previously frozen algorithm/field hashes before touching
        # outputs. Runner/reporting changes are deliberately independent.
        combined = {**source, **fields}
        for path, expected in prior.get('source_sha256', {}).items():
            if path in ('scripts/run_p3_paper.py', 'src/run_p3_paper.py'): continue
            if combined.get(path) != expected:
                raise RuntimeError(f'Frozen algorithm/field changed: {path}; use a new --out directory')
        for path, expected in prior.get('field_sha256', {}).items():
            if fields.get(path) != expected: raise RuntimeError(f'Frozen field changed: {path}')
        if prior.get('configs') != configs: raise RuntimeError('Frozen requested configurations changed')
        if prior.get('definition_fingerprint') not in (None, fingerprint):
            raise RuntimeError('Frozen experiment definition changed; use a new --out directory')
    return {'schema_version': SCHEMA_VERSION, 'experiment_id': (prior or {}).get('experiment_id', fingerprint[:20]),
            'definition_fingerprint': fingerprint, 'definition': definition,
            'source_sha256': source, 'field_sha256': fields, 'configs': configs,
            'seed_formula': '2026091200 + 100000*scenario_index + 97*i',
            'scenarios': clean(SCENARIOS), 'policies': [p[0] for p in POLICIES],
            'ablations': dict(ABLATIONS), 'python_at_audit': platform.python_version(),
            'numpy_at_audit': np.__version__, 'platform_at_audit': platform.platform(),
            'execution_runtime_recorded': (prior or {}).get('execution_runtime_recorded', {k: (prior or {}).get(k) for k in ('python', 'numpy')}),
            'audit_version': AUDIT_VERSION, 'audit_runner_sha256': sha256(HERE),
            'bootstrap': {'resamples': BOOTSTRAP_N, 'rng_seed': BOOTSTRAP_SEED,
                          'statistic': '1 - mean(comparison) / mean(reference)',
                          'sampling': 'paired scene indices, percentile interval; fresh seeded generator per statistic'},
            'physical_rules': {'start_xy': [0, 0], 'initial_measuring_channel': 1,
                               'speed_m_s': SPEED, 'measure_s': MEASURE_S, 'switch_s': SWITCH_S,
                               'clear_success_s': CLEAR_OK_S, 'clear_failure_s': CLEAR_FAIL_S,
                               'virtual_limit_s': 360000, 'clear_does_not_change_measuring_channel': True},
            'offline_real_time_policy': {'allowance_s': 1200,
                'historical_runs': 'Recorded policy-construction + run wall time; no real-time watchdog was installed. Virtual limit was 360000 s.',
                'future_runs': '1200 s POSIX timer around construction + run; timeouts retained as error rows.',
                'limitations': 'Wall time is local diagnostic timing, excludes source generation/base-map loading and independent audit; host load is uncontrolled. No formal online evaluation.'},
            'provenance_note': 'Algorithm/field hashes are checked against frozen records. Reporting-script revisions do not rerun algorithms. Earlier runner-hash metadata had been refreshed after reporting edits, so it is not treated as a byte-exact execution-script snapshot.',
            'primary_metric': 'complete exit_s, reported as mean and P90; all status values retained',
            'per_source_metric': 'mean across cases of exit_s / cleared; undefined when cleared=0; defined-value count disclosed',
            'coverage_note': '5 m sampled-disk audit is numerical evidence, not a continuous proof. The 50 m square-cell certificate is separately reconstructed from no_signal actions.'}


def scene_manifest(arena):
    sources = [{k: getattr(s, k) for k in ('channel', 'x', 'y', 'r_recv', 'cone_half', 'dir_deg')} for s in arena.sources]
    error = {'class': type(arena).__name__, 'waves': arena._waves, 'scale': arena._scale,
             'local_w': arena._local_w, 'local_seed': arena._local_seed,
             'probes': {f'{x},{y}': float(arena.env_error(x, y)) for x, y in
                        ((0., 0.), (500., -500.), (-500., 500.), (1200., 0.), (0., -1200.))}}
    scene = {'class': type(arena).__name__, 'seed': int(arena.seed), 'r_arena': float(arena.r_arena), 'sources': sources}
    return {'class': type(arena).__name__, 'seed': int(arena.seed), 'sources': sources,
            'scene_fingerprint': digest(scene), 'error_field': error, 'error_fingerprint': digest(error)}


class TruthGuard:
    def __init__(self, arena):
        self.arena = arena; self.exit_called = False; self.truth_calls = []
        self._truth, self._exit = arena.truth, arena.exit
        arena.truth, arena.exit = self.truth, self.exit
    def truth(self):
        self.truth_calls.append({'virtual_time_s': float(self.arena.virtual_time_s), 'exit_called': self.exit_called})
        if not self.exit_called: raise RuntimeError('truth() called before accepted exit()')
        return self._truth()
    def exit(self):
        result = self._exit()
        if result.get('accepted') is True: self.exit_called = True
        return result


def action_audit(actions, report=None, truth=None, scene=None):
    position, measuring_channel, elapsed, travel = (0., 0.), 1, 0., 0.
    counts = Counter(); records = []; errors = []; maxima = Counter(); last_clear = None
    cleared_channels = set(); entered = False; exited = False
    def compare(label, observed, expected, index=None):
        if observed is None:
            errors.append({'check': label, 'index': index, 'reason': 'missing value'}); return None
        delta = float(observed) - expected
        maxima[label] = max(maxima[label], abs(delta))
        if not math.isfinite(delta) or abs(delta) > TOL:
            errors.append({'check': label, 'index': index, 'observed': observed, 'expected': expected, 'error': delta})
        return delta
    for i, a in enumerate(actions):
        kind = a.get('kind'); before_channel = measuring_channel
        rec = {'index': i, 'kind': kind, 'independent_before_s': elapsed,
               'measuring_channel_before': before_channel}
        if kind == 'enter':
            if entered or exited: errors.append({'check': 'enter order', 'index': i})
            entered = True; counts['enter'] += 1
        elif kind == 'exit':
            if not entered or exited: errors.append({'check': 'exit order', 'index': i})
            exited = True; counts['exit'] += 1
        elif kind in ('measure', 'clear'):
            if not entered or exited: errors.append({'check': 'action outside session', 'index': i})
            xy = (float(a['x']), float(a['y'])); distance = math.dist(position, xy)
            channel = int(a['channel']); switch = 0.0
            if not 1 <= channel <= 20: errors.append({'check': 'channel range', 'index': i})
            if kind == 'measure':
                switch = SWITCH_S if channel != measuring_channel else 0.0
                measuring_channel = channel; service = MEASURE_S
                counts['measure'] += 1; counts['switch'] += int(switch > 0)
                counts[a.get('measure_result', 'invalid_measure_result')] += 1
                if a.get('measure_result') not in ('direction', 'near', 'no_signal'):
                    errors.append({'check': 'measurement result', 'index': i})
                rec['switch_error_s'] = compare('switch', a.get('switch_s'), switch, i)
            else:
                ok = a.get('clear_result') == 'success'
                service = CLEAR_OK_S if ok else CLEAR_FAIL_S
                counts['clear_success' if ok else 'clear_failure'] += 1
                if a.get('clear_result') not in ('success', 'no_target_in_range'):
                    errors.append({'check': 'clear result', 'index': i})
                if ok:
                    if channel in cleared_channels: errors.append({'check': 'duplicate clear success', 'index': i})
                    cleared_channels.add(channel)
            expected_dt = distance / SPEED + switch + service
            elapsed += expected_dt; travel += distance; position = xy
            if kind == 'clear' and a.get('clear_result') == 'success': last_clear = elapsed
            rec.update({'distance_m': distance, 'distance_error_m': compare('distance', a.get('dist_m'), distance, i),
                        'expected_switch_s': switch, 'expected_dt': expected_dt,
                        'dt_error': compare('dt', a.get('dt'), expected_dt, i)})
        else: errors.append({'check': 'unknown action kind', 'index': i, 'kind': kind})
        rec.update({'expected_cumulative': elapsed, 'reported_time': a.get('virtual_time_s'),
                    'cumulative_error': compare('cumulative', a.get('virtual_time_s'), elapsed, i),
                    'measuring_channel_after': measuring_channel})
        records.append(rec)
    if not entered or not exited: errors.append({'check': 'missing enter or exit', 'entered': entered, 'exited': exited})
    pieces = {'travel_s': travel / SPEED, 'measure_s': MEASURE_S * counts['measure'],
              'switch_s': SWITCH_S * counts['switch'],
              'clear_s': CLEAR_OK_S * counts['clear_success'] + CLEAR_FAIL_S * counts['clear_failure']}
    compare('cost_piece_sum', sum(pieces.values()), elapsed)
    checks = {}
    if report is not None:
        expected = {'virtual_time_s': elapsed, 'travel_m': travel, 'cleared': counts['clear_success'],
                    'n_clear': counts['clear_success'], 'n_clear_fail': counts['clear_failure']}
        # Policy measurement count can include rejected calls; accepted action
        # count is authoritative and report n_rejected is disclosed separately.
        if report.get('n_rejected', 0) == 0: expected['n_measure'] = counts['measure']
        if scene: expected['n_sources'] = len(scene['sources'])
        if counts['clear_success']: expected['avg_clear_time_s'] = last_clear / counts['clear_success']
        for name, value in expected.items():
            checks[name] = compare('report.' + name, report.get(name), value)
        expected_stats = {'travel_m': travel, 'n_measure': counts['measure'],
                          'n_clear': counts['clear_success'] + counts['clear_failure'],
                          'n_clear_fail': counts['clear_failure'], 'n_direction': counts['direction'],
                          'n_near': counts['near'], 'n_no_signal': counts['no_signal']}
        for name, value in expected_stats.items():
            if name in report.get('mock_stats', {}):
                checks['mock_stats.' + name] = compare('stats.' + name, report['mock_stats'][name], value)
    if truth is not None:
        compare('truth.n_cleared', truth.get('n_cleared'), len(cleared_channels))
        if scene: compare('truth.n_sources', truth.get('n_sources'), len(scene['sources']))
        actual = {int(s['channel']) for s in truth.get('sources', []) if s.get('cleared')}
        if actual != cleared_channels: errors.append({'check': 'truth cleared channels', 'truth': sorted(actual), 'actions': sorted(cleared_channels)})
    return {'version': AUDIT_VERSION, 'records': records, 'max_abs_dt_error_s': maxima['dt'],
            'max_abs_cumulative_error_s': maxima['cumulative'], 'max_abs_distance_error_m': maxima['distance'],
            'max_abs_switch_error_s': maxima['switch'], 'maxima': dict(maxima),
            'cost_pieces': pieces, 'independent_exit_s': elapsed, 'independent_last_clear_s': last_clear,
            'independent_travel_m': travel, 'counts': dict(counts), 'cleared_channels': sorted(cleared_channels),
            'accepted_action_count': counts['measure'] + counts['clear_success'] + counts['clear_failure'],
            'report_checks': checks, 'errors': errors, 'audit_ok': not errors}


def flat_metrics(report, audit, scene, wall_s, status='ok', error=None, completion_kind='unknown'):
    n = len(scene['sources']); counts = audit['counts']; cleared = counts.get('clear_success', 0)
    elapsed, last = audit['independent_exit_s'], audit['independent_last_clear_s']; p = audit['cost_pieces']
    return {'exit_s': elapsed, 'last_clear_s': last, 'tail_s': elapsed-last if last is not None else None,
            'exit_per_source_s': elapsed/cleared if cleared else None,
            'last_clear_per_source_s': last/cleared if cleared else None,
            'travel_m': audit['independent_travel_m'], 'n_measure': counts.get('measure', 0),
            'n_switch': counts.get('switch', 0), 'n_clear_success': cleared,
            'n_clear_failure': counts.get('clear_failure', 0),
            't_travel_s': p['travel_s'], 't_measure_s': p['measure_s'],
            't_switch_s': p['switch_s'], 't_clear_s': p['clear_s'],
            'cleared': cleared, 'n_sources': n, 'clear_ratio': cleared/n if n else None,
            'full_clear': bool(n > 0 and cleared == n), 'completion_kind': completion_kind,
            'complete_proof': report.get('complete_proof'), 'wall_s': wall_s,
            'wall_within_1200s': wall_s <= 1200 if wall_s is not None else None,
            'n_rejected': report.get('n_rejected', 0), 'status': status, 'error': error,
            'audit_ok': audit['audit_ok']}


def row_from_trace(t, path):
    return {k: t[k] for k in ('split', 'scenario', 'index', 'seed', 'policy')} | {'trace_file': str(path)} | t['metrics']


def run_case(kind, index, seed, policy, cls, cfg, base, outdir, manifest, split='main'):
    arena = make_scene(kind, seed); scene = scene_manifest(arena); guard = TruthGuard(arena)
    forced_exit = False; status = 'ok'; error = None; robot = None; report = {}
    def timeout_handler(_signum, _frame): raise TimeoutError('1200 s real-time offline allowance exceeded')
    old_handler = signal.signal(signal.SIGALRM, timeout_handler)
    start = time.perf_counter(); signal.setitimer(signal.ITIMER_REAL, 1200.0)
    try:
        robot = cls(arena, cfg, base=base, log=None); report = robot.run()
    except Exception as exc:
        status = 'error'; error = f'{type(exc).__name__}: {exc}'
    finally:
        signal.setitimer(signal.ITIMER_REAL, 0); signal.signal(signal.SIGALRM, old_handler)
        wall = time.perf_counter() - start
        if not guard.exit_called:
            forced_exit = True
            try: arena.exit()
            except Exception as exc: error = f'{error or ""}; exit: {exc}'
    truth = arena.truth() if guard.exit_called else None
    completion = policy if policy in ('field', 'sweep', 'tour') else 'continuous_cell'
    audit = action_audit(arena.actions, report or None, truth, scene)
    metrics = flat_metrics(report, audit, scene, wall, status, error, completion)
    name = f'{split}_{kind}_{index:03d}_{seed}_{policy}'
    path = Path('traces') / f'{name}.json'
    trace = {'schema_version': SCHEMA_VERSION, 'run_id': name, 'split': split, 'scenario': kind,
             'index': index, 'seed': seed, 'policy': policy, 'scene_manifest': scene,
             'requested_config': asdict(cfg), 'effective_config': asdict(robot.cfg if robot else cfg),
             'actions': arena.actions, 'post_exit_truth': truth, 'report': report,
             'action_audit': audit, 'guard': {'exit_called': guard.exit_called, 'truth_calls': guard.truth_calls,
                                           'forced_harness_exit': forced_exit}, 'metrics': metrics,
             'provenance': {'experiment_id': manifest['experiment_id'],
                 'definition_fingerprint': manifest['definition_fingerprint'],
                 'source_sha256': manifest['source_sha256'], 'field_sha256': manifest['field_sha256'],
                 'execution_runner_sha256': sha256(HERE)}}
    atomic_json(outdir/path, trace)
    return row_from_trace(trace, path)


class CoverageAuditor:
    """Grid verification using only per-channel raw no_signal observations."""
    def __init__(self, step=5.0):
        self.step = step
        xs = np.arange(-1800., 1800. + step/2, step)
        gx, gy = np.meshgrid(xs, xs); mask = gx*gx + gy*gy <= 1800.**2
        self.px, self.py = gx[mask], gy[mask]
        self.cache = {}; self.hits = 0; self.misses = 0; self.point_sets = {}
        # Independent reconstruction of the 50 m cell certificate also keeps
        # boundary-intersecting cells whose centers lie outside the disk.
        h = 50.; cx = -1800. + (np.arange(72) + .5)*h; qx, qy = np.meshgrid(cx, cx)
        m = np.maximum(np.abs(qx)-h/2, 0)**2 + np.maximum(np.abs(qy)-h/2, 0)**2 <= 1800.**2 + 1e-6
        self.qx, self.qy = qx[m], qy[m]; self.cell_radius = 1000.-h/math.sqrt(2)-1e-6
    def set_audit(self, points):
        # Sorting/deduplicating lets absent channels and unchanged variants
        # share a single distance computation without dropping any observation.
        points = tuple(sorted(set(points))); key = digest(points)
        if key in self.cache: self.hits += 1; return key, self.cache[key]
        self.misses += 1; self.point_sets[key] = [list(x) for x in points]
        dmin2 = np.full(len(self.px), np.inf); cell_min2 = np.full(len(self.qx), np.inf)
        for x, y in points:
            # Chunking bounds scratch memory; no SciPy nearest-neighbor code.
            for start in range(0, len(self.px), 32768):
                stop = start + 32768
                d2 = (self.px[start:stop]-x)**2 + (self.py[start:stop]-y)**2
                np.minimum(dmin2[start:stop], d2, out=dmin2[start:stop])
            np.minimum(cell_min2, (self.qx-x)**2 + (self.qy-y)**2, out=cell_min2)
        worst = int(np.argmax(dmin2)); unc = int(np.count_nonzero(dmin2 > 1000.**2))
        cunc = int(np.count_nonzero(cell_min2 > self.cell_radius**2))
        result = {'max_nearest_m': math.sqrt(dmin2[worst]) if points else None,
                  'worst_xy': [float(self.px[worst]), float(self.py[worst])] if points else None,
                  'uncovered_points': unc, 'no_signal_stations': len(points),
                  'uncovered_certificate_cells': cunc, 'cell_certificate': cunc == 0}
        self.cache[key] = result
        return key, result
    def audit(self, trace):
        acts = trace['actions']; cleared = {int(a['channel']) for a in acts if a.get('kind') == 'clear' and a.get('clear_result') == 'success'}
        channels = {}
        for ch in range(1, 21):
            if ch in cleared: channels[str(ch)] = {'status': 'cleared'}; continue
            obs = [a for a in acts if a.get('kind') == 'measure' and int(a['channel']) == ch]
            if any(a['measure_result'] in ('direction', 'near') for a in obs):
                channels[str(ch)] = {'status': 'positive_pending'}; continue
            points = [(float(a['x']), float(a['y'])) for a in obs if a['measure_result'] == 'no_signal']
            key, result = self.set_audit(points)
            channels[str(ch)] = {'status': 'covered' if result['uncovered_points'] == 0 else 'uncovered',
                                 'observation_set_sha256': key, **result}
        return {'grid_step_m': self.step, 'grid_points_inside': len(self.px), 'channels': channels,
                'claim_5m_grid': all(c['status'] in ('cleared', 'covered') for c in channels.values()),
                'independent_cell_certificate': all(c['status'] == 'cleared' or c.get('cell_certificate', False) for c in channels.values()),
                'certificate_grid_cells': len(self.qx), 'certificate_radius_m': self.cell_radius,
                'note': '5 m point-grid check is numerical evidence; the separate 50 m cell reconstruction implements a conservative geometric certificate.'}


def finite(v): return isinstance(v, (int, float)) and not isinstance(v, bool) and math.isfinite(v)


def bootstrap(a, b):
    if not a: return {'n': 0, 'saving': None, 'ci95': [None, None]}
    a, b = np.asarray(a, float), np.asarray(b, float)
    if a.mean() <= 0: return {'n': len(a), 'saving': None, 'ci95': [None, None]}
    rng = np.random.default_rng(BOOTSTRAP_SEED)
    idx = rng.integers(0, len(a), size=(BOOTSTRAP_N, len(a)))
    amean, bmean = a[idx].mean(axis=1), b[idx].mean(axis=1)
    if np.any(amean <= 0): raise ValueError('Undefined bootstrap relative saving due to zero reference mean')
    vals = 1-bmean/amean
    return {'n': len(a), 'saving': float(1-b.mean()/a.mean()), 'ci95': np.percentile(vals, [2.5, 97.5]).tolist()}


def aggregate(rows):
    def values(key): return [r[key] for r in rows if finite(r.get(key))]
    def mean(key):
        v = values(key); return float(np.mean(v)) if v else None
    times = values('exit_s'); sources = sum(r.get('n_sources', 0) for r in rows)
    out = {'cases': len(rows), 'completed': sum(r['status'] == 'ok' for r in rows),
           'errors': sum(r['status'] != 'ok' for r in rows), 'all_statuses_included': True,
           'full_clear_cases': sum(bool(r['full_clear']) for r in rows),
           'cleared': sum(r.get('cleared', 0) for r in rows), 'sources': sources,
           'clear_ratio_source_weighted': sum(r.get('cleared', 0) for r in rows)/sources if sources else None,
           'zero_cleared_cases': sum(r.get('cleared', 0) == 0 for r in rows),
           'primary_metric_defined_cases': len(times),
           'undefined_primary_cases': len(rows)-len(times),
           'per_source_metric_defined_cases': len(values('exit_per_source_s')),
           'undefined_per_source_cases': len(rows)-len(values('exit_per_source_s')),
           'p90_exit_s': float(np.percentile(times, 90)) if times else None,
           'max_exit_s': max(times) if times else None,
           'wall_max_s': max(values('wall_s'), default=None),
           'rejections': sum(r.get('n_rejected', 0) for r in rows),
           'completion_kinds': dict(Counter(r['completion_kind'] for r in rows)),
           'claimed_complete_cases': sum(r.get('complete_proof') is True for r in rows),
           'audit_failure_cases': sum(not r.get('audit_ok', False) for r in rows)}
    mapping = {'mean_exit_s': 'exit_s', 'mean_exit_per_source_s': 'exit_per_source_s',
               'mean_last_clear_s': 'last_clear_s', 'mean_last_clear_per_source_s': 'last_clear_per_source_s',
               'mean_tail_s': 'tail_s', 'mean_travel_m': 'travel_m', 'mean_measure': 'n_measure',
               'mean_switch': 'n_switch', 'mean_clear_failure': 'n_clear_failure', 'wall_mean_s': 'wall_s',
               'mean_t_travel_s': 't_travel_s', 'mean_t_measure_s': 't_measure_s',
               'mean_t_switch_s': 't_switch_s', 'mean_t_clear_s': 't_clear_s'}
    out.update({dst: mean(src) for dst, src in mapping.items()})
    out['metric_defined_counts'] = {src: len(values(src)) for src in mapping.values()}
    return out


def paired(reference, comparison):
    ref = {(r['scenario'], r['seed']): r for r in reference}
    cmp = {(r['scenario'], r['seed']): r for r in comparison}
    keys = sorted(ref.keys() & cmp.keys()); pairs = [(ref[k], cmp[k]) for k in keys]
    def one(metric):
        keep = [(a, b) for a, b in pairs if finite(a.get(metric)) and finite(b.get(metric))]
        result = bootstrap([a[metric] for a, b in keep], [b[metric] for a, b in keep])
        result.update({'candidate_pairs': len(pairs), 'undefined_metric_pairs': len(pairs)-len(keep),
                       'non_ok_pairs_included': sum(a['status'] != 'ok' or b['status'] != 'ok' for a, b in keep),
                       'incomplete_clear_pairs_included': sum(not a['full_clear'] or not b['full_clear'] for a, b in keep),
                       'comparison_faster_cases': sum(b[metric] < a[metric] for a, b in keep)})
        return result
    out = one('exit_s'); out['per_source'] = one('exit_per_source_s')
    out['reference_only_cases'] = len(ref.keys()-cmp.keys()); out['comparison_only_cases'] = len(cmp.keys()-ref.keys())
    out['convention'] = 'positive saving means comparison faster; include failures/incomplete-clear pairs when numeric metrics exist'
    return out


def summarize(rows):
    main = [r for r in rows if r['split'] == 'main']; variants = [r for r in rows if r['split'] == 'ablation']
    def group(rs):
        policies = sorted(set(r['policy'] for r in rs)); out = {p: aggregate([r for r in rs if r['policy'] == p]) for p in policies}
        out['paired_vs_field'] = {p: paired([r for r in rs if r['policy'] == 'field'], [r for r in rs if r['policy'] == p]) for p in policies if p != 'field'}
        out['joint_vs'] = {p: paired([r for r in rs if r['policy'] == p], [r for r in rs if r['policy'] == 'joint']) for p in policies if p != 'joint'}
        return out
    out = {'schema_version': SCHEMA_VERSION, 'counts': {'rows': len(rows), 'main': len(main), 'ablation': len(variants),
           'status': dict(Counter(r['status'] for r in rows))}, 'overall': group(main),
           'scenarios': {k: group([r for r in main if r['scenario'] == k]) for k, _ in SCENARIOS if any(r['scenario'] == k for r in main)},
           'ablations': {}, 'bootstrap': {'resamples': BOOTSTRAP_N, 'seed': BOOTSTRAP_SEED},
           'failure_policy': 'All runs contribute to aggregates, including error-status and incomplete-clear runs. Undefined values are null and counts are disclosed; never replace zero-cleared denominator with one.'}
    if variants:
        out['ablations']['random'] = {p: aggregate([r for r in variants if r['policy'] == p]) for p, _ in ABLATIONS}
        out['ablations']['random']['joint_baseline'] = aggregate([r for r in main if r['scenario'] == 'random' and r['index'] < 20 and r['policy'] == 'joint'])
        out['ablations']['random']['paired_vs_joint'] = {p: paired([r for r in main if r['scenario'] == 'random' and r['index'] < 20 and r['policy'] == 'joint'], [r for r in variants if r['policy'] == p]) for p, _ in ABLATIONS}
    return out


def static_policy_audit():
    accesses = []; suspicious = []; truth = []; dynamic = []
    blocked = {'sources', 'seed', 'rng', 'env_error', '_src_of', '_waves', '_scale', '_local_w', '_local_seed', '__dict__', '__getattribute__'}
    methods = []
    for filename in POLICY_FILES:
        root = ast.parse((SRC/filename).read_text(encoding='utf-8'))
        for cls in root.body:
            if not isinstance(cls, ast.ClassDef) or cls.name not in POLICY_CLASSES: continue
            for method in cls.body:
                if not isinstance(method, (ast.FunctionDef, ast.AsyncFunctionDef)): continue
                methods.append({'file': 'src/'+filename, 'class': cls.name, 'method': method.name, 'line': method.lineno})
                for node in ast.walk(method):
                    entry = {'file': 'src/'+filename, 'class': cls.name, 'method': method.name,
                             'line': getattr(node, 'lineno', None)}
                    if isinstance(node, ast.Attribute):
                        expr = ast.unparse(node)
                        if '.arena.' in expr:
                            accesses.append({**entry, 'expression': expr})
                            if node.attr in blocked: suspicious.append({**entry, 'expression': expr})
                            if node.attr == 'truth': truth.append({**entry, 'expression': expr})
                    if isinstance(node, ast.Call) and isinstance(node.func, ast.Name) and node.func.id in ('getattr', 'setattr', 'hasattr'):
                        if node.args and 'arena' in ast.unparse(node.args[0]):
                            expr = ast.unparse(node); dynamic.append({**entry, 'expression': expr})
                            if len(node.args) < 2 or not isinstance(node.args[1], ast.Constant) or node.args[1].value in blocked:
                                suspicious.append({**entry, 'expression': expr})
    allowed_truth = all(x['class'] == 'P3Robot' and x['method'] == 'report' for x in truth)
    return {'audit_ok': not suspicious and allowed_truth, 'files': list(POLICY_FILES),
            'classes': sorted(POLICY_CLASSES), 'method_count': len(methods),
            'decision_method_count': sum(x['method'] != 'report' for x in methods),
            'methods': methods, 'arena_attribute_accesses': accesses,
            'dynamic_attribute_accesses': dynamic, 'sensitive_hits': suspicious,
            'truth_accesses': truth, 'truth_only_in_report': allowed_truth,
            'manual_crosscheck': 'Independent reviewer checked 147 decision methods plus report, module-level benchmark imports, dynamic getattr/eval/exec/vars, and p3_expect_field/p1_intersection/p2_grid_expectation helpers. No protected arena state found in active decisions. LIVE1/LIVE2 exist as module imports in p3_frontier but are not referenced in policy methods.',
            'normal_exit_before_report_lines': {'p3_robot.py': [1189,1193], 'p3_sweep.py': [1086,1090], 'p3_tour.py': [1522,1527]},
            'scope': 'AST inspection of decision classes and all their methods; excludes module-level mock tests/scene generators. Runtime truth guard independently checks actual exit order.',
            'limitation': 'Static inspection and endpoint guard are not a language-level sandbox; reflective/dynamic evasion is not claimed impossible.'}


def audit_selfcheck():
    # A clear on channel 8 must leave the detector on its initial channel 1.
    actions = [{'kind': 'enter', 'virtual_time_s': 0.},
        {'kind': 'clear', 'x': 0., 'y': 0., 'channel': 8, 'clear_result': 'no_target_in_range', 'dt': 3., 'dist_m': 0., 'virtual_time_s': 3.},
        {'kind': 'measure', 'x': 15., 'y': 20., 'channel': 1, 'measure_result': 'no_signal', 'switch_s': 0., 'dt': 10., 'dist_m': 25., 'virtual_time_s': 13.},
        {'kind': 'measure', 'x': 15., 'y': 20., 'channel': 2, 'measure_result': 'no_signal', 'switch_s': 1., 'dt': 6., 'dist_m': 0., 'virtual_time_s': 19.},
        {'kind': 'exit', 'virtual_time_s': 19.}]
    checks = {'valid_clear_preserves_channel': action_audit(actions)['audit_ok']}
    for key in ('switch_s', 'dist_m', 'dt', 'virtual_time_s'):
        corrupt = [dict(a) for a in actions]; corrupt[2][key] += 1
        checks['detect_corrupted_'+key] = not action_audit(corrupt)['audit_ok']
    corrupt = [dict(a) for a in actions]
    for a in corrupt[2:]: a['virtual_time_s'] += 1
    checks['detect_cumulative_offset'] = not action_audit(corrupt)['audit_ok']
    m = flat_metrics({}, action_audit(actions), {'sources': [{}]}, .1, 'error', 'synthetic failure', 'field')
    checks['zero_clear_is_undefined'] = m['exit_per_source_s'] is None and m['last_clear_s'] is None
    agg = aggregate([m]); checks['failed_run_retained'] = agg['cases'] == 1 and agg['errors'] == 1 and agg['mean_exit_s'] == 19.
    if not all(checks.values()): raise AssertionError(checks)
    return checks


def raw_payload(t):
    keys = ('run_id', 'split', 'scenario', 'index', 'seed', 'policy', 'scene_manifest',
            'requested_config', 'effective_config', 'actions', 'post_exit_truth', 'truth_error', 'report', 'guard')
    return {k: t.get(k) for k in keys} | {'execution_wall_s': t['metrics'].get('wall_s')}


def archive_scratch(outdir, rows):
    scratch = Q3.parents[1] / 'tmp/q3-paper-scratch'
    scratch.mkdir(parents=True, exist_ok=True)
    refs = {(outdir/r['trace_file']).resolve() for r in rows}
    moved = []
    for path in sorted((outdir/'traces').glob('*.json')) + [outdir/'metadata.json']:
        if not path.exists() or path.resolve() in refs: continue
        relative = path.relative_to(outdir); dest = scratch/'previous-products'/relative
        dest.parent.mkdir(parents=True, exist_ok=True)
        if dest.exists(): dest = dest.with_name(dest.stem + '-' + sha256(path)[:12] + dest.suffix)
        if dest.exists(): raise RuntimeError(f'Archive destination already exists: {dest}')
        value = sha256(path); path.rename(dest)
        moved.append({'from': str(relative), 'to': str(dest), 'sha256': value})
    return moved


def snapshot_sources(outdir, manifest):
    hashes = {**manifest['source_sha256'], **manifest['field_sha256'], 'scripts/run_p3_paper.py': sha256(HERE)}
    for relative, expected in hashes.items():
        source = Q3/relative; target = outdir/'source_snapshot'/relative
        target.parent.mkdir(parents=True, exist_ok=True)
        if sha256(source) != expected: raise RuntimeError(f'Input changed during audit: {source}')
        if not target.exists() or sha256(target) != expected: shutil.copyfile(source, target)
    atomic_json(outdir/'source_snapshot/SHA256.json', hashes)
    return len(hashes)


def write_results(outdir, rows, manifest):
    atomic_json(outdir/'paired_results.json', {'schema_version': SCHEMA_VERSION, 'experiment_id': manifest['experiment_id'],
        'definition_fingerprint': manifest['definition_fingerprint'], 'manifest_file': 'manifest.json', 'rows': rows})
    columns = sorted({k for r in rows for k in r})
    with (outdir/'paired_results.csv').open('w', encoding='utf-8', newline='') as stream:
        writer = csv.DictWriter(stream, fieldnames=columns); writer.writeheader(); writer.writerows(rows)


def audit_all(outdir, rows, manifest):
    start = time.perf_counter(); auditor = CoverageAuditor(); audits = []; all_checks = []; new_rows = []
    fingerprints = defaultdict(set); scene_checks = []; configs_bad = []; max_errors = Counter(); raw_ok = True
    static = static_policy_audit(); selfcheck = audit_selfcheck(); prior = outdir/'audit_report.json'
    prior_report = read_json(prior) if prior.exists() else {}
    moved = archive_scratch(outdir, rows)
    for n, row in enumerate(rows):
        path = outdir/row['trace_file']; trace = read_json(path); raw_before = digest(raw_payload(trace))
        # Recompute from the raw action stream, never from the previous audit.
        audit = action_audit(trace['actions'], trace.get('report'), trace.get('post_exit_truth'), trace['scene_manifest'])
        old = trace['metrics']; trace['action_audit'] = audit
        trace['metrics'] = flat_metrics(trace.get('report') or {}, audit, trace['scene_manifest'], old.get('wall_s'),
                                       old.get('status', 'error'), old.get('error'), old['completion_kind'])
        scene = trace['scene_manifest']; fingerprints[(row['scenario'], row['seed'])].add((scene['scene_fingerprint'], scene['error_fingerprint']))
        scene_obj = {'class': scene['class'], 'seed': scene['seed'], 'r_arena': 1800.0, 'sources': scene['sources']}
        # Historical fingerprint serialization used ASCII default; values here
        # contain ASCII-only strings so this is byte-equivalent.
        scene_ok = digest(scene_obj) == scene['scene_fingerprint'] and digest(scene['error_field']) == scene['error_fingerprint']
        scene_checks.append(scene_ok)
        requested = trace['requested_config']; cfg_ok = requested == manifest['configs'][row['policy']]
        effective = dict(requested)
        if old['completion_kind'] == 'continuous_cell': effective.update({'cover_radius_m': 1000.-50./math.sqrt(2)-1e-6, 'cover_slack_m': 0.0})
        effective_ok = trace['effective_config'] == effective
        if not cfg_ok or not effective_ok: configs_bad.append({'trace_file': row['trace_file'], 'requested_ok': cfg_ok, 'effective_ok': effective_ok})
        exits = [a for a in trace['actions'] if a.get('kind') == 'exit']
        calls = trace.get('guard', {}).get('truth_calls', [])
        guard_ok = bool(exits and calls and trace.get('guard', {}).get('exit_called') and
                        all(c.get('exit_called') and c['virtual_time_s'] >= exits[0]['virtual_time_s']-TOL for c in calls))
        trace['audit_provenance'] = {'audit_version': AUDIT_VERSION, 'runner_sha256': sha256(HERE),
            'raw_payload_sha256': raw_before, 'definition_fingerprint': manifest['definition_fingerprint'],
            'execution_algorithms_unchanged': True,
            'verified_source_sha256': manifest['source_sha256'],
            'verified_field_sha256': manifest['field_sha256']}
        raw_after = digest(raw_payload(trace)); raw_ok &= raw_before == raw_after
        trace['audit_provenance']['raw_payload_unchanged'] = raw_before == raw_after
        check = {'trace_file': row['trace_file'], 'policy': row['policy'], 'scenario': row['scenario'], 'seed': row['seed'],
                 'raw_payload_sha256': raw_before, 'raw_unchanged': raw_before == raw_after,
                 'action_audit_ok': audit['audit_ok'], 'truth_guard_ok': guard_ok, 'scene_fingerprint_ok': scene_ok,
                 'requested_config_ok': cfg_ok, 'effective_config_ok': effective_ok,
                 'errors': audit['errors'], 'maxima': audit['maxima']}
        all_checks.append(check)
        for k, value in audit['maxima'].items(): max_errors[k] = max(max_errors[k], value)
        if old['completion_kind'] == 'continuous_cell' and trace['report'].get('complete_proof') is True:
            coverage = auditor.audit(trace)
            audits.append({'scenario': row['scenario'], 'index': row['index'], 'seed': row['seed'], 'policy': row['policy'],
                           'split': row['split'], 'trace_file': row['trace_file'], 'audit': coverage})
            trace['independent_coverage_audit'] = {'claim_5m_grid': coverage['claim_5m_grid'],
                'cell_certificate': coverage['independent_cell_certificate'], 'details_file': 'coverage_audit.json'}
        atomic_json(path, trace); new_rows.append(row_from_trace(trace, row['trace_file']))
        if (n+1) % 50 == 0: print(f'Audited {n+1}/{len(rows)} traces; {len(audits)} certificates, {auditor.misses} unique observation sets', flush=True)
    unpaired = [{'scenario': k[0], 'seed': k[1], 'different_fingerprints': len(v)} for k, v in fingerprints.items() if len(v) != 1]
    coverage_result = {'schema_version': SCHEMA_VERSION, 'experiment_id': manifest['experiment_id'], 'audit_version': AUDIT_VERSION,
                       'scope': 'All adaptive/joint/ablation runs claiming continuous-cell completion', 'audits': audits,
                       'counts': {'runs': len(audits), 'numerical_pass': sum(x['audit']['claim_5m_grid'] for x in audits),
                                  'cell_certificate_pass': sum(x['audit']['independent_cell_certificate'] for x in audits),
                                  'unique_observation_sets': auditor.misses, 'observation_cache_hits': auditor.hits},
                       'observation_sets': auditor.point_sets, 'observation_set_results': auditor.cache}
    atomic_json(outdir/'coverage_audit.json', coverage_result)
    summary = summarize(new_rows); summary['experiment_id'] = manifest['experiment_id']; summary['audit'] = coverage_result['counts']
    atomic_json(outdir/'summary.json', summary); write_results(outdir, new_rows, manifest)
    atomic_json(outdir/'manifest.json', manifest)
    snapshot_count = snapshot_sources(outdir, manifest)
    report = {'version': AUDIT_VERSION, 'experiment_id': manifest['experiment_id'], 'runs': len(new_rows),
        'all_raw_payloads_unchanged': raw_ok, 'policy_runs_executed_by_this_audit': 0,
        'action_audit_pass': sum(x['action_audit_ok'] for x in all_checks),
        'truth_guard_pass': sum(x['truth_guard_ok'] for x in all_checks), 'scene_fingerprint_pass': sum(scene_checks),
        'pairing_mismatches': unpaired, 'configuration_mismatches': configs_bad, 'max_errors': dict(max_errors),
        'coverage': coverage_result['counts'], 'static_policy_audit': static, 'audit_contract_checks': selfcheck,
        'source_snapshot_files': snapshot_count, 'source_snapshot_index': 'source_snapshot/SHA256.json',
        'archive_movements': prior_report.get('archive_movements', []) + moved,
        'per_run_checks': all_checks, 'audit_wall_s': time.perf_counter()-start}
    atomic_json(outdir/'audit_report.json', report)
    write_readme_and_evidence(outdir, summary, report, manifest)
    print(json.dumps({'rows': len(new_rows), 'action_pass': report['action_audit_pass'], 'truth_pass': report['truth_guard_pass'],
                      'coverage': report['coverage'], 'audit_wall_s': report['audit_wall_s']}), flush=True)
    return new_rows


def write_readme_and_evidence(outdir, summary, audit, manifest):
    evidence = {'experiment_id': manifest['experiment_id'], 'claims': [
        {'claim': 'Five policies compared on the same 120 scenes', 'evidence': ['manifest.json#/scenarios', 'paired_results.json#/rows', 'audit_report.json#/pairing_mismatches']},
        {'claim': 'Every reported action obeys independently recomputed physical timing', 'evidence': ['audit_report.json#/max_errors', 'traces/*#/action_audit']},
        {'claim': 'No truth call before accepted exit, and no sensitive arena attributes in policy classes', 'evidence': ['audit_report.json#/truth_guard_pass', 'audit_report.json#/static_policy_audit']},
        {'claim': 'Every continuous-cell completion claim passes 5 m numerical and independent 50 m cell verification', 'evidence': ['coverage_audit.json#/counts', 'coverage_audit.json#/audits'], 'limit': 'Numerical point verification alone is not a proof; finite-sample success is not a universal algorithm guarantee.'},
        {'claim': 'Joint efficiency improvements and uncertainty', 'evidence': ['summary.json#/scenarios', 'summary.json#/overall/joint_vs']},
        {'claim': 'Ablation modules have different effect sizes; after-service scan CI may include zero', 'evidence': ['summary.json#/ablations/random/paired_vs_joint'], 'limit': '20 paired random cases, one module removed at a time; effects need not add linearly.'}],
        'key_results': {'overall': summary.get('overall'), 'ablations': summary.get('ablations')},
        'remaining_limitations': ['Offline mock evidence only; formal online runs remain pending.',
            'Historical execution runner bytes were not preserved before metadata-only edits. Algorithm and field hashes are verified; current audit runner is snapshotted.',
            'Historical real-time allowance was measured but not enforced by a watchdog.',
            'The error field is a bounded synthetic model, with hash-noise stress scenes; these do not cover every physically possible error field.']}
    atomic_json(outdir/'evidence_map.json', evidence)
    text = f'''# Q3 frozen paper experiment

Experiment ID: `{manifest['experiment_id']}`. Definition fingerprint: `{manifest['definition_fingerprint']}`.

The final matrix contains 120 paired scenes × 5 default policies (600 runs), plus
20 random seeds × 4 single-module joint ablations (80 runs). Scenario order is
random 50, annulus 20, center 10, hash 20, worstrecv 20. Seed =
2026091200 + 100000 × scenario index + 97 × within-scenario index.
Every policy uses a fresh scene with the same exact source and error-field fingerprints.

## Reproduce or audit

From the repository root, with Python 3.12 and NumPy 2.3.5 (no SciPy needed):

```sh
python3 -B B_locator/Q3/scripts/run_p3_paper.py --audit-only
python3 -B B_locator/Q3/scripts/run_p3_paper.py --smoke
python3 -B B_locator/Q3/scripts/run_p3_paper.py --out /path/to/new-experiment
```

`--audit-only` only replays existing actions and never constructs/runs a robot.
`--smoke` writes to a separate smoke subdirectory. A new experiment runs serially;
existing finished traces are validated and skipped. Incompatible algorithm/config
hashes raise an error instead of overwriting data. Audit/report script changes do
not invalidate or rerun algorithms. `--max-new-runs N` can bound a partial new run.
The stored source snapshot contains all Q3 Python modules, the current runner,
nine loaded expectation maps, their manifest and the fallback map; SHA256.json
verifies their bytes. The immutable algorithm/field definition excludes reporting
script revisions. The discarded early src/run_p3_paper.py was not an imported
algorithm dependency and is excluded from the final manifest.

## Metric definitions

T = travel_m/5 + 5*N_measure + N_switch + 5*N_clear_success + 3*N_clear_failure.
The initial location is (0,0) and measuring channel is 1. Only a measurement
changes the measuring channel; clear and enter/exit do not. Every distance,
switch, dt and cumulative timestamp is independently recomputed from raw actions.
The cumulative sum never uses a reported timestamp as its next starting value.
Final metrics are derived from this reconstruction and checked against report,
mock statistics and post-exit truth. The paper primary metric is the complete T,
reported as mean and P90. The per-source statistic is mean(T/cleared), with null
for cleared=0 and a disclosed defined-value count. All status values,
including failed and incomplete-clear runs, remain in aggregates. Last-clear time
and tail after last clear are diagnostics and are undefined when no clear succeeds.
The bootstrap uses 4000 paired scene resamples, RNG seed 20260912, fresh for each
comparison. A positive saving means the comparison uses less time than the reference.
Both total time and per-source time have percentile 95% confidence intervals.

The simulator virtual limit is 360000 s. The 1200 s allowance is real computation
time. Historical policy construction/run wall times were recorded without a timer;
all are summarized, but local wall measurements are not official online validation.
Future new runs install a 1200 s timer and retain timeouts as failed cases.

## Coverage and information isolation

The independent 5 m grid audit uses only actual no_signal responses for each
uncleared channel and successful-clear actions to retire channels. Any positive
observation left uncleared blocks completion. Identical observation station sets
are deduplicated and cached. All {audit['coverage']['runs']} claimed continuous-cell
completions are audited ({audit['coverage']['numerical_pass']} numerical passes).
This numerical audit is not a continuous proof. A separate independent 50 m
closed-square-cell reconstruction checks all cells intersecting the disk at
radius 1000 − 50/sqrt(2) − 1e-6; its triangle-inequality argument certifies entire
cells. {audit['coverage']['cell_certificate_pass']} such reconstructions pass.

The runtime truth wrapper logs calls and permits them only after accepted exit.
AST inspection covers every method in the eight decision classes, with sensitive
arena properties and dynamic getattr calls separately listed. Module-level mock
tests and scene-generation helpers are outside decision scope. This is transparent
inspection plus a runtime guard, not a security sandbox.

## Files and provenance

- paired_results.json / .csv: the authoritative 680 rows, each pointing to one trace.
- traces/: only final referenced traces; raw actions/configuration/report/truth are preserved.
- summary.json: all cases, each scenario, all baselines and four paired ablations.
- coverage_audit.json: per-run/per-channel verification and deduplicated observation sets.
- audit_report.json: timing, truth, configuration, static-source, pairing and payload-integrity checks.
- evidence_map.json: paper claims linked to precise evidence and limitations.
- source_snapshot/: verified algorithm, field and current audit-runner bytes.

600 redundant early traces and obsolete metadata are moved (never deleted) to
`tmp/q3-paper-scratch/previous-products/`. Movements and hashes are listed in the
audit report. Historical provenance fields are retained in traces; the separate
audit_provenance section records this audit and canonical raw-payload hashes.
Earlier runner metadata was refreshed after reporting edits, so its hash is not
claimed to prove the exact historical runner bytes. Algorithm/field hashes remain
verified. This audit did not rerun any policy.
'''
    (outdir/'README.md').write_text(text, encoding='utf-8')


def main(argv=None):
    parser = argparse.ArgumentParser(description=__doc__)
    parser.add_argument('--out', default=str(OUT_DEFAULT)); parser.add_argument('--audit-only', action='store_true')
    parser.add_argument('--smoke', action='store_true'); parser.add_argument('--max-new-runs', type=int, default=0)
    args = parser.parse_args(argv); outdir = Path(args.out)
    if args.smoke: outdir = outdir/'smoke'
    (outdir/'traces').mkdir(parents=True, exist_ok=True)
    old_manifest = read_json(outdir/'manifest.json') if (outdir/'manifest.json').exists() else None
    manifest = make_manifest(old_manifest)
    result = read_json(outdir/'paired_results.json') if (outdir/'paired_results.json').exists() else {'rows': []}
    rows = result['rows']; by_key = {(r['split'], r['scenario'], r['index'], r['policy']): r for r in rows}
    if len(by_key) != len(rows): raise RuntimeError('Duplicate result key')
    if args.audit_only:
        if not rows: raise RuntimeError('--audit-only requires existing paired_results.json')
        audit_all(outdir, rows, manifest); return
    atomic_json(outdir/'manifest.json', manifest)
    registry = {p: (cls, cfg) for p, cls, cfg in POLICIES}; jobs = []
    for s, (kind, count) in enumerate(SCENARIOS):
        if args.smoke and s: continue
        for i in range(1 if args.smoke else count):
            for p, cls, cfg in POLICIES: jobs.append(('main', kind, i, SEED_BASE+100000*s+97*i, p, cls, cfg()))
    for i in range(1 if args.smoke else 20):
        for p, override in ABLATIONS: jobs.append(('ablation', 'random', i, SEED_BASE+97*i, p, JointRobot, replace(JointConfig(), **override)))
    new = 0; base = None
    for split, kind, i, seed, p, cls, cfg in jobs:
        key = split, kind, i, p
        if key in by_key:
            existing = read_json(outdir/by_key[key]['trace_file'])
            if existing['seed'] != seed or existing['requested_config'] != clean(asdict(cfg)):
                raise RuntimeError(f'Invalid resume metadata: {key}')
            continue
        if args.max_new_runs and new >= args.max_new_runs: break
        if base is None: base = load_base()
        row = run_case(kind, i, seed, p, cls, cfg, base, outdir, manifest, split)
        by_key[key] = row; new += 1; rows = list(by_key.values())
        write_results(outdir, rows, manifest)
        print(f'New run {new}: {kind} {i} {p} {row["status"]}', flush=True)
    audit_all(outdir, list(by_key.values()), manifest)
    print(f'Executed {new} new policy runs', flush=True)


if __name__ == '__main__': main()

\end{lstlisting}
\subsection{\texttt{scripts/make\_paper\_tables.py}}
% SHA256 73a2c4f0972210c8add0b065916ad1cc0f1aa0811593634d6dc0bc8b3e3ea4a1
\begin{lstlisting}[language=python,escapeinside={(*@}{@*)}]
"""Read frozen Q3 rows and emit paper tables, macros and supporting statistics."""
import json
import hashlib
from pathlib import Path
import numpy as np

Q3 = Path(__file__).resolve().parents[1]
DATA = Q3 / 'out/paper_q3'
SECTIONS = Q3.parent / 'paper/sections'
POLICIES = ('field', 'sweep', 'tour', 'adaptive', 'joint')
SCENARIOS = ('random', 'annulus', 'center', 'hash', 'worstrecv')

def main():
    obj = json.loads((DATA/'paired_results.json').read_text())
    rows = obj['rows']
    summary = json.loads((DATA/'summary.json').read_text())
    random = summary['scenarios']['random']
    body=[]
    for p in POLICIES:
        s=random[p]
        body.append(f"{p} & {s['clear_ratio_source_weighted']:.3f} & {s['full_clear_cases']}/50 & {s['mean_exit_s']:.0f} & {s['p90_exit_s']:.0f} & {s['mean_exit_per_source_s']:.0f}" + r' \\')
    (SECTIONS/'q3-core-rows.tex').write_text('\n'.join(body).rstrip().removesuffix(r'\\')+'\n')
    full={p:sum(r['full_clear'] for r in rows if r['split']=='main' and r['policy']==p) for p in POLICIES}
    pair=random['paired_vs_field']['joint']
    base20=[r for r in rows if r['split']=='main' and r['scenario']=='random' and r['policy']=='joint' and r['index']<20]
    values={'QThreeFieldFull':str(full['field']), 'QThreeSweepFull':str(full['sweep']),
            'QThreeSaving':f"{100*pair['saving']:.1f}", 'QThreeSavingLow':f"{100*pair['ci95'][0]:.1f}",
            'QThreeSavingHigh':f"{100*pair['ci95'][1]:.1f}", 'QThreeAblationBase':f"{np.mean([r['exit_s'] for r in base20]):.1f}"}
    for policy in POLICIES:
        title=policy.capitalize()
        first=next(r for r in rows if r['split']=='main' and r['scenario']=='random' and r['index']==0 and r['policy']==policy)
        values[f'QThreeFirst{title}Time']=f"{first['exit_s']:.0f}"
        values[f'QThreeFirst{title}Travel']=f"{first['travel_m']/1000:.2f}"
        values[f'QThreeFirst{title}Measure']=str(first['n_measure'])
        values[f'QThreeRandom{title}Travel']=f"{random[policy]['mean_travel_m']/1000:.2f}"
        values[f'QThreeRandom{title}Measure']=f"{random[policy]['mean_measure']:.1f}"
        values[f'QThreeCenter{title}Time']=f"{summary['scenarios']['center'][policy]['mean_exit_s']:.1f}"
    for suffix, policy in [('After','joint_no_after_service'),('Multi','joint_single_plan'),('Polish','joint_no_cover_polish'),('Prob','joint_no_probability_gate')]:
        values[f'QThreeAblation{suffix}Increase']=f"{-100*summary['ablations']['random']['paired_vs_joint'][policy]['saving']:.2f}"
    (SECTIONS/'q3-numbers.tex').write_text('% Generated from the frozen Q3 data; do not hand edit.\n'+''.join('\\newcommand{\\'+k+'}{'+v+'}\n' for k,v in values.items()))
    abstract=(r'{\heiti 针对问题三}\songti，建立{\heiti 联合覆盖、定位与清除模型}\songti，'
              '以保守定位控制清除风险，以连续覆盖判据验证任务完成，并逐步引入共用航路、沿途定位和逐频道检测调度。'
              '在 120 组场景开展五策略配对实验，另作 80 次单因素消融。随机 50 组中，'
              r'$\textit{joint}$ 全部清除，完整退出时间均值 '
              +f"{random['joint']['mean_exit_s']:.0f}"+r' s，{\heiti 较初代减少 '
              +values['QThreeSaving']+r'\%}\songti。结果表明，联合规划能同时减少行程与检测成本。'+'\n')
    (SECTIONS/'q3-abstract.tex').write_text('% Generated Q3 abstract from frozen data.\n'+abstract)
    lines=[]
    for kind in SCENARIOS:
        for p in POLICIES:
            rs=[r for r in rows if r['split']=='main' and r['scenario']==kind and r['policy']==p]
            s=summary['scenarios'][kind][p]
            lines.append(f"{kind}/{p} & {s['clear_ratio_source_weighted']:.3f} & {s['full_clear_cases']}/{len(rs)} & {s['mean_exit_s']:.1f} & {s['p90_exit_s']:.1f} & {s['mean_exit_per_source_s']:.1f} & {s['mean_travel_m']/1000:.2f} & {s['mean_measure']:.1f} & {s['wall_mean_s']:.3f}"+r' \\')
    (SECTIONS/'q3-supplement-rows.tex').write_text('\n'.join(lines).rstrip().removesuffix(r'\\')+'\n')
    lines=[]
    for p in ('joint_no_after_service','joint_single_plan','joint_no_cover_polish','joint_no_probability_gate'):
        s=summary['ablations']['random'][p]
        pair=summary['ablations']['random']['paired_vs_joint'][p]
        label={'joint_no_after_service':'取消清后扫描','joint_single_plan':'单方案规划','joint_no_cover_polish':'取消测站精修','joint_no_probability_gate':'取消概率门槛'}[p]
        lines.append(f"{label} & {s['mean_exit_s']:.1f} & {-100*pair['saving']:.2f} & [{-100*pair['ci95'][1]:.2f}, {-100*pair['ci95'][0]:.2f}] & {s['mean_travel_m']/1000:.2f} & {s['mean_measure']:.1f} & {s['wall_mean_s']:.3f}"+r' \\')
    (SECTIONS/'q3-ablation-rows.tex').write_text('\n'.join(lines).rstrip().removesuffix(r'\\')+'\n')
    lines=[]
    for p,label in [('joint_baseline','完整 joint'),('joint_no_after_service','取消清后扫描'),('joint_single_plan','单方案规划'),('joint_no_cover_polish','取消测站精修'),('joint_no_probability_gate','取消概率门槛')]:
        s=summary['ablations']['random'][p]
        lines.append(f"{label} & {s['clear_ratio_source_weighted']:.3f} & {s['full_clear_cases']}/20 & {s['mean_exit_s']:.1f} & {s['p90_exit_s']:.1f} & {s['mean_exit_per_source_s']:.1f}"+r' \\')
    (SECTIONS/'q3-ablation-complete-rows.tex').write_text('\n'.join(lines).rstrip().removesuffix(r'\\')+'\n')
    (DATA/'table_provenance.json').write_text(json.dumps({'experiment_id':obj['experiment_id'], 'data_sha256':hashlib.sha256((DATA/'paired_results.json').read_bytes()).hexdigest(), 'summary_sha256':hashlib.sha256((DATA/'summary.json').read_bytes()).hexdigest(), 'generated':values, 'all_main_full_clear':full},indent=2)+'\n')

if __name__=='__main__':main()

\end{lstlisting}
\subsection{\texttt{scripts/make\_paper\_figures.py}}
% SHA256 2fcc9481ab7a12d1de771f2692b15cf42c82691bcae2006554bdded328b498c3
\begin{lstlisting}[language=python,escapeinside={(*@}{@*)}]
"""Make the compact Q3 paper figures from frozen experiment records.

The script is intentionally a renderer only: it never runs a policy and never
rewrites the frozen JSON.  The runner writes ``out/paper_q3`` using the schema
documented in the module docstring of the evidence runner; this script reads
that directory and writes ``paper/figures/q3-*.pdf`` and ``.png``.

Usage (from the repository root)::

    python B_locator/Q3/scripts/make_paper_figures.py

The five figures are sized for the 160 mm text width of the paper.  Full trajectories remain in the frozen evidence directory.
"""
from __future__ import annotations

import argparse
import json
import math
from pathlib import Path
from typing import Any, Iterable

import numpy as np

import matplotlib
matplotlib.use("Agg")
import matplotlib.pyplot as plt
from matplotlib.lines import Line2D
from matplotlib.patches import Circle, Rectangle


MM = 1.0 / 25.4
FIG_W = 160 * MM
DPI = 330
ROOT = Path(__file__).resolve().parents[2]
DATA = ROOT / "Q3" / "out" / "paper_q3"
OUT = ROOT / "paper" / "figures"

POLICIES = ["field", "sweep", "tour", "adaptive", "joint"]
POLICY_LABEL = {
    "field": "field",
    "sweep": "sweep",
    "tour": "tour",
    "adaptive": "adaptive",
    "joint": "joint",
}
POLICY_COLOR = {
    "field": "#667085", "sweep": "#b7791f", "tour": "#3b82a0",
    "adaptive": "#7656a6", "joint": "#118a8c",
}
SCENE_LABEL = {
    "random": "随机", "annulus": "外圈", "center": "中心",
    "hash": "非平滑误差", "worstrecv": "最小接收半径",
    "live1": "重建布局1", "live2": "重建布局2",
}


def setup_style() -> None:
    """Use paper-safe sizes; all visible text stays >= 8.5 pt."""
    from matplotlib import font_manager
    font = ROOT / "paper" / "fonts" / "simsun.ttc"
    font_manager.fontManager.addfont(str(font))
    matplotlib.rcParams["font.family"] = [font_manager.FontProperties(fname=font).get_name(), "Times New Roman"]
    matplotlib.rcParams["mathtext.fontset"] = "stix"
    matplotlib.rcParams.update({
        "font.size": 9.0,
        "axes.titlesize": 9.2,
        "axes.labelsize": 9.0,
        "xtick.labelsize": 8.5,
        "ytick.labelsize": 8.5,
        "legend.fontsize": 8.5,
        "axes.spines.top": False,
        "axes.spines.right": False,
        "axes.unicode_minus": False,
        "pdf.fonttype": 42,
        "ps.fonttype": 42,
    })


def load_data(data_dir: Path) -> tuple[dict[str, Any], list[dict[str, Any]]]:
    path = data_dir / "paired_results.json"
    if not path.exists():
        raise FileNotFoundError(f"冻结数据不存在：{path}")
    obj = json.loads(path.read_text(encoding="utf-8"))
    rows = obj.get("rows")
    if rows is None:
        # Backward-compatible reader for the older nested benchmark files.
        rows = []
        for scenario, group in obj.get("scenarios", {}).items():
            for old in group.get("rows", []):
                seed = old.get("seed")
                for policy in POLICIES:
                    if policy in old:
                        r = dict(old[policy])
                        r.update({"scenario": scenario, "seed": seed,
                                  "policy": policy, "split": "main"})
                        rows.append(r)
    if not isinstance(rows, list) or not rows:
        raise ValueError("paired_results.json 中没有可绘图的 rows")
    return obj, rows


def val(row: dict[str, Any], *names: str, default: Any = None) -> Any:
    """Look up a scalar in metrics/report/audit and then at row level."""
    buckets: list[Any] = [row.get("metrics"), row.get("report"),
                          row.get("audit"), row]
    for name in names:
        for bucket in buckets:
            if isinstance(bucket, dict) and name in bucket:
                return bucket[name]
    return default


def num(row: dict[str, Any], *names: str, default: float = float("nan")) -> float:
    x = val(row, *names, default=default)
    try:
        return float(x)
    except (TypeError, ValueError):
        return default


def policy_of(row: dict[str, Any]) -> str:
    return str(row.get("policy", row.get("strategy", ""))).lower()


def scene_of(row: dict[str, Any]) -> str:
    return str(row.get("scenario", row.get("scene", "random"))).lower()


def is_ok(row: dict[str, Any]) -> bool:
    return str(row.get("status", "ok")).lower() in {"ok", "complete", "success", ""} and not row.get("error")


def time_value(row: dict[str, Any]) -> float:
    return num(row, "virtual_time_s", "exit_s", "total_time_s", "exit_time_s", "T_s")


def trace_path(row: dict[str, Any], data_dir: Path) -> Path | None:
    p = row.get("trace_file")
    if not p:
        return None
    q = Path(str(p))
    return q if q.is_absolute() else data_dir / q


def first_trace(rows: list[dict[str, Any]], policies: Iterable[str], data_dir: Path,
                scenario: str = "random") -> dict[str, dict[str, Any]]:
    out: dict[str, dict[str, Any]] = {}
    policies = list(policies)
    candidates = [r for r in rows if is_ok(r) and scene_of(r) == scenario and
                  str(r.get("split", "main")) == "main" and int(r.get("index", -1)) == 0]
    # Select one key for which every requested policy has a saved trace.  This
    # prevents silently comparing different seeds when one policy failed.
    groups: dict[tuple[Any, Any], dict[str, dict[str, Any]]] = {}
    for row in candidates:
        p = trace_path(row, data_dir)
        if not p or not p.exists():
            continue
        key = (row.get("index", 0), row.get("seed", 0))
        groups.setdefault(key, {})[policy_of(row)] = row
    for key in sorted(groups, key=lambda z: (int(z[0]), int(z[1]))):
        if all(pol in groups[key] for pol in policies):
            for pol in policies:
                row = groups[key][pol]; p = trace_path(row, data_dir)
                out[pol] = json.loads(p.read_text(encoding="utf-8"))
            break
    return out


def actions(trace: dict[str, Any]) -> list[dict[str, Any]]:
    return [a for a in trace.get("actions", [])
            if isinstance(a, dict) and "x" in a and "y" in a]


def action_points(trace: dict[str, Any]) -> np.ndarray:
    aa = actions(trace)
    return np.asarray([[0.0, 0.0]] + [[float(a["x"]), float(a["y"])] for a in aa], float)


def source_points(trace: dict[str, Any]) -> np.ndarray:
    truth = trace.get("post_exit_truth", trace.get("truth", {})) or {}
    src = truth.get("sources", []) if isinstance(truth, dict) else []
    return np.asarray([[float(s["x_m"]), float(s["y_m"])] for s in src
                      if "x_m" in s and "y_m" in s], float)


def route_legend():
    return [Line2D([], [], color="#d7a06b", marker="o", ls="", ms=3, label="检测点"),
            Line2D([], [], color="#b33e36", marker="x", ls="", ms=4.5, label="成功清除"),
            Line2D([], [], color="#b7791f", marker="^", mfc="none", ls="", ms=4, label="清除失败"),
            Line2D([], [], color="#424b54", marker="*", ls="", ms=5, label="源位（事后）"),
            Line2D([], [], color="#252a35", marker="s", ls="", ms=3.5, label="起点")]


def draw_route(ax, trace: dict[str, Any], color: str, title: str,
               annotate_sources: bool = True, show_ylabel=True) -> None:
    p = action_points(trace)
    if len(p) > 1:
        ax.plot(p[:, 0], p[:, 1], color=color, lw=1.05, zorder=3)
    aa = actions(trace)
    m = np.asarray([[a['x'], a['y']] for a in aa if a.get('kind') == 'measure'])
    if len(m):
        m = np.unique(np.round(m, 2), axis=0)
        ax.scatter(m[:, 0], m[:, 1], s=6, color='#d7a06b', alpha=.6, zorder=4)
    src = source_points(trace)
    if annotate_sources and len(src):
        ax.scatter(src[:, 0], src[:, 1], marker='*', s=26, color='#424b54', zorder=6)
    for success, marker, size in ((True, 'x', 20), (False, '^', 16)):
        q = np.asarray([[a['x'], a['y']] for a in aa if a.get('kind') == 'clear' and
                       (a.get('clear_result') == 'success') == success])
        if len(q):
            if success:
                ax.scatter(q[:, 0], q[:, 1], marker=marker, s=size,
                           color='#b33e36', linewidth=.8, zorder=7)
            else:
                ax.scatter(q[:, 0], q[:, 1], marker=marker, s=size,
                           facecolors='none', edgecolors='#b7791f', linewidth=.8, zorder=8)
    th = np.linspace(0, 2 * math.pi, 400)
    ax.plot(1800*np.cos(th), 1800*np.sin(th), '--', color='#aeb5bd', lw=.65)
    ax.scatter([0], [0], marker='s', s=13, color='#252a35', zorder=9)
    met = trace.get('metrics', {})
    t = met.get('exit_s', trace.get('report', {}).get('virtual_time_s', 0))
    length = met.get('travel_m', trace.get('report', {}).get('travel_m', 0))
    n = met.get('cleared', 0); total = met.get('n_sources', 0)
    ax.set_title(f'{title}\nT={t:.0f} s; L={length/1000:.2f} km\n清除 {n}/{total}',
                 fontsize=9, linespacing=1.25, pad=5)
    ax.set(xlim=(-1950, 1950), ylim=(-1950, 1950), xlabel='x / m')
    if show_ylabel: ax.set_ylabel('y / m')
    else: ax.tick_params(labelleft=False)
    ax.set_xticks([-1000, 0, 1000]); ax.set_yticks([-1000, 0, 1000])
    ax.set_aspect('equal'); ax.grid(alpha=.17, lw=.4)


def save_figure(fig: plt.Figure, stem: str) -> None:
    OUT.mkdir(parents=True, exist_ok=True)
    # Exact print canvas; individual builders reserve their own text margins.
    fig.savefig(OUT / f'{stem}.pdf')
    fig.savefig(OUT / f'{stem}.png', dpi=DPI)
    plt.close(fig)


def fig1_geometry(rows, data_dir) -> None:
    """Two deterministic geometric diagrams, explicitly independent of data."""
    fig, (ax, bx) = plt.subplots(1, 2, figsize=(FIG_W, 57*MM))
    fig.subplots_adjust(left=.02, right=.98, bottom=.04, top=.84, wspace=.16)
    ax.set_title('(a) 单元整体排除：三角不等式', fontsize=9.5)
    ax.set_xlim(-3.9, 1.7); ax.set_ylim(-1.4, 2.05); ax.set_aspect('equal'); ax.axis('off')
    s = np.array([-3., 0.]); q = np.array([0., 0.]); x = np.array([.65, .65])
    ax.add_patch(Circle(s, 3.98, fill=False, lw=1.2, ec='#3b82a0'))
    ax.add_patch(Rectangle((-.65, -.65), 1.3, 1.3, fc='#e3f0f1', ec='#118a8c'))
    ax.plot([s[0], q[0], x[0]], [s[1], q[1], x[1]], color='#118a8c', lw=1.2)
    ax.plot([s[0], x[0]], [s[1], x[1]], ls='--', color='#b7791f', lw=1)
    ax.plot(s[0], s[1], 's', color='#252a35', ms=4)
    ax.plot(q[0], q[1], 'o', color='#118a8c', ms=3)
    ax.plot(x[0], x[1], 'o', color='#b7791f', ms=3)
    ax.text(-2.9, -.60, '$s$：无信号测点', ha='center')
    ax.text(.03, -.34, '$q$', ha='center'); ax.text(.75, .74, '$x$')
    ax.text(-1.55, -.22, r'$\|q-s\|\leq r_{\rm cert}$', ha='center')
    ax.text(.20, 1.13, r'$\|x-q\|\leq h/\sqrt{2}$', ha='center')
    ax.text(-2.95, 1.5, r'$R_{\min}=1000$ m', color='#3b82a0')
    ax.text(-.95, -1.13, r'$r_{\rm cert}=1000-h/\sqrt{2}-\varepsilon$', ha='center')
    bx.set_title('(b) 测站滑移：保留独占覆盖职责', fontsize=9.5)
    bx.set_xlim(-.4, 5.2); bx.set_ylim(-.95, 3.8); bx.set_aspect('equal'); bx.axis('off')
    a=np.array([0., 0.]); b=np.array([4.8, 0.]); old=np.array([2.2, 2.0])
    centers=[np.array([1.4,2.4]),np.array([2.95,2.0])]; radius=1.7
    v=np.array([0.,-2.]); delta=np.array(centers)-old
    bb=-2*delta@v; cc=np.sum(delta**2,axis=1)-radius**2
    tt=min(1., float(np.min((-bb+np.sqrt(bb**2-4*(v@v)*cc))/(2*(v@v)))))
    new=old+tt*v
    for c in centers:
        bx.add_patch(Circle(c,radius,fill=False,ec='#cbd6dc',lw=.7,ls=':'))
        bx.add_patch(Rectangle(c-.12,.24,.24,fc='#bededb',ec='#118a8c',lw=.7))
    bx.plot([a[0], b[0]],[0,0],color='#aeb5bd',lw=.8)
    bx.plot([a[0],old[0],b[0]],[0,old[1],0],'--',color='#b7791f',lw=1)
    bx.plot([a[0],new[0],b[0]],[0,new[1],0],color='#118a8c',lw=1.6)
    bx.annotate('',xy=new,xytext=old,arrowprops=dict(arrowstyle='->',color='#252a35',lw=1))
    bx.plot(old[0],old[1],'o',mfc='white',mec='#b7791f',ms=5)
    bx.plot(new[0],new[1],'o',color='#118a8c',ms=4)
    bx.text(old[0]+.50,old[1]+.50,'原站 $p_0$')
    bx.text(new[0]+.18,new[1]+.14,'新站 $p(t)$')
    bx.text(.25,3.1,'方格：必须保留的单元中心 $Q$')
    bx.text(2.4,-.6,r'$\forall q\in Q:\ \|p(t)-q\|\leq r_{\rm cert}$',ha='center')
    bx.text(0,-.30,'$a$'); bx.text(4.8,-.30,'$b$')
    save_figure(fig,'q3-geometry-certificate')


def route_comparison(rows, data_dir, selected, stem):
    trs=first_trace(rows,POLICIES,data_dir)
    if len(trs)!=5: raise ValueError('固定首局的五策略轨迹必须齐全')
    fig,axes=plt.subplots(1,3,figsize=(FIG_W,65*MM))
    fig.subplots_adjust(left=.09,right=.985,bottom=.25,top=.70,wspace=.12)
    for i,pol in enumerate(selected):
        draw_route(axes[i],trs[pol],POLICY_COLOR[pol],f'({chr(97+i)}) {pol}',show_ylabel=i==0)
    fig.legend(handles=route_legend(),loc='lower center',ncol=5,frameon=False,
               bbox_to_anchor=(.5,.005),columnspacing=.7,handletextpad=.4,handlelength=1.0)
    save_figure(fig,stem)


def fig2_field_sweep_tour(rows,data_dir):
    route_comparison(rows,data_dir,['field','sweep','tour'],'q3-route-field-sweep-tour')


def fig3_tour_adaptive_joint(rows,data_dir):
    route_comparison(rows,data_dir,['tour','adaptive','joint'],'q3-route-tour-adaptive-joint')


def audit_from_actions(trace: dict[str, Any]) -> dict[str, float]:
    aa = actions(trace)
    if not aa:
        return {"travel": 0., "measure": 0., "switch": 0., "clear": 0., "time": 0., "cleared": 0.}
    prev = np.array([0., 0.]); channel = 1; parts = {"travel": 0., "measure": 0., "switch": 0., "clear": 0.}
    cleared = 0
    for a in aa:
        q = np.array([float(a["x"]), float(a["y"])])
        d = float(np.linalg.norm(q - prev)); parts["travel"] += d / 5.; prev = q
        if a.get("kind") == "measure":
            parts["measure"] += 5.; parts["switch"] += float(int(a.get("channel", channel)) != channel)
            channel = int(a.get("channel", channel))
        elif a.get("kind") == "clear":
            ok = a.get("clear_result") == "success"; parts["clear"] += 5. if ok else 3.; cleared += int(ok)
    parts["time"] = sum(parts[k] for k in ("travel", "measure", "switch", "clear"))
    parts["cleared"] = float(cleared)
    return parts


def fig4_cost_and_clear(rows: list[dict[str, Any]], data_dir: Path) -> None:
    main=[r for r in rows if str(r.get('split','main'))=='main' and scene_of(r)=='random' and policy_of(r) in POLICIES]
    fig,(ax,bx)=plt.subplots(1,2,figsize=(FIG_W,70*MM))
    fig.subplots_adjust(left=.13,right=.985,bottom=.25,top=.88,wspace=.45)
    parts=['t_travel_s','t_measure_s','t_switch_s','t_clear_s']
    cols=['#3b82a0','#d97736','#7656a6','#55a868']; labels=['移动','检测','切频','清除']
    means=np.array([[np.mean([num(r,key) for r in main if policy_of(r)==pol]) for key in parts] for pol in POLICIES])
    y=np.arange(5); left=np.zeros(5)
    for k in range(4):
        ax.barh(y,means[:,k],left=left,color=cols[k],label=labels[k],height=.60); left+=means[:,k]
    for i,t in enumerate(left):ax.text(t+60,i,f'{t:.0f}',va='center',fontsize=8.5)
    ax.set_yticks(y,POLICIES);ax.invert_yaxis();ax.set_xlim(0,left.max()*1.18)
    ax.set_xlabel('平均完整时间 / s');ax.set_title('(a) 随机50局：动作时间分解',fontsize=9.5)
    ax.grid(axis='x',alpha=.16,lw=.4)
    fig.legend(*ax.get_legend_handles_labels(),frameon=False,ncol=4,loc='lower center',
               bbox_to_anchor=(.31,.01),columnspacing=.65,handlelength=.8,handletextpad=.35)
    trs=first_trace(rows,POLICIES,data_dir)
    for i,pol in enumerate(POLICIES):
        tr=trs[pol];aa=actions(tr)
        ts=[float(a['virtual_time_s']) for a in aa if a.get('kind')=='clear' and a.get('clear_result')=='success']
        tend=float(tr.get('metrics',{}).get('exit_s',aa[-1]['virtual_time_s']))
        bx.step([0.]+ts+[tend],[0]+list(range(1,len(ts)+1))+[len(ts)],where='post',
                color=POLICY_COLOR[pol],lw=1.2,ls=['-','--','-.',':','-'][i],label=pol)
        bx.plot(tend,len(ts),'s',ms=3.5,color=POLICY_COLOR[pol],zorder=6)
    bx.set(xlabel='虚拟时间 / s',ylabel='累计清除数',ylim=(-.3,13.8))
    bx.set_title('(b) 固定首局：清除进度',fontsize=9.5)
    bx.text(.97,.96,'(*@\ensuremath{\blacksquare}@*) 为退出时刻',transform=bx.transAxes,ha='right',va='top',fontsize=8.5)
    bx.set_xticks([0,2000,4000]);bx.set_yticks([0,4,8,12]);bx.grid(alpha=.16,lw=.4)
    bx.legend(frameon=False,ncol=2,loc='lower right',handlelength=1.2,columnspacing=.6,handletextpad=.35)
    save_figure(fig,'q3-time-cost-clearance')


def fig5_scenarios_ablation(rows: list[dict[str,Any]]) -> None:
    main=[r for r in rows if str(r.get('split','main'))=='main' and policy_of(r) in POLICIES]
    scenes=['random','annulus','center','hash','worstrecv']
    mat=np.array([[np.mean([time_value(r) for r in main if scene_of(r)==sc and policy_of(r)==p]) for p in POLICIES] for sc in scenes])
    bad=np.array([[any(not val(r,'full_clear',default=False) for r in main if scene_of(r)==sc and policy_of(r)==p) for p in POLICIES] for sc in scenes])
    ratio=mat/mat[:,2,None]
    fig,(ax,bx)=plt.subplots(2,1,figsize=(FIG_W,92*MM),gridspec_kw={'height_ratios':[1.1,1]})
    fig.subplots_adjust(left=.19,right=.89,bottom=.14,top=.92,hspace=.85)
    from matplotlib.colors import TwoSlopeNorm
    norm=TwoSlopeNorm(vmin=min(.7,float(ratio.min())),vcenter=1.,vmax=max(1.5,float(ratio.max())))
    cmap=plt.get_cmap('RdYlBu_r');im=ax.imshow(ratio,cmap=cmap,norm=norm,aspect='auto')
    ax.set_xticks(range(5),POLICIES);ax.set_yticks(range(5),[SCENE_LABEL[s] for s in scenes])
    ax.set_title('(a) 五类场景：平均完整时间（s）',fontsize=9.5,pad=6)
    for i in range(5):
        for j in range(5):
            rgba=cmap(norm(ratio[i,j]));lum=.299*rgba[0]+.587*rgba[1]+.114*rgba[2]
            ax.text(j,i,f'{mat[i,j]:.0f}'+('*' if bad[i,j] else ''),ha='center',va='center',fontsize=9,color='white' if lum<.48 else '#1f2933')
    cax=fig.add_axes([.91,.632,.015,.28]); cb=fig.colorbar(im,cax=cax,ticks=[.75,1,1.5]);cb.ax.tick_params(labelsize=8.5)
    ax.text(0,-.39,'颜色：相对 tour 耗时；*：该组存在未全清局',transform=ax.transAxes,fontsize=8.5)
    base={(scene_of(r),r['seed']):r for r in main if policy_of(r)=='joint'}
    variants=['joint_no_after_service','joint_single_plan','joint_no_cover_polish','joint_no_probability_gate']
    labels=['取消先清后测','单一规划权重','取消测站精修','取消概率筛选']
    values=[];ci=[]
    for p in variants:
        ab=[r for r in rows if r.get('split')=='ablation' and policy_of(r)==p and scene_of(r)=='random']
        pair=np.array([[time_value(base[(scene_of(r),r['seed'])]),time_value(r)] for r in ab])
        if len(pair)!=20:raise ValueError(f'{p}: requires 20 paired ablations, got {len(pair)}')
        values.append(100*(pair[:,1].mean()/pair[:,0].mean()-1))
        rng=np.random.default_rng(918120);idx=rng.integers(0,len(pair),(4000,len(pair)))
        means=pair[idx].mean(axis=1);ci.append(np.percentile(100*(means[:,1]/means[:,0]-1),[2.5,97.5]))
    values=np.asarray(values);ci=np.asarray(ci)
    bx.bar(range(4),values,color='#5b7f99',width=.50)
    bx.errorbar(range(4),values,yerr=[values-ci[:,0],ci[:,1]-values],fmt='none',ecolor='#263846',capsize=3,lw=1)
    bx.axhline(0,color='#767e87',lw=.65)
    for i,v in enumerate(values):bx.annotate(f'{v:+.2f}%',(i,ci[i,1]),xytext=(0,3),textcoords='offset points',ha='center',fontsize=8.5)
    bx.set_xticks(range(4),labels);bx.set_ylabel('耗时增加 / %')
    bx.set_title('(b) 配对20局消融：均值及95%区间',fontsize=9.5,pad=6)
    bx.set_ylim(min(-.8,float(ci.min())-.6),float(ci.max())+1.0)
    bx.grid(axis='y',alpha=.16,lw=.4)
    save_figure(fig,'q3-scenarios-ablation')


def main() -> None:
    global OUT
    ap = argparse.ArgumentParser()
    ap.add_argument("--data", type=Path, default=DATA)
    ap.add_argument("--out", type=Path, default=OUT)
    args = ap.parse_args()
    OUT = args.out
    setup_style()
    obj, rows = load_data(args.data)
    fig1_geometry(rows, args.data)
    fig2_field_sweep_tour(rows, args.data)
    fig3_tour_adaptive_joint(rows, args.data)
    fig4_cost_and_clear(rows, args.data)
    fig5_scenarios_ablation(rows)
    # Full raw trajectories remain available in the frozen evidence directory.
    print(f"wrote five Q3 paper figures to {OUT} (width={160} mm, dpi={DPI})")


if __name__ == "__main__":
    main()

\end{lstlisting}
\subsection{\texttt{scripts/make\_paper\_trajectory\_atlas.py}}
% SHA256 83fe78c797c47a8da0b08b6ee5bfd520b2ccb8db86be95f8453a033dc78944bd
\begin{lstlisting}[language=python,escapeinside={(*@}{@*)}]
"""Render every frozen main-experiment scene as a five-policy atlas page.

This command reads paired_results.json and its existing trace files. It never
loads a policy class, reruns a simulation, or rewrites experiment records.

    python B_locator/Q3/scripts/make_paper_trajectory_atlas.py
"""
from __future__ import annotations

import argparse
import json
from collections import defaultdict
from pathlib import Path

import matplotlib
matplotlib.use("Agg")
import matplotlib.pyplot as plt
from matplotlib.backends.backend_pdf import PdfPages

from make_paper_figures import (DATA, MM, POLICIES, POLICY_COLOR, SCENE_LABEL,
                               draw_route, route_legend, setup_style)

SCENES = ["random", "annulus", "center", "hash", "worstrecv"]


def load_groups(data: Path) -> list[tuple[tuple[str, int, int], dict]]:
    records = json.loads((data / "paired_results.json").read_text(encoding="utf-8"))
    groups = defaultdict(dict)
    for row in records["rows"]:
        if row["split"] != "main":
            continue
        key = (row["scenario"], int(row["index"]), int(row["seed"]))
        policy = row["policy"]
        if policy in groups[key]:
            raise ValueError(f"Duplicate record for {key}/{policy}")
        groups[key][policy] = row
    for key, rows in groups.items():
        if set(rows) != set(POLICIES):
            raise ValueError(f"Incomplete five-policy group: {key}")
    return sorted(groups.items(), key=lambda x: (SCENES.index(x[0][0]), x[0][1], x[0][2]))


def load_traces(data: Path, rows: dict) -> dict:
    traces = {}
    for policy in POLICIES:
        source = Path(rows[policy]["trace_file"])
        if not source.is_absolute():
            source = data / source
        trace = json.loads(source.read_text(encoding="utf-8"))
        if trace["policy"] != policy:
            raise ValueError(f"Trace policy mismatch: {source}")
        traces[policy] = trace
    for field in ("scene_fingerprint", "error_fingerprint"):
        if len({t["scene_manifest"][field] for t in traces.values()}) != 1:
            raise ValueError(f"The five policies do not share {field}")
    return traces


def atlas_page(key: tuple[str, int, int], rows: dict, traces: dict, page: int, total: int):
    scenario, index, seed = key
    fig, axes = plt.subplots(2, 3, figsize=(297 * MM, 210 * MM))
    fig.subplots_adjust(left=.07, right=.97, bottom=.12, top=.82, wspace=.10, hspace=.72)
    for i, policy in enumerate(POLICIES):
        ax = axes.flat[i]
        draw_route(ax, traces[policy], POLICY_COLOR[policy], f"({chr(97+i)}) {policy}",
                   show_ylabel=(i % 3 == 0))
    ax = axes[1, 2]
    ax.axis("off")
    ax.text(.04, 1.12, "图例与读取说明", transform=ax.transAxes, fontsize=11, va="top")
    ax.legend(handles=route_legend(), loc="upper left", bbox_to_anchor=(.015, 1.02),
              frameon=False, fontsize=10, handlelength=1.4, labelspacing=.8)
    ax.text(.04, .33, "同页共用源位、接收半径及误差场。\n"
            "T 为从进入到退出的完整虚拟时间；\n"
            "L 为动作账本记录的总行程。\n"
            "源位仅在退出之后用于绘图。\n"
            "未全清的运行照常保留。", transform=ax.transAxes,
            fontsize=9.5, va="top", linespacing=1.6)
    fig.text(.5, .955, "问题三补充材料：五策略同场景轨迹对照", ha="center", fontsize=16)
    fig.text(.5, .908, f"{SCENE_LABEL.get(scenario, scenario)} · 第 {index+1} 组 · seed {seed}",
             ha="center", fontsize=12)
    fig.text(.07, .025, f"冻结记录：main_{scenario}_{index:03d}_{seed}_<policy>.json",
             fontsize=8.5, color="#626d78")
    fig.text(.97, .025, f"{page} / {total}", ha="right", fontsize=9.5)
    return fig


def main() -> None:
    parser = argparse.ArgumentParser(description=__doc__)
    parser.add_argument("--data", type=Path, default=DATA)
    parser.add_argument("--out", type=Path, default=DATA / "supplementary_trajectories.pdf")
    args = parser.parse_args()
    setup_style()
    groups = load_groups(args.data)
    if len(groups) != 120:
        raise ValueError(f"Expected 120 main-experiment scenes, found {len(groups)}")
    args.out.parent.mkdir(parents=True, exist_ok=True)
    with PdfPages(args.out, metadata={"Title": "Q3 Supplementary Trajectories",
                                     "Subject": "120 frozen paired scenes; five policies per page",
                                     "Author": "", "CreationDate": None, "ModDate": None}) as pdf:
        for page, (key, rows) in enumerate(groups, 1):
            traces = load_traces(args.data, rows)
            fig = atlas_page(key, rows, traces, page, len(groups))
            pdf.savefig(fig)
            plt.close(fig)
            if page % 20 == 0:
                print(f"Rendered {page}/{len(groups)} frozen scenes", flush=True)
    print(args.out)


if __name__ == "__main__":
    main()

\end{lstlisting}
\endgroup
